% 备赛文档：双方共用。
\documentclass[prep,both]{debate}
\motion{辩题：我厌蠢，我错了 / 我厌蠢，我没错}
\meta{赛制}{中国科学技术大学校赛}
\meta{生成日期}{2026-10-12}
\begin{document}
\debatetitle

\begin{summary}
\begin{fields}
\field{持方优劣评估}{反方略劣。五个维度里四个偏正方，只有论据可得性到「偏」，其余都是「略」，定义空间均衡。劣势主要来自论据可得性和论证义务，两者是同一件事的两面：决胜点上的实证（评判无能时起作用的是轻蔑、轻蔑贬品性多于行为、轻蔑可能从一次次的恼火里长出来）在正方手里，反方「划得来吗」上的实证又塌在材料上，没有抵账的路，只能在「看准了吗」和「对他公道吗·尊重」两条线上都正面否认「系统性」。反方的胜机也在同一处：正方的「平常」同样只有媒体描述，论点二最后一步「回头不讲理」是推论，核心实验估的都是陌生人。反方的打法是把「平常」钉在熟人之间、记录在手的场景，逼正方承认它的实验够不着那里；正方要防的是「平常」的举证塌在两头，以及被打成霸凌或要求圣人。}
\thesis{正方}{常人的厌蠢平常凭比较下判断，尺子只往一头偏、自己看不见；就算看准了，劝了还犯，也会从「说了那么多遍」走到「他就是那么蠢」，回头不再讲理：问完三句，这份厌平常的样子该改。}
\thesis{反方}{常人的厌蠢平常冲着教过了还犯，看得有据、怪得公道；它恼他不弃他，收回的是对这一次的评价，作为人的尊重不动，这回放下、下回照样说：问完三句，这份厌平常的样子站得住。}
\end{fields}
\end{summary}

\section{一、赛制要点}
\begin{flowtable}
\block{A}
\flow{1 正方一辩开篇立论}{180 秒}{正一}{定稿：定义、中立判准、两个论点，抛出后场反复追问的问题；判准由正方先说出口}
\flow{2 反方二辩质询正方一辩}{120 秒 双边}{反二 问，正一 答}{判准的第一个战场：反方开取证链，正方当场答出已让出的五处，再说清只争两处}
\flow{3 反方一辩开篇立论}{180 秒}{反一}{套件：开头接住质询；正方用了中立版，一句「大体一致」，省下的时间讲门槛}
\flow{4 正方二辩质询反方一辩}{120 秒 双边}{正二 问，反一 答}{取「看准了就能看低吗」「分寸调到不看低，是不是改了平常的厌」两处表态}
\flow{5 评委打分}{30 秒}{—}{—}
\block{B}
\flow{6 反方二辩驳论}{120 秒}{反二}{先驳：只打正方两处软肋（论点一估的都是陌生人，论点二最后一步是推论），不铺新战场}
\flow{7 正方二辩驳论}{120 秒}{正二}{后驳：套件里有回应反二的选择组；拆「教过了还犯」只赢到准与归责}
\flow{8 正反方四辩对辩}{各 90 秒 单边}{正四、反四}{四辩中场就进入交锋，比每次发言的密度；每一句都要能在口径表里找到}
\flow{9 评委打分}{30 秒}{—}{—}
\block{C}
\flow{10 正方三辩盘问反方二辩、四辩}{90 秒 单边}{正三 问}{抓口径：「先问」是不是平常，「一时」有没有换成「一事」}
\flow{11 反方三辩盘问正方二辩、四辩}{90 秒 单边}{反三 问}{抓「平常」的材料与推论；正方二辩、四辩答案必须一致}
\flow{12 正方三辩盘问小结}{120 秒}{正三}{套件：把反方在盘问里认下的东西落到判准上}
\flow{13 反方三辩盘问小结}{120 秒}{反三}{C 块最后一句：把正方认下的列成清单，只列台上确实认了的}
\flow{14 评委打分}{30 秒}{—}{—}
\block{D}
\flow{15 自由辩论}{双方各 240 秒}{全员}{围绕决胜点：这份厌平常停在事上，还是落到看低人、回头不讲理}
\flow{16 评委打分}{30 秒}{—}{—}
\block{E}
\flow{17 反方四辩结辩}{210 秒}{反四}{先结辩：封路式，按正方全场主打的路线分条件逐条先封，再收价值，不引新论据}
\flow{18 正方四辩结辩}{210 秒}{正四}{最后一发：回应反四的选择组，另留约 60 字临场位；预留时间正面回答对方的奇袭}
\flow{奇袭（任意大环节结束后，每方一次）}{质询 150 秒双边 / 申论 120 秒}{见下}{两种都要备，上场再选；插入所在块计分}
\end{flowtable}
质询双边计时：被质询方拖延就是在吃质询方的时间，所以问题全部写成封闭式，每问不超过 25 字，对方说超过 15 秒还没回到问题就礼貌打断。对辩、盘问单边计时，比的是每次发言的密度；盘问时被盘问方回答是否计时，以当天公布为准，赛前向组委会确认。二辩同时负责质询与驳论，四辩同时负责对辩与结辩，是负荷最重的两个位置。

\begin{fields}
\field{三票制}{评委小组赛与淘汰赛第一轮 3 人、半决赛 5 人、决赛 7 人。每位评委三票：\sub{环节票}{每个环节单独打分，总分高的一方得票；没有可以放弃的环节，盘问、对辩都是净得分点。}\sub{印象票}{交完环节票立刻投，凭全场整体印象，可分 0.5:0.5；全场姿态、口径一致、是否正面回应都被算进去。}\sub{决胜票}{述票之后投，不可分；总比分 4.5:4.5 时由它定胜负。述票从少数派开始，每人不超过 7 分钟，只陈述理由。每场 1 名最佳辩手。}}
\field{奇袭}{双方可在任意大环节结束后发动，每方一次，二选一：奇袭质询 150 秒双边，或奇袭申论 120 秒单方。\sub{怎么选}{对方刚说出一个可以钉死的表态，用质询当场取证；我方论点被打塌或有一个对方一直没碰的战场，用申论，稿子提前写好、与既有框架接得上。}\sub{四个窗口}{①对方一辩立论之后，框架最脆，适合质询钉定义或判准；②对辩之后，抓四辩临场的口径不一致；③盘问小结之后、自由辩之前，用申论给自由辩定议题；④自由辩之后、结辩之前，最狠：对方必须在结辩里回答。不要在开场用掉，也别留到反四结辩之后。}\sub{本题}{反方首选窗口④发动奇袭质询，逼正四把最后一发的一部分时间花在回答「平常还有什么材料」「回头不讲理是不是推论」上（链首问见第九节）；正四要为这两问预留时间。}}
\field{论据红线}{章程第十九条第四款：对 AI 生成的假数据零容忍，一旦出现，印象票判负；第十七条第三款：虚假论据直接判该场负。对方可在评委述票结束前当场提出伪证申诉，评委有权要求出示来源；假证申诉赛后书面提交，每队每赛季一次。所以：只有「已核实」进文稿与速查；「有把握」不进文稿，【待核实】连速查也不进；每条数据都要能当场报出机构、年份、报告名；不报任何厌蠢比例（调研里没有这类数据，\#102）；标〔智〕的材料测的是智力或认知能力，用之前先搭桥到「蠢」。}
\field{赛程纪律}{辩题在比赛当周周一公布，备赛窗口只有几天；弃赛与改期须在当周周三 24 时前提交。章程第九条正文的日期与流程表对不上，赛前请领队向组委会确认一次。}
\end{fields}

\section{二、辩题性质与争议焦点}
\begin{fields}
\claim[辩题性质]{价值判断题：评的是一种态度对不对，里面夹着两个事实子问题；不是利弊题、政策题，也不是比较题。「我厌蠢」是双方共同承认的前提，要争的只有后半句；「错了」与「没错」是矛盾关系，没有中间项，谁也不占推定，胜负几乎全看判准和「错」的门槛。}
\field{两个事实子问题}{①「我觉得他蠢」准不准，属于认知层；②这份态度平常停在事上，还是落在看低人那一头，以及它稳定地带来什么。双方都会把它们拉进来当论据，题目性质不变。学理上这正是态度恰当性问题：情绪对不对得上对象（契合）是一问，通盘考虑下该不该怀着它是另一问，两问互相推不出（\#164、\#59、\#60）。}
\field{正方倾向把题目解释成}{规范题：这种态度照它平常的样子，该不该改。看厌蠢平常的样子与内在倾向（整体化、贴标签、把人划出讲理的范围），门槛放在「平常的样子有应当修正的缺陷」。\sub{理由}{中文「我错了」的日常语义是认错、该改，不是自认坏人；第一人称题里，厌蠢者自己也认可的理由最有说服力。}\sub{好处}{看错人的认知证据、伤关系伤合作的后果证据，都能算进「错」里。}\sub{风险}{一旦滑成「厌蠢弊大于利」的利弊题，就得算净效果，反方一句「有弊不等于错」就能顶回来。}}
\field{反方倾向解释成}{三条路，学理强度依次递减。\sub{①恰当性题}{「我的厌对不对得上真实存在的蠢」，只看契合，并把评价对象收窄到那个有分寸、针对屡教不改、不说出口的「我」。最强的一条。}\sub{②权利题}{「我有没有权利厌、有没有权利远离」。毛病是「有权」回答的是别人能不能干涉，回答不了这份态度本身对不对。}\sub{③不可评价题}{「情绪不经选择，无所谓对错」。毛病是「无所谓对错」不等于「没错」，被引用的这批作者自己也没有推出「所以无可指摘」。}\sub{风险}{评委会觉得反方换了评价对象；走第③条，还和「我没错」的正面姿态自相矛盾。}}
\field{争议焦点}{\sub{战场一：厌的是事还是人}{是对一件蠢事的恼火，促使人去纠正、交涉；还是对一个人的看低，让人走开、不再讲理？最清楚的界线是 Roberts 转述的 Bell：轻蔑针对人、促使退出，怨恨针对行为、促使交涉（\#160）。「走开」双方共用 Strawson 这条线：一时抽身躲开投入之累，正常人可以用，但不能长久（\#34）；长久下去，就是把人当作只能回避、不能与之讲理的对象（\#33）。正方要证厌蠢平常会落到对人，反方要证它平常停在事上、走开只是一时。}\sub{战场二：被厌的蠢该不该由他负责}{是不会、不懂，还是想得到却不想、听得进却不听？缺努力的失败比缺能力的失败招来更多惩罚（\#58）。还连着一个认知问题：厌的人当场分不分得清这两种情况。}\sub{战场三：「错」的层次与门槛}{看人不准（认知）、待人不公（道德）、对事对己划不来（实用与对己），哪一层算数？门槛是「该改」还是「该受谴责」？「平常的样子该改」和「个别场合调分寸」怎么分？双方共用的分界规则见第四节。}}
\field{评委直觉}{主流媒体几乎一边倒地批评厌蠢（\#10、\#11、\#13、\#117、\#118），正方有舆论顺风，论点也最容易被对方预判。本题在 2026 年新国辩复赛打过一场，战报记 4:5，反方 5 票胜（\#113）；一场比赛的票数说明不了持方优劣，只说明反方并非打不赢。}
\end{fields}

\section{三、关键词定义}
\subsection{3.1 厌（厌蠢）}
\begin{fields}
\field{调研}{\sub{词典义}{汉典「厌」现代义「嫌恶，憎恶」，本义「吃饱」（\#1）；「厌恶」：对人或事物产生强烈的反感（\#2）；「嫌弃」：厌恶而不愿接近（\#4）；「厌烦」：对某事感到厌恶不耐烦（\#6）；「鄙视」：轻视，看不起（\#5）。台湾地区《重編國語辭典修訂本》释「厭」为「憎惡、嫌棄」，与汉典并引（\#3）。词典义里没有「看不起」，那是「鄙视」「轻蔑」的释义。《现代汉语词典》原文没有核到（\#7，【待核实】），场上只说「汉典释……」。「学而不厌」的厌是满足义，与本题无关（\#150）。}\sub{热词义}{央视网：「看不得别人犯蠢、没常识、太笨」（\#10）；新京报：「看不惯别人没常识、太笨」（\#118）；中国网·心理中国：不是正式医学术语，是对他人「相比自己而言显得无知的言行感到烦躁、难以忍受、厌恶的情绪状态」，有的人或许不说出口，「但会用行动离对方远远的」（\#16）。使用者多以自己为主语，用「症」字自嘲，还常自我限定，比如「我只厌打扰到我的蠢」「遇见不会但努力的人，并不会激发她的厌蠢症」（\#14、\#9）。}\sub{学术义}{中文心理学把 disgust 译作「厌恶」：它的倾向是把自己与污染隔开，contempt 更像向下看（\#21）。contempt 的核心是看低对方，并试图忽视或排斥对方（\#19）；它把对方贬为愚蠢、无能，贬的是品性多于行为（\#20）；评判「无能行为」时起作用的是 contempt（\#26）；contempt 只指向人，含优越感、不看好对方改善、疏离（\#27）；Bell 笔下的轻蔑有整体性、轻慢、比较、退出四个特征（\#159）。两者在日常体验中交织，表情识别里常被混淆（\#24）。}\sub{情绪与立场}{亚里士多德说「我们感到愤怒和恐惧不经选择」，但受褒贬的是「以某种方式」去感受（\#151）；重评降低了厌恶体验（\#154）；换位思考是调节效果最大的一项，压住感受几乎无效（\#155）；Adams 承认愤怒不能像坐下、站起那样选，仍认为「对他人的轻蔑」可以是过错（\#156）；康德说比较之下在心里看低他人有时不可避免，但表现出来就是冒犯（\#48，中文据德文自译，赛前对照李秋零译本）。}\sub{小结}{字面上的「厌」是嫌恶、反感加上不愿接近；「厌蠢」在实际用法和心理机制上，常常夹带看低。学理上它是一条光谱：一端是对蠢事的反感，另一端是对蠢人的看低。定义之争的关键是「看低」算不算在里面。辩题写的是「厌蠢」，没有「症」字：正方别拿「症是要治的」当论据，反方也别拿「症」字论证「不由自主」。}}
\begin{procard}{正方有利定义}\definition{厌蠢，是因为认定一个人蠢，就看低他这个人的态度。}轻蔑版。\sub{有利之处}{直接接上轻蔑的结构：看低、整体化（\#19、\#20、\#159），对无能行为的评判本来就有轻蔑参与（\#26）；轻蔑面临的四条道德反对都能用上：整体化几乎总会错误呈现对象、违背应得的尊重、向下看损害关系、退出交流（\#161）。它也让「厌」天然成为立场，和「我厌蠢」作为习惯性陈述的语法一致。}\sub{为什么不算定义杀}{看低一个人对不对仍然可以争：轻蔑与承认尊重相容（\#37），按品质、成就区别评价人不违背尊严（\#44），康德把「比较之下在心里看低」和「拒绝给人应得尊重的蔑视」分开（\#48）。不能再往里加：写成「看低他、不愿再认真对待他」，后半句就是拒绝基本尊重（\#48、\#33），那才是定义杀。}\sub{反方如何反驳它}{词典里「厌」只有反感、嫌恶、嫌弃，「看不起」属于「鄙视」（\#1、\#2、\#4、\#5）：「题目写的是厌，不是鄙视」；正方爱引的 Hutcherson 与 Gross，摘要写的是「无能的行为」（\#26）；网友自己会划界限（\#14、\#9），正方把最坏的版本当成了常态。}\sub{正方退守}{不把「看低」写死在定义里，退回中立定义的「可以只冲着那件事，也可以冲着这个人、看低他」，再在论证里证明厌蠢平常落在看低人那一头。}\end{procard}
\begin{concard}{反方有利定义}\definition{厌蠢，是碰上别人犯蠢时自然生起的那股烦：反感、不耐烦、想离远一点，并不等于看不起他。}反感—回避版。\sub{有利之处}{把厌钉在「自然生起」的情绪上，接上「情绪不经选择」（\#151 前半句）；明确排除「看不起」，对应 disgust「把自己隔开」的结构（\#21），正方那一套轻蔑材料就落了空；接上媒体里偏功能的一路解读：「厌蠢」很多时候是「厌麻烦」（\#12），是「对潜在风险的厌烦」（\#15，只引明确写作医生原话的这一句）。它仍以「别人犯蠢」为对象，不是「只厌蠢事」那种定义杀。}\sub{正方如何反驳它}{题目是「厌蠢」，不是「厌烦」，也不是「厌麻烦」：「厌烦」的对象是「某事」（\#6），通行说法落在人身上（\#10、\#118），「讨厌麻烦叫怕麻烦，用不着『厌蠢』这个词」。排除「看不起」是硬切：对能力的评判有轻蔑参与（\#26、\#20），两者在体验里本来就混在一起（\#24），中国网写厌蠢者会「用行动离对方远远的」（\#16）。「不经选择」只引了半句：同一段亚里士多德紧接着说受责备的是以某种方式感到愤怒的人（\#151），Adams 的结论是存在「不自愿的罪」（\#156），康德那句后半是「表现出来就是冒犯」（\#48）。新京报书评那篇后半转向批判：这个标签「给了我们合理化贬低他人的契机」（\#13）。disgust 指向人时，去人性化的色彩更重（\#21）。}\sub{反方退守}{不说「情绪无对错」，改说「情绪可以被评价，而这份厌经得起评价」；退回中立定义，主打「厌而不看低，一时走开也不把人划出讲理的范围」。}\end{concard}
\begin{neutralcard}{中立定义}口播版：「厌蠢，是觉得别人犯蠢时生出的嫌恶、反感，常带着想躲开的心；它可以只冲着那件事，也可以冲着这个人、看低他；它是一种态度，不等于骂人或欺负人。」\sub{合理性}{贴近词典：「嫌恶」「反感」「想躲开」分别对应汉典「厌」「厌恶」「嫌弃」（\#1、\#2、\#4）与台湾辞书的「憎惡、嫌棄」（\#3）；不写「强烈」，强弱留给场上；不写「不耐烦」，它来自「厌烦」，对象是「某事」，放在前半句等于替反方选了对事那一端。用「觉得」不用「看到」：写成「看到别人犯蠢」，等于先认定对方真蠢，误判被挤出题目，「看准了吗」一问落空；写成「觉得」，误判留在厌蠢里，看得准的厌同样在里面。贴近热词：评委一听就知道是网上那句「我厌蠢症犯了」（\#8、\#10），「想躲开」对应「用行动离对方远远的」（\#16）。两端并列、都用「可以」，不暗示哪一端是本体；「常带着」说明想躲开是常见倾向，不是必带成分，更不是已经走开。骂人、欺负人是行为层面的伤害，写进定义就把「错」预埋进题目，所以排除；这份态度稳定带来什么，判准里照样要算。}\sub{双方都有得打}{正方打「这份厌平常落在看低人那一头，从想躲开滑到把人划出讲理的范围」；反方打「它平常停在那件事上，走开也只是一时」。第一秒的厌不由自主，双方都承认，所以定义里既不写「自然生起」，也不写「选择」。}\sub{质询守法}{反二问「也可以冲着这个人看低他，是不是把鄙视塞进了厌」：「两头都写了，都用『可以』，没替任何一方选。平常落在哪一头、落在那里错不错，是今天要辩的。」问「觉得犯蠢，那看错了也算厌蠢」：「算。厌蠢说的是你的厌；看没看准，是第一问要称的，看得准的厌也在里面。」正二问「把骂人切出去，剩下的厌还有什么可评」：「切出去的只是霸凌。想躲开的心、冲着人的看低都还在定义里，欢迎来评。」双方都会被问「只在心里厌、从不表现出来，算不算厌蠢」，统一答：「算，心里那份就是厌蠢本身；它错不错，看判准。」正方千万别答「不算」。}\sub{退守线}{正方：「我们不要求你第一秒不厌。我们说的是，你第二秒以后还认同、放任那份看低人的厌，就错了。」再退：「就算只在心里，轻蔑也会从表情和线索里漏出来，对方感觉得到。」（\#156，一派看法）反方：「我们不说情绪无对错。我们说这份厌经得起评价。」再退：「心里看低别人，连康德都承认有时免不了，表现出来才是冒犯，而我们不表现。」（\#48 必须整句引；康德另把拒绝给人作为人应得尊重的蔑视说成无论如何违背义务，反方的厌必须停在「比较之下的看低」这一边。）}\end{neutralcard}
\end{fields}

\subsection{3.2 蠢}
\begin{fields}
\field{调研}{\sub{词典义}{汉典「蠢」释「愚笨，笨拙」，古义还有「不逊」（《尔雅·释训》「蠢，不逊也」）（\#66）；「蠢人」：缺乏判断力或谋虑的人，「愚蠢」：愚笨无知，「无知」：缺乏知识或不明事理（\#67）；「笨」偏理解、记忆能力不佳，比「蠢」更中性；《荀子》有「非是是非谓之愚」。蠢偏判断与行为，本身带骂人色彩。《现代汉语词典》原文没有核到（\#68，【待核实】）。}\sub{热词义}{「看不得别人犯蠢、没常识、太笨」（\#10）；「看不惯别人没常识、太笨」（\#118）；网友的自我限定：讨厌的是「蠢得理直气壮还喜欢帮倒忙的人」（\#9）；吉林工商学院一场院赛里，主张「厌蠢没错」的一方说「厌蠢是对认知惰性的理性批判」（\#115）；中大附中的国旗下讲话说「『蠢』并没有明确的标准」（\#117）。}\sub{学术义}{人们把三类情境叫作蠢：信心与能力失衡、注意失误、缺乏控制；行为人被认为责任越大、后果越严重，就越显得蠢（\#93；各类比例与样本数只到「有把握」，不报，\#69）。理性与智力是两回事，「理性障碍」指智力足够却不能理性地思考和行动（\#70）；大量思维偏差与认知能力不相关（\#71、\#72）；聪明人也会做惊人的蠢事（\#78）。朋霍费尔说愚蠢「不是智力上的缺陷，而是人的缺陷」，人在某些条件下会「被弄蠢，或任由自己被弄蠢」（\#75；中文流传本作「道德缺陷」是意译）。穆齐尔说骂一句「蠢」是「回应的开始，也是它的结束」（\#77）。}\sub{归因}{缺努力导致的失败比缺能力导致的失败招致更多惩罚（\#58）；反应态度回应的是他人意志的好坏，不是能力高低（\#32）；从完全可以由情境解释的行为推出持久的特质，叫对应偏差（\#61）。}\sub{成因}{〔智〕智力遗传率随年龄上升（\#57，不要说「八成由基因决定」）；可遗传不等于不可改变（\#80）；每多受一年教育，认知能力约提高 1 到 5 个 IQ 分（\#82）；具备科学素质的比例随受教育程度陡升（\#83，相关不是因果）。这一组测的是智力或科学素质，用来说「蠢」之前要先搭桥。}\sub{小结}{日常说的「蠢」，是对一个人在某件事上言行的评价，不是智商鉴定；既包括不会、不懂，也包括没用心、自以为是。核心争点：不会还是不肯，对事还是对人。}}
\begin{procard}{正方有利定义}\definition{蠢，是在「我」看来，别人在知识、能力或判断上的明显不足；不会、不懂、没想到都算，原因不限。}\sub{有利之处}{「不会、不懂」都在范围里，热词的原始场景（「坐高铁不是谁都会吗？」，\#119）完整进场，信息差（\#120、\#121）、知识的诅咒（\#253）、对应偏差（\#61）、角色优势造成的误判（\#254）都能用；「在『我』看来」把判断权落在厌者身上，「蠢」成了一张贴上去的标签；「原因不限」让运气与应得进场：没有人应得他更高的自然能力（\#52），没人责备天生的盲人（\#41），「自然原谅缺陷时，没有人发怒」（\#46）。}\sub{反方如何反驳它}{把「蠢」等同于「无知」「笨」是偷换：汉典里「无知」重知识，「蠢人」重判断力与谋虑（\#67），「不知道叫无知，不肯想才叫蠢」；「在我看来」把结论预埋进定义，正方赢下的只是「你可能看错了」；日常说的蠢也包括不肯学，「困而不学，民斯为下矣」（\#49）；蠢也不等于智商低（\#70、\#75），拿遗传和 IQ 讨论「蠢」是跑题。}\sub{正方退守}{放弃「在我看来」，退回中立定义的「一般人看来本该」，再在论证里证明厌的人常把「我这样的人」当成了一般人（\#254、\#253），放到「看准了吗」去打。}\end{procard}
\begin{concard}{反方有利定义}\definition{蠢，是在本可以想明白的事上没去想，或者不懂又自以为懂，以致言行明显失当。}\sub{有利之处}{落在汉典「蠢人：缺乏判断力或谋虑的人」（\#67）和《荀子》「非是是非谓之愚」上，把单纯的「笨」和「无知」分了出去；把蠢放进可以归责的范围：缺努力招致惩罚（\#58），因疏忽而来的无知可以受责备（\#41），反应态度本来就回应「冷漠、缺乏关切」这种意志品质（\#32），朋霍费尔说人会「任由自己被弄蠢」（\#75），人们最常叫作蠢的一类是信心与能力失衡（\#93），认知反思越低越容易信假新闻，原因是「懒」而不是「偏」（\#107）。它没有写成「理直气壮、损人利己」，没有把蠢偷换成坏，不进定义杀。}\sub{正方如何反驳它}{日常义的蠢是「愚笨」，汉典「愚蠢」就释为「愚笨无知」（\#66、\#67），无知本来就在蠢里面；厌蠢症的原始对象正是「没常识、太笨」（\#10、\#118），出圈事件是教人坐高铁（\#8、\#119）。「不肯想」和「想不到」当下分不清：你以为的常识常常是知识的诅咒（\#253），从完全可以由情境解释的行为推出持久特质是对应偏差（\#61），很多人转发错误信息是因为注意力没放在准不准上（\#108）。常识带着阶层和地域的印记（\#120）。反方引朋霍费尔，正方补下一节：他写完〈论愚蠢〉，紧接着就是「蔑视人？」，说蔑视一个人就永远不能使他有所成就（\#76）。顺势要对方让步：「那请承认，对真正的笨，厌是错的。」}\sub{反方退守}{不把「不会」整个排除，退回中立定义的「一般人看来本该想得到」，第一次、没人教的不会自然不在其中；对剩下那些本该会的，再主张厌是对得上的。}\end{concard}
\begin{neutralcard}{中立定义}口播版：「蠢，是在一般人看来本该想得到、做得对的事上，明显想错、做错；可能出于不会，也可能出于不肯，原因不在定义里预设。」\sub{合理性}{贴近词典：汉典「蠢」兼有「愚笨」（能力）和「缺乏判断力或谋虑」（判断）两层（\#66、\#67），「不会」「不肯」都不能单独拿掉；三类情境（信心与能力失衡、注意失误、缺乏控制）都在常人所说的蠢里，每一类都可能出于不会，也可能出于不肯（\#93）。贴近常识：尺度用「一般人看来本该」，就是媒体说的「没常识」（\#10、\#118）；这里的一般人指了解他处境、讲道理的一般人，不是评论区的多数票，口播不加字，被问到时这样答。不用 IQ（蠢不等于低智商，\#70、\#75），也不用「我看来」。「明显」二字排除鸡毛蒜皮。不写「天生」（正方的定义杀），也不写「明知故犯」（反方的定义杀）；成因留给论证。}\sub{对事还是对人}{定义里的蠢描述的是想法和做法，也就是事；把它说成「这个人蠢」，是一种推断。这个推断准不准、该不该，是辩论要打的内容，不在定义里决定。}\sub{双方都有得打}{正方打「厌的人常把『我这样的人』当成一般人」、常识带着阶层和地域的印记、蠢常常出自对方没选过的处境、厌会从事滑向人；反方打「本该」意味着可以归责、「明显」意味着不冤枉人、常识是讲道理的一般人的公共标准。}\sub{质询守法}{正方陈述后，反二问「你们也承认不肯是蠢，那厌不肯错在哪」：「错不在你对那件事有反感，而在这份厌平常就从事滑到人，从这一次滑到这一辈子。」反方陈述后，正二问「第一次不会坐高铁算不算蠢」：「不算。定义里的一般人，是了解他处境、讲道理的一般人；第一次、没人教，这样的人不觉得他本该会。我们的定义不冤枉人。」追问「那评论区为什么说『这不是谁都会吗』」：「评论区的多数票不是定义里的一般人，那是评论区看错了。他们那份也算厌蠢，看错是第一问要称的；但这是个别场合的误判，不是这份厌平常的样子。」这一步把战场引到「平常还是个别」，正方接着要证明这类误判是普遍倾向（\#61、\#254）。}\sub{退守线}{正方：「就算蠢里有不肯，我们也不必证明每一次厌都错；只要证明这份厌平常就对人不对事、把人一概而论、冲着对方没选过的东西，就够了。」反方：「就算一部分蠢是不会，对不会我们同样主张耐心；我们守的是对『本该会却明显做错』的那份反感。」}\end{neutralcard}
\end{fields}

\subsection{3.3 错了 / 没错}
\begin{fields}
\field{调研}{\sub{词典}{调研没有收「错」的词典条目。按日常语感，「我错了」是承认「这样不对、该改」，「没错」是「站得住、不用改」；这一条凭语感，赛前查《现代汉语词典》「错」字条，查到之前不说「词典释」。「我错了」不等于「我是坏人」。}\sub{「错」的类型}{恰当性之错（呈现失准：对方其实不蠢，或者蠢不是该厌的那类性质）、道德之错（违背尊重、对运气不公平）、审慎之错（代价大于收益）、对自己的错（情绪本身的代价）。SEP「Emotion」区分情绪的契合与策略、道德上的恰当（\#164）；D'Arms 与 Jacobson 把从「感到它是错的」推出「它不恰当」称为道德主义谬误（\#59）；「错误种类的理由」：恶魔可以逼你「鄙视一个并不可鄙的人」，这个理由与对方是否真的可鄙无关（\#60）。}\sub{情绪可以有对错}{在恰当的时候、对恰当的对象、以恰当的方式去感受，才是德性（\#152）；「发而皆中节，谓之和」（\#158）；该怒而不怒被视为愚蠢（\#39）；Adams 认为不自愿的轻蔑仍可以是过错（\#156）。「情绪无对错」在哲学上不是共识。}\sub{改不是硬压}{重评降低了厌恶体验，压抑反而提高交感神经激活（\#154）；换位思考 d=0.45，压住感受 d=−0.04（\#155）。}\sub{门槛}{正方要低门槛（错＝应当修正），反方要高门槛（错＝应受谴责）。}}
\begin{procard}{正方有利定义}\definition{「我错了」，是说这份厌有该改的毛病：常常看错人，或者亏待了被厌的人，占一条就该改。}低门槛，任一层即可。\sub{有利之处}{门槛定在「该改」，贴合中文认错的语义，正方不必证明厌蠢者品格恶劣；只要赢下认知或公平其中一层就够，反方「厌有用」的好处救不了「看错人」；对应偏差（\#61）、知识的诅咒（\#253）、问答者实验（\#254）、运气与应得（\#52、\#58）都能直接用上。}\sub{反方如何反驳它}{门槛太低，接近定义杀：「照这个标准，世上没有不错的情绪」；道德主义谬误反过来同样成立，某一层有瑕疵推不出通盘该改（\#59、\#164）；「常常看错人」说的是误判的个案，看得准的厌又怎么算；该怒而不怒会被看作愚蠢（\#39）。}\sub{正方退守}{接受通盘衡量，但主张「看准了吗」「对他公道吗」两问分量最重：只要这份厌平常的样子在这两处系统性地站不住，而它在「划得来吗」上的好处换一种态度也拿得到，它就该改。}\end{procard}
\begin{concard}{反方有利定义}\definition{「我错了」，是说这份厌理亏、该受指责：要么厌了根本不该厌的人，要么因此伤了人；只是心里不舒服，谈不上错。}高门槛，只认道德层和行为。\sub{有利之处}{心里的厌几乎不可能「该受指责」：据 SEP 概括，康德认为我们没有「感到」尊重他人的道德义务（\#43），比较之下在心里看低他人有时不可避免（\#48），亚里士多德说我们不因情感本身受褒贬（\#151 前半句）；把正方的审慎材料挡在「错」之外；保留了「厌了根本不该厌的人」，不是定义杀。}\sub{正方如何反驳它}{门槛太高：日常说「我错了」只到「该改」，「看错了人，也叫错」；「伤了人」把评价对象从「厌」换成了「霸凌」；心里的态度也有对错：亚里士多德那句后半是受责备的是以某种方式去感受的人（\#151），Adams 点名对他人的轻蔑（\#156），康德同一句后半是「表现出来就是冒犯」、同一节还说拒绝给人作为人应得尊重的蔑视无论如何违背义务（\#48），尊重义务虽不要求「感到」尊重，却包括不持有轻蔑、嘲笑这类态度（\#43）；在心里冤枉人也是冤枉（\#60）。}\sub{反方退守}{接受通盘衡量，但主张「错」要落在这份厌平常的样子，而不是它冒出来的个别场合；误判和过头只在个别场合出现，调分寸就好。}\end{concard}
\begin{neutralcard}{中立定义}口播版（与第四节中立判准口播版是同一句，全队只背这一句）：「判断这份厌对不对，就看一个清醒的我问完三句：看准了吗？对他公道吗？划得来吗？问完觉得这份厌平常的样子就该改，就是错了；站得住，只在个别场合调调分寸，就是没错。」\sub{多义怎么处理}{不挑某一种「错」，把四种放进同一组三问：认知之错在「看准了吗」；道德之错在「对他公道吗」（前一小问管归责，后一小问管尊重，「他不肯学」只回答了前一小问；就算看准了这个人，看低他还守不守得住作为人应得的尊重，要在这里再问一次，\#48、\#37、\#44）；实用之错和对自己的错在「划得来吗」。只赢一问不自动赢下整场：契合不等于道德上恰当（\#164 的例子：性别歧视的笑话也许确实好笑，被它逗乐却不恰当），有代价也不等于不契合（\#59）。}\sub{门槛分界规则（双方共用）}{「要改的如果是这份厌平常的样子，就是这份厌本身该改，哪怕改完还剩一点反感；要改的如果只在个别场合冒出来，那叫调分寸。平常就这样、只是改得掉，照样算错了。」最后一句堵住「什么毛病都改得掉，所以都只是校准」：可以校准不能当分界，分界只看平常还是个别。这是本题最实战的一条线。}\sub{修正不等于压抑}{要求的是换一种看法（换位思考 d=0.45），不是「不许有感觉」（压住感受 d=−0.04）（\#155）。反方不能用「我控制不了」免于评价，正方也不能要求对方「别有感觉」。}\sub{合理性}{「我错了」就是该改，「我没错」不等于「我完美」；门槛在中间：比「有点瑕疵」高，比「该受谴责」低；分界看「平常的样子」，正对应「我厌蠢」是一句习惯性陈述。情绪对不对，看是否中节、是否以恰当的方式发出（\#158、\#152），本来就是通盘衡量；Bell 的七条恰当性条件也是把契合、证据、道德、公平放在一起看（\#36）。既不是「占一条就算错」，也不是「不伤人就没错」。}\sub{质询守法}{正方陈述后，反二问「调整分寸也是改，那岂不是一点不完美就算错」：「不是。只在个别场合冒出来的，调分寸就好；要改的是这份厌平常的样子，才叫错。今天我们要证明的是后者。」反方陈述后，正二问「分寸调到不看低他，这是改了你平常的厌，不是调分寸吧」：「看低不是这份厌平常的样子，只在个别场合冒出来，所以是调分寸。平常那份反感冲着事，它告诉我标准被踩了，我们守的就是它。」}\sub{退守线}{正方：「就算不是每一问都赢，只要这份厌平常的样子在看人和待人上系统性地站不住，而它带来的好处换一种态度也拿得到，它就该改。」反方：「就算这份厌有时看错、有时过头，那也只在个别场合，不是它平常的样子；该做的是调分寸，不是放下。」}\end{neutralcard}
\end{fields}

\subsection{3.4 隐含要素}
\begin{fields}
\field{主体}{评价者与被评价者是同一个「我」。取「一个普通的、会说『我厌蠢』的人，回头审视自己这份态度」：不是圣人，也不是网暴者，也不预设他是亚里士多德笔下的大度者；他给自己划的界限守不守得住，由双方去证。质询里被问「你说的是哪个我」，统一答：「一个普通的、会说出这句话的我。」受影响的是被厌的人、我自己，以及我所在的关系和团队。}
\field{「我」带来的后果}{\sub{习惯性陈述}{「我厌蠢」像「我怕狗」，说的是一种稳定的倾向，评价对象天然是这份倾向平常的样子，不只是一次情绪。}\sub{内在理由}{第一人称题里内在理由最有说服力：正方要证这份厌违背了你自己看重的东西，反方要证这份厌正是你的标准在起作用。}\sub{对称优势}{正方评的是你这份厌，不是你这个人，这正是正方要你对「蠢人」做到的；反方反咬「你们说我错，不也是在厌我」，正方用 Roberts 的区分回应：针对行为、促使交涉的是批评，针对整个人、促使退出的是轻蔑（\#160）。}\sub{中科大语境}{台下多是高学历者。Kuppens 等发现高学历者对低学历者存在基于学历的群际偏见（\#54，西方样本，测的是学历不是蠢）；新京报书评周刊认为被认为有「精英」身份的群体反而最容易有厌蠢症（\#13，作者论断，只能说「新京报书评周刊认为」）。正方可以借此把「我」拉近评委，反方不能有一个字的精英腔。}}
\field{范围}{当代中国的日常人际：同学、同事、家人和网上。「厌蠢症」一词约在 2023 年元旦前后出现（\#8）。两类情形不当主战场：医疗事故、航空安全这类专业失误（\#104、\#105、\#106）只作「错在系统，不在蠢人」的机制例证；教师、医生这类负有额外义务的职业角色（\#147、\#148）可以举例，但对方一句「那是职业要求」就能切掉。}
\field{时间尺度}{一次厌（当下几秒）和一份厌（平常的样子）不同。双方都承认第一秒不由自主，真正被评价的是第二秒以后常人实际怀着的那份厌：他认同它、放任它，还是修正它。正方要把战场放在这份厌平常的样子，反方要把毛病落回一次次具体场合；两者怎么分，见门槛分界规则。}
\field{比较对象}{「没错」不等于「最好」：存在更好的态度，不等于这份厌错了。替代态度只回答「改得了吗、改了亏不亏」，本身不构成「错」：正方先证出这份厌平常的样子有系统性缺陷，替代才登场；这时反方只需证明替代做不到，或者换了会丢掉这份厌换不来的好处。替代不能是「压着不厌」，双方都排除。可选的替代有四种：只对事的批评（怨恨型，促使交涉，\#160）；换位与重评（\#154、\#155）；对不会但态度好的人的怜悯（刻板印象内容模型的假设，\#29）；一时抽身、回头还肯讲理（\#34、\#33）。}
\end{fields}

\begin{warnbox}{3.5 定义杀陷阱：双方都要背}
\begin{fields}
\begin{procard}{正方的定义杀陷阱}最容易犯的是把伤害写进定义（「厌蠢＝蔑视并伤害蠢人」）；其次是把「不愿再认真对待他」写进定义，那就是拒绝基本尊重（\#48），等于把「错」预埋进去；再次是把「蠢」限定为「天生能力不足」，反方一句「困而不学，民斯为下矣」（\#49）就能拆掉；最后是把门槛压成「有任何不完美就算错」，让对方怎么打都输。评委会觉得正方不敢面对对方最好的版本：只在心里厌、从不表现出来的我，错了吗？反方一句「今天讨论的是厌，不是霸凌；我只远离，不伤害」就能切掉。\sub{实战规矩}{任何环节都不说「厌蠢就是霸凌」「厌蠢就是网暴」；伤害、退出、漏出来的轻蔑只当作厌蠢稳定带来的后果去论证（\#156、\#160、\#25），不当作定义去宣布。被问「只在心里算不算」，答：「算。心里看低他，问完三句同样站不住；而心里的厌，平常就是冲着人去的。」不要答成「心里对事的反感也错」，那和「对那件事生气、批评都可以」自相矛盾。}\end{procard}
\begin{concard}{反方的定义杀陷阱}①「厌只是生理反应，情绪无对错」：技术上「无法评价」也算没错，但评委会觉得反方在躲，还和「我没错」的正面姿态自相矛盾。②「蠢＝明知故犯、拒不改正、损人不利己」（「又菜又犟」）：把蠢偷换成坏，厌坏人叫义愤，不叫厌蠢；正方用「坏人不是蠢人」（\#111）和英语里常说的汉隆剃刀（\#130）指出偷换，还能顺势要反方承认「对真正的笨，厌是错的」。③「厌蠢＝只讨厌蠢事，与人无关」：拿掉了指向人的那一面，而那正是全部争议所在，讨厌错误本来就没人反对。门槛杀是「错＝违法或伤人」。隐蔽的一种：把厌蠢写成「看到别人犯蠢时」的反应，等于先认定对方真蠢，「看准了吗」一问被整个拿走。评委会觉得反方是在为「厌坏」或「讨厌错误」辩护，换了题。\sub{实战规矩}{可以拿「又菜又犟」当例子，但定义里要保留「不会」，「多数被厌的蠢是想得到却不想」放到论证里去证；可以主打「厌而不伤人」，但不说「情绪无对错」。}\end{concard}
\end{fields}
\end{warnbox}

\section{四、判准与论证义务}
\begin{fields}
\claim[判准是什么]{评委凭什么认定「这份厌错了」或「没错」。要回答三件事：评哪一份厌（常人说「我厌蠢」时那份厌平常的样子，连同它稳定带来的东西；不是霸凌式的最坏版本，也不是圣人式的最好版本，也不是反思修正过的版本）；用什么尺子（三问）；门槛在哪（这份厌平常的样子该改，还是平常站得住、只在个别场合调分寸）。推定谁也不占：「我没错」本身也是一句正面的自我辩护，双方都要正面举证。}
\begin{procard}{正方有利判准}\definition{这份厌，是让「我」更接近、还是更远离我自己也认可的目标：看清事实、把事做好、与人合作；如果换一种态度能更好地达到这些目标，这份厌就该改。}内在目标判准，加上与替代态度的比较。\sub{有利之处}{第一人称题里内在理由最有说服力，厌蠢者通常自认重视理性和效率，正方就用他自己的尺子去量他。三个目标里正方占两个：「看清事实」有对应偏差、知识的诅咒、问答者实验（\#61、\#253、\#254）；「与人合作」有轻蔑导向排斥（\#25、\#86）、人们怕显得无能就不敢认错求助（\#64）。「替代态度」把正方的方案摆上桌面：怨恨针对行为、促使交涉（\#160），换位思考是最有效的调节方式（\#155），事实纠正通常有效（\#109）。}\sub{反方如何反驳它}{把情绪当工具来评，按「有没有用」评情绪是把情绪当成行为；三个目标分量不一样，对厌蠢者来说「把事做好」是首要目标，对造成损失者的负面情绪是维持合作的机制（\#110，对象是故意搭便车者），不尽本分的成员会拖垮团队（\#65）；「有更好的」不等于「这个错了」，正方的替代也未必做得到：共情费力，人们稳定地回避它（\#98，同一研究也说提高共情效能感后回避就消失），长期耐心会疲劳（\#99，护理人员样本）。一句话：「存在更好的，不等于这个就错了。」}\sub{拉回中立}{正方提出这个判准后，反方指出它把「没错」偷换成了「最优」，把评委拉回中立判准的「这份厌平常的样子该不该改」；中立判准同样认「存在更好的不等于错了」。}\end{procard}
\begin{concard}{反方有利判准}\definition{只看这份厌对不对得上它的对象：对方确实明显犯蠢，而且蠢在他本可以负责的地方；对得上就没错，至于表达成什么行为，是另一个问题。}恰当性判准。\sub{有利之处}{学理出处最硬：Bell 论证了轻蔑在某些情况下是恰当的，并给出七条恰当性条件（\#35、\#36）；亚里士多德说骄傲的人轻视别人是正当的，「因为他想得对」（\#162）。结构上偏反方：只要说明存在一个满足条件、又融贯的「核心厌蠢」（针对屡教不改、把代价转嫁给别人、本可以避免的蠢），就可以主张没错。用道德主义谬误挡掉正方的后果材料（\#59）。在「看准了吗」上也有材料：〔智〕陌生人看过有声录像后对智力的评分与实测相关（\#88，不报系数），从行为片段推断的智力强烈预测测验分数（\#89）；教师期望能预测学生表现，更多是因为期望本来就准（\#255、\#257，测的是学业表现）；这几条都要先搭桥到「蠢」。}\sub{正方如何反驳它}{个案化，换了评价对象：满足七条的厌已经不是常态的厌蠢，而是「恰当的批评」；连 Bell 也没把蠢列为恰当对象，她说的是傲慢、伪善这类「优越之恶」（\#35），按她的说法，自恃聪明而看低别人的人自己倒成了恰当轻蔑的对象；七条里常态的厌蠢至少「有证据」「公平」两条过不了（\#36）；亚里士多德同一句后半是「多数人的轻视则是随意的」（\#162）；道德主义谬误是双向的，「对得上」同样推不出「没错」（\#164）；行为切不干净，轻蔑会漏在当事人察觉不到的线索里（\#156），表现出来就是冒犯（\#48）。}\sub{拉回中立}{反方提出这个判准后，正方指出它只回答了「看准了吗」和归责那一半，没有回答厌了以后还给不给他作为人应得的尊重，也没有回答「划得来吗」，把评委拉回三问通盘衡量。}\end{concard}
\begin{neutralcard}{中立判准}口播版（与 3.3「错了 / 没错」中立定义口播版是同一句）：「判断这份厌对不对，就看一个清醒的我问完三句：看准了吗？对他公道吗？划得来吗？问完觉得这份厌平常的样子就该改，就是错了；站得住，只在个别场合调调分寸，就是没错。」\sub{看准了吗（准，认知之错）}{对方真的蠢吗？我了解他的处境吗？从这件事推到这个人，推得准吗？我凭的证据够不够？}\sub{对他公道吗（公，道德之错）}{这份蠢该不该由他负责，是不会还是不肯？厌了以后，我还给不给他作为人应得的尊重？前一小问管归责，后一小问管尊重；「他不肯学」只回答了前一小问。就算看准了这个人，看低他还守不守得住作为人应得的尊重，仍要在这里再问一次（\#48、\#37、\#44）。}\sub{划得来吗（值，实用之错和对自己的错）}{这份厌让事情、关系和我自己，离我在乎的东西（把事做对、守住标准、与人共事、自己过得好）更近还是更远？改掉它要付什么代价、换上什么态度，换上的态度能不能拿到同样的好处？后一组只在已证出缺陷时才问。}\sub{「清醒」是什么}{双方会摆出来的事实他都知道：蠢的成因（归因偏差、知识的诅咒、天赋与机会）和代价（损失、风险、别人的不负责），厌的后果（关系、合作、身心）和功能（守住标准、及时止损）。「清醒的我」是评价者，不是被评价的那份厌。}\sub{合理性}{与辩题同构：「我厌蠢，我错了 / 我没错」本身就是一句第一人称的反思，人会想要、也会不想要自己的某些欲望和动机（\#157）。是普通人实际的判断方式：失败该同情还是该恼火看可不可控（\#58），一件事蠢不蠢看责任和后果（\#93），「发而皆中节」与「对恰当的人、以恰当的方式」问的都是情绪中不中节（\#158、\#152）。先分开问契合与通盘恰当，再合起来判（\#164、\#59、\#60）。吸收了情绪与立场之争：第一秒不在评价范围内。修正了「了解成因」会偏向正方的问题：「清醒」里同时写进蠢的代价与厌的功能；「看清事实」只在「准」里算一次；评价对象不取网暴者，也不取圣人；「公」同样计入「不肯学」。学界本身把「轻蔑应当作为守护规范的情绪来培养，还是应当尽量远离」列为开放问题（\#165），只能说是开放问题。}\sub{正方需要证明}{常人说的厌蠢，平常的样子在三问的某一处有系统性缺陷，而且没有被它在「值」上换不来的好处抵掉。①评价对象是常人说的厌蠢平常的样子，不是霸凌版本；网友的自我限定代表不了常态（有间大学写「我只厌打扰到我的蠢」这个前缀是为了不显得傲慢才加的，\#14）。②平常的样子有系统性缺陷：至少在「准」或「公」一处系统性失守。一是倾向于从一次犯蠢推出「这个人蠢」（\#26、\#20、\#61、\#159；\#256 只作调节变量）；二是常常冲着对方没选过的东西（\#32、\#52、\#58；〔智〕材料要搭桥）；三是看低一旦冲着人，就碰到拒绝基本尊重那条线（\#48、\#43、\#33）。③存在做得到的替代：针对行为的批评与交涉、重评与换位、一时抽身而回头还肯讲理（\#160、\#154、\#155、\#34）；替代绝不能是「压住不厌」。\par 不必证明：厌蠢者是坏人；每一次厌都错；第一秒的情绪可以控制。不够的：只证明「有时厌错了」或「厌有代价」；只证明存在更好的态度；把厌等同于伤害；要求「不许有感觉」。必须回答「对屡教不改、损害别人的蠢怎么办」：「对那件事生气、批评都可以，一时走开也可以；要改的是看低他，把他划到不用再讲理那边。」}\sub{反方需要证明}{这份厌平常的样子在三问上都站得住，毛病只在个别场合；或者就算某一处有系统性的毛病，它在「值」上有换不来的好处抵得过。①常人说的厌蠢平常对准「本该会却明显做错」，尤其是可归责的部分（\#93、\#58、\#41、\#49）；判断大体有据，靠「明显」和「屡次」；不能缩到只厌「坏」。②平常停在那件事上，不否认对方作为人应得的基本尊重（\#37、\#44、\#48）。③（条件项，正方证出系统性缺陷后才需要，实战上仍要备）这份态度本身有正当功能，至少不必然伤人（\#165、\#110、\#65、\#34）。④（条件项）正方的替代要么做不到（\#98、\#99），要么丢掉了守护标准的功能，要么只是换了名字的厌。⑤误判和过头只在个别场合出现。\par 不必证明：厌蠢是美德；每一次厌都厌得准；厌要说出口；厌蠢对社会有益。不够的：只证明「存在一个厌得对的个案」或「这份厌可以停在事上」；只证明毛病「改得掉」；把蠢换成坏；说「情绪控制不了，所以无所谓对错」。必须回答「厌了之后，是去交涉还是走开」：「厌，可以一时走开，但不把他划出讲理的范围：回头照样教、照样合作。」答成「从此走开」，就要面对「把人划出讲理范围」的质疑（\#33、\#160、\#161）。}\end{neutralcard}
\field{一辩开头合并版}{约 32 秒，只在正一开头用（反一在正方用了偏向版时重述）：「厌蠢，是觉得别人犯蠢时的嫌恶、反感，常带着想躲开的心；可以只冲着那件事，也可以冲着这个人、看低他；它是态度，不等于骂人或欺负人。蠢，是在一般人看来本该想到、做对的事上明显想错、做错，不论出于不会还是不肯。错没错，看清醒的我问完三句：看准了吗？对他公道吗？划得来吗？这份厌平常的样子该改，就是错了；只在个别场合调调分寸，就是没错。」}
\field{两条约定}{①评价对象取一个普通的、会说「我厌蠢」的人，不是网暴者（\#124），也不是圣人；他给自己划的界限守不守得住，由双方去证。②「对他公道吗」既计入对方没选过的成因（\#52、\#82），也计入「不肯学」这类可以归责的成因（\#49、\#93、\#58）。}
\field{怎么称}{\sub{三问都进场}{只赢一问不自动赢下整场，双方都要说明自己赢的那一问为什么分量最重。}\sub{分歧裁决}{一方在某一问上证出系统性的东西（平常就有的缺陷，或平常就有、换不来的好处），另一方只证出个别场合的东西，前者胜；双方都证出系统性的东西，就比分量，并看替代态度能不能拿到那份好处。}\sub{门槛}{用 3.3 的门槛分界规则：任何毛病都改得掉，「可以校准」不能当分界，分界只看平常还是个别。}\sub{「对自己」分量偏低}{普通人不会因为「生气伤身」就说一份该有的义愤是错的（\#63，只是相关），也不会因为「体谅费力」就说该有的体谅可以不给（\#98）；只在其他几处胶着时加分量。}\sub{替代后登场}{替代态度只在正方证出系统性缺陷之后登场，回答「改得了吗、改了亏不亏」；存在更好的，不等于这个就错了。}\sub{不报比例}{调研里没有任何厌蠢比例数据（\#102 是检索未果的记录）。「多数时候」「平常」只能靠机制和典型情形论证，双方都不得报百分比。}\sub{修正不等于压抑}{「改」指换一种看法，不是「别有感觉」；「压住不厌」双方都排除，压抑的生理代价（\#154）不算任何一方在「值」上的好处，只用来挡对方立的稻草人。}}
\field{三问弹药索引}{每一问双方都有已核实的材料，判准受质疑时用它证明中立；标〔智〕的先搭桥。\sub{准·正方}{对应偏差（\#61）；知识的诅咒（\#253）；问答者实验（\#254，测的是学识评价）；整体化几乎总会错误呈现对象（\#161）；实体论者从单次行为推出更强的特质判断（\#256，调节变量）；〔智〕自评与实测只相关 r=.33（\#92）。}\sub{准·反方}{「明显」加上「屡次」：屡次犯同样的错就不是情境能解释的；人们叫作蠢的是具体行为，跟着责任和后果走（\#93）；〔智〕看到言行之后对智力的判断有一定效度（\#88、\#89，不报系数）；教师期望的预测力更多来自期望准确（\#255、\#257，测的是学业表现）。}\sub{公·正方}{罪与病（\#58）；反应态度回应意志而非能力（\#32）；自然彩票（\#52）；拒绝基本尊重的蔑视违背义务（\#48）；尊重义务包括不持有轻蔑、嘲笑这类态度（\#43）；高学历者对低学历者的偏见（\#54、\#56，西方样本，测的是学历）；〔智〕智力的环境与教育来源（\#82、\#81）。}\sub{公·反方}{因疏忽而来的无知可以受责备（\#41）；困而不学（\#49）；任由自己被弄蠢（\#75）；缺努力招致惩罚（\#58，双方共用）；轻蔑与承认尊重相容（\#37）；区别评价不违背尊严（\#44）；比较之下在心里看低有时不可避免（\#48，必须整句引）。}\sub{值·正方}{轻蔑促使退出（\#160、\#25）；心理安全（\#64）；事实纠正没有引发逆火（\#109）；系统取向（\#104、\#105，只作机制例证）；愤怒、敌意与冠心病（\#63，相关，分量低）；〔智〕协作过程比个体技能更能预测集体智力（\#95）。}\sub{值·反方}{守住规范的功能是开放问题（\#165）；对搭便车者的负面情绪维持合作（\#110，对象是故意背叛者）；不尽本分者拖垮团队（\#65）；被轻蔑者任务表现提升（\#87，测的是被「表达」的轻蔑，以内隐自尊下降为中介，攻击性也上升，引它等于自认表达了轻蔑）；〔智〕个体 IQ 解释约八成群体差异（\#94，与 \#95 结论相反）；共情费力（\#98，分量低）。}}
\field{决胜点}{全场最后多半落在一处：常人说「我厌蠢」时，那份厌平常停在事上，还是落在看低人那一头；去掉看低，是改了这份厌平常的样子，还是只在个别场合调了分寸。按门槛规则，这是一个关于「平常」的事实之争。正方要证明看低和把人划出讲理范围是这份厌平常的样子（\#26、\#20、\#61、\#25；\#256 只作调节变量，对方可以说「我是渐进论者」，所以它不能打头阵），并说清换成什么。反方要证明那份反感平常停在事上，看低和把人划出讲理范围只在个别场合冒出来；调完分寸之后剩下的仍是厌，而且问完三句站得住。评委看哪一边对常人所说的厌蠢平常的样子描述得更准。}
\field{两边共用的口径}{「本可以」归归责小问。「明显」挡不掉「看准了吗」，也不等于罕见。接触回来不等于讲理回来。越线按时间划：这回先不说、一时走开，不越线；回头他开口也不听、合作只下指令不给理由、认定「他这个人说什么都没用」，越线；不再主动提醒，不算不讲理。「教过、被指出、又做错」这串记录，在「准」上双方都按它判。引每日商报那一案（\#248），必须把「多次劝告」「一度降至冰点」「很微妙」放在一起说。}
\field{对方还可能提的判准}{\sub{道德尊重判准（反方拆）}{「一份对人的负面态度，只要因为对方难以自主的缺陷看低他，就是错的。」拆：它跳过了看准没看准，等于说对方确实蠢、确实造成了损失也不能厌；「难以自主」正是要争的，困而不学不是难以自主（\#49、\#41）；在它之下任何对能力的负面评价都算错，是判准杀。}\sub{后果判准（反方拆）}{看这份厌对被厌者、自己和共同体的净影响。拆：「把情绪当行为来评；有弊不等于错。」}\sub{权利判准（正方拆）}{「我有权选择和谁亲近。」拆：「有权这样做，不等于这样做没错。」}\sub{行为判准（正方拆）}{「情绪本身没有对错，只看它是否转化为伤害。」拆：「题目问的是厌，不是骂；而且『无所谓对错』不等于『没错』。」}}
\field{判准质询守法}{反二问正一「『清醒的我』谁说了算」：「清醒只是说双方摆出来的事实都知道：蠢的成因和代价，厌的后果和功能。你们的事实也在里面。」问「划得来吗，那有用就是对的」：「有用要算，但它是三问之一，不是全部；只赢一问不自动赢下整场，对你们对我们都一样。」问「他不肯学，厌他就公道了吧」：「公道问两件事：该不该由他负责，厌了以后还给不给他作为人应得的尊重。不肯学，只回答了前一件。」正二问反一「看准了就够了吗」：「不够，所以我们同意三问都问；我们要证明的是，问完三句，这份厌平常的样子站得住。」问「看准了这个人蠢，就能看低他吗」：「看低他的能力可以，拒绝给他作为人的尊重不行；我们的厌停在前一边。」}
\field{赛制提示}{正方先立论，判准由正一先提出，反二质询正一是判准的第一个战场。正方用了中立定义和中立判准，反一一句「定义与判准和对方大体一致」就够，省下的时间讲门槛；正方用了偏向性版本，反一用一句话点出偏在哪里，再重述中立版。反方先结辩，反四要把判准封死；正四最后发言，留临场位。三辩盘问专抓口径不一致，全队对「我们的判准是什么」必须背出同一句话。}
\field{来源提醒}{反方引蔡康永「交朋友需要付出心力」（\#111）要小心：同一篇报道里有对反方「蠢不等于坏」不利的说法，那一句只到「有把握」，场上不说；「150 人」是随口引的数，不能当论据。切割「蠢就是坏」时用同条已核实的高晓松「坏人不是蠢人」。Torvalds 原话是「为了更好的补丁，这在我看来说得通，现在我知道这不 OK」（\#139），反方只能说他道歉的是邮件里的人身攻击，不能说他至今仍认为说得通。康德要整句引，中文是据德文自译，赛前对照李秋零译本。\#87 要连同「以内隐自尊下降为中介、攻击性上升」一起说；\#94 与 \#95 结论相反，谁引都要预判对方引另一篇。每日商报医生那段只引明确写作医生原话的一句（\#15）；新京报不是央媒（\#118）；Aczél 2015 不报样本数和比例（\#69 只到有把握），三类情境引 \#93。}
\field{双方退守线}{正方：「我们不要你第一秒不厌；我们要你问完三句之后，别再认同那份看低人的厌。」反方：「我们不说情绪无对错；我们说问完三句，那份平常冲着明显犯蠢的事去的反感，仍然站得住。」}
\end{fields}

\prosection{五、正方论点}
\prosubsection{论点一：拿我当一般人}
\begin{fields}
\claim{常人的厌蠢平常凭比较（「这都不会？」）而不凭记录；尺子握在已经会了的人手里，只往「他本该会」那头偏，自己还看不见。这份厌平常的判断方式，该改。}
\begin{mechanism}\step 说他「本该会」，依据只有记录或比较。凭记录的（教过、被指出、又做错）不争；厌蠢的典型说法是比较：「这都不会？」「我都可以做到你为什么做不到」（\#10、\#212、\#213）。\step 知道答案的人，会高估不知情的人本来知道多少；替一个不认识的同龄人估也一样；自己察觉不到，提醒也压不下（\#204、\#176、\#205）。\step 要改的是「凭比较、不问就判」：先开口问他的处境，再下判断（归因信息一换，情绪就跟着变，\#188 研究 2）；不是「别有感觉」。\end{mechanism}
\begin{evidence}{学理}[已核实]后见之明偏差（Fischhoff，1975）\sub{核心主张}{知道结果的人，高估不知情者实际知道多少，而且多半察觉不到。}\sub{支撑}{机制第二、三步。}\sub{边界}{从「知道结果」迁移到「这都不会？」是推论；元分析发现减偏操纵没让效应变小（\#246），所以改法押在「开口问」，不押在提醒。}\sub{出处}{\#204（只核摘要）。}\end{evidence}
\begin{evidence}{数据}[已核实]「1975 年的经典实验发现：知道结果的人会高估不知道的人本来能知道多少，自己还几乎察觉不到。」\sub{来源}{Fischhoff, 1975, “Hindsight ≠ foresight: The effect of outcome knowledge on judgment under uncertainty”, Journal of Experimental Psychology: Human Perception and Performance 1(3): 288–299, doi:10.1037/0096-1523.1.3.288（\#204）。}\sub{方法}{3 项实验，实验 3 估「别人」实际知道多少；摘要未报样本量与效应量。}\sub{局限}{只核摘要；迁移到日常的「这都不会？」是推论；不描述实验材料。}\end{evidence}
\begin{evidence}{补充数据}[已核实]「2023 年的实验里，知道字谜答案的大学生替不认识的同龄同学估难易，偏差和替自己估基本一样。」\sub{来源}{Hom, 2023, “Perspective-taking and hindsight bias: When the target is oneself and/or a peer”, Current Psychology 42(16): 13987–13998, doi:10.1007/s12144-021-02413-z（\#176，PMC 开放全文）。}\sub{方法}{字谜任务，随机分组；招募 193 人、排除 10 人；原文 “essentially the same”；要求忽略答案也没用。}\sub{局限}{美国大学生；研究指示被试设想同伴「经历相同」，这条指示本身推着人投射，被问时必须认。}\end{evidence}
\begin{evidence}{案例}[已核实]「厌蠢症」出圈的小张视频下，评论「怎么会有人不会坐飞机啊？我厌蠢症犯了」（\#167、\#119）；新京报的文章写这类言论「评论区总有」，写到同学同事时概括为「我都可以做到你为什么做不到」（\#213、\#212）；出圈一年后，每日商报开头照样列出同类说法（\#214）。\sub{代表性}{中：作者与记者的描述，不是统计；上游新闻与新京报所引可能是同一批评论，不当两份独立证据。}\sub{贴合度}{高：凭比较下判断。}\end{evidence}
\closing{回扣「看准了吗」。分量最重：看准是归责与尊重的前提；偏差只往一头偏，错落在被厌的人身上，是平常而非个别；归责只作连带，不在「对他公道吗」里重复计算。}
\begin{clash}
\pair{「记录就是他的处境，你们只剩网上陌生人。」}{凭记录的我们不争；比较也在同学同事之间：新京报的文章写到被吐槽的同学、同事，写的也是「我都可以做到你为什么做不到」（\#212）。}
\pair{「系统地偏不等于平常判错，后见之明元分析平均才零点三九个标准差（\#246）。」}{我们不说每次判错，也不评偏多少；方向只往一头偏、自己看不见、每回都用的尺子，按门槛规则就是平常。同一篇还写，设法减少偏差的研究，效应量没有更低。}
\end{clash}
\field{残余风险}{Hom 的指示（设想同伴经历相同）可能推着人投射，已写进数据局限；「平常凭比较还是凭记录」两边都只有媒体描述，是全场一处事实之争。攻击 2 的回应只按门槛规则答，不说偏得大。}
\end{fields}

\prosubsection{论点二：说了也没用}
\begin{fields}
\claim{就算看准了，熟人劝了还犯的时候，常人的厌平常会从「说了那么多遍」走到「他就是那么蠢」，回头不再把他当值得说理的人。这就越过了作为人应得的尊重。}
\begin{mechanism}\step「蠢」字落在人身上：轻蔑把对方贬为愚蠢、无能，贬的是品性多于行为（\#20）。恼火不剔，看它往哪走。\step 劝了还犯，现成供上轻蔑的预测因素：「说了那么多遍」是觉得管不了他，「还是那么蠢」是品性归因（\#185）；轻蔑可能从对同一个人的一次次恼火里长出来，而且不看好他会变好（\#86、\#27）。\step 越线看回头：这回先不说、一时走开，不越线；回头他开口先认定不值得听、合作只下指令不给理由，越线（\#48、\#45；走到这一步是推论）。\end{mechanism}
\begin{evidence}{学理}[已核实]愤怒与轻蔑的社会功能区分（Fischer 与 Roseman，2007）\sub{核心主张}{愤怒短期攻击、长期和解；轻蔑短期长期都导向拒绝和排斥，可能从先前的愤怒发展而来。}\sub{支撑}{机制第一、二步。}\sub{边界}{自我报告；原文是 may（可能）；预测因素不是必要条件。}\sub{出处}{\#86、\#185、\#170（只核摘要）。}\end{evidence}
\begin{evidence}{数据}[已核实]「2007 年的三项研究发现：生气过后会和解，轻蔑把人推出去，还可能从一次次生气里长出来；觉得管不了他、把错归到他这个人，都预示轻蔑。」\sub{来源}{Fischer, A. H., \& Roseman, I. J., 2007, “Beat them or ban them: The characteristics and social functions of anger and contempt”, Journal of Personality and Social Psychology 93(1): 103–115, doi:10.1037/0022-3514.93.1.103（\#86、\#185、\#170）。}\sub{方法}{3 项研究，回忆式情绪报告与情境研究；研究 2 检验能否从对同一个人的先前愤怒预测轻蔑；摘要未报样本量与效应量。}\sub{局限}{自我报告；必须主动认 “anger and contempt often occur together”。}\end{evidence}
\begin{evidence}{案例}[已核实]每日商报 2024-01-13：同办公室的同事小刘工作超过两年，仍常填错表格；小陈「也曾多次劝告」，同样的错接二连三，同事间的关系「一度降至冰点」，现在「陷入一种很微妙的状态」。报道紧接着这个案例，单独写了一句：「说了那么多遍，怎么还是那么蠢。」（\#248）\sub{代表性}{中：化名个案，记者挑选。}\sub{贴合度}{中高：回头讲不讲理，报道没写。}\sub{用法}{三段一起引，不标说话人；不说「小陈从此不理小刘」。}\end{evidence}
\closing{回扣「对他公道吗」的尊重小问。分量最重：它就压在全场决胜点上；尊重是底线，「划得来吗」上的好处抵不掉。}
\begin{clash}
\pair{「小陈多次劝告，讲理一直在；『一度』只是一时。」}{劝告在前，冰点和「还是那么蠢」在后，顺序就是从要求走到判词；报道写的现状是「很微妙」，不是照常教、照常听。}
\pair{「事情跟自己相关，推高的是愤怒（\#26）。」}{恼火我们不反对；同一篇摘要写着，评判无能的行为时起作用的是轻蔑。麻烦落到我身上，和「蠢」字落到人身上，是同时发生的（推论）。}
\end{clash}
\field{残余风险}{最后一步「回头不讲理」报道没写，只有 \#27 与康德撑。机制第三步标「走到这步为推论」，案例写「回头讲不讲理没写」，被问「是不是推论」直接认，再给推论的依据：轻蔑不看好对方改善；认定人改不了，人就少教（\#218）；康德问：在一个判断上否认他有理智，又怎能让他明白自己错了（\#45）。}
\end{fields}

\begin{fields}
\field{两个论点的分工}{第一点管这份厌凭什么起，第二点管它往哪走；第二点从「就算看准了」说起，场上把「就算」挂在嘴上，不说「屡次也救不了」。两个论点不共用主论证来源：每日商报一案只在论点二作主案例。}
\field{正方禁止说法}{①「厌蠢只在我知道、他不知道的时候出现」「屡次也救不了」：注意失误、缺乏控制这两类，我方不争「准」。②「厌蠢每一次都判错」「研究证明厌蠢偏了 0.39」；也不说「提醒压不下」是防洪那项研究证的（\#206 只说去偏程序没有显著效果）。③「小陈说：说了那么多遍，怎么还是那么蠢」「小陈从此不理小刘」：那句话没有标说话人，报道写的现状是「很微妙」。④「朋霍费尔只是暂停说服」「朋霍费尔是我们改法的榜样」：原文是「再也不」；只引「被推向蔑视人的危险很大」一句，说「1942 年底写的」。⑤「你们第一秒就不该厌」「别有感觉」「躲开的才叫厌」：我方不剔恼火，不要求压抑。}
\end{fields}

\subsection{（被淘汰的正方候选论点，交给反驳库与自由辩备用）}
\begin{enumerate}[leftmargin=1.8em]
\item 厌蠢逼人藏拙（值）：厌让露怯变贵，人藏拙不补拙。淘汰原因：要求厌冲着「不会」本身，与论点一同生共死；证据的自变量不是厌（团队氛围、专业系统、公开骂人）；求助反显能干（\#227、\#181）直打「藏拙」；只赢「值」。残值是「纠正换得来」模块（\#183、\#109、\#189、\#64），只在反方打「值」时用。
\item 蠢多不由己（公·归责）：蠢多出于天赋、教育、处境。淘汰原因：对方一句「我们只厌本该会却不肯的」就切掉（\#49、\#41）；核心数据 \#82、\#57、\#80 标〔智〕。
\item 标签止于骂：一个「蠢」字让思考到此为止。淘汰原因：与论点二同一机制，已并入；穆齐尔那句（\#77）只作修辞。
\item 尺子量回自己：人人都会蠢，厌蠢者自认在蠢之外。淘汰原因：人身反诘，一句「我蠢时也厌自己」打掉；数据 \#91、\#92 标〔智〕。
\item 小厌汇成偏见：个人之厌汇成对低学历者等群体的偏见。淘汰原因：判准评一个普通的「我」那份厌，不是社会现象；Kuppens 测的是学历不是蠢（\#55）；「汇聚」没有数据。
\item 身份要求耐心：教者、医者的耐心是义务。淘汰原因：职业角色不当主战场，「那是职业要求」一句切掉。
\item 坏印象改不回：厌让误判锁死。淘汰原因：前提是「看不准」，单列就与论点一同生共死；Denrell 说经常接触者的坏印象更容易改（\#193）。
\item 事后才显得蠢：知道结果的人把「本该想到」算到别人头上。淘汰原因：与论点一机制第二步同一机制，已并入；只覆盖「想得到」，覆盖不了「做得对」。
\item 认定他听不进：分歧一持续就判他有偏，不再讲理。淘汰原因：「意见分歧不是犯蠢」一刀切掉；自变量是觉得对方有偏，分不开厌与对事批评（\#208）。
\end{enumerate}

\consection{六、反方论点}
\consubsection{论点一：教过了还犯}
\begin{fields}
\claim{常人的厌蠢平常冲着「教过了还犯」：位置要他会的基本功，他被教过、提醒过，仍一再出错、不问不改。这样厌，看得有据，怪得公道。}
\begin{mechanism}\step 情绪按「能不能控制」分流：判断他该负责才生气，控制不了就同情（\#219、\#58）。\step「本该」出自位置的公共要求；「本可以」看他被教过没有、待了多久、教过之后改不改（\#41、\#231），查得到，正方在「准」上也不争。新人先教，在问在改还学不会的给耐心，不凭「一贯做不到」就厌（\#241）。\step 这正是厌蠢平常的场景：新京报的文章写，现实里说厌蠢更多冲着拖后腿的同学同事（\#243、\#249，作者描述）；评论区凭比较的，反方认作误判、个别场合。\end{mechanism}
\begin{evidence}{学理}[已核实]归因—责任理论（「罪与病」，Weiner，1993；另见 1986、1995 两部专著）\sub{核心主张}{缺努力的失败招来愤怒，缺能力的失败招来同情。}\sub{支撑}{机制第一步。}\sub{边界}{只描述人怎么反应，不管判断准不准；没有线索时人默认偏向「可控」（\#188），先认。}\sub{出处}{\#58、\#179。}\end{evidence}
\begin{evidence}{数据}[已核实]「2004 年汇总六十四项研究、一万二千多名被试的元分析发现：判断对方该负责就生气，控制不了就同情，跨文化高度一致。」\sub{来源}{Rudolph, U., Roesch, S. C., Greitemeyer, T., \& Weiner, B., 2004, “A meta-analytic review of help giving and aggression from an attributional perspective”, Cognition and Emotion 18(6): 815–848, doi:10.1080/02699930341000248（\#219）。}\sub{方法}{元分析：64 项研究，12,000 余名被试；摘要未报效应量。}\sub{局限}{只核摘要；对象是求助者与冒犯者，不是「蠢」。}\end{evidence}
\begin{evidence}{案例}[已核实]每日商报 2024-01-13：同事小刘工作超过两年仍常填错表格，小陈「也曾多次劝告」，同样的错「接二连三地再犯」（\#248）。\sub{代表性}{中：化名个案。}\sub{贴合度}{高：本职、有时间、劝后再犯。}\sub{局限}{小陈一方的叙述；后半的冰点与「说了那么多遍，怎么还是那么蠢」交给论点二，三句（「多次劝告」「一度降至冰点」「很微妙」）必须一起引。}\end{evidence}
\closing{回扣「看准了吗」与归责小问，两段分开称。分量最重：常人判断一份恼火对不对，本来就看该不该他负责。}
\begin{clash}
\pair{「新京报读全，同学同事那段是『你能不能有点常识』『我都可以做到你为什么做不到』（\#212）。」}{这几句全认：只凭「我都可以」就是比较，那里反方先问；新京报只取场景，记录由每日商报那一案补。}
\pair{「教爸妈付款，教三遍还按错，是不会还是不肯？」}{看教过之后问不问、改不改：在问、在改还学不会的，给耐心；分不清时，先问他卡在哪。}
\end{clash}
\field{残余风险}{「本可以看教过之后的做法」没有直接实证，只有亚里士多德、孟子的规范判别线（\#41、\#231）；「转眼间就将嘱咐抛诸脑后」也是小陈一方的叙述，不能当事后可查的记录。被追问时答「分不清时先问」。}
\end{fields}

\consubsection{论点二：恼他不弃他}
\begin{fields}
\claim{厌蠢收回的是对他这一次的评价，作为人的尊重不动；恼里的「你本可以」，是对能负责的人提要求，讲理由它带着：这回放下，下回照样说、照样听。}
\begin{mechanism}\step 区别评价不违背平等尊严，只要受承认尊重约束（\#44）；心里比较着看低有时不可避免，拒绝基本尊重才违背义务（\#48）。\step 恼他不上心，是把他当要为自己意志负责的参与者（\#32）；态度完全客观时才「不能讲理」（\#174），所以恼里带着讲理，不等他做对（推论）；恼还在，只证明不是纯粹的轻蔑（\#38）。\step 越界按时间判：一时躲开可以，但作为人不能长久这样（\#34）；这回放下、下回照样说照样听（\#142）；认定「他说什么都没用」才越线。\end{mechanism}
\begin{evidence}{学理}[已核实]反应态度与客观态度（Strawson，1962，《自由与怨恨》）\sub{核心主张}{反应态度回应对方意志里的善意、恶意或冷漠，是把他当参与者；态度完全客观时，可以跟他谈、甚至谈判，却不能跟他讲理。}\sub{支撑}{机制第二、三步。}\sub{边界}{「还在恼就还在讲理」是推论；必带「完全」二字；原刊卷页属有把握，只引全文转载页。}\sub{出处}{\#32、\#34、\#174。}\end{evidence}
\begin{evidence}{数据}[已核实]「2007 年的三项研究发现：生气当下冲上去、长远会和好；轻蔑才把人推出去。」\sub{来源}{Fischer \& Roseman, 2007, Journal of Personality and Social Psychology 93(1): 103–115, doi:10.1037/0022-3514.93.1.103（\#86、\#170）。}\sub{方法}{3 项研究，回忆式情绪报告与情境研究；摘要未报样本量与效应量。}\sub{局限}{自我报告；必须主动认 “often occur together” 与 “contempt may develop out of previously experienced anger”：可能，不是平常的走向。}\end{evidence}
\begin{evidence}{案例}[已核实]孔子与子路论正名：子路当面说「子之迂也」，孔子斥一句「野哉，由也」，接着把正名的道理讲完；朱熹集注记胡氏「子路終不喻也」（\#226）。\sub{代表性}{中低：圣人师生，「平常」靠机制。}\sub{贴合度}{高：恼与讲理同在当场，讲理也不等他做对。}\end{evidence}
\closing{回扣「对他公道吗」的尊重小问。分量最重：决胜点就落在「平常停在事上，还是落到看低人」。}
\begin{clash}
\pair{「轻蔑的预测因素在你们的场景里现成：劝了还犯是管不了，『还是那么蠢』是品性归因。」}{预测因素不是必要条件；熟人之间亲密在，他别处做对的记录逆着品性归因（推论）；真到「他说什么都没用」，反方认越线、要调分寸。}
\pair{「讲理挂在他的表现上；你们把『一时』换成了『一事』。」}{讲理不等他做对：子路「終不喻」，孔子照样讲完。这回先不劝，下回照说；这件事上从此不听他，反方认越线。}
\end{clash}
\field{残余风险}{三处推论：「还在恼就还在讲理」；主案例是圣人；「熟人间亲密在、别处做对的记录逆着品性归因」。三处场上都照说是推论，案例代表性说「中低」。}
\end{fields}

\begin{fields}
\field{两个论点的分工}{论点一管归责，论点二管尊重；每日商报一案只在论点一作主案例，Fischer 与 Roseman 只在论点二。「平常」只有一条主线：厌跟着代价走，代价反复落在一起做事的人身上，那里记录现成。}
\field{反方禁止说法}{①「情绪控制不了，所以无所谓对错」「我有权远离谁」。②「新京报认为厌蠢没错」「多数厌蠢冲着拖后腿的同事」：只说「新京报的文章写道」，同篇的不利句要认。③「小刘就是无所谓」「抛诸脑后是事后可查的记录」：这些都是小陈一方的叙述。④「恼还在，所以里面没有轻蔑」「轻蔑要三个前因凑齐才出现」。⑤「郑人宁信尺码也不信自己的脚」：原文是「無自信也」，而且那一问在集市散了之后；也不说 Torvalds 至今认为骂人说得通。}
\end{fields}

\subsection{（被淘汰的反方候选论点，交给反驳库与自由辩备用）}
\begin{enumerate}[leftmargin=1.8em]
\item 看见「本可以」才厌：认「本可以」靠他自己的话和记录（做到过、推开检验、理直气壮）。淘汰原因：新京报读全倒向正方（\#212）；Frieze 与 Weiner 把一贯做不到读成不会（\#241）；「做到过」与「位置」成了两把尺子；郑人买履时序不合（\#202）。
\item 守住本该（值）：反感让不肯想有代价，标准才立得住。淘汰原因：主数据 \#110、\#220 测的是故意背叛、能做不做，是坏不是蠢；要让对方感到代价，与论点二「照样讲理」冲突；修好的版本就是正方的替代态度；Linux 案例可被反用（\#189）。残值是「换不掉的那点反感」模块，只在正方拿出替代态度后登场。
\item 明显屡次，看得有据：凭「明显」「屡次」「被提醒过」下判断。淘汰原因：主数据 \#88、\#89 标〔智〕；思路已并入论点一。
\item 有要求才会厌：「你本可以」的恼是把他当成年人。淘汰原因：与论点二机制第二步同一机制，已并入；单列会被 \#33、\#38 两面夹击。
\item 一时抽身，止损留力：躲开是有限心力下的止损。淘汰原因：「对自己」的代价分量偏低；\#98 同一摘要说提高效能感后回避消失；\#99 是护理人员样本。
\item 情绪无罪：厌不经选择，谈不上对错。淘汰原因：判准写明「控制不了」不能免于评价，「无所谓对错」不等于「没错」，还与「我没错」的正面姿态自相矛盾。
\item 远离是我的权利：只远离不伤害。淘汰原因：有权这样做不等于这样做没错；「150 人」不能当论据（\#111）；「从此远离」撞上「把人划出讲理范围」。
\item 躲开会自己退：想躲开会随时间退去。淘汰原因：原谅研究的回头靠关系价值与对方示好，是外加的（\#228、\#199、\#200）；对象是人际冒犯，不是蠢；测的是回避，不是讲理。
\item 认错就消气：他一问、一认错，恼就转成耐心。淘汰原因：只证明「改得掉」，门槛规则「平常就这样、只是改得掉，照样算错了」打掉。
\item 熟人有记录：熟人之间凭记录，印象会更新。淘汰原因：一身撑两问，违反论点互不重叠；Denrell 前一句（坏印象让人少接触）利正方。
\item 同一把尺子先量自己：厌蠢是拿要求自己的期待要求别人。淘汰原因：没有测过这一点的数据（\#239 是作者动机）；中国网的投射说偏正方（\#16）；「对己严不等于对人可以厌」一句打掉。
\item 诚实之德：压住厌恶是伪善。淘汰原因：打稻草人，「压住不厌」双方都排除。
\item 动力之用：被厌的人表现会提升。淘汰原因：\#87 测的是被表达的轻蔑，以内隐自尊下降为中介，攻击性也上升，引它等于自认表达了轻蔑。
\item 结构反用：根在效率社会与优绩主义。淘汰原因：没回答这份厌对不对，还把「蠢多出于处境」送给正方。
\item 后果之厌 / 理性之守：厌的是麻烦与风险。淘汰原因：单引 \#12、\#15 会被补出同一来源 \#13 的批判与「慕强」一句。
\item 限定之蠢：只厌又菜又犟的人。淘汰原因：写进定义就是把蠢偷换成坏（3.5 反方定义杀②），只能留在论证里。
\item 愚蠢比恶更危险（朋霍费尔）：淘汰原因：下一节「蔑视人？」反打（\#76），语境是纳粹治下。
\item 开口问反而显得聪明：淘汰原因：\#227 写明任务难时效应更明显，容易任务的结果只有报道转述（\#181）。
\end{enumerate}

\prosection{七、正方反驳库（针对反方全部可能论点）}
\begin{fields}
\field{全队底线}{\sub{三问与门槛}{看准了吗；对他公道吗（前一小问归责，后一小问尊重）；划得来吗。只赢一问不自动赢下整场。看这份厌平常的样子：「平常就这样、只是改得掉，照样算错了」。替代态度只在证出系统性缺陷之后登场，「存在更好的不等于这个错了」我方认。}\sub{我方让出的}{凭记录（教过、被指出、又做错）的「准」；注意失误、缺乏控制两类的「准」；对那件事的恼火和批评；一时走开、这回先不说、不再主动提醒。}\sub{越线}{按时间划：回头他开口也不听、合作只下指令不给理由、认定「他这个人说什么都没用」，才越线；接触回来不等于讲理回来。}\sub{论据纪律}{不报任何厌蠢比例，「许多」「总有」只说「某某文章写道」；研究规模只当研究本身的规模；〔智〕先搭桥；康德整句引；每日商报一案第一次引用三段一起说（「多次劝告」「一度降至冰点」「很微妙」），之后才回指「还是那么蠢」；Strawson 不报卷页；朋霍费尔只引「被推向蔑视人的危险很大」；Torvalds 只说他曾经觉得说得通。}\sub{条数}{共 45 条：反方最终论点拆成的子论点 20 条（1–19、44），淘汰候选 17 条（20–35、45），公开讨论常见说法 8 条（36–43）。}}
\end{fields}
\begin{rebuttals}
\rebuttal{1}{平常就在身边：「新京报的文章写道：现实里说厌蠢的人，更多是在说拖后腿的同学同事；厌蠢平常在有记录的人身上，评论区那种是个别。」}{同一篇读全：写同学同事那段是「你能不能有点常识」「我都可以做到你为什么做不到」。身边不等于有记录，平常凭的还是比较。}{回扣：准｜平常。反方从这篇取了场景，我方取同一篇写下的判断方式：同学同事被吐槽，是「『你能不能有点常识』的烦躁」，这种愤怒背后的思维「更像是一种『我都可以做到你为什么做不到』的抱怨」（\#212）。反方引的那段没写当事人被教过没有（\#249）。「最基本」由已经会了的人判，人估别人知道什么，起点就是自己知道什么，而且常常不加批判（\#216）。凭记录的那一格我们让，凭比较的这一格才是平常。}{①「作者也写了，更可怕的是理直气壮。」→「更可怕」是比较级，先被厌的是无知；作者的判断是被叫作蠢的人「有许多却是努力了但也没有学会的人」（\#212）。反方只能说「新京报的文章写道」。②「天天在一起，教没教过一问就知道。」→一问就知道，那就是先问；反方要证的是常人平常问了。记录要攒，「这都不会？」那一下只要比较（推论）。}{最终}
\rebuttal{2}{情绪按「能不能控制」分流：「汇总六十四项研究、一万二千多名被试的元分析发现：判断他该负责就生气，控制不了就同情。常人的厌冲着该负责的人，怪得公道。」}{元分析管的是判断之后生出什么，不管判断准不准；没人告诉你他卡在哪时，人默认把他读成「不肯」。}{回扣：准｜公·归责。情绪跟着责任判断走，我方全认；可它只管判断之后（\#219 局限），对象是求助者和冒犯者，不是蠢。毛病在判断这一步：1988 年的实验里，「心理—行为」类的毛病默认被归为起因可控，招来更多愤怒、更少同情、倾向不管（\#188）；反方立论自己也先认了这一点。厌蠢那一下恰恰线索最少，分流越灵，读错的火就越准地落在被读成「不肯」的人身上（推论）。「公」的归责要先过「准」。}{①「Weiner 测的是污名。」→对，「蠢属于这一类」是推论，我方照这样说；反方的元分析对象也不是蠢，两边同一个跳步。②「给了信息情绪就变，你们证的只是改得掉。」→平常就这样、只是改得掉，照样算错了；我们要的改法正是先开口问，把信息问出来。}{最终}
\rebuttal{3}{「本该」出自位置的公共要求：「本该会什么，是他所处位置的公共要求，写得出、教得出；亚里士多德说，对应当知道又不难知道的东西无知，要担。」}{位置给的是「该会」，不是「已经会」；他被没被教过，位置答不了，答的还是会了的人的估计。}{回扣：准｜平常。教过、交代过的，是记录，我方不争；没教过的，反方自己说先教。剩下的才是战场：不知道他被教过没有，「人人都该会」就是我的估计。研究测的正是这一下：判断自己知道的东西在同龄人里有多普遍，熟悉带来的流畅感就足以造成知识的诅咒（\#215，3 项实验、359 名本科生）；越内行越低估新手的难处，去偏办法也压不下（\#203）。亚里士多德讲的是立法者对成文法规的处理（\#230），借得来判别线，借不来「哪些算基本」的清单。}{①「Hinds 说中等水平的人预测新手更准。」→说「这都不会」的人在这件事上恰恰是已经会了的那个（推论）。②「Birch 摘要里没有『高估』。」→迁移是推论；方向据 Hom 2023 引言转述（\#217），场上只说「往『大家都知道』那头偏」。}{最终}
\rebuttal{4}{「本可以」看记录：「被教过没有、待了多久、教过之后改不改，查得到；这串记录，正方自己承认在『看准了吗』上不争。」}{有记录的我们让，可让的是「准」；「教过了还犯」赢不下「平常」，更回答不了厌了以后还给不给他尊重。}{回扣：准｜公·尊重。让的只是「准」和归责那一半：凭记录看准了他不上心，恼他、批评他，我方本来允许。可「对他公道吗」还有尊重小问，「他不肯学」只回答了前一小问；反方论点一赢到头，也只是「冲着不肯的恼火公道」，判准把「只证明这份厌可以停在事上」列为不够的。决胜点在看准以后：同一个小刘，报道写下的是多次劝告之后关系一度降至冰点、现状很微妙，和那句「说了那么多遍，怎么还是那么蠢」（\#248）。}{①「你们让了准，就只剩网上陌生人。」→比较也在同学同事之间（\#212）；就算只剩陌生人，论点二照样在：报道里从要求走到了判词。②「劝了多次还犯，还不能怪？」→能怪，怪那件事；要改的是「怪」走到看低他、回头不再跟他讲理那一步。}{最终}
\rebuttal{5}{厌跟着代价走：「与自己相关推高的是愤怒；新京报的文章写，对蠢的厌恶更多源自它可能造成的后果；代价落在一起做事的人身上，他们的记录本来就在我手里。」}{代价证明的是厌在哪儿起得重，不证明起的时候手里有记录；返工的麻烦当场就有，记录要一次次攒。}{回扣：准｜平常。链条断在「代价落在我身上」到「我手里有他的记录」这一步：代价第一次落下来，就是「这都不会？」冒出来的时候（推论）。那句摘要前半是与自己相关推高愤怒，后半是评判无能的行为时起作用的是轻蔑（\#224、\#26）：代价让人更恼，不让人更准。「我只厌打扰到我的蠢」，有间大学写这个前缀是为了不显得傲慢才加的（\#14）。讨厌麻烦叫怕麻烦，用不着「厌蠢」这个词。}{①「天天在一起，教没教过当然知道。」→知道「他又错了」，不等于知道他卡在哪；反方主案例里劝的原话是「工作上的事情要多上心」，劝的是态度（\#248）。②「你们也说对事生气不越线。」→对，争的是这份气往哪走：「说了那么多遍」正是轻蔑的现成预测因素（\#185），见第 17 条。}{最终／淘汰}
\rebuttal{6}{比较的口气说的是「本该」：「『我都可以做到』只是本该说出口的样子；拿一般人比是中立定义自带的。」}{「我都可以做到」比的不是一般人，是「我」；定义里的一般人了解他的处境，比较句恰恰把处境省掉了。}{回扣：准｜平常。中立定义里的一般人，是了解他处境、讲道理的一般人，不是评论区的多数票。「我都可以做到」量的是我的经验，「你能不能有点常识」量的是一条知识在一般人里有多普遍，这两下正是知识的诅咒发作的地方（\#215、\#216）。真凭记录，说出口的该是「上次教过你」（推论）。作者给这种思维的定性是「抱怨」，还写把失败归咎于不够努力「显然是不客观」（\#212）。}{①「嘴上说的和心里的依据可以不一样。」→请拿出心里依据是记录的材料；调研到的文字里，教过劝过又犯的记录只出现在一个化名个案（\#248）。②「说话的人明明知道他被教过。」→那一格我们让；争的是不知道教没教过、却一样说「这都不会」的那一格。}{最终}
\rebuttal{7}{评论区是最响的地方：「网上的愤怒表达得到点赞，人以后就更常这样表达；评论区量不出常人平常的样子。」}{被点赞强化出来的厌蠢也是厌蠢，强化正是一份态度变成习惯的过程；何况比较句也在同学同事之间。}{回扣：平常。Brady 等的研究说愤怒表达得到正反馈，以后更可能再这样表达（\#101）：放大的不是假象，是习惯（推论）；「我厌蠢」本来就是习惯性陈述。而且它测的是推特上的道德愤怒。我方论点一也不靠评论区独撑：同一篇新京报写同学同事，同样是「你能不能有点常识」（\#212）；出圈一年后，每日商报开头照列「怎么会有人不会打车啊？我厌蠢症犯了」（\#214，报道描述，未注出处）。平常看判断方式是不是一贯，不看平台。}{①「你们也承认量不出比例。」→承认，双方都不报比例；所以「平常」靠机制：比较那一下人人随手可用（推论）。②「它说的是表达，不是心里。」→那反方也不能拿它说评论区的不是真厌蠢；那份也算厌蠢，看错了，是第一问要称的错。}{最终}
\rebuttal{8}{教过还不改，不是「看不出」：「提高技能，人就认得出自己的局限；被一再提醒过，『看不出』就换成了『被告知了』，再不改就是不上心。」}{被告知「你错了」，不等于被教会「怎样才对」；研究里认得出局限，靠的是提高技能，不是听几遍劝。}{回扣：公·归责。Kruger 与 Dunning 原句：提高被试的技能、从而提高元认知，帮他们认出了自己能力的局限（\#251）。条件是研究者给的训练。反方主案例里劝的是「工作上的事情要多上心」（\#248），没写教他怎样才对；技能没提上去，「看不出」还在（推论）。分不清时反方自己的答法是「先问」，可案例里没写有人问过他卡在哪。「教过了还犯」里有一截其实是「劝过了还不会」，按反方自己的分流该得耐心。}{①「他自己说会了。」→那项研究的被试只写出步骤、没动手做（\#201）；说会了做出来不对，是信心与能力失衡，「准」我们不争，归责上以为自己会不是「本可以」。②「这样谁都能说自己不会。」→所以要问；问了他说不会又不学，归责就到反方那边。}{最终}
\rebuttal{9}{没记录先问，替代本来就有：「没有记录时我们先问，新人先教；正方的替代『先开口问』，在我们的厌里本来就有。」}{「先问」是反方替常人许的愿；厌蠢的典型句子全是反问句，反问是判完了，询问才是在问。}{回扣：平常｜门槛。判准评的是常人那份厌平常的样子，不是反思修正过的版本。请反方拿出常人先问的样子。调研到的典型说法：「怎么会有人不会坐飞机啊？我厌蠢症犯了」（\#167）；「你能不能有点常识」（\#212）；「怎么还有人到医院不会挂号呀？太蠢了吧」（\#214），全是反问。「先问」要靠立论承诺才补得上，恰恰说明这份厌自己不往那里走；平常就这样、只是改得掉，照样算错了。替代也不是「别有感觉」：把「这都不会？」换成「你卡在哪？」，换位思考是调节效果最大的一项，d=0.45（\#155）。}{①「替代只在证出缺陷之后登场。」→所以先证：凭比较的判断只往一头偏、知道答案的人察觉不到（\#204），替不认识的同龄人估也一样（\#176）；反方说它「本来就有」，等于先认了它做得到。②「这只是我们的厌换了名字。」→差别在有没有「看低他」那一截、回头还讲不讲理。}{最终}
\rebuttal{10}{主案例小刘：「同事工作两年，表格常填错，劝了多次，同样的错接二连三；本该、本可以都在，这就是厌蠢平常的样子。」}{三段一起念：劝告之后，关系「一度降至冰点」，现状「很微妙」，报道紧接着一句「说了那么多遍，怎么还是那么蠢」。这一案证的是我方论点二。}{回扣：准｜公·尊重。「准」上我方让：两年多、多次劝告、接二连三。可报道没停在这里：同一段写关系「一度降至冰点」，开头写两人「陷入一种很微妙的状态」，案例之后单独一句「说了那么多遍，怎么还是那么蠢」（\#248，未标说话人）。从「多上心」的要求，走到「那么蠢」的判词，走到冰点，这是从对事走到对人。这是化名个案，「无所谓」「抛诸脑后」都是小陈一方的说法，证明不了平常。三问分开称：准，反方赢；归责，反方赢一半；尊重，报道写下的是我方那条路。}{①「多次劝告就是讲理一直在。」→劝告在前，判词和冰点在后；我方不说「从此不理」，只说现状很微妙，不是照常教、照常听。②「那句没标说话人。」→我们不说是谁说的，只说报道把它紧接在案例后面；「蠢」字落在人身上（\#20；推论）。}{最终}
\rebuttal{11}{看低这一方面不等于不当人：「平等的尊严不排除区别评价；康德也说比较之下在心里看低有时不可避免。我们下调的只是对他这一方面的评价。」}{两句都要念全：区别评价必须受承认尊重约束；康德下半句是「表现出来就是冒犯」，同一节还说拒绝作为人应得的尊重无论如何违背义务。}{回扣：公·尊重｜平常。SEP 那句后半：区别评价必须始终受「给人承认尊重」这条要求约束（\#44）；康德整句见第四节（\#48）。所以问题不在能不能评价，在这份评价停不停得住。心理学说它平常停不住：轻蔑把对方贬为愚蠢、无能，贬的是品性多于行为（\#20），「蠢」字本身就落在人身上；评判无能时起作用的正是轻蔑（\#26），而轻蔑不看好对方改善（\#27），看低一个人改不了，就到了「不值得再说」的门口（推论）。每日商报那句「还是那么蠢」里，没有「这一方面」。}{①「你们没证明越到了承认层。」→越不越看回头；康德附释说，在某一个判断上否认他有理智，又怎能让他明白自己错了（\#45），越线不必等到整个人。②「表现出来才冒犯，我们不表现。」→心里那份就是厌蠢本身；Adams 一派还认为轻蔑会从察觉不到的线索里漏出来（\#156）。}{最终}
\rebuttal{12}{恼还在就不是纯粹的轻蔑：「客观态度装不下愤怒；去人性化的是不掺愤怒的纯粹轻蔑。我们的厌带着恼，到不了轻蔑那一头。」}{「不纯」不等于「没有」：两种情绪常常一起出现，轻蔑还可能就从对同一个人的恼里长出来。}{回扣：公·尊重。反方自己写明，「恼还在」只证明不是纯粹的轻蔑。Fischer 与 Roseman 原话：愤怒和轻蔑常常一起出现，轻蔑可能从先前经历的愤怒里发展出来（\#170、\#86）；恼不是挡住轻蔑的盾，是它生长的地方。Strawson 也只划了两头：态度完全客观，才可以谈判却不能讲理（\#174）；他没说有恼就一定在讲理。判准问的是这份厌落不落在看低人那一头，不问纯不纯；夹着恼的轻蔑照样是轻蔑。}{①「恼退了是和解。」→恼退了，盾牌跟着退；剩下的想躲开和调低的评价，正是 Strawson 说客观态度可以带着的「排斥性的反感」（\#197）。②「may 只是可能。」→我方只说可能；研究 2 测的正是对同一个人的多次愤怒能否预测轻蔑（\#185），条件就是反方挑的「教过了还犯」。}{最终}
\rebuttal{13}{恼传过去的是「你没用心」：「1984 年的实验里，孩子做砸后大人冲他生气，孩子更多想的是『我没用心』；恼里的『你本可以』，只对做得到的人说。」}{那项实验测的是孩子听到大人生气后怎么想，不是生气的人心里还看不看得起他；恼平常停不停在「你本可以」，它答不了。}{回扣：公·尊重。三处。第一，Graham 测的是接收方的归因，不测生气的人怎么想，也不测讲不讲理；被试是儿童，一次实验室任务，摘要自己说愤怒那条联系弱于同情那条（\#250）。第二，劝了还犯之后说出口的是什么？报道写下的是「说了那么多遍，怎么还是那么蠢」（\#248），那是判词，按反方自己引的逻辑，传过去的就是「你就这样」（推论）。第三，「你本可以」和「他就是那么蠢」可以前后脚：前一句是要求，后一句是品性判决（\#20）。}{①「我们只借它说明恼的内容。」→那「愤怒里有你本可以」我方不争；争的是它平常停不停在那里：「说了那么多遍」是觉得管不了，「还是那么蠢」是品性归因（\#185）。②「提要求本身就是尊重。」→所以对事批评我们支持（\#160）；要求变成判词，就从交涉走到了退出。}{最终／淘汰}
\rebuttal{14}{愤怒短期攻击、长期和解：「2007 年的三项研究发现，生气长远会和好，轻蔑才把人推出去；厌蠢里的是恼，不是推开。」}{同一篇摘要往下读：轻蔑可能从对同一个人的愤怒里长出来；劝了还犯的熟人场景，正是两种情绪交接的地方。}{回扣：公·尊重｜平常。这篇双方共用，我方整篇引：愤怒长期和解，轻蔑长短期都排斥；两者常常一起出现；研究 2 发现轻蔑可能由先前的愤怒发展而来；觉得控制不了对方、作负面品性归因，都预测轻蔑（\#86、\#185、\#170，自我报告）。反方挑的「教过了还犯」正好把这几个预测因素凑齐。「愤怒会和解」管的是一次的恼，厌蠢说的是「说了那么多遍」之后那一份。说愤怒是趋近型的综述比较的是愤怒和焦虑，「趋近」不等于讲理，更不等于和好（\#244）。}{①「预测因素不是必要条件。」→对，我方不说必然；说的是反方挑的典型场景自带这几个因素（推论）。②「熟人之间亲密在。」→同事未必亲密（推论）；「还是那么蠢」已经是品性归因，见第 17 条。}{最终}
\rebuttal{15}{越界按时间判，这回放下、下回再来：「Strawson 说可以一时躲开投入之累但不能长久；莎莉文绝望时放下、又重新开始；孟子说不屑之教诲也是教诲。」}{线我们同意按时间划；可反方的样子全是老师和圣人，评的却是常人，常人那份厌的「下回」谁来保证？}{回扣：公·尊重｜平常。越线标准双方一致，分歧在常人平常落在哪边。莎莉文是老师，海伦是聋盲不是蠢（\#142）；孟子讲的是老师的教法，朱熹注是条件句，对方能「退自修省」才算教诲（\#232）；Strawson 同一处还说投入之累持续太大，就得做别的事，比如断绝关系（\#197）：撑不久的出口不止回头讲理一个。反面证据：轻蔑不看好对方改善（\#27）；认定人的特质改不了，经理就不愿花力气辅导表现差的下属，看法一改，辅导意愿和建议质量都上升（\#218）。}{①「Heslin 测的是经理，那是职业要求。」→我们拿它证机制：「改不了」的看法让人少教；有职责的人尚且如此，没职责的同事只会更少（推论）。②「压力太大才断绝。」→反方的典型场景恰是压力持续：同样的错接二连三（\#248），那正是 Strawson 写到断绝的那一格（推论）。}{最终}
\rebuttal{16}{回头在熟人之间是常态：「经常接触的人坏印象更容易改；做对的信息在能力评价里分量重；被冒犯之后，原谅随时间上升。想躲开的心会自己退。」}{接触回来、原谅回来，都不等于讲理回来；Denrell 前一句写着：坏印象让人少接触，所以坏印象比好印象更稳定。}{回扣：公·尊重｜门槛。Denrell 先说负面经历降低再接触的概率，所以负面初始印象更稳定（\#193，模型推导）；「经常接触的人才更容易改」说的是躲不开的接触。原谅随时间上升，与「这段关系仍有价值」的知觉相关（\#228），道歉补偿能让它加快（\#199）；可原谅测的是回避和报复的念头下降、善意回升（\#229），对象是人际冒犯，不测讲不讲理。频率那条，同一原则说越少见的行为分量越重（\#240），往哪边推双方都只当推论。三条材料撑不起「平常自己走回讲理」。}{①「回头是自然走向。」→回来的是善意和来往，讲理回不回来没人测过；就算回来了，平常就这样、只是改得掉，照样算错了。②「他做对一次，评价就改写。」→被看成不会就分不到难的活：那位扮蠢的大学生自述小组作业「大家不会把最难的部分分给你」（\#237）；能改写评价的表现要先有机会（推论）。}{最终／淘汰}
\rebuttal{17}{轻蔑的预测因素在熟人间偏弱：「熟人之间亲密在；还在劝，就是还觉得话管得到他；『怎么还是』还在要求他。」}{「怎么还是」后面接的是「那么蠢」：要求的是一件事，判的是这个人。}{回扣：公·尊重｜平常。这句话在报道里跟在「多次劝告」「一度降至冰点」之后，现状「很微妙」（\#248）。「说了那么多遍」说的是话管不到他了，就是「觉得控制不了对方的行为」；「还是那么蠢」把错落到他这个人身上，就是负面品性归因（\#185）。整句的落点是「那么蠢」，是对人的品性判决（\#20）。「熟人亲密在」：同事未必亲密（推论）；「别处做对的记录逆着品性归因」，反方自己标了推论。预测因素不是必要条件，我方照认；我方说的是反方挑的场景自己就把它们供上了。判词冲着人，就碰到拒绝基本尊重那条线（\#48）。}{①「还在劝，就还觉得管得到他。」→「说了那么多遍」是一边说一边宣布说了没用（推论）；劝的原话是「多上心」，怎么做没人教。②「真到说了白说，我们认越线。」→那反方要证常人平常停在「这回先不说」；报道写下的那句是后者。}{最终}
\rebuttal{18}{孔子骂了子路照样讲完：「子路顶撞『子之迂也』，孔子骂他『野哉』，同一口气把正名讲完；恼和讲理同在当场。」}{孔子是圣人、是老师，这正是我们要常人学的改法，不是常人厌蠢平常的样子。}{回扣：平常。判准评的是一个普通的「我」，不是圣人。孔子骂了还讲完，是我方替代态度的样子：对事批评、回头讲理（\#160、\#34）。反方自己写这一案代表性「中低」。「野」朱熹释为「鄙俗」，可以读成对人的评语；同卷紧接着一章，樊迟出去后孔子说「小人哉，樊須也」（\#226）。子路鼓瑟那一则：门人听孔子一句批评就不敬子路，有记录的同门一句话就从事滑到了人，拦住他们的是圣人（\#245）。这是要改出来的样子，不是平常就有的样子。}{①「鼓瑟那条，门人拿的是别人的判词。」→常人「他蠢」这句判词是自己下的，比听来的更难收回（推论）；朱熹「未可以一事之失而遽忽之也」是告诫，告诫写给人容易犯的毛病。②「我们只说这个样子常人身上也有。」→「也有」就是个别场合，门槛看平常。}{最终}
\rebuttal{19}{误判过头只在个别场合：「就算这份厌有时看错、有时过头，也只在个别场合，调分寸就好；什么毛病都能校准，照你们的标准世上没有不错的情绪。」}{门槛规则写着：平常就这样、只是改得掉，照样算错了；「个别」要证，不能靠喊。}{回扣：平常｜个别。我方两处都证在平常。判断方式：知道答案的人高估不知情的人本来能知道多少，自己察觉不到（\#204）；替不认识的同龄同学估，偏差和替自己估「基本相同」（\#176）：方向单一、看不见、每回都用的尺子，不是个别。走向：「蠢」字落在人身上（\#20）；劝了还犯的场景自带两个轻蔑预测因素（\#185）。反方要把它们降成个别，就得证常人平常有记录、平常先问、平常回头讲理；反方拿出的是一则化名个案、一位圣人、一位老师。}{①「平均才零点三九个标准差。」→我们不说每次判错，也不说偏了多少；同一篇还写，设法减少偏差的研究效应量没有更低（\#246）。②「Hom 的指示推着人投射。」→局限照认；另一项没有这条指示：熟悉感就足以造成知识的诅咒（\#215）。③「世上没有不错的情绪。」→不是，只在个别场合冒出来的调分寸就好；我们证的是平常。}{最终}
\rebuttal{20}{看见「本可以」才厌：「认本可以，靠他自己的话和记录：做到过这回没做到、推开了现成的检验、理直气壮说自己会。郑人买履就是蠢的原型。」}{Frieze 与 Weiner 说：一贯做不到读成不会，偶尔没做到读成运气或努力，别人会他不会分不出；没有一格稳稳读出「本可以」。}{回扣：准｜公·归责。与他过去一致的结果被归于能力和难度，不一致的引出运气和努力，与别人不一致的三样都有（\#241，1971 年实验室研究）；「屡教不改」按这条读出的是「不会」。Fernbach 的实验里，是研究者要求被试解释机制，没有人「选择」接不接检验（\#242）。郑人那一问在「及反，市罷」之后，集市散了，没有鞋可试（\#202）。理直气壮是以为自己会，在「准」上明显想错我们不争，在归责上不是「本可以」；同一份记录读出不会，「准」站住了，「公」就倒了（推论）。}{①「理直气壮本身就可恶。」→可恶的是让别人替他买单，那是厌坏不是厌蠢；英语里常说的汉隆剃刀讲的就是能用愚蠢解释的别归于恶意（\#130）。②「做到过这回没做到总能怪。」→那一格读出的是运气或努力，挑哪一样又是看的人在读；郑人原文是「無自信也」。}{淘汰}
\rebuttal{21}{明显、屡次，看得有据：「屡次犯同样的错不是情境能解释的；陌生人看过一个人的言行，对他智力的判断就有一定效度。」}{后半先搭桥：那两项测的是智力，中立定义写明蠢不等于低智商；前半「屡次」在准上我们让，可归责上读出来的是「不会」。}{回扣：准｜公·归责。\#88、\#89 测的是对智力的判断与测验分数的关联，系数不报；不先搭桥，一句「你们自己说蠢不是智商」就切掉，桥也难搭：大量思维偏差与认知能力不相关（\#71）。「屡次」我们在准上让，可与他过去一致的失败被归于能力和难度（\#241），按反方自己的分流该得同情。「明显」说的是错得清楚，不说他被教过，也不等于罕见。注意失误、缺乏控制两类我们不在准上争，那一类看厌了以后往哪走。}{①「搭桥：蠢是判断失当。」→理性与智力是两回事，智力足够却不能理性行事叫理性障碍（\#70），桥搭不过去。②「屡次读成不改。」→不改要看教过之后问不问、改不改，那是第 8 条；只凭屡次，读出的是不会。}{淘汰}
\rebuttal{22}{守住本该，换不掉的那点反感：「这份反感让不肯想有代价，标准才立得住；对搭便车者的负面情绪维持合作，不尽本分的人拖垮团队。」}{三条数据的对象都是「能做却不做」，那是坏，不是蠢；守标准，对事的恼就守得住，用不着看低人。}{回扣：值。Fehr 与 Gächter 的对象是故意搭便车的背叛者（\#110）；Kerr 的傻瓜效应是搭档能出力却不肯，实验也没操纵「有没有人反感他」（\#220、\#190）；坏苹果是不做自己那份、欺凌他人的人（\#65，大学新闻稿）。轻蔑该培养还是该远离，那本文集只说是开放问题（\#165）。守标准靠怨恨那一种：针对行为、促使交涉（\#160）。厌在场的代价也在「值」上：人怕显得无能，就不敢承认错误、不敢求助（\#64）。守标准这份好处，对事批评拿得到，看低他那一截拿不到，也不该拿。}{①「存在更好的，不等于这个错了。」→同意，替代只在证出缺陷后登场；我们先在准、公两处证缺陷。②「对事批评没力量。」→力量在恼，不在看低；事实纠正通常有效，没有引发逆火（\#109）；反方「要让他感到代价」，又和自己「照样讲理」冲突。}{淘汰}
\rebuttal{23}{一时抽身，止损留力：「想躲开是有限心力下的止损；共情费力，长期耐心会疲劳；你不是讨厌蠢，你只是太累了。」}{一时走开我们不反对；累是原因，不是理由，「对自己」的代价双方分量都低，共情费力那项研究也说，提高效能感后回避就消失。}{回扣：值｜门槛。判准写明「对自己」的代价分量偏低。一时抽身正方本来允许，Strawson 说可以一时借用、不能长久（\#34）。给人选择时人倾向回避共情，因为觉得费力；同一研究提高共情效能感后，回避就消失（\#98）。护理人员的共情疲劳来自创伤和重负荷，不来自别人蠢（\#99）。「你只是太累了」是自媒体的说法，学术上的共情疲劳特指助人者（\#18）。蔡康永那期「150 人」是随口引的数，不能当论据（\#111）。累了一时走开，是个别场合的调分寸；累成了「他就是那么蠢」，就是平常的样子要改。}{①「累了躲开你们也说可以。」→一时躲开可以；躲开之后回头还讲不讲理，是判准问的。②「把耐心留给肯学的人。」→分配的前提是分得清谁肯学；分不清时凭比较判成「不肯」，耐心就给错了人（推论）。}{淘汰／常见}
\rebuttal{24}{认错就消气：「他一问、一认错，恼就转成耐心；给了新的归因信息，情绪就跟着变。我们的厌本来就会回头。」}{这只证明改得掉；门槛规则：平常就这样、只是改得掉，照样算错了。}{回扣：平常｜值。转的条件是对方给的，不是这份厌自己走到的。Weiner 1988 那篇第一句：心理—行为类的毛病默认被归为起因可控，招来愤怒、少同情、倾向不管；第二句才是归因一变情绪跟着变（\#188）。第一句是平常，第二句是改得掉。「他一问」的前提是他敢问：承认错误或求助的人可能显得无能（\#64）；心理安全量表第一题就是「在这个团队里犯了错，常常会被记在账上」（\#222）。厌在场时，「他一问」更难发生（推论）。}{①「认错就消气，说明本来讲理。」→消的是气，判词消不消研究没测；要他先认错才回头，讲理就挂在他的表现上，和反方论点二「讲理不等他做对」矛盾。②「不敢问是他自己的问题。」→代价落在团队：组织里的人常常不愿披露错误、不愿求助（\#184）。}{淘汰}
\rebuttal{25}{同一把尺子先量自己：「把行为叫作蠢，反映的是人要求自己也期待别人遵守的期待；对己对人一把尺子，是公平，我蠢时也厌自己。」}{对己严，不等于对人可以看低；何况那句是作者写的研究动机，不是研究发现。}{回扣：准｜公·尊重。\#239 那句是 Aczél 等写在摘要开头的立论动机；调研里没有材料测过说「我厌蠢」的人是否同样要求自己。就算一把尺子，问题在尺子准不准、量完以后怎么待人：尺子是我的经验，量别人时只往「他本该会」一头偏（\#204、\#216）；量完说「他就是那么蠢」，是看低人（\#20）。新京报那篇写，这种思维更像「我都可以做到你为什么做不到」的抱怨（\#212）：尺子是「我」，不是「一把」。尊重小问问的是厌了以后还给不给他尊重。}{①「你们说我错，不也是在厌我？」→我们评的是你这份厌，不是你这个人：针对行为、促使交涉的是批评，针对整个人、促使退出的是轻蔑（\#160）。②「严是美德。」→严我们不反对，对事批评就是严；要改的是看低和不再讲理。}{淘汰}
\rebuttal{26}{开口问反而显得聪明：「研究说求助的人被看作更有能力；厌只落在不问也不改的人身上。」}{容易的任务上，求助不加分；厌蠢说的恰恰是「这么简单都不会」，那一格问了也不显聪明。}{回扣：公·归责｜值。\#227 写明任务难时效应更明显；容易任务的结果只有《MIT 斯隆管理评论》的报道：求助既不损害、也不提升能力评价（\#181）。厌蠢的对象按定义正是容易那一格。同一摘要开头一句更要紧：人们怕显得无能，所以不愿求助（\#227）；这份怕从哪来，承认错误或求助的人可能显得无能（\#64）。厌蠢越在场，「不问」越可能是被厌出来的，反方却拿「不问」当归责的线索（推论）；同一条线索不能两头都记在他身上。}{①「研究说那份怕放错了地方。」→放错地方的怕照样让人不问；把被厌出来的不问再算到他头上，等于厌了两次（推论）。②「我们只说在问的人不厌。」→小张评论区那句「坐高铁不是谁都会吗？这个也要教？」（\#119、\#225），冲着的恰恰是来学的人。}{淘汰}
\rebuttal{27}{愚蠢比恶更危险：「朋霍费尔说，对于善，愚蠢是比恶意更危险的敌人，不要再三努力同蠢人论理；厌蠢是正当的警觉。」}{朋霍费尔写完〈论愚蠢〉，下一节的标题就是「蔑视人？」，开头一句：被推向蔑视人的危险很大。}{回扣：公·尊重。〈论愚蠢〉是 1942 年底写的《十年之后》里的一节，紧接着「蔑视人？」开头说被推向蔑视人的危险很大（\#211，据德文自译）：他论完愚蠢，紧接着提防的就是蔑视人。他说的是强大权力扩张下的群体愚蠢，克服它靠解放，不靠说教（\#75），语境是纳粹治下，不是日常同事。「道德上的缺陷」是中文流传本的意译，原文是「人的缺陷」（\#133、\#75）。奇波拉那篇是戏谑文（\#73）。汉隆剃刀的本意是别把愚蠢当恶意（\#130）：反方先回答，你厌的是蠢还是恶？}{①「他原话是再也不试图用理由说服蠢人。」→原文是「再也不」，我方不改成「暂停」；不再说服和不蔑视，在他那里是两件事，反方只取了前一件（\#247、\#211）。②「蠢人才是最坏的恶人。」→那是自媒体的「有人认为」（\#124）；高晓松那句「坏人不是蠢人」正好回它（\#111）。}{淘汰／常见}
\rebuttal{28}{情绪无罪：「厌是自然生起的反应，不是我选的；亚里士多德说感到愤怒不经选择；『应当』以『能够』为前提。」}{亚里士多德下一句就是：受责备的是以某种方式去感受的人。第一秒我们不评，评的是第二秒以后你认不认同它。}{回扣：三问｜门槛。同一段亚里士多德紧接着说，被责备的是「以某种方式」感到愤怒的人（\#151）；Adams 承认愤怒不能随意选，结论却是存在「不自愿的罪」，点名了对他人的轻蔑（\#156）；康德那句后半是「表现出来就是冒犯」（\#48）。被引来说「不经选择」的作者，没有一位推出「所以无可指摘」。James 的说法受到有力批评，不能当「情绪不可控」的定论（\#153）。改的是换一种看法：换位思考 d=0.45，压住感受 d=−0.04（\#155）。「无所谓对错」也不等于「没错」，跟「我没错」的正面姿态矛盾。}{①「换看法也是要人压住。」→不是压，是换：压住体验无效，换位有效；重评降低了厌恶体验本身（\#154）。②「没证明错就是没错。」→「我没错」本身是正面的自我辩护，推定谁也不占。}{淘汰／常见}
\rebuttal{29}{远离是我的权利：「我只是远离，不嘲讽、不伤害；和谁深交是我的自由。」}{有权这样做，不等于这样做没错；题目问的是这份厌对不对，不是别人能不能管你。}{回扣：公·尊重。权利回答的是别人能不能干涉，判准问的是这份厌平常的样子该不该改。「远离蠢人是自己的权利」也不是网友原话，是记者「由此看出」的概括；那位网友原话里还有「她甚至可以耐着性子教对方」（\#168）。远离分两种：一时抽身可以；从此远离，就是把人当作只能回避、不能与之讲理的对象（\#33、\#34）。反方论点二自己说躲开只是一时，「从此远离」和它冲突。「不伤害」也要打折：刻板印象研究里的被动伤害，正是忽视、不与之交往（\#30，群体层面）。}{①「远离不等于不讲理，下次碰到照样讲。」→那就是一时抽身，我们不反对，但那就不叫「远离蠢人」。②「交朋友的心力有限，总得挑。」→凭兴趣、凭投缘挑可以；凭「他蠢」挑，判断又回到准和公。}{淘汰／常见}
\rebuttal{30}{诚实之德：「压抑厌恶、假装包容是伪善；承认自己厌、管住行为才是成熟；莎莉文也曾绝望。」}{我们从没要你压住不厌，「压住不厌」双方都排除；要改的是看法，不是表情。}{回扣：门槛。判准写明修正不等于压抑，「压住不厌」双方都排除，只用来挡稻草人。我方要的改是换一种看法（\#155）。莎莉文那句全文是：她在绝望中把这个话题暂时放下，一有机会又重新开始（\#142）：绝望之后是重新开始，这正是我方的改法；她是老师，海伦是聋盲。承认自己厌，就是题目前半句；人还会想要或不想要拥有某种欲望和动机（\#157），题目问的是承认之后认不认同。诚实地承认「我刚才觉得他蠢」，再问一句「他卡在哪」。}{①「换看法也是装。」→重评降低的是厌恶体验本身（\#154），不是嘴上一套。②「你们的改法太理想。」→换位思考是那项元分析里效果最大的一项；反方第 9 条还说「先问」本来就有，做得到。}{淘汰／常见}
\rebuttal{31}{动力之用：「厌蠢推动标准提高；实验里被轻蔑的人任务表现更好；Torvalds 也说为了更好的补丁这说得通。」}{被轻蔑的人表现提升，靠的是内隐自尊下降，攻击性也上升；Torvalds 原句后半是「现在我知道这不 OK」。}{回扣：值。Melwani 与 Barsade：被轻蔑的人任务表现更好，对表达者的攻击性也显著更高，中介是内隐自尊下降（\#31、\#87，实验室短期任务、虚拟伙伴）；它测的是被「表达」的轻蔑，反方主张不看低、不表现，引它等于自认表达了轻蔑。Torvalds 原话：在追求更好补丁的过程中这在我看来说得通，现在我知道这不 OK（\#139）；同一封信说他的行为可能把一些人彻底赶出了内核开发（\#186）。2023 年那道辩题只查到公示（\#112）。守质量的好处，对事批评拿得到。}{①「他道歉的是攻击，不是厌劣质的情绪。」→我们只说他道歉的是攻击；但他也说要学「理解别人的情绪并恰当回应」（\#172），内核照样拒绝不合格的补丁（\#221），守质量用不着看低人。②「高标准本来就让人难受。」→难受可以来自对事批评；内隐自尊下降、攻击性上升是轻蔑的代价。}{淘汰／常见}
\rebuttal{32}{结构反用：「媒体自己说，厌蠢源于效率社会、优绩主义；根在结构，错的不是『我』。」}{根在结构，说的是这份厌从哪来，不回答它对不对；而且「蠢多出于处境」这半句，正好送给正方。}{回扣：准｜公·归责。判准问的是这份厌平常的样子该不该改，结构造出来的态度照样可以该改，辩题本身就是第一人称的反思（\#157）。同一批文章还写，这个标签「给了我们合理化贬低他人的契机」（\#13），调侃成病「都不能掩盖背后浓浓的傲慢和尖刻」（\#14）。反方借「结构」，就得承认结构也造出了被厌的「蠢」：新京报写常识本是由个体社会经验组成的，要求乡村的人立刻拥有城市常识「本就是一种苛责」（\#213），那反方的「本可以」就松了。结构越强，越说明这份厌是平常，不是个别人的脾气。}{①「我们只是说我不是坏人。」→我们也不说；「我错了」是该改，不是自认坏人。②「结构的错该由社会改。」→社会改社会的，个人改自己那份，两件事不冲突（推论）。}{淘汰／常见}
\rebuttal{33}{后果之厌：「厌蠢本质上是厌麻烦、厌风险；新京报写厌蠢很多时候是厌麻烦，医生说背后是对潜在风险的厌烦。」}{讨厌麻烦叫怕麻烦，用不着「厌蠢」这个词；厌麻烦冲着事，厌蠢冲着人。}{回扣：准｜公。「厌烦」的对象是「某事」（\#6），「厌蠢症」的通行说法是「看不得别人犯蠢、没常识、太笨」（\#10），落在人身上。同一篇新京报后文转向：这个标签给了我们合理化贬低他人的契机（\#13）。每日商报那篇，明确写作医生原话的只有「对潜在风险的厌烦」一句，标题写的是专家说厌蠢的实质是「慕强」心态（\#15）。咨询圈那篇是自媒体，只能算「有人认为」，原句前还有「很多情况下」（\#17）。真只厌麻烦，对麻烦生气、批评都可以，那就不是辩题里的厌蠢；把对象挪到事，题目也跟着换了。}{①「蠢就是麻烦的来源。」→来源是一件事还是一个人？厌那件事我们不反对；把来源定成「这个人蠢」，就要过准、公两问。②「风险是真的。」→风险管理那头说，怪个人在情绪上比追究机构更痛快，最好的人也会犯最糟的错（\#183，只作机制例证）。}{淘汰／常见}
\rebuttal{34}{理性之守：「对拒绝思考保持不耐烦，是对常识与理性的坚守；厌蠢是对认知惰性的理性批判，是觉醒，不是傲慢。」}{守理性，先过理性这一关：看准了吗？把处境读成蠢、把不会读成不肯，恰恰是理性没守住。}{回扣：准｜值。用反方看重的理性来量这份厌：对应偏差就是从完全可以由情境解释的行为推出持久特质（\#61）；知道答案的人高估不知情的人本来能知道多少，自己察觉不到（\#204）。理性的厌蠢者，第一件事该是怀疑自己这一下。治反智靠纠正和提醒：人分享错误信息常常是因为注意力不在准确性上，轻轻提醒就能改善（\#108）；事实纠正没有引发逆火（\#109）。北大战报那句「是觉醒还是傲慢」是一个问句，不是答案（\#113）。}{①「对谣言都不反感，那是没有是非。」→对谣言反感是对事，我们不剔；对信谣言的人，先问他从哪看来的（推论）。②「信假新闻是懒不是偏。」→那项研究测的是认知反思得分与信假新闻的相关（\#107）；同组作者另一项说注意力一转就能改善（\#108）。}{淘汰／常见}
\rebuttal{35}{限定之蠢：「我厌的不是不会，是又菜又犟、理直气壮还帮倒忙、学了也不改的人；朋霍费尔说人会任由自己被弄蠢。」}{两条路都输：只厌「犟」和「帮倒忙」，那是厌坏，不叫厌蠢；要还算蠢，你当场怎么分清「不想知道」和「不能知道」？}{回扣：准｜公·归责。把「理直气壮、损人、帮倒忙」写进蠢，是把蠢偷换成坏（3.5 反方定义杀②）；高晓松那句「坏人不是蠢人」（\#111）。退一步只留「不想知道」：当场看到的是同一个动作，「别人都会、他不会」那一格，人归于能力、努力、运气三样，分不出（\#241）。苏格拉底那句是「他什么都不知道却以为自己知道」（\#146）：自以为知的，也可以是觉得「这都不会」的人。分不清时先问一句「你卡在哪」，问出来是不肯，你怪那件事，我们不拦；把人判成「就是蠢」，是另一回事。}{①「我们不写进定义，只说平常厌的是这一类。」→那就要证平常：写同学同事的那段是「你能不能有点常识」，不是「你怎么那么犟」；作者还写被叫作蠢的人「有许多却是努力了但也没有学会的人」（\#212）。②「朋霍费尔说人会任由自己被弄蠢。」→同一节说权力扩张让许多人变蠢、克服靠解放（\#75）；下一节就提防蔑视人（\#211）。}{淘汰／常见}
\rebuttal{36}{切厌：「汉典里厌只有嫌恶、反感、嫌弃，看不起是鄙视的释义；厌恶对译 disgust，是把自己隔开。厌蠢就是烦，正方的轻蔑材料全落空。」}{定义里两头都写了、都用「可以」，平常落在哪头是今天要辩的；而心理学说，评判「无能」时起作用的，恰恰是轻蔑。}{回扣：平常。中立定义不替任何一方选端点，我方不把看低写进定义，在论证里证它是平常的样子。证据：轻蔑在评判无能的行为时起作用（\#26）；轻蔑把对方贬为愚蠢、无能，贬品性多于行为（\#20）；轻蔑与厌恶在表情识别里常被混淆（\#24）。退一步，就算是 disgust：对人产生 disgust，去人性化的色彩更重（\#21）。中国网写有的人「会用行动离对方远远的」（\#16）。这一条不是定义之争，是事实之争；只留对事那一端，是反方定义杀③。}{①「Hutcherson 摘要写的是无能的行为。」→起作用的是轻蔑这种情绪，而轻蔑只指向人（\#27）。②「那你们承认看低只是可能。」→定义里两头都在，平常落在哪头靠论证，见第 11、17 条和正方论点二。}{常见}
\rebuttal{37}{切蠢：「第一次不会坐高铁按中立定义不算蠢；正方的知识诅咒、信息差打的都是误判；我厌的是第三次还问同一个问题。」}{那份厌照样算厌蠢：定义写的是「觉得」别人犯蠢；看错，正是第一问要称的错，而且不是哪一次眼花，是平常那把尺子造成的。}{回扣：准｜平常。中立定义用「觉得」不用「看到」，就是为了把误判留在题目里：说「这都不会？」的人当时就厌了，那份厌就是厌蠢。误判来自一把平常就用的尺子：人估别人知道什么，起点就是自己知道什么，常常不加批判（\#216）；知道答案的人察觉不到自己高估了别人（\#204）。媒体批的也正是这把尺子：每日商报的小标题是「不要用自己的『常识』去攻击别人的『无知』」（\#214）。既然第一次不会不算蠢，厌它就是看错，而看错用的尺子平常就在用。}{①「我第三次问才厌。」→第三次问是有记录那一格，准我们让；接下来看那份厌往哪走（第 10、17 条）。②「误判只是少数时候。」→不报比例，双方一样；我方说的是方向单一、察觉不到，提醒也压不下（\#205，认脸任务）。}{常见}
\rebuttal{38}{只看行为：「情绪本身没有对错，只看是否转化为伤害；我只在心里厌，不骂不欺负；冯小刚错在公开骂观众。」}{题目问的是厌，不是骂；心里那份就是厌蠢本身，它错不错看三问，不看有没有说出口。}{回扣：三问｜门槛。「无所谓对错」不等于「没错」；只看伤害，是把评价对象从厌换成了霸凌。心里的态度也有对错：Adams 点名对他人的轻蔑是不自愿但可受责备的过错（\#156）；SEP 概括康德的尊重义务，虽不要求「感到」尊重，却包括不持有轻蔑、嘲笑这类态度（\#43）；康德那句整句念：心里看低有时不可避免，但表现出来就是冒犯（\#48）。心里那份厌决定你回头还跟不跟他讲理，越线看的正是回头。反方自己退守说「这份厌经得起评价」，那就请按三问来评。}{①「SEP 明说没有感到尊重的义务。」→同一段又说这些义务包括不持有某些态度，两句要一起念。②「Torvalds 道歉的是人身攻击。」→我们也只这样说；可他同时说要去学理解别人的情绪（\#172）。}{常见}
\rebuttal{39}{被无知伤害过：「新国辩反方结辩的标题说：厌蠢的人，大多真的被无知伤害过；受伤后的厌是正当反应。」}{对伤害你的那件事愤怒，正当；从一次伤害推到「这一类人」，是泛化。你怎么确定伤你的是蠢，而不是坏、不是流程？}{回扣：准｜公。对造成损害的那个行为的愤怒，我们不争；对「蠢」这一类人的泛化厌恶，要过「准」：从完全可以由情境解释的行为推出持久特质是人普遍的倾向（\#61）；还要过「公」：最好的人也会犯最糟的错（\#183），有些不安全行为确实恶劣，但绝大多数不是（\#169，专业领域只作机制例证）。那个标题是账号编辑起的，结辩内容没有文字稿（\#114），「大多」不是统计。厌冲着那件事，三问都过；冲着「蠢人」这一类，准和公都失守。受害者的愤怒值得尊重，这正是我方不剔恼火的原因。}{①「被伤害过的人有权厌。」→有权，见第 29 条；问题是厌冲着谁。②「受害者的情绪你们也要评，太冷血。」→我们不评受伤那一下，评的是第二秒以后你认不认同「他们都蠢」（推论）。}{常见}
\rebuttal{40}{理由判准：「一份厌只看对不对得上对象；Bell 论证轻蔑有时恰当；亚里士多德说骄傲的人轻视别人正当；该怒不怒是愚蠢；从坏结果推出不恰当是道德主义谬误。」}{对得上只回答了一问；道德主义谬误是双向的，对得上同样推不出没错。Bell 自己也没把蠢列为恰当对象。}{回扣：三问通盘。三问都进场，只赢一问不自动赢下整场；契合不等于道德上恰当：性别歧视的笑话也许确实好笑，被它逗乐却不恰当（\#164）。Bell 说轻蔑有时是回应傲慢、伪善这类优越之恶的最好方式（\#35），自恃聪明看低别人的人倒成了恰当轻蔑的对象；她的七条里，常态的厌蠢至少「有证据」「公平」两条过不了（\#36）。亚里士多德同一句后半：多数人的轻视是随意的（\#162）。「该怒不怒是愚蠢」说的是愤怒，主要针对受到轻慢（\#39），恼火我们本来不剔。}{①「针对品性某一方面、保留改变可能的轻蔑可以正当。」→条件就是我方的改法（\#22）；「他就是那么蠢」两条都不满足（\#20、\#27）。②「我们赢了准和归责就够了。」→「公」还有尊重小问。}{常见}
\rebuttal{41}{古今都说别跟蠢人纠缠：「《箴言》说不要照愚昧人的愚妄话回答他；席勒说与愚蠢作战众神也徒劳；还有别和傻子争论、可以原谅坏人不能原谅蠢人……」}{先查出处：后面几句是伪名言或找不到出处，校赛可当场申诉；前面几句读全文，都不站在反方那边。}{回扣：防守条，这些引句一问都答不了。「别和傻子争论」找不到出自马克·吐温的实质证据（\#127）；「宇宙和人类的愚蠢」大概不是爱因斯坦说的（\#132）；颜回「三八二十三」《论语》里没有（\#135）；「可以原谅坏人不能原谅蠢人」找不到可靠出处（\#134，【待核实】），只请对方出示，不说它是假的。《箴言》第 5 节紧接着说「要照愚昧人的愚妄话回答他，免得他自以为有智慧」（\#128），14:7 也没有「离开」的意思（\#129）。席勒那句是剧中濒死败将的诅咒，他骂作愚蠢的正是全剧歌颂的贞德（\#131）。「suffer fools gladly」源自保罗的反讽（\#136）。}{①「名言真假不影响道理。」→道理要自己论证；拿伪名言撑道理，章程规定判负。②「写进讣告说明是被接受的性格。」→维基原文只说它带着各种含义，没有说是褒义（\#136）。}{常见}
\rebuttal{42}{连蠢都能装：「职场流行装蠢；连蠢都能装，说明场上谈的蠢大多是策略和态度，厌一种态度有什么错？」}{装出来的蠢是偷懒，冲着它的是对坏的不满，不叫厌蠢；能装成功，恰恰说明旁人分不清真不会和装不会。}{回扣：准｜值。装蠢是能做却不做，按中立定义不是「明显想错、做错」，是偷懒。同一篇写那位扮蠢的大学生挂在嘴边的是「完了我不会搞」，同学「不会把最难的部分分给你」（\#237）：蠢的标签一贴，机会就少了（推论）。「大多是装的」也说不得：那是非虚构写作，受访者化名，没有数据（\#123）。装得成，说明从外面看是同一个动作（推论），这正是「准」的难处。真不会的人被当成装的，就是平常的误判。}{①「真有人装，厌他不对吗？」→对装的人生气是冲他偷懒，我们不反对；辩题是厌蠢。②「分不清就宁可当他是装的。」→宁可错判，代价落在真不会的人身上：默认读成可控，就是多愤怒、少同情、倾向不管（\#188），这正是要改的平常。}{常见}
\rebuttal{43}{厌蠢症多是自嘲玩梗：「用个『症』字自嘲，本来就是给刻薄打折；说它错是上纲上线。」}{「症」字是给刻薄裹的糖衣；题目写的是「我厌蠢」，没有「症」，评的是那份厌，不是那个梗。}{回扣：平常。有间大学那篇正是说，用病症的说法自嘲来中和刻薄，但「无论怎样把它调侃为一种疾病，都不能掩盖背后浓浓的傲慢和尖刻」，「我只厌打扰到我的蠢」这个前缀是为了不显得傲慢才加的（\#14）。评价对象是常人说「我厌蠢」时那份厌平常的样子：轻的那头是对一件蠢事的烦，我们不反对；要改的是从「这都不会？」滑到「他就是那么蠢」那一截。「智商税」那条，受访者说别人用开玩笑的语气说，「我会觉得非常不礼貌且傲慢」（\#125）：说的人觉得轻，听的人未必。}{①「玩梗没人当真。」→不当真的不在题目里；题目问的是当真的那份。②「一点烦就说它错，太重了。」→「我错了」是该改，不是该受罚；我们也不改那点烦，改的是滑过去的那一截。}{常见}
\rebuttal{44}{反方分量：「三问里归责小问分量最重，常人判断一份恼火对不对，本来就看该不该他负责；有记录时，看准和归责一起过关。」}{那句称的是恼火，对那件事的恼火我们本来不反对；我们要改的是看低他、回头不讲理那一截，称这一截的是尊重，不是归责。}{回扣：公·归责｜公·尊重。先认：情绪按责任分流，常人确实这样称恼火（\#58、\#219）。可判准写这句，是说三问合乎常人的判断方式，没说归责压过其余几问。分量要落在争点上：全场决胜点是这份厌平常停在事上，还是落到看低人、回头不讲理；那一截不是恼火，称它的是尊重小问：拒绝给人作为人应得的尊重，无论如何违背义务（\#48），「无论如何」就包括他该负责的时候。Weiner 1988 那项实验里，被看作起因可控的那一路配的是多愤怒、少同情、倾向不管（\#188），往外推的倾向和我方论点二同向（推论）。}{①「看低是你们加上去的。」→中立定义两头都写了，平常落在哪头靠论证（\#26、\#20）。②「倾向不管说的是帮不帮。」→对，所以标了推论（原文 “judgments to neglect”）。③「各说各的。」→我方理由落在争点上：看不准，后面讲公道是在错的判断上讲；看准了，也不等于可以看低这个人。}{最终}
\rebuttal{45}{教师期望更多因为准：「三十多年的教师期望研究说，期望能预测学生表现，更多是因为期望本来就准；人看别人行不行，大体看得准。」}{先搭桥：那篇测的是老师对学生学业表现的预期，不是对「蠢」的判断，原文还是「可能」。老师有一学期的作业可看，「这都不会？」那一下手里什么也没有。}{回扣：准。Jussim 与 Harber 的综述，对象是教师期望与学生表现（\#255、\#257），原文是 may。第一，测的是学业表现，不是中立定义里的蠢；同一篇还写影不影响智力仍不清楚。第二，老师的期望通常是看过作业、考试之后才形成的（推论）：那是记录，凭记录的准我方本来不争。第三，同一篇写着，强有力的自我实现预言可能选择性地发生在受污名群体的学生身上（\#257）：看低的预期最可能在已被看低的人身上应验（推论）。手里有证据的人大体看得准，恰好说明证据要先有。}{①「同事天天见，也有记录。」→有记录的那一格我们让（第 4、5 条）；争的是平常下判断时用的是不是记录（第 1、6 条）。②「效应很小，厌了也无害。」→我方论点不靠「贴标签把人变蠢」；何况受污名群体里可能更强。③「陌生人看录像判断智力也有效度。」→那两项测的是智力，同样要先搭桥（第 21 条）。}{淘汰}
\end{rebuttals}

\consection{八、反方反驳库（针对正方全部可能论点）}
\begin{fields}
\field{全队共守的六条线}{\sub{①尊重线}{看低他这件事上的表现可以，拒绝给他作为人应得的尊重不行；我们的厌停在前一边（\#44、\#48，康德两句必须整句引）。}\sub{②越界按时间划}{这回先不劝、一时走开、不再主动提醒，不越线；他开口也不听、合作只下指令不给理由、认定「他这个人说什么都没用」，越线。越线我们认要调分寸，那是个别场合。}\sub{③记录线}{没有记录先问；新来的、没人教过的先教；在问、在试、在改还学不会的给耐心；只凭「一贯做不到」不厌（\#241）。「教过、被指出、又做错」这串记录，在「准」上双方都按它判。}\sub{④每日商报一案三句一起说}{「多次劝告」「一度降至冰点」「很微妙」；「说了那么多遍，怎么还是那么蠢」没有标说话人；「无所谓」「转眼间就将嘱咐抛诸脑后」是小陈一方的叙述（\#248）。}\sub{⑤「平常」只有一条主线}{厌跟着代价走，代价反复落在一起做事的人身上，那里他被教过没有、改没改是现成的记录（第 1 条）。其余各条说「那是个别场合」，都回指这一条；双方都只有描述，没有比例。}\sub{⑥禁说}{「情绪控制不了，所以无所谓对错」「我有权远离谁」「新京报认为厌蠢没错」「多数厌蠢冲着拖后腿的同事」「恼还在，所以里面没有轻蔑」「轻蔑要三个前因凑齐才出现」「Torvalds 至今认为骂人说得通」；个别场合不说「那一份我们认错」，只说「那一次是看错了，要调分寸」。}\sub{条数}{共 43 条：正方最终论点及其单独抛出的子说法 23 条（1–4、24、26–43），淘汰候选 9 条（5–13），公开讨论常见说法 11 条（14–23、25）。}}
\end{fields}
\begin{rebuttals}
\rebuttal{1}{拿我当一般人（上）：「本该只有两种依据，记录或比较；写厌蠢的文字给出的平常说法落在比较：『这都不会？』『我都可以做到你为什么做不到』。所以这份厌平常凭的是比较。」}{厌跟着代价走，蠢的代价天天落在一起做事的人身上；他被教过没有、改没改，我手里本来就有。评论区那句最响，不是厌蠢最常住的地方。}{回扣：准｜门槛。研究发现，事情与自己相关推高的是愤怒（\#224）；新京报的文章写道，对蠢的厌恶更多源自它可能造成的后果，想的是「生活中有这样一个人要怎么办」（\#252）。蠢的代价反复落到我身上，是在同一个组、同一间办公室：每日商报那位同事工作超过两年、被多次劝告、同样的错接二连三，关系一度降至冰点，现在很微妙（\#248），下判断时记录是现成的。网上的愤怒得到点赞就更常表达（\#101，测的是推特，迁移是推论），评论区是被放大的声音；新京报写完评论区，下一段就是「当然，在现实生活中」（\#243）。评论区那句我们认：算厌蠢，也看错了，是个别场合。}{①「新京报写同学同事，原话就是『我都可以做到』（\#212）。」→全认；比较的口气说的是「本该」，没说出口的是他被教过没有，熟人之间那一步手里有（推论）；只凭「我都可以」的，我们认该先问。②「记录只出现在一个化名个案里。」→新京报那段只取场景，记录由每日商报补；我们的「平常」压在机制上。③「出圈现场就是定义这个词的场景。」→这个词元旦前后就出现了，小张事件是出圈不是出生（\#8、\#167）；出圈靠的是好批，好批不等于常见（推论）。}{最终}
\rebuttal{2}{拿我当一般人（下）：「知道答案的人高估不知情的人、自己察觉不到、提醒压不下；替不认识的同龄人估也一样；所以要改成先开口问。」}{这些实验估的都是一个你一无所知的人。被教过、劝过多次的同事不是不知情者，他学过没有，查得到，用不着估；对方的药只治得到没有记录的那一种。}{回扣：准｜门槛。只看估的是谁：Fischhoff 的被试估不知道结果的「别人」（\#204，只核摘要）；Hom 的被试替一个不认识的同学估，还被要求设想他和自己经历相同（\#176），这条指示本身推着人投射；Birch 估的是同龄人里有多少人知道一道知识题（\#215，方向据 Hom 转述）。偏差长在「对他一无所知」的地方。厌蠢平常冲着天天一起做事的人，这串记录对方在准上不争。对方的药是开口问；问是为了拿到记录，对共事的人记录本来在手，没有记录的我们也认该先问。}{①「位置要求的『人人』还是你估的。」→位置上的基本功写得出、教得出；现实里说的蠢，是缺乏所处位置需要的基本知识和技能（\#243）。②「方向固定、提醒压不下就是平常。」→我们不评偏多少；这把尺子只在没有记录时才是全部依据。③「Nickerson 说起点是自己。」→同一句前半说这总体上对人有用，出错在不加批判（\#216）。④「教过也是你以为，说话的人总高估自己说清楚了。」→知道本意的旁听者没有系统性高估（\#234）；判断看他之后改没改。}{最终}
\rebuttal{3}{说了也没用（上）：「就算看准了；轻蔑贬品性多于行为；劝了还犯把轻蔑的预测因素现成供上，轻蔑可能从一次次的恼火里长出来。」}{还在劝，就是还觉得话管得到他。恼里那句「你本可以」，只对做得到的人说，和轻蔑「不看好他改」正好相反。}{回扣：公·尊重。逐个看预测因素。觉得管不了他（\#185）：劝了不听，这回放下、下回照样说，就是还觉得话管得到他。品性归因：生气本身跟着「他该负责、他做得到」的判断走（\#219）；1984 年的实验里，大人冲做砸的孩子生气，孩子更多想的是「我没用心」，表示同情反倒更多觉得「我能力不行」（\#250，儿童样本，生气这条联系较弱）。没用心是下回改得了的事；轻蔑的特点恰恰是不看好对方改善（\#27）。同一篇研究说生气长远会和解（\#86）；轻蔑「may」从生气里长出来，是可能，不是平常的走向。真硬化成「说什么都没用」，我们认越线，那是个别场合。}{①「愤怒和轻蔑常一起出现。」→我们不说恼里没有轻蔑，看混着的那份往哪走；去人性化的是不掺愤怒的纯粹轻蔑（\#38），客观态度装不下愤怒（\#197）。②「你们辩护的只是生气。」→中立定义不写「强烈」，烦躁、难以忍受都在厌蠢里（\#16）。③「怎么还是那么蠢，宾语是人。」→这句没标说话人；「怎么还是」还在要求他（推论）；平等尊严不排除区别评价（\#44）。}{最终}
\rebuttal{4}{说了也没用（下）：「认定他改不了，人就不再教；拒绝基本尊重违背义务；每日商报一案劝了多次、关系冰点，接着是『还是那么蠢』；回头他开口先不听，就越线。」}{报道写的是劝了多次、冰点是「一度」、现在「很微妙」，没写从此不讲；回头讲不讲理，对方自己也标了推论。}{回扣：公·尊重｜门槛。三句一起读：多次劝告，讲理没等他改好；关系「一度」降至冰点，是降到过；现在「很微妙」，不是断了往来（\#248）；那句话没标说话人，之后讲不讲理报道没写。撑推论的经理研究（\#218，只核摘要）测的是对人性的一般看法，不是厌，不是对一个人的判断，辅导下属是经理的职务；它还替我们说话：决定教不教的是「人能不能改」，恼里的「你本可以」站在相信他改得了的那边（推论）。康德附释管的是指出错误时别斥为荒谬（\#45），「这回先不劝」没有否认他有理智。越线要证是平常，对方只有一个没写结尾的报道。}{①「完全客观时能谈判、不能讲理（\#174）。」→方向反了：完全客观，所以不能讲理；推不出没讲理就是完全客观；混着恼的厌到不了「完全」（\#197）。②「只在心里看低也会漏出来（\#156）。」→康德整句：表现出来就是冒犯（\#48），所以我们管住表现、守住尊重；Adams 是一派哲学论证，不是实证。③「这件事上从此不听他算不算越线？」→算；它是个别场合，正常人不能长久这样（\#34）。}{最终}
\rebuttal{5}{厌蠢逼人藏拙：「这份厌让露怯变贵，人怕显得无能就不认错、不求助；心理安全越低，学习越少；厌蠢的氛围逼人装蠢自保。」}{开口问的人不在我们厌里，我们厌的是不问就乱做、教过还犯。这份厌让「不问」变贵，让「问」变便宜。}{回扣：值。这条藏着一个前提：厌冲着「不会」本身。上游新闻引的一位早期自称厌蠢症的网友说，遇见不会但努力的人不会激发她的厌蠢症，她甚至可以耐着性子教对方（\#168，二手转述，代表性低）。人怕求助显得无能，这种担心放错了地方：求助的人被看作更有能力（\#227）。Edmondson 的心理安全，原文定义是相信团队不会因为某人开口而羞辱、排斥或惩罚他，并特意说它不是随意的放任（\#222）；怕显得无能、不愿认错是组织里常见的倾向（\#184），没有厌蠢者也在，没有一条数据说厌蠢让它更多。这条只打「值」，单赢一问赢不下整场。}{①「对方只看见你不耐烦，于是不敢问。」→那要改的是表达，是个别场合的分寸；表现出来的看低我们不辩护（\#48）。②「量表第一题：犯了错常被记在账上。」→题目说犯一个错，我们厌的是提醒过还犯同一个错。③「装蠢流行。」→装蠢的大学生自述好处是少分到难活、大家更照顾她，文章说「本质也算一种自私自利」（\#237）。④「问容易的事加不了分。」→据报道，任务容易时求助既不损害也不提升能力评价（\#181）。}{淘汰}
\rebuttal{6}{蠢多不由己：「判断力出于天赋、教育、处境；没有人应得他更高的自然能力；厌他是把运气当应得。」}{不会的，我们先教、给耐心；我们厌的是被教过、有时间、还不上心的那一份。孔子说「困而不学，民斯为下矣」：困不怪，不学才怪。}{回扣：公·归责。中立定义里的蠢，是了解他处境的一般人看来本该会；第一次、没人教的，本来不在里面。我们守的是被教过、有时间、还不问不改的那一份，古人划得清楚：没人责备天生的盲人，人人责备醉酒致盲的人，对应当知道又不难知道的事无知要担（\#41、\#230，只借判别线）；孔子按困了学不学排等次（\#49）；孟子说为长者折枝「是不為也，非不能也」（\#231）。汇总六十四项研究的元分析发现，该负责就生气，控制不了就同情（\#219）。对方遗传、教育那两条测的是智力（\#57、\#82），不搭桥用不上。}{①「当下分得清吗？人默认读成可控（\#188）。」→分不清时先问，有记录看记录；只凭一贯做不到读出的多是能力和难度（\#241）；同一研究说给了信息情绪就变。②「不肯也有成因。」→照这个逻辑世上没有可以怪的事；判准约定写明「公」也计入不肯学；人会「任由自己被弄蠢」（\#75）。③「请承认对真正的笨，厌是错的。」→对努力了还学不会的人厌他，那一次是看错了，要调分寸；它不是常人厌蠢平常的样子。④「何不食肉糜。」→惠帝错在不知道饥民的处境（\#137）；我们没记录先问，问的正是处境。}{淘汰}
\rebuttal{7}{标签止于骂：「骂一句『蠢』是回应的开始也是结束；厌蠢的逻辑下焦点从事情转到愚蠢；污名把人贬为被打折扣的人。」}{骂出口的「蠢」才是短路，我们不骂；心里那份恼想的是「他这回又没上心，下回怎么办」，焦点就在事上。}{回扣：准｜公。穆齐尔说的是「蠢」作为骂人话的用法（\#77），不等于一切厌蠢都是短路；中立定义写明厌不等于骂人，骂出口的看低，康德说就是冒犯（\#48）。心里那份厌想的是什么？新京报的文章写道，人会延伸到「生活中有这样一个人要怎么办」，对蠢的厌恶更多源自它可能造成的后果（\#252）：焦点在后果、在下回怎么办。硬核读书会那篇和有间大学那篇是同一作者同一篇（\#122、\#14），只能引一次，说焦点转移是描述，不是机制。Goffman 谈的是污名（\#163），那是对整个人的判决；我们下调的是这一次、这一方面的评价（\#44）。}{①「不说出口，心里也贴上了。」→标签贴在人身上才止于骂；「他这一处又没上心」是对一件事的评价，下回照样说（\#34）。②「焦点在后果，那是厌麻烦。」→后果只说厌在哪里冒出来，对象仍是本该会却明显做错；只有麻烦没有蠢的场合不叫厌蠢，我们也不替它说话。}{淘汰}
\rebuttal{8}{尺子量回自己：「人人都会蠢，聪明人也会做蠢事；厌蠢可能是『恐弱』，讨厌的往往是愚蠢的自己；智慧始于知道自己无知。」}{「你也会蠢」推不出「你的厌错了」。我蠢的时候也厌我自己，一把尺子量两个人，这是公平，不是傲慢。}{回扣：准。这是人身反诘：就算我也会犯蠢，推出的是我也该被这样要求。研究人们管什么叫蠢的那篇论文，作者说明动机时写，把行为叫作蠢反映的是人拿来要求自己、也期待别人遵守的理性期待（\#239，作者的说明，不是发现）。觉得自己比一般人聪明、自评与实测相关不高（\#91、\#92）测的是智力自评，标〔智〕。「恐弱」是新京报书评作者的论断（\#13）；「讨厌的是愚蠢的自己」是中国网特约稿的说法，投射没有经过实证（\#16）。苏格拉底那句我们接得住：不懂又自以为懂，正是人们叫作蠢的第一类（\#146、\#93）。}{①「朋霍费尔：我们在别人身上蔑视的，没有一样是自己全然陌生的（\#211）。」→说的是对人的蔑视；我们守的是对事，不蔑视人。②「人对别人总比对自己苛刻。」→请拿出常人厌蠢平常就双标的材料；就算有人双标，那是分寸。③「鲁迅笔下的傻子是唯一肯开窗的人。」→那是反讽的称呼，傻子不是真蠢（\#145）。}{淘汰}
\rebuttal{9}{小厌汇成偏见：「个人的厌蠢汇成对低学历者、老人、进城务工者的偏见；精英主义信念越强越归咎；精英身份的群体最容易厌蠢。」}{学历、年龄、出身是没选过的处境，我们先教；我们恼的是教过还犯的一件件事，不是一群人。题目评的是一个普通的「我」，不是社会现象。}{回扣：公。判准约定写明评价对象是一个普通的「我」那份厌；从个人的厌走到群体偏见，中间那一步「汇聚」没有一条数据。Kuppens 测的是对「学历」这个群体身份的态度（\#54，西方样本），不是对一个本该会却不改的人的反应。新京报的文章写道，现实里说的蠢，是缺乏所处位置需要的基本知识和技能（\#243）：是位置上的要求，不是出身上的标签。把一个群体的处境算成活该，我们和对方一样反对。智商税、嘲笑受骗老人是嘲笑，是表达，我们不辩护。「精英最容易厌蠢」是作者论断，没有数据（\#13）。}{①「研究六：愤怒由责任归因中介（\#55）。」→对一个群体做责任归因才是错配；同一研究只报了愤怒的中介，没报轻蔑（\#171）。②「精英主义信念让人把处境怪到他们头上（\#56）。」→那是相关；我们归的是一件被教过、又做错的事。③「你们厌的同事多是条件差的人。」→没有这样的数据，我们也不说。}{淘汰}
\rebuttal{10}{身份要求耐心：「教师准则要求诲人不倦，大医精诚要求普同一等；费曼、莎莉文、特教老师都是榜样。厌蠢就是失职。」}{那是职业要求，不是普通人的门槛；而且准则同一条写着「严慈相济」「严格要求学生」，禁止的是歧视侮辱，不是心里的恼。}{回扣：公。职业角色不当本题主战场，它最多证明某些身份下要求更高。就算拿来看：教师准则第五条同一句写着严慈相济、严格要求学生，禁止的是歧视、侮辱学生（\#147）；大医精诚管的是救不救、怎么救，没有要求医者心里不恼（\#148）。莎莉文那段，原话是她在绝望中暂时放下这个话题，又在下一个机会重新开始（\#142）：有绝望、有放下、有再来，正是我们说的这份厌；海伦是聋盲，不是蠢。费曼那句说的是讲不清是讲者自己的问题（\#140）；特教老师面对的是有障碍的孩子，那是不会，我们本来就给耐心。}{①「连老师都要求耐心。」→职业把门槛抬高，恰恰说明普通人的门槛不在那里。②「你的讥笑中有爱迪生。」→讥笑是表达，我们不辩护；这句原诗我方没有打开原文（\#144，有把握），请对方说明出处。③「常善救人，故无弃人（\#51）。」→我们不弃人：这回放下，下回照样说。}{淘汰}
\rebuttal{11}{坏印象改不回：「人只愿意和印象好的人来往，坏印象更稳；被判为坏的东西被回避了，错误的负面预期抵抗改变。厌了就躲，误判锁死。」}{锁死的前提是判错了；判对了，少接触不丢任何信息。同一篇还写着，对经常接触的人，坏印象更容易改，厌蠢的典型对象正是天天见面的同事。}{回扣：准。Eiser 那篇模拟的「偏见」，定义就是对其实是好的对象抱着错误的负面预期（\#236，计算机模拟）：咬人的前提是看错了，那是「准」的事，没有记录时先问。Denrell 的机制里没有情绪变量，只有当前印象（\#235）：任何坏印象都让人少来往，对方要定罪的不是厌，是有坏印象本身。同一篇摘要写着，对经常接触的人坏的初始印象更容易改（\#193，模型推导）；常人厌蠢的典型场景是同一间办公室，见面躲不开。「改不回」只落在少来往的陌生人身上，那是个别场合。}{①「接触回来不等于讲理回来。」→同意；本条打的是误判锁死，讲理见第 4 条。②「同一摘要前一句：坏印象更稳。」→整句认；它说的是印象稳不稳，判对了的印象稳，不是错。③「对陌生人看错了就改不回。」→对陌生人只凭一句评论下判断，我们认是误判、个别场合。}{淘汰}
\rebuttal{12}{事后才显得蠢：「同一份防洪材料，事前看的人认为不值得再防，事后看的人判疏忽；知道结果的人把『本该想到』算到别人头上。」}{我们的「本可以」不看结果：他被教过没有、教过之后改没改，都在结果出来之前就摆着。事后诸葛只管「想得到」，管不了「做得对」。}{回扣：准｜公·归责。Kamin 与 Rachlinski 判的是一场洪水当初该不该防：事前本来就是不确定的风险，又是法律上的过失认定（\#206，只核摘要），专业情境只作机制例证。常人厌蠢平常判的是位置上的基本功，教过就会，不需要预测；判断他本可以，靠的是之前教过、劝过。对方自己的淘汰记录也承认，「事后」只覆盖想得到，覆盖不了做得对；注意失误这一类也不靠事后才知道（\#93）。「知道了就觉得早该料到」（\#207）管的是结果的可预见性，不管一个人被教过之后改没改。记录在结果之前，就不是事后诸葛。}{①「察觉不到，说明你冤枉了人还不知道。」→冤枉要先判错；我们凭的是他被教过、被提醒过。②「教过他未必听懂（\#191）。」→知道本意的旁听者没有系统性高估（\#234）；判断看他改没改。③「法律也是事后判疏忽。」→刑法第十五条写的是「应当预见」而疏忽大意没有预见（\#180）：看的是事前该不该预见；我们只借这条判别线，不借它的严厉。}{淘汰}
\rebuttal{13}{认定他听不进：「人把意见不合的人看成有偏，就更对抗地回应，形成冲突螺旋；厌蠢就这样从『他错了』走到『他不可理喻』。」}{意见分歧不是犯蠢。同一处基本功错了又错，没有立场之争；这组研究管的是观点不合，管不到教过还犯。}{回扣：准｜公。Kennedy 与 Pronin 的自变量是「觉得对方有偏」，管的是意见分歧（\#208，只核摘要），不问有没有厌：对方允许的对事批评被顶回来，一样进螺旋，这条数据分不开厌和对方开的药。摘要的词是「更对抗地回应」，还在回应；冲上去是愤怒的样子，愤怒长期走向和解（\#86）。我们厌的典型情形是同一处错接二连三，那是做没做对，不是观点对不对。对方自己的淘汰记录写着，这条一句「意见分歧不是犯蠢」就能切掉。}{①「朴素实在论是 Ross 与 Ward 1996 年提出的。」→我方只核到 Pronin 等 2004 年的综述（\#209）；「1996 年提出」没有核到（\#178，【待核实】），请对方报原文。②「劝了不听不就是分歧？」→劝的是「多上心」，同样的错再犯，不是他另有主张；他真有主张，那就讲理。}{淘汰}
\rebuttal{14}{认知之错：「所谓蠢多半是你没看见的信息差：坐高铁的评论、进城农民工的知识盲区、职场新人不懂不成文规定；你当下怎么分辨不想与不能？」}{信息差不是蠢：了解他处境的一般人，不觉得第一次坐高铁的人本该会。我们没记录先问、新人先教；评论区那句，我们认是看错了。}{回扣：准｜门槛。这一条我们大半同意。中立定义里的一般人是了解他处境的一般人；工人日报说得对，不是每个人一出生就有坐高铁的机会（\#120）。评论区那句「这也要教」，算厌蠢，看错了，是个别场合。职场新人那项调查没有测厌蠢，也不是随机样本（\#121）。川观新闻那篇，同一篇写着受访老人多数在用智能手机，「太贵」只来自少数不用的老人，两代人都说难在操作太复杂，年轻人愿意主动教父母（\#175、\#238）：读偏了，生出的是教，不是厌。怎么分辨？有记录看记录，没记录先问，问到了情绪就跟着变（\#188）。}{①「第三次问同一个问题，也可能是没人好好教。」→那就先问他卡在哪；教过、交代过还错在同一处，才是我们恼的。②对方报新人调查的百分比。→「那是新人遇到尴尬的情形，不是厌蠢的比例。」③「评论区总有、一年后还在列。」→量的是评论区，被点赞的表达会被放大（\#101）。④「要求乡村的人立刻有城市常识是苛责（\#213）。」→同意，这正是我们先教、先问的理由。}{常见}
\rebuttal{15}{傲慢之错：「厌蠢预设我在上你在下；央视网说优越感从屏幕里飞出来，红网说自以为是、傲慢自大；轻蔑含着优越感，违背人格平等。」}{评价有高下，不等于人格不平等。看低他这件事上的表现可以，拒绝给他作为人应得的尊重不行，我们的厌停在前一边。}{回扣：公·尊重。平等的尊严不排除按品质、成就区别评价人，只要受承认尊重约束（\#44）；康德整句：比较之下在心里看低别人有时不可避免，但表现出来就是冒犯，拒绝给人作为人应得的尊重无论如何违背义务（\#48）。同事同一处错了又错，我看低他在这件事上的用心，照样讲理、照样合作，这是区别评价。舆论批的是什么：央视网那篇批的是把别人的不了解当成蠢（\#10）；国旗下讲话自己说，没有经历过而不会做并不是蠢（\#117）。标准也不是我私定的：现实里说的蠢，是缺乏所处位置需要的基本知识和技能（\#243）。舆论批的傲慢是表现出来的看低和看错了人，我们不辩护。}{①「P17 后半句就是受承认尊重约束。」→整句认，那正是我们守的线；破没破看回头还讲不讲理。②「新京报整篇是批判的（\#13）。」→我们只说「新京报的文章写道」；贬低人不对，我们同意。③「子张：人将拒我，如之何其拒人。」→同章说「嘉善而矜不能」（\#50、\#182），矜的是不能，我们不拒人。④「评委直觉在我们这边。」→新京报不是央媒（\#118）；舆论批的两样我们都不辩护。}{常见}
\rebuttal{16}{厌就是轻蔑：「对无能的情绪反应心理学上归轻蔑；轻蔑以整个人为对象，几乎总是错误呈现对象；厌蠢就是轻蔑。」}{题目写的是厌，不是鄙视。同一篇摘要说，事情落到自己头上推高的是愤怒；轻蔑只是「起作用」。厌蠢平常起在蠢的代价落到我身上的时候。}{回扣：准｜公。汉典里「厌」是嫌恶、憎恶，「看不起」是「鄙视」的释义（\#1、\#5）；中立定义两端都写了，平常落在哪头要证。综述说轻蔑会把对方看成愚蠢、无能（\#20），推不出反过来，把一件事看成蠢就是轻蔑。Hutcherson 与 Gross 的摘要只说轻蔑在评判无能的行为时起作用；同一句前一项写着，与自己相关推高的是愤怒（\#26、\#224）。厌蠢平常在蠢的后果落到我身上时冒出来（\#252），推高的是愤怒（推论）。愤怒与轻蔑常一起出现，但区分清楚（\#170）；去人性化的是不掺愤怒的纯粹轻蔑（\#38）。整体性的那条反对（\#159、\#161），针对的是以整个人为对象的轻蔑。}{①「混着恼也还是有轻蔑。」→我们不说没有，看它往哪走：恼传过去的是「你没用心」（\#250），和轻蔑不看好改善（\#27）方向相反。②「用法里就是看不起（\#10、\#16）。」→中国网原话带着「有的人」「或许」；用法里有看不起的一端，定义里也写了。③「特质性轻蔑的人对低能力目标反应尤其强（\#28）。」→那是人格特质研究，相关，说的是一部分人。}{常见}
\rebuttal{17}{被伤害可以怨那件事：「就算你真被无知伤害过，也只该对那个行为生气；『我厌蠢』是对一类人的泛化厌恶，从一件事滑到一类人，正是错。」}{「我厌蠢」说的是我反复碰到才认得出来的反应，每一回都冲着一件具体的、本该会却做错的事；我们不厌一类人。}{回扣：准｜公。「我厌蠢」像「我怕狗」，是习惯性陈述，说的是一种稳定的反应，不是对一类人的判决。它稳定的是每一回都冲着具体的人做具体的事：本该会、被教过、还明显做错。中立定义写明，蠢描述的是想法和做法，说成「这个人蠢」是推断，要另证。对方允许的「对造成损害的那个行为生气」，正是我们这份厌的内容。蠢、坏、制度怎么分：坏是存心损人；反应态度回应的是对方意志里的冷漠或缺乏关切（\#32），不上心就在这里，用不着恶意（推论）；制度造成的不会，我们先教。那个视频标题是账号起的（\#114）。泛化到一类人是个别场合的过头。}{①「稳定的反应，说明平常就这样。」→平常冲着什么才是问题：平常冲着教过还犯的事，这份平常站得住。②「见了他就先烦，这就是对人。」→先烦不等于不讲理；这回放下、下回照样说。}{常见}
\rebuttal{18}{错在系统，不在蠢人：「戴锦华说责任不在观众在资本垄断；医疗差错报告说不是烂苹果问题；怪个人比追究机构痛快。」}{系统取向自己也分流：Reason 要的公正文化，是大家对「无可指责」和「应受指责」的界线有共识。体制造成的不会，我们先教；教过还犯，才是他那一份。}{回扣：公·归责｜值。对方引的 Reason，同一篇写着报告文化靠信任，信任要求公正文化，也就是大家共同清楚无可指责与应受指责的界线；他还写有些不安全行为确实恶劣（\#223、\#169）。系统取向要的是分流，不是什么都不怪。医疗、航空这类专业系统只作机制例证（\#104、\#105）；日常同事之间，同一个人同一处错教过三遍，不是一个系统没拦住。冯小刚错在把一个群体公开骂成「垃圾」（\#138），那是群体加羞辱，我们两样都不做。结构让很多人没机会会，我们先教；教过之后还不改的那一部分，才是我们恼的。}{①「同样的情境会让不同的人犯同样的错（\#183）。」→那是不同的人在同一情境犯同一个错，该改情境；我们说的是同一个人被教过之后还在同一处错。②「绝大多数不安全行为并不恶劣。」→他要的是划线，不是不划。③「费曼说巴西学生背下一切却不懂（\#141）。」→体制造出的不会不该怪学生，我们同意；这条只到有把握，对方报出处时请说明。}{常见}
\rebuttal{19}{Torvalds 都认错了：「最有名的厌蠢型天才 2018 年公开道歉：邮件里的攻击既不专业也没必要，尤其是人身攻击；他要学会理解别人的情绪。」}{他道歉的是邮件里的攻击，尤其是人身攻击；社区随后的准则删掉的是辱骂，留下的是严格审查、拒绝不合格的贡献。骂人去掉，标准留下，这叫调分寸。}{回扣：公｜值。他道歉的是邮件里轻率的攻击，尤其是在把它变成人身攻击的时候（\#139、\#186）：攻击是行为，中立定义写明厌不等于骂人。社区随后的规矩：贡献会被仔细审查、常招来批评，拒绝不合格的贡献不算违反准则，禁止的是辱骂、贬损和人身攻击（\#221）。对坏补丁的不满留下了，变成审查和拒绝；人身攻击去掉了。他说要学会理解别人的情绪、恰当回应（\#172），改的是怎么回应。这是专业社区里的名人个例，伴随新准则、有外部压力，证明不了常人平常的样子。}{①「他说自己多年误读别人，是态度问题。」→他给的办法是学会恰当回应；该不该拒绝不合格的贡献，规章写明拒绝不算违规。②「准则写着对其他成员表现出共情（\#189）。」→准则约束的是行为，不是感受；同一套规矩也写着贡献会被仔细审查。③「他说为了更好的补丁这说得通，现在知道不 OK。」→我们不说他现在还认为骂人说得通。}{常见}
\rebuttal{20}{朋霍费尔自己警告蔑视人：「写完〈论愚蠢〉，下一节就是『蔑视人？』：我们无权这样；蔑视一个人就永远不能使他有所成就。」}{不蔑视人，正是我们守的线。「蔑视一个人就永远不能使他有所成就」，而恼还在指望他成，那恰恰不是蔑视。}{回扣：公·尊重。我们不靠朋霍费尔立论，但这一节接得住。他警告的是对「人」整体的蔑视，比厌重得多，写在 1942 年底，语境是纳粹治下的十年（\#211、\#247）。他说我们无权蔑视人，我们同意：拒绝给人作为人应得的尊重无论如何违背义务（\#48）。他说蔑视一个人就永远不能使他有所成就（\#76）；恼里那句「你本可以」恰恰还在指望他成，冲着失败生气对方读出的是「我没用心」（\#250）。同一个人在上一节写的是再也不试图用理由说服蠢人（\#247）；那一步我们不学，我们回头还讲。}{①「他说要少看人做了什么，多看人承受了什么。」→那是对被权力弄蠢的群体说的（\#75）；对新人、没人教过的，我们也先看他承受了什么。②「他说再也不说服，你们引他就是承认走到头不讲理。」→那是政权之下对群体愚蠢的结论，期限是外部解放；我们不拿它说日常。③「流传版说是道德缺陷、狱中写的（\#133）。」→原文是「人的缺陷」，写在被捕之前。}{常见}
\rebuttal{21}{厌蠢伤自己、伤关系：「愤怒敌意与冠心病风险相关；积怨难消的人是自己和至交的大麻烦；轻蔑在婚姻里预测离婚。」}{普通人不会因为「生气伤身」就说一份该有的恼是错的；那几项测的是长期的敌意、夫妻间表达出来的轻蔑，不是同事填错表时的那一下恼。}{回扣：值。判准写明「对自己」的代价双方分量都偏低。冠心病那项元分析测的是特质性的愤怒与敌意，是相关不是因果（\#97）；Gottman 那项是夫妻冲突讨论里的轻蔑，79 对夫妻，轻蔑只是几个负面预测项之一（\#84）。我们说的是同事教过还犯时冲着那件事的恼，这回放下、下回再来，不是积怨，也不是对伴侣表达轻蔑。亚里士多德同一部书说，该生气时生气受称赞，该怒而不怒被看作愚蠢（\#39）；积怨难消（\#40）是他说的过头那一端。生气的社会功能是长远会和解（\#86）；关系的代价落在轻蔑那一头，不在恼这一头。}{①「Gottman 研究所说轻蔑是离婚最强预测因子（\#85）。」→那是研究所博客的说法；「研究发现轻蔑是最佳预测因子」在原始摘要里找不到（\#23 那句只到【待核实】）；那是对伴侣表达出来的轻蔑。②「厌蠢的人活得累。」→累能解释情绪为什么发生，不说明它错；压住感受本身无效（\#155），对方也不要求压。③「身边的人更难开口。」→见第 5 条。}{常见}
\rebuttal{22}{判准之争：占一条就错，有更好的就该换：「认知、尊重、成长占一条就该改；换成对事批评、换位、事实纠正、公正文化，既拿到守标准的好处又没有代价。」}{只赢一问，不自动赢下整场；存在更好的，不等于这个就错了。替代态度只在对方证出平常的缺陷之后才登场。}{回扣：三问｜门槛。共用判准写明三问都进场，只赢一问不自动赢下整场；「占一条就算错」，世上就没有不错的情绪。内在目标判准把「没错」偷换成了「最优」。再看对方开的替代：怨恨针对行为、促使交涉（\#160），我们的恼冲着这件事、回头还讲，本来就在怨恨这一格；事实纠正有效（\#109，美国政治议题），纠正治的是不知道，我们本来先教，治不了教过还不上心；公正文化要分清该怪与不该怪（\#183），我们划的正是这条线。对方的替代，一半在我们这份厌里，另一半管不到我们恼的那一类。}{①「那你们承认三问里有输的？」→每一问我们都守：准，记录在手；公，只怪不上心的那一份、尊重不动；值，关系的代价落在轻蔑那头。②「换位效果最大，为什么不换？」→对没有记录的人先问，就是在换位；对方自己的研究说设身处地替同伴想偏差照旧（\#176）。③「协作过程比个人技能更能预测团队表现（\#95）。」→标〔智〕；另一项结论相反（\#94），学界有争议。}{常见／淘汰}
\rebuttal{23}{情绪也有对错：「情绪有认知内容，可以不恰当；受责备的是以某种方式去感受的人；不自愿的轻蔑仍可以是过错。『控制不了』不是理由。」}{同意，情绪可以评，所以我们才问完三句；可以错，不等于错了。}{回扣：三问。这一条我们不打，我们接：反方不说情绪无对错，也不拿「控制不了」免于评价。亚里士多德说，在恰当的时候、对恰当的对象、以恰当的方式去感受才是德性（\#152）；《中庸》说发而皆中节，谓之和（\#158）。中节的厌：对象是本该会、被教过、还明显做错的事；方式是冲着这件事、回头还讲理；时候是这回放下、下回再来。Adams 举的是无端的愤怒、自以为是的轻蔑，「不自愿也可责」是一派立场（\#156）；我们的厌有端可循，记录在手。问题不是能不能评，是问完三句站不站得住。}{①「那你们认同这份厌？」→认同冲着教过还犯的那份恼；不认同把人划出讲理的那一步，那一步出现时我们自己也改。②「中节是改过之后的样子，你们描述的是最好版本。」→最好版本是圣人；常人会恼、会一时走开、会再来，正常人不能长久躲开（\#34）。}{常见}
\rebuttal{24}{厌蠢话语自己写下的处方：「别和蠢人争论；蠢就是天生无可救药的坏，请远离你身边的蠢人；有的人会用行动离对方远远的；不要再三努力同蠢人论理。」}{这几句，一句找不到出处，一句是自媒体的极端说法，一句带着「有的人或许」；越线的那几句我们和它切割，不替网暴者辩护。}{回扣：公·尊重｜门槛。「别和蠢人争论」找不到出自马克·吐温的实质证据，而且说的是别争论（\#127）。「蠢就是天生无可救药的坏」出自自媒体，只能当「有人认为」的话语样本（\#124）；评价对象是普通人，不是网暴者。中国网那句原话是「有的人或许不会用语言来表达，但会用行动离对方远远的」（\#16）。朋霍费尔流传版那句，语境是纳粹治下的群体愚蠢（\#133、\#74）。越线的我们认越线；可它们是网上最响的几句，不是常人平常的样子：同一篇上游新闻引的网友，遇到不会但努力的人还能耐着性子教（\#168）。}{①「最响的就是最常见的。」→被点赞的表达会被放大（\#101）；好批不等于常见（推论）。②「远离蠢人是自己的权利（\#9）。」→那是记者「由此看出」的概括，不是网友原话（\#168）；我们也不说「我有权远离」。}{最终／常见}
\rebuttal{25}{偷换：你们厌的是坏，不是蠢：「不上心、又菜又犟、理直气壮是态度，是坏；坏人不是蠢人；请承认对真正的笨，厌是错的。」}{我们用的是中立定义：蠢不论出于不会还是不肯。不上心不是恶意，是对本该会的事没用心；汉隆剃刀说的正是别把它当成坏。}{回扣：定义｜公·归责。我们没有改定义：中立定义写明蠢可能出于不会，也可能出于不肯。汉典释「蠢人」是缺乏判断力或谋虑的人（\#67）；人们叫作蠢的情境里有一类是注意失误（\#93），常常就是心思没放在该放的地方（推论）。坏是存心损人；不上心是对该会的事没用心、把代价甩给旁人，反应态度回应的正是意志里的冷漠或缺乏关切（\#32），用不着恶意（推论）。汉隆剃刀说别把蠢当成恶意（\#130），我们同意，所以只恼他不上心，回头照样讲理。我们也不拿「理直气壮」来定义蠢。对努力了还学不会的人，我们给耐心；有人厌他，那一次是看错了。}{①「你怎么知道是不肯？」→看教过之后的做法：被教过、提醒过还不问不改的才是我们恼的（\#41、\#230）；分不清时先问。②「屡教不改是品行问题。」→落在做事上：同一处基本功错了又错，是本该会的事做错，正是蠢；它没有损人利己的存心。③「第三类『缺乏控制』就是控制不了。」→摘要没有定义这一类，双方都不该说它具体指什么。}{常见}
\rebuttal{26}{从一次推到这个人：「心理学叫对应偏差，从情境能解释的行为推出持久特质，是人普遍的倾向；信特质固定的人推得更狠。看低这个人，是这份厌平常的样子。」}{对应偏差咬的是处境没看见、就推到人的那一下。教过、劝过，处境已经摆出来一回；我们恼的是这件事又没做好，不是「他这个人就蠢」。}{回扣：准｜公·尊重。对应偏差是从完全可以由情境解释的行为推出持久特质；摘要列的第一个机制，就是没觉察到情境的约束（\#61，只核摘要）：前提是处境本来解释得了，看的人却没看见。所以我们的第一条线是记录：没有记录先问、新来的先教、只凭一贯做不到不厌（\#241）。教过、被指出、同一处又做错，是看着一串事说话。我们下调的是他在这件事上的评价（\#44）；推成「他这个人就是蠢」，那一步越线了，要调分寸。信特质固定的人推得更狠（\#256，只核摘要），说的是人和人不一样；「你本可以」站在相信他改得了的那边（推论）。}{①「屡次也可以是处境。」→那就先问他卡在哪；分不清时我们不厌。②「『他这回又没上心』也是对人的判断。」→说的是这件事上的用心，下回照样跟他说；分不分得清看回头还讲不讲理。③「你怎么证明自己不是信特质固定的人？」→我们不靠自称；真到「他就这样」，我们认越线。④「蠢字本身落在人身上（\#20）。」→推不出反过来：把一件事看成蠢不等于轻蔑。}{最终}
\rebuttal{27}{正方分量（上）：「看准是归责与尊重的前提，『看准了吗』分量最重；归责只作连带，判断平常就偏向『不肯』。」}{排在前面，不等于分量最重。「准」这一问，对方自己让了一半：凭记录的不争；没记录的那一格，我们也认先问。}{回扣：准｜分量。前提说的是先问哪一问，不说哪一问最重（推论）；判准原文是只赢一问不自动赢下整场。偏差只往一头偏：这把尺子只在没有记录时才是全部依据，有记录时依据是他的记录（第 2 条）。对方说最重的这一问，有记录那一格对方让了，没记录那一格我们认先问；真正的分歧在平常落在哪一格（第 1 条）。「本可以」归归责小问，这是两边共用的口径；不会还是不肯，看教过之后问不问、改不改，不是「准不准」能顺带答掉的。常人判断一份恼火对不对，本来就看可不可控（\#58、\#219）。}{①「你们认没有记录时会看错，就是认系统性缺陷。」→我们认的是那一次看错了；是不是平常，看常人厌蠢平常手里有没有记录。②「判断偏了，归责跟着偏。」→教没教过、改没改是查得到的事，归责靠这串记录。③「那你们说哪一问最重？」→归责小问；尊重我们另守，见第 28 条。}{最终}
\rebuttal{28}{正方分量（下）：「尊重是底线，拒绝基本尊重无论如何违背义务；『划得来吗』上的好处都抵不掉。」}{底线我们同意，我们本来就没拿「值」去抵它。尊重是一条线，问的是这份厌平常越没越线；越线我们认要调分寸，对方要证的是平常越线。}{回扣：公·尊重｜分量。康德说拒绝给人作为人应得的尊重无论如何违背义务（\#48），我们同意；我们不开「值」的进攻，也不说厌蠢守住了社会标准。底线的分量要落到事实上：常人这份厌平常越没越过这条线。康德同一段还说比较之下在心里看低有时不可避免，表现出来才是冒犯；平等尊严不排除区别评价，只要受承认尊重约束（\#44）。线划在拒绝他作为人应得的尊重，不划在对他这件事的评价；越没越线按时间看。对方论点二最后一步「回头不讲理」报道没写，对方自己标了推论：底线越重，越要拿出平常越线的材料。}{①「心里看低有时免不了，就是越线了。」→康德自己分了两件事，我们守后一条线。②「底线不讲平常，越过一次就是越过。」→越过的那一次要调分寸；门槛看平常，这条规则两边共用。③「守标准在尊重面前一文不值。」→我们本来没把它放上秤。}{最终}
\rebuttal{29}{常态与最好版本：「判准评的是常人那份厌，不是反思修好的版本；『我只厌打扰到我的蠢』是为了不显得傲慢才加的前缀；网友那句只是一个人的转述。」}{我们的限定不是前缀，是双方共用的定义：一般人看来本该会、明显做错；第一次、没人教的，本来不算蠢。我们不靠谁的自我声明，靠记录。}{回扣：门槛｜评价对象。「我只厌打扰到我的蠢」是网友给自己加的一句话，我们一个字没靠它，也不拿一位网友证常人。我们的限定写在中立定义里。没记录先问、新来的先教，不是立论时许的愿：对这一类厌了，那一次是看错了，要调分寸。评价对象照判准约定，不是网暴者，也不是圣人；常人会恼、会一时走开、下回会再说，正常人可以一时躲开投入之累但不能长久（\#34）。有间大学那篇是评论，没有数据（\#14）；它批的傲慢是表现出来的贬低，我们不辩护。前缀管不了「平常」，记录管得了。}{①「先问要靠立论许诺。」→先问是记录线的一半：没记录就还不知道他本该会没有（推论）；熟人之间劝本身就是开口（\#248）。②「典型句子全是反问。」→反问是口气，说的是本该；问没问看下一步。③「常人有分寸，证据就是一位网友。」→那位网友只看用法（\#168）；平常压在机制上（第 1 条）。④「有间大学说掩盖不了傲慢。」→同一作者同一篇（\#122），只算一次。}{最终}
\rebuttal{30}{多数人的轻视是随意的：「连为轻视辩护的亚里士多德都说，骄傲的人轻视别人正当，因为他想得对；多数人的轻视则是随意的。常人就是这里的多数人。」}{那句讲的是轻视人，我们不为轻视人辩护；「随意」是说没根据，同一句说有根据才正当。我们的恼有根据：他被教过没有、改没改。}{回扣：准｜门槛。整句念：骄傲的人轻视得正当，因为他想得对；多数人的轻视则是随意的（\#162）。第一，这句的对象是看不起人；我们的线正好划在这里。第二，分界不在多数少数，在想得对不对：有根据的才正当。我们的厌根据是记录；没有记录先问，就是不让它随意（推论）。第三，「多数人」是描述大度之人时的对照，不是统计；原文没说多数人看不准，也没说多数人的恼（推论）。对方要拿这句证平常，得证常人厌蠢平常没有根据，判准不许报比例。}{①「你们不是大度者，凭什么说想得对？」→不靠自称，交给「看准了吗」去称。②「他说的轻视就是厌蠢。」→他说的是看不起人；厌可以只冲着那件事，平常落在哪头要证。}{最终}
\rebuttal{31}{知识的诅咒：「成年人自己知道结果，会损害他推断别人怎么想的能力；你眼里的蠢常常是自己的知识造成的误判。」}{知识的诅咒，是知道答案的人猜不准别人「以为」什么。对一个教过三遍的同事，我们用不着猜他以为什么：他学过没有，记录在手。没记录的，先问。}{回扣：准。Birch 与 Bloom 2007 测的是错误信念推理：自己知道事情的结果，再推断另一个人对这件事以为什么（\#253，只核摘要）；它研究的是信念推理的偏差，不是情绪，拿它解释厌蠢是媒体的延伸。偏差长在对他一无所知、只能拿自己去推的地方。我们恼的那一种不用推：看得见的是他的做法，不是猜出来的信念。知识的诅咒正是我们没有记录先问、新来的先教的理由（推论）。评论区里对陌生人的那句「这都不会」，我们认：算厌蠢，看错了，是个别场合。}{①「教的人也总高估自己说清楚了。」→判断看他之后改没改（\#234）。②「媒体都拿知识的诅咒解释厌蠢。」→那是延伸；它解释得了把信息差当成蠢，那一种我们认看错；解释不了教过还犯。③「提醒也压不下。」→压不下的是估计；有记录时我们不靠估计（第 36 条）。}{最终}
\rebuttal{32}{问答者实验：「随机分配的出题者拿自己知道的冷门知识出题，被一致评为更有学识；说『你怎么连这都不知道』的人，就站在出题人的位置上。」}{出题人拿的是自己的冷门知识，测的是谁显得更有学识；我们说的本该会，是位置上教过、交代过的基本功，不是我出的题。}{回扣：准。出题人是随机分的，题是自己编的，可以拿冷门知识来出（\#254）；偏差出在旁人没把「题是他挑的」这层角色优势算进去。它咬的是两样：题目冷门，而且只有出题人知道；我们恼的那一种两样都不沾。中立定义里的本该会，不是我会什么；现实里说的蠢，是缺乏所处位置需要的基本知识和技能（\#243），表格怎么填、资料怎么分类，带教交代过。这项实验测的是对学识的评价，不是蠢。拿冷门知识说「这都不知道」，那一次是看错了。校正角色偏差的办法就是把角色说破：题是谁出的，他学过没有。}{①「日常说『这都不会』的人就是拿自己会的当题。」→只凭「我都会」的，我们认该先问。②「答题者自己都觉得不如出题人。」→同一篇摘要写明偏差在答题者和旁观者身上更强（\#254 方法栏）；被厌的人怎么看自己，回答不了我们恼得准不准。}{最终}
\rebuttal{33}{康德的尊重义务包括不持有轻蔑、嘲笑：「据 SEP，康德说没有感到尊重的义务，但这些义务包括不持有某些态度；心里的轻蔑本身就违背尊重的义务。」}{两句一起念，我们都认，心里那份也要称。可百科写的轻蔑，是否认他有任何价值；我们看低的是他这件事上的表现，不否认他作为人的价值。}{回扣：公·尊重。心里那份就是厌蠢本身，它错不错看判准。百科两句：没有感到尊重的义务；这些义务同时包括不持有某些态度（\#43）。康德原文里，蔑视就是拒绝给人作为人应得的尊重，无论如何违背义务；同一段又说比较之下在心里看低别人有时不可避免（\#48）：他自己分了两件事，否认一个人的价值，和对他某一方面的评价。同一个百科条目也说平等尊严不排除区别评价，只要受承认尊重约束（\#44）。我们的厌停在后一件：回头照样讲理、照样合作。嘲笑是表达，我们不辩护。}{①「不尊重的主因是想自我感觉良好，『我比你聪明』就是这种欲望。」→后半句是对方的引申，不是百科原话。②「蠢就是对人的评价。」→中立定义写明蠢描述的是想法和做法。③「心里否没否认他的价值，谁看得见？」→看回头：下回照样说、照样听，就没有否认。}{最终}
\rebuttal{34}{新京报写同学同事也是比较：「写到被吐槽的同学同事，写的是『你能不能有点常识』，背后是『我都可以做到你为什么做不到』；被叫作蠢的人有许多是努力了也没学会的。」}{这几句我们全认：只凭「我都可以」、不看他学过没有的，该先问。那段我们只取场景，记录由每日商报那一案补；努力了还学不会的，给耐心。}{回扣：准｜门槛。同一篇读全，三处。第一，「我都可以做到」是作者剖析出来的思维，是评论（\#212）；比较的口气说的是本该，没说出口的是他被教过没有（推论）。第二，作者自己写了下一层：「当然，在现实生活中」，说的是拖后腿的同学同事，现实里说的蠢更多是指缺乏所处位置需要的基本知识和技能（\#243）。第三，「努力了也没学会」「归咎于不够努力不客观」，我们同意：在问在改还学不会的给耐心，只凭一贯做不到不厌（\#241）。新京报没写当事人被教过没有；记录由每日商报那一案补（\#248）。}{①「厌蠢与厌烦没有常识近乎同义，这就是比较。」→常识说的是本该；分歧在本可以，要看记录。②「记录要攒，比较不挑场合。」→记录攒在同一个组里，厌也跟着代价在那里起；「比较不挑场合」对方自己标了推论。③「袁湘琴努力了还是笨，观众照样厌她。」→那是偶像剧角色、观众代入时的反应；有人厌她，那一次是看错了。}{最终}
\rebuttal{35}{不是一次：「出圈现场『我厌蠢症犯了』，评论区『总有』，一年后的报道开头照列；说法不变，这就是平常。」}{持续出现在评论区，量的是评论区。一年后那篇报道，第 2 页写的就是工作两年、劝了多次的同事。评论区那种我们认：算厌蠢，看错了。}{回扣：准｜门槛。先数清是几份：新京报所引的评论和上游新闻所引可能出自同一批，不能当两份独立证据（\#213）；「总有」是作者描述，不是统计。每日商报开头那两句没注出处，只能说是报道开头列出的说法（\#214）。网上的道德愤怒得到点赞以后就更常表达（\#101，推特，迁移是推论）。对方没念的一处：同一篇每日商报第 2 页写的是一位工作超过两年的同事，多次劝告，同样的错接二连三，关系一度降至冰点，现在很微妙（\#248）；同一篇新京报写完评论区，下一段也是「当然，在现实生活中」（\#243）。两篇文字一写到身边，写的都是一起做事的人。}{①「被强化出来的厌蠢也是厌蠢。」→是厌蠢，也看错了；被强化的是网上的表达，推到同事之间怎么判断是推论。②「开头是记者挑的典型说法。」→第 2 页的案例也是记者挑的，双方各取一半，都不是比例。③「两篇写身边也是『你能不能有点常识』。」→见第 34 条。}{最终}
\rebuttal{36}{提醒也压不下、越内行越低估新手：「先认出答案的人总高估不知情的人，越熟越高估，指示避免也压不下；越内行越低估新手的难处。」}{压不下的是估计；有记录时我们不靠估计，他学过没有、改没改，查得到。认脸是知觉，不是做事；越熟，是对那张脸越熟，不是对那个人。}{回扣：准｜门槛。Harley 的任务是辨认由模糊变清晰的名人面孔，先认出的人再估一位设想中的不知情同伴（\#205，只核摘要）：是知觉，不是做事；「越熟越高估」熟的是那张脸。Hinds 测的是内行预测新手完成一项任务要多久；同一句后半还说，中等水平的人预测新手反而可能更准（\#203）。拿自己去估一个不知情的人会偏，提醒也压不下，我们不争；所以不靠提醒，靠记录。我们不评偏多少，问的是平常还是个别：这把尺子只在没有记录、只能靠它时才是全部依据，对方还要证常人厌蠢平常手里没有记录。}{①「说他本该会，本身就是估。」→本该会看位置上写得出、教得出的基本功；本可以看他被教过没有，都不是拿自己去估。②「你天天做这件事，就是内行。」→所以没人教过的新手我们先教，不拿内行的手感判他。③「设想中的同伴就是被厌的人。」→同伴是被试设想出来的；同事不是，问一句就知道。}{最终}
\rebuttal{37}{替不认识的同龄人估也一样：「知道字谜答案的大学生替不认识的同龄同学估难易，偏差和替自己估基本一样，要求忽略答案也没用；还更相信自己估得准。」}{Hom 的被试替不认识的同学估，还被要求设想他和自己经历相同，这本身就推着人投射。对一个共事两年的人，我们用不着设想，也用不着估。}{回扣：准｜门槛。被试替一个不认识的同龄同学估字谜，还被要求设想这位同伴的经历和自己相同（\#176，PMC 全文，美国大学生）；作者在讨论里也写，被试是以自己的知识为锚，再投射到同伴身上（\#217）。对一个陌生人，除了自己没有别的锚，这条指示又把锚定在了自己身上。「更相信自己估得准」那一项，作者自注是事后分析，只在一项指标上一致（\#217）。我们恼的那一种，锚不在我身上：同事被教过没有、教过之后改没改，是他的记录。「不认识」三个字，正是这项研究和我们说的场景分开的地方。}{①「要求忽略答案也没用，压不住。」→压不住的是拿自己当锚的估计；锚换成他的记录，就不是这项实验估的那一下（推论）。②「你们说设身处地偏差照旧，那先问也没用。」→设身处地是闭门替他想；先问是开口问他卡在哪（推论）。③「美国大学生也是普通人。」→我们不拿样本挑刺，只看任务。}{最终}
\rebuttal{38}{替代态度，先开口问：「要改的是凭比较、不问就判：先问一句『你卡在哪』；心理和行为一类的毛病默认被看作可控，归因信息一变情绪就跟着变；厌蠢的典型句子全是反问。」}{没有记录先问，本来就在我们的厌里。对一个共事两年、劝过多次的同事，先开口问，问什么？}{回扣：值｜门槛。替代态度只在对方证出平常的缺陷之后才登场；存在更好的，不等于这个就错了。就算现在谈，对方的替代一半本来就在我们这份厌里：没记录先问、新来的先教、分不清时先问他卡在哪。Weiner 那项研究两句我们都认：默认被看作可控；给了归因信息情绪就跟着变（\#188，只核摘要）；所以没有线索时我们不凭默认判。另一半管不到我们恼的那一类：被教过、被劝过、同一处又错，问过劝过之后还不改的那一份，开口问换不掉什么（推论）。「这都不会？」是口气，熟人之间的下一步是劝。}{①「那项研究第一句：不会就默认读成了不肯。」→那是污名研究，「蠢属于这一类」是推论；没有线索时我们不判，先问。②「说先问本来就有，等于承认做得到。」→它本来就是平常的一半；对方要证的是常人平常不问。③「问什么？问他为什么还不会。」→他说不会，我们再教；他说会了又不改，归责就落到他那一份上。}{最终}
\rebuttal{39}{轻蔑可能从一次次恼火里长出来：「研究 2 专门检验从对同一个人的先前愤怒预测轻蔑，结论是轻蔑可能由愤怒发展而来；恼不是盾，是轻蔑生长的地方。」}{研究写的是「可能」，不是平常的走向；同一篇说，生气长远会和解。我们不说恼里没有轻蔑，看它往哪走：回头还讲，就没走到那一头。}{回扣：公·尊重。这篇双方共用，我们整句认：愤怒和轻蔑常常一起出现，轻蔑可能从先前的生气里长出来（\#86、\#170，只核摘要，自我报告）。三个词分清：「可能」，原文是 may（\#185）；「预测因素」，不是必要条件；「一起出现」，所以我们从不说恼里没有轻蔑。这篇的主干是另一句：生气短期冲上去、长远会和解，轻蔑短期长期都把人推出去（\#86）。恼还在，只说明还不是纯粹的轻蔑（\#38）；往哪走，看回头。真硬化成「他这个人说什么都没用」，那一步越线了；「可能长出来」证的是有，没证平常。}{①「恼退了，剩下的反感正是客观态度装得下的排斥（\#197）。」→Strawson 只说客观态度可以带着排斥性的反感，不说反感就是客观态度；想躲开按时间划。②「研究 2 的条件就是教过了还犯。」→研究 2 测的是对同一个人的多次愤怒，没测教没教过、回头讲没讲理。③「你们只会说预测因素不是必要条件。」→还有一句：「你本可以」和不看好改善方向相反（\#27；推论）。}{最终}
\rebuttal{40}{「说了那么多遍」是管不了，「还是那么蠢」是品性归因：「每日商报那一句自己就把轻蔑的两个预测因素供上了：从『多上心』的要求，走到『那么蠢』的判词。」}{这句话报道没标说话人；照字面读，「说了那么多遍」是还在说，「怎么还是」还在要求他（推论）。还在说，就是还觉得话管得到他。}{回扣：公·尊重。三句一起读：多次劝告，同样的错接二连三，关系一度降至冰点，现在很微妙；报道紧接着写了那一句，没有标说话人（\#248）。逐个对预测因素（\#185）：「说了那么多遍」，前提是一遍遍说了，这回放下、下回照样说，就是还觉得话管得到他；「怎么还是」是失望，失望的前提是还在期待（推论）；蠢字落在人身上，下调的是对他这一方面的评价（\#44）。同一句摘要里还有一项预测因素是与对方缺乏亲密，天天一起做事的同事未必缺这一项（推论）。真到不再说、他开口也不听，我们认越线。}{①「一边说一边宣布说了没用。」→宣布没用是不再说；报道写的现状是很微妙，不是不说了（推论）。②「同事未必亲密。」→那这一项双方都只有推论；剩下两项都还没走到不看好他改（\#27）。③「这句就是小陈说的。」→报道没标，对方也不该替报道补说话人。}{最终}
\rebuttal{41}{认定改不了就不再教：「经理信不信人能改，预测他之后辅导下属多少；看法一改，辅导的意愿和建议就上来。有职责的经理尚且如此，同事只会更少。」}{那项研究测的是经理对人性能不能改的一般看法，不是厌，也不是对一个人的判断；它反倒替我们说话：「你本可以」，站在相信他改得了的那一边。}{回扣：公·尊重。Heslin 的自变量是经理对「人的特质能不能改」的一般看法，因变量是后续辅导和对一名表现差员工的辅导意愿、改进建议（\#218，三项研究，只核摘要）；不测厌，不测对某一个人的判断，也不测听不听他的理由；辅导下属是经理的职务；摘要第一句是作者对已有研究的概括。决定教不教的是「人能不能改」；恼里那句「你本可以」说的正是你能，冲着失败生气对方读出的更多是没用心（\#250）。这项研究要定罪的是「他改不了」，不是恼；那一步我们和对方站在一起。「同事只会更少」对方自己标了推论。}{①「劝了还犯，人会改信他改不了。」→会不会改信研究没测；改信那一步越线了，要调分寸。②「信特质固定的人推得更狠（\#256）。」→推得狠不狠因人而异，正说明推到「改不了」不是这份厌本身带的（推论）。③「不再主动提醒就是放弃他。」→两边共用的口径：不算越线。}{最终}
\rebuttal{42}{康德附释：「在某个判断上否认对方有任何理智，又怎能让他明白自己错了？越线不必等到整个人。」}{康德说的是指出错误时别斥为荒谬，要说明他为什么错：那正是讲理。我们恼他这回没上心，回头照样跟他说错在哪，没有否认他有理智。}{回扣：公·尊重。这一节我们整段认（\#45，据德文原文的意思，中文为自译）。康德要的三件事：别把对方的失误斥为荒谬；先假定其中有真的东西，把它找出来；说明他为什么会错，好让他明白。它管的是指出错误的方式和言辞，不是禁止心里不以为然；同一节正文说比较之下在心里看低别人有时不可避免（\#48）。恼他没上心，下回照样跟他说哪里错、为什么错，就是让他明白自己错了。骂他荒谬、不再解释，是表现出来的看低，我们不辩护。在一个判断上从此不再听他的理由，我们认越线；它是个别场合，正常人不能长久这样（\#34、\#197）。}{①「劝的原话是多上心，没告诉他为什么错。」→劝态度和讲事可以都有；报道只录了劝的那一句，双方都不该替报道补。②「康德要你先假定对方有真的东西，厌蠢先认定他蠢。」→蠢说的是这一处明显做错，不是这个人没有理智；先问他卡在哪，就是在找他那么做的理由（推论）。③「这件事上不再听他算不算越线？」→算，那是个别场合（第 4 条）。}{最终}
\rebuttal{43}{正方的门槛主张：「平常就这样、只是改得掉，照样算错了；方向只往一头偏、自己察觉不到、提醒压不下、每回都用的尺子，就是平常。」}{方向固定说的是偏向哪边，不是用得多不多。有记录时凭记录，这把尺子不上场。对方只证明了有，没证明平常。}{回扣：门槛。门槛规则我们一个字不改，从不拿改得掉来挡。我们说的是这把尺子平常不是判断的依据。「方向固定」说的是偏的时候往哪边偏，回答不了偏的时候多不多；我们也不评偏多少。「自己看不见、提醒压不下」说的是估计本身改不过来；有记录时判断的依据不是估计。「每回都用」是对方没证的一步：这些实验里，被试估的都是一个他一无所知的人（第 2 条）；常人厌蠢平常面对的是天天一起做事、记录在手的人（第 1 条）。按分歧裁决，一方证出平常就有的东西、另一方只证出个别场合的东西，前者胜。}{①「什么毛病都能推给有记录。」→不是：没有记录、只凭比较下的判断，我们认那一次看错了；是平常还是个别，看常人平常手里有没有记录。②「第一次碰上时手里只有尺子。」→第一次碰上我们先问，不凭尺子下判断。③「减少偏差的研究效应量没有更低（\#246）。」→那说明估计压不下，我们不争；我们不靠压它，靠不用它。}{最终}
\end{rebuttals}

\section{九、持方优劣评估}
\begin{assessment}
\assess{定义空间}{三个词方向不一，互相抵消。「厌」词典义偏反（汉典只有嫌恶、憎恶，「看不起」是「鄙视」的释义），热词义和心理学偏正（「看不得别人犯蠢」，评判无能时起作用的是轻蔑）；中立定义两端都写、都用「可以」，打平。「蠢」略偏反：「一般人看来本该」「明显」让第一次、没人教的不会不算蠢，正方的信息差材料从「蠢」移到了「看准了吗」的误判，反方可以认作个别场合；反方的代价是定义里要留着「不会」。「错了 / 没错」略偏正：「我错了」的日常语义是该改，中立定义取了「该改」；分界规则「平常就这样、只是改得掉，照样算错了」「哪怕改完还剩一点反感」堵住了「校准就好」。谁也不能靠定义赢，残余倾向全部转进「这份厌平常落在哪头」的事实之争。}{均衡}
\assess{论证义务}{推定谁也不占，字面上双方镜像。结构上反方义务更宽：正方破一问即可，反方要守每一问，「看准了」不能把「看低这个人」带过关，所以正方论点二可以从「就算看准了」起步，反方论点一赢下准和归责之后，还得靠论点二单独守住尊重；抵账的路也基本断了，「守住本该」已淘汰（主数据测的是故意背叛、能做不做），反方只能在准和公上正面否认「系统性」。正方也有硬义务：要证「平常」却不能报比例；替代必须做得到且不能是「压住不厌」；论点一已让出凭记录的准，论点二最后一步是推论。}{略偏正}
\assess{论据可得性}{正方的实证成族：准上有对应偏差、后见之明与知识的诅咒一组实验、问答者实验；决胜点上有评判无能时起作用的是轻蔑、轻蔑贬品性、轻蔑导向排斥且可能从先前的愤怒中发展、不看好对方改善。反方的实证集中在归因—情绪一族，其余靠规范哲学和案例；试过的另外三块都塌在材料上：看人准不准的 \#88、\#89 标〔智〕，守住标准的 \#110、\#220 测的是坏，动力之用的 \#87 测的是被表达的轻蔑。正方的软处正是反方的突破口：核心几条只核了摘要；迁移到「这都不会？」是推论；Hom 的指示本身推着人投射；“may” 是可能，不是平常；这些实验估的都是陌生人。两边共同的空白是没有任何厌蠢比例数据，「平常」都只有媒体描述。}{偏正}
\assess{评委直觉}{没有评委信息，按一般判断：校赛评委多为高年级辩手与教练，偏好定义清楚、正面回应、不玩定义杀。常人直觉分两半：对屡教不改的同事恼火，多数人觉得说得过去（偏反）；看不起人，多数人觉得不对（偏正）。反方开场要先背「傲慢」的第一印象：「蠢」字本身带骂人色彩，「我没错」听着像自我辩护，主流媒体又几乎一边倒；反方最顺手的三句里，「情绪控制不了」「我有权远离」被判准挡掉，「只厌又菜又犟」不能进定义，剩下能争直觉的是「恼他不弃他」这种要解释两步的话。正方的风险在另一头：一旦被打成「要求圣人」「不许有感觉」或把厌说成霸凌，评委会觉得正方在回避真正难的问题。}{略偏正}
\assess{赛制顺序}{正方的好处：正一先立，定义、判准、门槛由正方先说出口，用中立版时反一只需一句「大体一致」；正二驳论在反二之后，可以回应；正四最后结辩，是全场唯一不会被回应的发言，而本题的决胜点是一个没有比例数据可查的「平常」之争，最后由谁替评委称分量，分量很重。反方的好处：反一在正一立论和反二质询之后发言，可以按正方实际的说法调整；反三盘问小结拿到 C 块最后一句；环节票每个环节都计分。中性但关键的是奇袭：反方在自由辩之后、结辩之前发动奇袭质询，能逼正四把最后一发的一部分时间花在回答上，部分抵消正方的末位优势。}{略偏正}
\end{assessment}
\begin{fields}
\claim[结论]{反方较劣势，程度是「略劣」。劣势来自论据可得性和论证义务，两者是同一件事的两面：决胜点上的实证在正方手里，反方「值」上的实证又塌了，于是反方没有抵账的路，只能在「看准了吗」和「对他公道吗·尊重」两条线上都正面否认「系统性」；评委直觉（傲慢的第一印象、最顺手的说法被挡）和赛制顺序（正四末位）各添一点。反方的胜机在同一处：正方的「平常」同样只有媒体描述，论点二最后一步是推论，核心实验估的都是陌生人。反方的赢法，就是把「平常」钉在熟人之间、记录在手的场景，逼正方承认它的实验够不着那里。}
\begin{concard}{劣势方（反方）打法}\sub{收窄战场}{只打两个战场：①熟人之间、教过了还犯（论点一）：每日商报一案三句一起说，新京报只说「新京报的文章写道」，事情与自己相关推高的是愤怒；②恼与讲理同在（论点二）：Strawson、SEP、康德整句，Graham，孔子与子路。主动让掉、不接的：评论区里对陌生人的比较（认：算厌蠢，看错了，个别场合）；第一次、没人教的不会（先教、先问）；越线的行为（认要调分寸）；专业系统里的失误（只作机制例证）；群体偏见（判准评的是一个普通的「我」）。不开「值」的进攻，「换不掉的那点反感」只在正方摆出替代态度之后登场。全队用同一个问题给正方的每个例子分格：「这个例子里，他被教过没有？」没有，是我们先问、先教的那一种；有，是我们守的那一种。}\sub{抢先定义}{反方不能先立论，「抢先」落在反二质询正一和反一的前几句。反二开取证链，让正方当场确认既定口径里已经让出的东西：「教过、指出过、又做错，这种准你们争不争？」「走神、手滑这类犯蠢，你们争不争准？」「对那件事本身的恼火，你们反不反对？」「这回先不说、一时走开，算不算越线？」「『一般人』是了解他处境的一般人吗？」正方都不会否认，价值在于让评委第一个回合就看到：正方的战场只剩「没有记录的比较」和「回头不讲理」两块。反一一句「定义与判准和对方大体一致」，省下的时间讲门槛和「一般人看来本该」「明显」的意思；正一用了偏向版，一句话点出偏在哪，再重述中立版。反方自己不碰任何定义杀。}\sub{判准之争}{判准不能改，争的是三问里的分量和「平常」的举证。核心一句：「对方只证明了有，没证明平常。」依据：后见之明一组实验估的都是对其一无所知的陌生人；后见之明元分析只说明有偏差，没说明这份厌平常判错，大小不评；轻蔑是 may 从愤怒里长出来，预测因素不是必要条件；正方论点二最后一步自己标了推论。分量：主张归责小问最重，常人判断恼火对不对本来就看该不该他负责；正方自己说「归责只作连带」，反方就把准和归责绑在一起讲。两道共用的盾随时举：「存在更好的，不等于这个就错了」「只赢一问，不自动赢下整场」。正方拿出「先开口问」时回一句：没有记录先问，本来就在我们的厌里。}\sub{价值层}{把「我没错」从「我比你强」换成「我把你当成能负责的人」：反应态度回应的是对方意志里的冷漠或缺乏关切，是把他当参与者；把人当作要管理、回避、不能与之讲理的对象，才是客观态度（\#32、\#33）。所以「你本可以」的恼是一种看得起；对教过还犯的人只给怜悯、只管不问，反倒靠近客观态度（推论）。旁证：做砸之后冲他生气，对方更多读成「我没用心」（\#250，儿童样本，只说方向）。第二层：在意标准，也在意一起做事的人，「我没错」是不必为在意这件事认错；这是价值升华，不报任何效果数据。结辩收在全场退守线。禁区：一个字的精英腔都不能有，不说「我们这种人」「聪明人」。}\sub{顺序利用}{反一按正一实际的说法调整，不照背预写稿。反二驳论在正二之前，只打正方两处软肋（论点一估的都是陌生人，论点二最后一步是推论），不铺新战场。反三盘问小结是 C 块最后一句：把正方在质询、盘问里确实认下的东西列成清单，再说「剩下的只有评论区和一则没写结尾的报道」。奇袭首选自由辩之后、结辩之前发动奇袭质询，链首问：「你们说的平常，除了评论区和一则没写结尾的报道，还有什么？」「回头不讲理，有直接材料吗？还是推论？」「对一个共事两年、劝过多次的同事，『先开口问』问什么？」备用：「恼火你们不反对，去掉看低后剩下的那份，错没错？」反四先结辩，写封路式：①「平常就是凭比较」→记录主线，新京报写完评论区下一段就是「当然，在现实生活中」；②「就算看准了也会滑到看低」→对方自标推论，报道写的现状是「很微妙」；③「替代做得到，所以该改」→存在更好的不等于这个就错了，「先问」「讲理」本来就在我们的厌里；④「反方自己承认越线要调」→只在个别场合冒出来，那叫调分寸。封完再收价值，结辩不引新论据。}\end{concard}
\begin{procard}{优势方（正方）要防的}最容易翻车的是第 1 条。\sub{1 「平常」的举证塌在两头}{论点一已经让出凭记录的准，只剩没有记录的比较；论点二最后一步是推论。反方只要把「平常」钉在熟人之间、记录在手，论点一就被分进「个别场合」。防法：用新京报写同学同事的「我都可以做到你为什么做不到」拆「熟人＝有记录」；把决胜点拉回机制：「蠢」字一落到人身上，贬的就是品性（\#20），从情境可以解释的行为推出持久特质是人普遍的倾向（\#61），反方会接「对应偏差咬的是处境没看见，有记录时处境已经摆出来一回」，要预判；按门槛规则正方不必证明恼火消失，只要证明看低是平常那份厌的一部分；越线照共用口径，不往外推，不说「小陈从此不理小刘」。}\sub{2 打成霸凌或要求圣人}{「厌蠢就是霸凌」「你们第一秒就不该厌」「别有感觉」一出口，反方一句「我只远离，不伤害」「正方要我们做圣人」就能切掉。被问「对屡教不改的怎么办」，只答：「对那件事生气、批评都可以，一时走开也可以；要改的是看低他、把他划到不用再讲理那边。」}\sub{3 证据越界}{虚假论据判该场负，对方可以当场提伪证申诉。不说「研究证明厌蠢偏了 0.39」「厌蠢每一次都判错」；不报任何厌蠢比例；「说了那么多遍，怎么还是那么蠢」不标说话人，和「多次劝告」「一度降至冰点」「很微妙」一起说；康德整句引；朋霍费尔只说「1942 年底写的」，只引「被推向蔑视人的危险很大」；〔智〕材料不搭桥不引。}\sub{4 把舆论当证据}{媒体的批评是评论，不是数据，反方会用同一篇补全（新京报那篇对反方有利的句子要认）。只说「某某文章写道」，不说「新京报证明」，也别说新京报是央媒。}\sub{5 两个论点自相矛盾}{论点一说「看不准」，论点二说「就算看准了」，场上把「就算」挂在嘴上；不说「屡次也救不了」，注意失误、缺乏控制两类不争准。}\sub{6 滑成利弊题}{一旦把「值」上的代价当主论证，反方一句「有弊不等于错」就顶回来。被淘汰的「藏拙」只作「纠正换得来」模块，反方打「值」时才用。}\sub{7 末位松懈}{正四最后发言，但反方很可能在自由辩之后发动奇袭质询；正四要预留时间正面回答「平常还有什么材料」「回头不讲理是不是推论」，临场位留给反四结辩的封路点，结辩里不引新论据。}\end{procard}
\end{fields}

\section{十、口径表}
全队唯一真源：立论、质询、驳论、对辩、盘问、小结、自由辩、结辩、奇袭、速查，一律从这里取词；可以在「允许的换说法」里换着说，不许换意思、换数字。「标准表述」是稿件里说的原话，数字写口语形式；精确值、机构、年份、题名在出处栏与第十一节，被要求出示来源时照出处报。表外的数据、案例、学理不上场，要加先进表。出处栏的 \#N 是第十一节出处清单的编号。

\prosubsection{10.1 正方口径表}
\begin{canon}{正方}
\canongroup{定义}
\canonrow{厌（厌蠢）·口播版}{「厌蠢，是觉得别人犯蠢时生出的嫌恶、反感，常带着想躲开的心；它可以只冲着那件事，也可以冲着这个人、看低他；它是一种态度，不等于骂人或欺负人。」}{①开头合并版用「厌蠢，是觉得别人犯蠢时的嫌恶、反感，常带着想躲开的心；可以只冲着那件事，也可以冲着这个人、看低他；它是态度，不等于骂人或欺负人。」②被问「是不是把鄙视塞进了厌」：「两头都写了，都用『可以』，没替任何一方选。平常落在哪一头、落在那里错不错，是今天要辩的。」③被问「看错了也算厌蠢？」：「算。厌蠢说的是你的厌；看没看准，是第一问要称的，看得准的厌也在里面。」④讲词典义只说「汉典释『厌』为嫌恶、憎恶；『嫌弃』是厌恶而不愿接近」}{①「厌蠢就是看低人」「厌蠢就是鄙视」②「看低他、不愿再认真对待他」③「厌蠢就是霸凌」「厌蠢就是网暴」④把「觉得」换成「看到」⑤「厌蠢症是病，要治」⑥「《现代汉语词典》释『厌』为……」}{汉典「厌」（\#1）；汉典「厌恶」（\#2）；汉典「嫌弃」（\#4）；《重編國語辭典修訂本》2021（\#3）；央视网 2023-02-23（\#10）；中国网·心理中国 2026-02-27（\#16）；Miceli \& Castelfranchi 2018（\#27）；Roberts 2013 书评（\#160）}
\canonrow{只在心里厌，算不算}{「算。心里那份就是厌蠢本身；它错不错，看判准。」}{正方可接一句：「心里看低他，问完三句同样站不住；而且有哲学家说过，轻蔑会从当事人自己察觉不到的表情和线索里漏出来，对方感觉得到。」（Adams 一派的看法）}{①「不算」②「心里对那件事的反感也错」③「研究证明心里的轻蔑一定会被看出来」}{第三节；Adams 1985（\#156）}
\canonrow{蠢·口播版}{「蠢，是在一般人看来本该想得到、做得对的事上，明显想错、做错；可能出于不会，也可能出于不肯，原因不在定义里预设。」}{①开头合并版用「蠢，是在一般人看来本该想到、做对的事上明显想错、做错，不论出于不会还是不肯。」②被问「一般人是谁」：「了解他处境、讲道理的一般人，不是评论区的多数票。」③「第一次、没人教的不会，按定义不算蠢；可觉得他蠢、厌了，那份厌照样是厌蠢，看错是第一问要称的。」④「蠢不等于智商低。」}{①「蠢就是智商低」「蠢是天生能力不足」②「在我看来就算蠢」③「第一次不会坐高铁也是蠢」④把「又菜又犟」「明知故犯」写进蠢⑤「蠢百分之八十由基因决定」⑥「《现代汉语词典》释『蠢』为……」}{汉典「蠢」（\#66）；汉典「蠢人」（\#67）；Aczél 等 2015（\#93）；Stanovich 2009（\#70）；Bonhoeffer〈论愚蠢〉（\#75）}
\canonrow{错了 / 没错}{与判准第 1 行是同一句（「判断这份厌对不对，就看一个清醒的我问完三句……」），全队只背这一句。}{①「『我错了』是该改，不是自认坏人。」②「错不是罪；看错了人，也叫错。」③「我们不说你是坏人，说的是这份厌平常的样子该改。」}{①「我错了就是该受谴责」「厌蠢的人是坏人」②「有任何一点不完美就算错」「占一条就算错」③「《现代汉语词典》释『错』为……」}{第三节；Scarantino \& de Sousa 2018（SEP）（\#164）；Frankfurt 1971（\#157）}
\canonrow{我}{「一个普通的、会说出这句话的我。」}{①「不是网暴者，也不是圣人。」②「我们评的是你这份厌，不是你这个人，这正是我们要你对『蠢人』做到的。」③退守：「『看得准』的人，亚里士多德自己都说是少数：多数人的轻视是随意的。」}{①「厌蠢的人都傲慢」「你们这些精英」②「新京报证明精英最容易厌蠢」（\#13 是作者论断，只能说「新京报书评周刊认为」）③「我们说的就是网上的喷子」}{第三节；亚里士多德《尼各马可伦理学》IV.3（\#162）；Roberts 2013 书评（\#160）；新京报书评周刊 2023-11-01（\#13）}
\canongroup{判准}
\canonrow{中立判准·口播版}{「判断这份厌对不对，就看一个清醒的我问完三句：看准了吗？对他公道吗？划得来吗？问完觉得这份厌平常的样子就该改，就是错了；站得住，只在个别场合调调分寸，就是没错。」}{简称「三问」：准、公、值；质询里可说「问完三句」。}{①丢掉「平常的样子」，说成「只要有时厌错就算错」②「只看看准了没有」「只看伤没伤人」③「三问都要赢才算赢」「赢一问就赢了」}{第三节；第四节；Frankfurt 1971（\#157）}
\canonrow{一辩开头合并版（约 32 秒）}{「厌蠢，是觉得别人犯蠢时的嫌恶、反感，常带着想躲开的心；可以只冲着那件事，也可以冲着这个人、看低他；它是态度，不等于骂人或欺负人。蠢，是在一般人看来本该想到、做对的事上明显想错、做错，不论出于不会还是不肯。错没错，看清醒的我问完三句：看准了吗？对他公道吗？划得来吗？这份厌平常的样子该改，就是错了；只在个别场合调调分寸，就是没错。」}{只在正一开头用；后场引定义、判准，用定义组各行与判准组第一行。}{删掉「可以只冲着那件事」半句；改动任何一个词}{第四节}
\canonrow{三问·看准了吗（准）}{「看准了吗：对方真的蠢吗？我了解他的处境吗？从这件事推到这个人，推得准吗？我凭的证据够不够？」}{「准，就是先看清他的处境再下判断。」}{①「厌蠢永远看不准」②把「看准了吗」说成「智商测得准不准」}{第四节}
\canonrow{三问·对他公道吗（公）}{「对他公道吗：这份蠢该不该由他负责，是不会还是不肯？厌了以后，我还给不给他作为人应得的尊重？」}{被问「他不肯学，厌他就公道了吧」：「公道问两件事：该不该由他负责，厌了以后还给不给他作为人应得的尊重。不肯学，只回答了前一件。」}{①「蠢都不该由他负责」「不肯也不能怪」②「看准了就能看低」③「心里不许有一点看低」}{第四节；Kant 1797，§39（\#48）}
\canonrow{三问·划得来吗（值）}{「划得来吗：这份厌让事情、关系和我自己，离我在乎的东西更近还是更远？改掉它要付什么代价，换上的态度能不能拿到同样的好处？」}{「替代态度只在先证出缺陷之后登场。」}{①「厌蠢弊大于利，所以错」②「生气伤身，所以错」③把「值」当主论证}{第四节}
\canonrow{「清醒」是什么}{「清醒只是说双方摆出来的事实都知道：蠢的成因和代价，厌的后果和功能。你们的事实也在里面。」}{「清醒的我是评价者，不是被评价的那份厌。」}{「清醒就是认同我方」「清醒的人不会厌」}{第四节}
\canonrow{门槛分界规则}{「要改的如果是这份厌平常的样子，就是这份厌本身该改，哪怕改完还剩一点反感；要改的如果只在个别场合冒出来，那叫调分寸。平常就这样、只是改得掉，照样算错了。」}{①短版：「平常就这样、只是改得掉，照样算错了。」②「任何毛病都改得掉，『可以校准』不能当分界，分界只看平常还是个别。」③被问「岂不是一点不完美就算错」：「不是。只在个别场合冒出来的，调分寸就好；要改的是这份厌平常的样子，才叫错。」}{①「厌蠢每一次都判错」②「有一点毛病就算错」③「我们要连恼火一起去掉」}{第三节；第四节}
\canonrow{评哪一份厌·全场退守线}{「我们不要你第一秒不厌；我们要你问完三句之后，别再认同那份看低人的厌。」}{①「第一秒的厌不由自主，双方都认；评的是第二秒以后，你认同它、放任它，还是修正它。」②「我们评的是常人实际怀着的那份厌，不是反思修好的版本，也不是网暴版本。」}{①「你们第一秒就不该厌」②「别有感觉」「压住不厌」③「厌可以随手关掉」}{第四节；亚里士多德《尼各马可伦理学》II.5（\#151）}
\canonrow{修正不等于压抑}{「改，是换一种看法，不是压住感受。」}{「恼照样恼，只是把『这都不会？』换成『你卡在哪？』」}{「你们要忍着别生气」「要压住」}{第四节；Webb 等 2012（\#155）；Gross 1998（\#154）}
\canonrow{只赢一问不自动赢；我方分量主张}{「三问都进场，只赢一问不自动赢下整场。我方说准和公分量最重：看不准，后面讲公道也是在错的判断上讲；尊重是底线，『值』上的好处抵不掉。」}{「这条对你们对我们都一样。」}{「我们赢了准就赢了」「值无所谓」}{第四节}
\canonrow{存在更好的不等于错了（我方认）}{「存在更好的，不等于这个就错了，这一点我们认。所以我们先证这份厌平常的样子有缺陷，再说换成什么。」}{「替代只回答改得了吗、改了亏不亏。」}{「有更好的态度，所以这份厌错了」}{第四节}
\canonrow{推定谁也不占}{「『我没错』本身是一句正面的自我辩护，推定谁也不占，双方都要正面举证。」}{「我方证平常该改，对方证平常站得住。」}{「对方没证明没错，所以就是错」}{第四节}
\canonrow{对屡教不改的怎么办（必须回答）}{「对那件事生气、批评都可以，一时走开也可以；要改的是看低他，把他划到不用再讲理那边。」}{「怪那件事，我们不拦。」}{①「连一点不满都不该有」②「你就该一直耐心教下去」}{第四节；第九节}
\canonrow{越线按时间划（双方共用口径）}{「越线按时间划：这回先不说、一时走开，不越线；回头他开口先认定不值得听，合作只下指令不给理由，认定『他这个人说什么都没用』，才越线。」}{①「接触回来不等于讲理回来。」②「不再主动提醒，不算不讲理。」③「正常人可以一时躲开投入之累，只是不能长久。」}{①「走开就是越线」「躲开的才叫厌」②「一辈子不来往才算越线」}{Strawson 1962（\#34）}
\canonrow{我方让出的（全队同一答案）}{「教过、被指出、又做错，这种凭记录的看准，我们不争；走神、手滑这类，我们也不争准；对那件事的恼火和批评，我们不反对；这回先不说、一时走开，不算越线。」}{反二取证五问逐条答「不争」「不争」「不反对」「不算」「是，了解他处境的一般人」，再补一句：「我们争的是两处：没有记录、凭比较就判的那一下，和看准之后回头还讲不讲理。」}{①「屡次也救不了」「有记录也看不准」②「厌蠢只在我知道、他不知道的时候出现」③当场改口不认}{第九节}
\canonrow{决胜点}{「全场落在一处：常人说『我厌蠢』时，那份厌平常停在事上，还是会从『这都不会？』走到『他就是那么蠢』，回头不再跟他讲理。」}{「比的是谁描述的平常更像常人。」}{「决胜点是厌蠢有没有坏处」}{第四节}
\canonrow{拆·恰当性判准（对得上就没错）}{「对得上只回答了一问。它没回答厌了以后还给不给他作为人的尊重，也没回答划得来吗；契合不等于道德上恰当。」}{①「性别歧视的笑话也许确实好笑，被它逗乐却不恰当。」②「为轻蔑辩护的 Bell 自己说，轻蔑该回应的是傲慢、伪善这类优越之恶。」}{「对得上也是错」}{Scarantino \& de Sousa 2018（SEP）（\#164）；Bell 2013（\#35）；Roberts 2013 书评（\#36）}
\canonrow{拆·权利判准（我有权远离）}{「有权这样做，不等于这样做没错；题目问的是这份厌对不对，不是别人能不能管你。」}{「那篇报道里，『远离蠢人是自己的权利』是记者的概括；那位网友原话里还有『她甚至可以耐着性子教对方』。」}{「你没有权利远离谁」}{第四节；上游新闻 2023-02-24（\#168）}
\canonrow{拆·行为判准、情绪无对错}{「题目问的是厌，不是骂；『无所谓对错』也不等于『没错』。」}{「亚里士多德说我们感到愤怒不经选择，紧接着就说，受责备的是以某种方式去感受的人。」}{「情绪都该被谴责」}{第四节；亚里士多德《尼各马可伦理学》II.5（\#151）；Adams 1985（\#156）}
\canongroup{论点}
\canonrow{论点一「拿我当一般人」·主张}{「常人的厌蠢，平常凭比较下判断，不凭记录。『这都不会？』那一下，尺子握在已经会了的人手里，只往『他本该会』那头偏，自己还看不见。这份厌平常的判断方式，该改。」}{①「拿自己当尺子。」②「凭比较、不问就判。」③「尺子在会了的人手里。」}{①「厌蠢每一次都判错」②「研究证明厌蠢偏了零点三九」③「厌蠢的人是故意冤枉人」④「后见之明研究测的就是厌蠢」⑤「熟人之间也从来不凭记录」}{第五节论点一；回扣「看准了吗」}
\canonrow{论点一·机制 1：本该凭什么}{「说他『本该会』，依据只有两种：记录，或者比较。凭记录的，教过、被指出、又做错，我们不争；厌蠢的典型说法是比较：『这都不会？』『我都可以做到你为什么做不到』。」}{「真凭记录，说出口的该是『上次教过你』；说出口的是『我都会』，用到的就是我。」（推论）}{①「新京报证明厌蠢都凭比较」（只说「新京报的文章写道」）②报任何「多少人凭比较」的比例}{央视网 2023-02-23（\#10）；新京报 2023-11-01（\#212、\#213）；上游新闻 2023-02-24（\#167）}
\canonrow{论点一·机制 2：尺子往一头偏}{「知道答案的人，会高估不知道的人本来能知道多少；替一个不认识的同龄人估，也一样偏；自己察觉不到，提醒也压不下。」}{①「越内行，越低估新手的难处。」②「熟悉的东西，会让人觉得『大家都知道』。」}{①「提醒压不下是那项防洪研究证的」（\#206 只说去偏程序没有显著效果，它的比例也不报；「提醒压不下」只挂 \#205、\#246）②「设身处地想一想就能估准」③「察觉不到，所以没责任」}{Fischhoff 1975（\#204）；Hom 2023（\#176）；Birch 等 2017（\#215）；Hinds 1999（\#203）；Harley 等 2004（\#205）；Guilbault 等 2004（\#246）}
\canonrow{论点一·机制 3：要改的是什么}{「要改的不是『别有感觉』，是『凭比较、不问就判』：先开口问他的处境，再下判断。」}{①「把『这都不会？』换成『你卡在哪？』」②「没有记录，先问。」}{①「闭门换位就能看准」②「你们不该觉得他蠢」}{Weiner 等 1988（\#188）}
\canonrow{论点一·回扣与分量}{「这一点回扣『看准了吗』。它分量最重：看准是归责和尊重的前提；偏差只往一头偏，错落在被厌的人身上；这不是哪一次眼花，是平常那把尺子。」}{归责只作连带：「判断这一步平常就偏向『不肯』，后面分流得再整齐，也是在错的判断上讲公道。」}{①「准错了，公和值就全错了」②把归责在「公」里再算一遍}{第四节}
\canonrow{论点一·预答「记录就是处境，你们只剩网上陌生人」}{「凭记录的我们不争；可比较也在同学同事之间：新京报的文章写到被吐槽的同学、同事，写的是『你能不能有点常识』的烦躁，是『我都可以做到你为什么做不到』的抱怨。」}{①「身边不等于有记录，平常凭的还是比较。」②「记录要一次次攒，『这都不会？』那一下只要比较就够。」（推论）}{①「新京报认为厌蠢错了」②「同学同事之间全凭比较」}{正驳 1、5、6；新京报 2023-11-01（\#212）}
\canonrow{论点一·预答「平均才偏一截，不等于平常判错」}{「我们不说每次判错，也不说厌蠢偏了多少；说的是方向只往一头偏、自己看不见、每回都用。这样的尺子，就是平常。」}{「同一篇元分析还写，设法减少偏差的研究，效应量没有更低：看不见，也压不下。」}{①「偏得很大」②「厌蠢偏了零点三九」}{Guilbault 等 2004（\#246）}
\canonrow{论点二「说了也没用」·主张}{「就算看准了，熟人劝了还犯的时候，常人的厌平常会从『说了那么多遍』走到『他就是那么蠢』，回头不再把他当值得说理的人。这就越过了作为人应得的尊重。」}{①「从要求走到判词。」②「从对事走到对人。」③「就算看准了，也会越线。」}{①「恼火也是错的」「生气就是轻蔑」②「轻蔑一定会从愤怒里长出来」（原文是 may，只说「可能」）③「厌蠢的人都会从此不理人」④「厌蠢就是霸凌」⑤论点二里又说「其实也看不准」}{第五节论点二；回扣「对他公道吗」的尊重小问}
\canonrow{论点二·机制 1：「蠢」字落在人身上}{「『蠢』字一落下来，就落在人身上：心理学说，轻蔑把对方贬为愚蠢、无能，贬的是品性多于行为。恼火我们不剔，要看它往哪走。」}{「评判无能的时候，起作用的正是轻蔑。」}{「研究证明厌蠢就是轻蔑」}{Fischer \& Giner-Sorolla 2016（\#20）；Hutcherson \& Gross 2011（\#26）}
\canonrow{论点二·机制 2：劝了还犯，供上轻蔑}{「劝了还犯，正好把轻蔑的预测因素供上：『说了那么多遍』，是觉得管不了他；『还是那么蠢』，是把错归到他这个人。轻蔑可能就从对同一个人的一次次恼火里长出来，而且不看好他会变好。」}{「恼不是挡住轻蔑的盾，是它生长的地方。」}{①「预测因素凑齐就必然轻蔑」②「愤怒和轻蔑分不开」}{Fischer \& Roseman 2007（\#86、\#185、\#170）；Miceli \& Castelfranchi 2018（\#27）}
\canonrow{论点二·机制 3：越线看回头}{「越线看回头：这回先不说、一时走开，不越线；回头他开口，先认定不值得听，合作只下指令、不给理由，就越线了。康德说得直白：在一个判断上否认他有理智，又怎能让他明白自己错了？」}{①「接触回来不等于讲理回来。」②「认定人改不了，人就不愿再教：经理研究里，看法一改，辅导的意愿就上来了。」}{①「报道写了小陈不再跟小刘讲理」②「研究证明厌蠢的人回头不讲理」③「一走开就是不讲理」}{Kant 1797，§39 附释（\#45）；Kant 1797，§39（\#48）；Heslin 等 2006（\#218）}
\canonrow{论点二·回扣与分量}{「这一点回扣『对他公道吗』里的尊重。它分量最重：它就压在全场决胜点上；尊重是底线，『划得来吗』上的好处抵不掉。」}{「看准了，也不等于可以看低这个人。」}{「尊重受损，所以厌蠢的人该受罚」}{第五节论点二回扣}
\canonrow{论点二·预答「多次劝告就是讲理一直在，『一度』只是一时」}{「劝告在前，冰点和『还是那么蠢』在后，顺序就是从要求走到判词；报道写的现状是『很微妙』，不是照常教、照常听。」}{「我们不说从此不理，只说现状很微妙。」}{①「小陈从此不理小刘」②「冰点到现在还没化」}{每日商报 2024-01-13（\#248）}
\canonrow{论点二·预答「麻烦落到自己头上，推高的是愤怒」}{「对，事情跟自己相关，推高的是愤怒，这条我们不争；争的是这份气往哪走。同一篇摘要写着，评判无能的行为时，起作用的是轻蔑。」}{「麻烦落到我身上，和『蠢』字落到人身上，是同时发生的。」（推论）}{「与自己相关推高的是轻蔑」}{Hutcherson \& Gross 2011（\#224、\#26）}
\canonrow{两个论点的分工}{「第一点管这份厌凭什么起，第二点管它往哪走；第二点从『就算看准了』说起。」}{场上把「就算」挂在嘴上。}{①「两个论点都在说看不准」②「屡次也救不了」}{第五节}
\canonrow{替代态度（证出缺陷后才登场）}{「换成什么？恼照样恼，对那件事照样批评；先开口问一句『你卡在哪』；一时走开也行，回头照样跟他讲理。」}{①「对事批评、先问再判、一时抽身回头还讲理。」②「孔子骂了子路，还把道理讲完，这正是我们要常人学的样子。」}{①「压住不厌」「别有感觉」②「要像老师一样一直耐心教」③「设身处地想想就行」}{第四节；Roberts 2013 书评（\#160）；Strawson 1962（\#34）；Webb 等 2012（\#155）}
\canonrow{模块「纠正换得来」（只在反方打「值」时用）}{「守标准，对那件事的恼和纠正就守得住，用不着看低人：事实纠正通常有效，没有引发逆火；可厌一在场，人怕显得无能，就不敢认错、不敢求助。」}{「Linux 内核照样拒绝不合格的补丁，守质量用不着看低人。」}{①主动开「值」的进攻②把「藏拙」当独立论点}{Wood \& Porter 2019（\#109）；Edmondson 1999（\#64、\#184）；Linux 内核行为准则 2018（\#221、\#189）}
\canonrow{学理·后见之明偏差（论点一）}{「心理学叫后见之明偏差：知道结果的人，会高估不知道的人本来能知道多少，而且多半察觉不到。」}{通俗说「知识的诅咒」，学理出处仍说 Fischhoff 1975。}{①「后见之明偏差证明了厌蠢是错的」②描述 Fischhoff 的实验材料}{Fischhoff 1975（\#204）}
\canonrow{学理·愤怒与轻蔑的功能区分（论点二）}{「2007 年 Fischer 和 Roseman 的研究把愤怒和轻蔑分开：生气过后会和解，轻蔑才把人推出去。」}{「愤怒冲着事、长远和好；轻蔑冲着人、长短期都排斥。」（前半「冲着事」是 Roberts 的区分，不说成这项研究）}{同第四节该行}{Fischer \& Roseman 2007（\#86、\#185、\#170）}
\canonrow{学理·对应偏差}{「从完全可以由情境解释的行为，推出一个人持久的特质，心理学叫对应偏差，是人普遍的倾向。」}{「从这件事推到这个人。」}{「对应偏差在哪种文化里都一样强」}{Gilbert \& Malone 1995（\#61）}
\canonrow{学理·Strawson：一时躲开，不能长久}{「Strawson 说，正常人可以一时躲开投入之累，但作为人，不能长久这样；态度完全客观的时候，你可以跟他谈、甚至谈判，却不能跟他讲理。」}{「可以一时，不能长久。」}{①「Strawson 说厌蠢就是客观态度」②去掉「完全」二字③报原刊卷页}{Strawson 1962（\#34、\#174、\#197、\#33）}
\canonrow{学理·康德两句（必须整句引）}{「康德说，比较之下在心里看低别人，有时免不了；可表现出来就是冒犯。他还说，拒绝给人作为人应得的尊重，无论如何都违背义务。」}{附释一句：「在一个判断上否认他有理智，又怎能让他明白自己错了？」}{①只引前半句或只引后半句②「康德说心里看低就违背义务」③把中文说成正式译本原文}{Kant 1797，§39（\#48）；Kant 1797，§39 附释（\#45）}
\canonrow{学理·区别评价受承认尊重约束}{「平等的尊严不排除按品质、成就区别评价人；但这种评价，必须始终受『给人承认尊重』这条要求约束。」}{「能评价，但评价要停得住。」}{只引前半句}{SEP「Respect」（\#44）}
\canonrow{学理·尊重义务包括不持有某些态度}{「据《斯坦福哲学百科》的概括，康德说我们没有『感到』尊重的义务，可这些义务包括不持有某些态度，比如不轻蔑、不嘲笑。」}{}{只引「没有感到尊重的义务」，或只引后半}{SEP「Respect」（\#43）}
\canonrow{学理·轻蔑与怨恨（对人与对事）}{「轻蔑针对人、促使退出；怨恨针对行为、促使交涉。对事的批评我们支持，要改的是对人的那一截。」}{被反咬「你们说我错，不也是在厌我」：「我们评的是你这份厌，不是你这个人：针对行为、促使交涉的是批评，针对整个人、促使退出的是轻蔑。」}{「Bell 认为轻蔑一律不对」}{Roberts 2013 书评（\#160）}
\canonrow{学理·不自愿的罪}{「Adams 承认愤怒不能像坐下、站起那样选，可他说存在『不自愿的罪』，点名了对他人的轻蔑。」}{}{说成学界共识}{Adams 1985（\#156）}
\canonrow{学理·罪与病（双方共用）}{「Weiner 的归因理论说，缺努力的失败比缺能力的失败招来更多惩罚：情绪跟着『能不能控制』走。它描述人怎么反应，不管判断准不准。」}{「这条我们全认，问题在判断那一步。」}{「Weiner 证明厌不肯的人没错」}{Weiner 1993（\#58）}
\canonrow{对方学理·应当知道又不难知道}{「亚里士多德讲的是立法者对成文法规的处理，借得来判别线，借不来『实验室里哪些算基本』这张清单。」}{「位置给的是『该会』，不是『已经会』。」}{「亚里士多德说无知都不该怪」}{亚里士多德《尼各马可伦理学》III.5（\#230、\#41）}
\canonrow{对方学理·轻蔑守护规范}{「那本文集只说：轻蔑该当作守护规范的情绪来培养，还是该尽量远离，是开放问题。」}{}{「学界认为轻蔑该远离」}{Mason 编 2018（\#165）}
\canongroup{数据}
\canonrow{后见之明·Fischhoff 1975（论点一主数据）}{「1975 年的经典实验发现：知道结果的人，会高估不知道的人本来能知道多少，自己还几乎察觉不到。」}{「知道答案的人估不准不知道答案的人，而且自己看不见。」}{①报样本量、效应量②描述实验材料③「Fischhoff 证明了厌蠢会偏」④「被试是故意高估」}{Fischhoff 1975（\#204）}
\canonrow{替陌生同伴估·Hom 2023（论点一辅数据）}{「2023 年的实验里，知道字谜答案的大学生，替一个不认识的同龄同学估难易，偏差和替自己估基本一样；要求他们忽略答案，也没用。」}{①「替陌生的同龄人估，照样按自己知道的估。」②被追问细节时可加：「他们还更相信自己替同学估得准」（作者自注是事后分析，只在一项指标上一致）}{①「偏差完全一样」②被追问时隐去「研究要被试设想同伴经历和自己相同」这条指示（必须认）③报效应量④「美国大学生的结果就是中国常人」}{Hom 2023（\#176、\#217）}
\canonrow{知识有多普遍·Birch 2017}{「2017 年的三项实验、三百五十多名本科生发现：人判断『自己知道的东西，同龄人里有多普遍』的时候，光是熟悉带来的顺手感，就足以把判断带偏。」}{「会被自己的熟悉感带偏，往『大家都知道』那头偏。」（方向据 Hom 2023 引言转述，被追问时说明）}{①把「高估」说成摘要原话②「熟悉感是唯一原因」}{Birch 等 2017（\#215）；Hom 2023（\#217）；3 项实验，359 名本科生}
\canonrow{拿自己估别人·Nickerson 1999（整句引）}{「1999 年的一篇综述说：人估别人知道什么，起点就是自己知道什么；这总体上管用，可人常常不加批判地这样做，结果误以为别人知道得和自己一样。」}{「起点是自己，常常不加批判。」}{①只引后半句②把「常常」说成多数或比例}{Nickerson 1999（\#216）}
\canonrow{专家之咒·Hinds 1999}{「1999 年的研究发现：越是内行，越会低估新手完成任务的难处，专门的去偏办法也压不下。」}{「越内行，越看不见新手卡在哪。」}{①说成对「蠢」的判断②被追问时隐去「中等水平的人预测可能更准」}{Hinds 1999（\#203）；2 项研究}
\canonrow{提醒压不下·Harley 2004}{「2004 年的三项实验：先认出模糊面孔的人，总高估不知情的人能认出多少；越熟悉越高估，提醒他们避免，偏差也没减小。」}{「提醒也压不下。」（只挂这一条和 \#246）}{「认脸实验证明常识判断会偏」}{Harley 等 2004（\#205）}
\canonrow{后见之明元分析·Guilbault 2004（对方提时才说）}{「对方说的那项元分析，汇总了九十多项研究，平均偏差是零点三九个标准差，这个数我们认。我们不说偏得大，说的是：方向只往一头偏，自己看不见；同一篇还写，设法减少偏差的研究，效应量没有更低。」}{「用常识问答题的研究，效应比用真实事件的更大。」（迁移到「这都不会？」是推论）}{①主动拿它证「偏得大」②「研究证明厌蠢偏了零点三九」③换算成「百分之几的人判错」或任何比例④评「大」「小」}{Guilbault 等 2004（\#246）}
\canonrow{愤怒与轻蔑·Fischer \& Roseman 2007（论点二主数据）}{「2007 年的三项研究发现：生气过后会和解，轻蔑才把人推出去，还可能从一次次生气里长出来；觉得管不了他、把错归到他这个人，都预示轻蔑。」}{①「愤怒短期冲上去、长远会和好；轻蔑长短期都排斥。」②被问时主动认：「两种情绪常常一起出现，但区分清楚。」}{①「轻蔑常常、一定从愤怒里长出来」②「研究证明厌蠢就是轻蔑」③报样本量、效应量④隐去「常常一起出现」}{Fischer \& Roseman 2007（\#86、\#185、\#170、\#25）}
\canonrow{轻蔑贬品性·Fischer \& Giner-Sorolla 2016}{「心理学的综述说：轻蔑把对方贬为愚蠢、无能，贬的是这个人的品性，多过他做的那件事。」}{「轻蔑的核心是看低对方，想忽视、排斥他。」}{①「研究证明厌蠢的人都看低人」②说成新实验}{Fischer \& Giner-Sorolla 2016（\#20、\#19）}
\canonrow{轻蔑管无能·Hutcherson \& Gross 2011}{「2011 年的五项研究发现：评判无能的行为时，起作用的是轻蔑；事情跟自己相关，推高的是愤怒。」}{「评判无能，靠的是轻蔑。」}{①「研究发现三种情绪分不开」②只引「与自己相关推高愤怒」不引后半}{Hutcherson \& Gross 2011（\#26、\#224）}
\canonrow{轻蔑不看好改善·Miceli \& Castelfranchi 2018}{「轻蔑只指向人，带着优越感，不看好对方会变好，心里跟他疏远。」}{「不看好他会变好。」}{说成实证发现}{Miceli \& Castelfranchi 2018（\#27）}
\canonrow{默认读成可控·Weiner 1988}{「1988 年的实验发现：被归为『心理和行为』一类的毛病，人们把它的起因看作当事人自己能控制的，于是更多愤怒、更少同情、倾向不管；归因信息一换，情绪就跟着变。」}{「没人告诉你他卡在哪时，人默认把他读成『不肯』。」（「蠢属于这一类」是推论，被追问时认）}{①举具体病名②「研究对象就是蠢」③只引第二句不引第一句}{Weiner 等 1988（\#188）；2 项实验，10 种污名}
\canonrow{认定改不了就少教·Heslin 2006}{「2006 年的三项研究发现：经理怎么看『人能不能改』，预示他后来辅导下属的多少；让认定人改不了的经理转而相信人能改，他辅导一名表现差的下属的意愿、给的改进建议的数量和质量，都上来了。」}{「『改不了』的看法让人少教。」}{①「经理研究证明同事有义务辅导」②「这项研究测的是厌」③把「持实体论的人不愿投入帮人」说成本研究的发现}{Heslin 等 2006（\#218）；3 项研究}
\canonrow{换位与压抑·Webb 2012}{「一项汇总三百多个实验比较的元分析发现：换位思考是调节情绪最有效的一种办法；光压住感受，几乎没有效果。」}{「换一种看法有用，硬压没用。」}{①「换位思考能让判断变准」②「专门测了厌蠢」③念小数}{Webb 等 2012（\#155）；306 个实验比较；换位思考 d=0.45；压住情绪体验 d=−0.04}
\canonrow{重评与压抑·Gross 1998}{「1998 年的实验里，一百二十名被试看令人厌恶的影片：换个想法去看，厌恶感真的降了；硬压着不表现，交感神经反而更兴奋。」}{「换看法降的是感受本身，不是嘴上一套。」}{「研究证明厌蠢能被重评消除」}{Gross 1998（\#154）；120 名被试}
\canonrow{一贯做不到读成不会·Frieze \& Weiner 1971}{「1971 年的实验发现：一个人这次的结果跟他自己过去一贯一致，人们归于能力和难度；跟过去不一致，才归于运气和努力。」}{「屡次做不到，读出来的是『不会』，不是『不肯』。」（迁移到厌蠢场景是推论）}{①「研究证明屡教不改的都是不会」②说成智力}{Frieze \& Weiner 1971（\#241）；实验室研究}
\canonrow{提高技能才认出局限·Kruger \& Dunning 1999}{「1999 年那项研究说：提高被试的技能，他们才认出了自己能力的局限。条件是提高技能，不是听几遍劝。」}{「被告知『你错了』，不等于被教会『怎样才对』。」}{①报「第六十二、第十二百分位」②「邓宁克鲁格效应证明蠢人不知道自己蠢」}{Kruger \& Dunning 1999（\#251）}
\canonrow{只写步骤没动手·Rozenblit \& Keil 2002 研究 8}{「那项研究的被试只是写出步骤，没有动手去做，研究者自己写明了。」}{「说会了，和做出来对，是两回事。」}{「研究证明人估自己会不会做事估不准」}{Rozenblit \& Keil 2002（\#201）；研究 8，37 名被试}
\canonrow{情绪按责任分流·Rudolph 2004（复述对方）}{「2004 年汇总六十四项研究、一万两千多名被试的元分析发现：判断他该负责就生气，判断他控制不了就同情。这条我们全认；可它只管判断之后生出什么，不管判断本身准不准。」}{「元分析管的是判断之后，不管判断这一步。」}{①「元分析证明厌不肯的人没错」②报效应量③说成研究对象是蠢}{Rudolph 等 2004（\#219）；64 项研究，12,000 余名被试}
\canonrow{恼传过去的是什么·Graham 1984（复述对方）}{「Graham 那项实验测的是孩子听到大人生气以后怎么想，不是生气的人心里还看不看得起他；摘要还写，生气那条联系比同情那条弱。」}{「被试是孩子，表达情绪的是实验者，一次实验室任务。」}{「被生气反而更有信心」}{Graham 1984（\#250）}
\canonrow{坏印象更稳·Denrell 2005}{「2005 年的模型研究说：坏印象让人少接触，所以坏印象比好印象更稳；只有经常被迫接触的人，坏印象才更容易改。」}{「接触回来，不等于讲理回来。」}{①只引后半句②说成实测}{Denrell 2005（\#193）}
\canonrow{原谅研究·McCullough 三篇（复述对方）}{「原谅研究里，原谅会随时间上升，但跟『这段关系还有没有价值』的感觉绑在一起；道歉、补偿能让它加快。原谅测的是回避和报复的念头下降，不测还讲不讲理，对象也是人际冒犯，不是蠢。」}{「回来的是善意和来往，讲理回不回来没人测过。」}{①把「两个对数优势单位」换算成百分比或任何比例②「原谅研究证明厌蠢会自己消退」③「回来的就是讲理」}{McCullough 等 2010（\#228）；McCullough 等 2003（\#229）；McCullough 等 2014（\#199）；337 名大学生}
\canonrow{少见行为分量重·Mende-Siedlecki 2013（复述对方）}{「越少见的行为，在评价里分量越重；放到『蠢』上往哪边推，双方都只能推论。」}{}{「脑成像证明……」}{Mende-Siedlecki 等 2013（\#240）；研究 1 80 人，研究 2 fMRI 23 人}
\canonrow{求助显能干·Brooks 2015 与 MIT 斯隆报道}{「那项研究说，人怕求助显得无能，所以不愿求助；任务难的时候，求助反被看作更能干。任务容易时，据《MIT 斯隆管理评论》2015 年的报道，求助既不加分也不减分。」}{「厌蠢说的恰恰是『这么简单都不会』，那一格问了也不显聪明。」}{①「论文发现容易任务求助不加分」②「求助总会被看不起」}{Brooks 等 2015（\#227）；MIT 斯隆管理评论 2015（\#181）}
\canonrow{怕显得无能·Edmondson 1999}{「1999 年对一家企业五十一个团队的研究写道：承认错误、求助的人可能显得无能，所以人不敢提；组织里的人常常不愿披露错误、不愿求助。心理安全量表第一题就是：『在这个团队里犯了错，常常会被记在账上。』」}{「心理安全越高，团队越会承认错误、求助、讨论问题。」}{①「心理安全决定三成学习」②「研究测了厌蠢」③说量表「反向计分」}{Edmondson 1999（\#64、\#184、\#222）；一家制造企业 51 个工作团队，横截面、相关；表 6 观察者评定学习行为 B=.46、R²=.33}
\canonrow{纠正没有逆火·Wood \& Porter 2019}{「五个实验、一万多名被试、五十二个议题里，没有一种事实纠正引发逆火：纠正通常有效。」}{「纠正，人是听得进的。」}{「纠正能改变行为」}{Wood \& Porter 2019（\#109）；5 个实验，10,100 余名被试，52 个议题，美国样本}
\canonrow{提醒准确性·Pennycook 2021}{「人分享错误信息，常常是因为注意力不在准不准上；轻轻提醒一下准确性，分享的新闻质量就上来了。」}{「拉回注意力，比看低他管用。」（后半推论）}{①报效应量②「提醒能消灭谣言」}{Pennycook 等 2021（\#108）；4 项调查实验与 1 项推特现场实验}
\canonrow{懒得想·Pennycook \& Rand 2019（复述对方）}{「那项研究测的是认知反思测验得分和信假新闻的相关，是相关研究；同一组作者的另一项研究说，注意力一转就能改善。」}{「懒得想是注意力放错了，提醒就能拉回。」}{「认知反思测验测的是智力」}{Pennycook \& Rand 2019（\#107）；N=3,446}
\canonrow{愤怒被点赞就更常表达·Brady 2021（对方会引，我方反用）}{「愤怒的表达被点赞，人以后就更常这样表达。被奖赏的态度会更常出现，放大的不是假象，是习惯。」（第二句是推论）}{「被强化出来的厌蠢也是厌蠢。」}{「这项研究测的是厌蠢」}{Brady 等 2021（\#101）；7,331 名用户、1,270 万条推文与行为实验}
\canonrow{共情费力·Cameron 2019（整句引）}{「共情研究说，给人选择时，人倾向回避共情，因为觉得费力；可同一研究提高了共情的效能感以后，回避就消失了。」}{「累是原因，不是理由。」}{①只引前半句②「共情做不到」}{Cameron 等 2019（\#98）；11 项研究，1,204 人}
\canonrow{共情疲劳的边界（复述对方）}{「护理人员的共情疲劳，来自接触创伤和重负荷，不是来自别人蠢；国内的综述也说，非助人行业有没有共情疲劳，还有待探讨。」}{「『你只是太累了』是自媒体的说法。」}{「共情疲劳人人都有」}{Yi 等 2024（\#99）；江语童等 2024；滑洋心理 2025（\#18）；196 项研究}
\canonrow{被轻蔑的代价·Melwani \& Barsade 2011（复述对方）}{「被轻蔑的人任务表现是好了，可对表达者的攻击性也明显更高，中间靠的是内隐自尊下降；而且它测的是被『表达』出来的轻蔑。」}{「对方主张不看低、不表现，引它就等于自认表达了轻蔑。」}{只引「表现更好」}{Melwani \& Barsade 2011（\#87、\#31）；三项实验，程序扮演的虚拟伙伴}
\canonrow{守标准三条的对象（复述对方）}{「这三条的对象都是能做却不做：故意搭便车的人，能出力却不肯的搭档，不做自己那份的成员。那是坏，不是蠢；而且没有一条操纵过『有没有人反感他』。」}{「守标准，对事的恼就守得住。」}{①网传「坏苹果」拖低绩效的比例②「研究证明反感不能维持合作」}{Fehr \& Gächter 2002（\#110）；Kerr 1983（\#220、\#190）；华盛顿大学新闻稿 2007（\#65）}
\canonrow{愤怒是趋近型·Carver \& Harmon-Jones 2009（复述对方）}{「说愤怒是趋近型的那篇综述，比较的是愤怒和焦虑，不是愤怒和轻蔑；趋近不等于讲理，更不等于和好。」}{}{「愤怒都会导向和解」}{Carver \& Harmon-Jones 2009（\#244）}
\canonrow{被动伤害·BIAS map 2007}{「刻板印象研究里说的被动伤害，正是忽视、不与之交往。」}{「『我只远离』，在这项研究里叫被动伤害。」（群体层面，迁移是推论）}{「远离就是伤害，研究证明了」}{Cuddy 等 2007（\#30）；美国样本，群体层面}
\canonrow{错在系统·Reason 2000 等（只作机制例证）}{「风险管理那头的说法是：怪个人在情绪上比追究机构更痛快；最好的人也常犯最糟的错，错误不是少数倒霉人的专利。」}{①「有些不安全行为确实恶劣，但绝大多数不是。」②「美国医学研究所 1999 年的报告说，多数医疗差错不是『坏苹果』问题。」}{①当主战场②「医疗差错每年死二十五万人」③报「约九成质量失误被判无可指责」}{Reason 2000（\#183、\#169、\#105）；IOM 1999（\#104）；NASA 1999（\#106）}
\canonrow{理性不等于智力·Stanovich（用来要求对方搭桥）}{「理性和智力是两回事：斯坦诺维奇把『智力足够、却不能理性地思考和行动』叫作理性障碍；大量思维偏差，跟认知能力不相关。」}{「你们自己的定义说，蠢不等于智商。」}{①「1993 年论文写道……」②「聪明人和普通人犯一样多的错」}{Stanovich 2009（\#70）；Stanovich \& West 2008（\#71）；7 项研究}
\canonrow{人们把什么叫蠢·Aczél 2015}{「人们把三类情境叫作蠢：信心和能力失衡、注意失误、缺乏控制；行为人责任越大、后果越重，就越显得蠢。」}{被对方引「把行为叫作蠢，反映人们要求自己也要求别人的期待」时：「那是作者写在摘要开头的研究动机，不是研究发现。」}{①报样本数和各类比例②「研究证明蠢都可以归责」}{Aczél 等 2015（\#93、\#239）}
\canonrow{想别的原因能减偏·Roese \& Vohs 2012（备用）}{「后见之明的综述说，鼓励人去考虑别的原因，能减少这种偏差。所以我们要的改法是开口问，把别的原因问出来。」（第二句是推论）}{}{①「能消除偏差」②「能减少很多」}{Roese \& Vohs 2012（\#207）}
\canonrow{教师期望更多因为准·Jussim \& Harber 2005（复述对方；回应先搭桥）}{「三十多年的教师期望研究，综述说：课堂里的自我实现预言确实有，但通常很小；教师期望能预测学生表现，可能更多是因为期望本来就准。」}{回应先搭桥：「那篇测的是老师对学生学业表现的预期，不是对蠢的判断；老师有一学期的作业、考试可看，那是记录，凭记录的准我们本来不争。」（后半是推论）；「同一篇还写，强有力的自我实现预言，可能选择性地发生在受污名群体的学生身上。」}{①「研究证明老师看学生都准」②把学业表现说成智力或蠢③「贴标签会把人变蠢」④报效应量}{正驳 45；Jussim \& Harber 2005（\#255、\#257）}
\canongroup{案例}
\canonrow{小张评论区（论点一主案例之一）}{「2023 年初，一位博主发视频，教第一次去高铁站的人怎么坐高铁。评论区里有人说：『怎么会有人不会坐飞机啊？我厌蠢症犯了。』还有人留言：『坐高铁不是谁都会吗？这个也要教？』」}{①「『厌蠢症』这个词，就是在这个评论区出圈的。」（上游新闻的判断）②「记者把那条留言写作『调侃』。」③「它冲着的，恰恰是来学的人。」}{①「满屏都是，说明多数网友厌蠢」②报涨粉的精确数（两家报道口径不同，只说「上百万」）③「厌蠢症就是从这里诞生的」}{上游新闻 2023-02-24（\#167）；扬子晚报 2023-02-15（\#119、\#225）；新京报 2023-11-01（\#213）}
\canonrow{新京报：同学同事那段（论点一主案例之二）}{「新京报 2023 年的一篇文章，写到被吐槽的同学、同事，写的是『你能不能有点常识』的烦躁；作者说这种愤怒背后的思维，『更像是一种我都可以做到你为什么做不到的抱怨』。作者还写，被叫作蠢的人，『有许多却是努力了但也没有学会的人』。」}{①「作者还说，把失败归咎于不够努力，『显然是不客观』的。」②「同一篇写：常识本是由个体社会经验组成的。」}{①「新京报认为厌蠢错了」「新京报证明」②「多数被厌的蠢是不会」③「新京报是央媒」④对方补出「比缺乏常识更可怕的是理直气壮」时否认（要认：那是加码，先被厌的仍是无知与无能）}{新京报 2023-11-01（\#212、\#213、\#243、\#249）}
\canonrow{每日商报：一年后照列}{「出圈一年后，每日商报 2024 年 1 月的报道，开头小标题是『不要用自己的常识去攻击别人的无知』，下面列着两句：『怎么会有人不会打车啊？我厌蠢症犯了。』『怎么还有人到医院不会挂号呀？太蠢了吧。』」}{「报道开头列出的说法。」}{①「网友原话」②把「常识本就是由个体社会经验组成的」说成医生原话}{每日商报 2024-01-13（\#214）}
\canonrow{每日商报：程序员与同事（论点二主案例，三段一起说）}{「还是这篇每日商报：一位程序员说，同办公室的同事工作两年多，还常把表格数据填错；他也曾多次劝告对方工作上要多上心，可同样的错误还是接二连三地再犯，同事间的关系一度降至冰点，现在『陷入一种很微妙的状态』。报道紧接着这个案例，单独写了一句：『说了那么多遍，怎么还是那么蠢。』」}{①三段已说过之后可以回指：「说了那么多遍，怎么还是那么蠢。」②「从『多上心』的要求，走到『那么蠢』的判词。」③「这一案在准上我们让：两年多、劝了多次、接二连三。」}{①「小陈说：说了那么多遍……」②「小陈从此不理小刘」「关系到现在还是冰点」③「小刘就是无所谓」「转眼就忘是事后可查的记录」④「职场里的厌蠢多半是这样」「厌蠢一词被用得极为频繁」⑤第一次引只说其中一段⑥「报道写了回头不讲理」}{每日商报 2024-01-13（\#248）}
\canonrow{扮蠢的大三学生}{「一位扮蠢的大三学生自己说，好处之一是小组作业时『大家不会把最难的部分分给你』。她是装的，同学的反应是真的：蠢的标签一贴，机会就少了。」（后半是推论）}{「能改写评价的表现，要先有机会。」（推论）}{①「职场里大多是装蠢」②「她就是蠢」}{真故研究室 2024-07-07（\#237、\#123）}
\canonrow{Torvalds 道歉}{「Linux 的创始人 Torvalds 2018 年公开道歉：他说自己邮件里那些轻率的攻击，既不专业也没有必要，尤其是上升到人身的时候；『在追求更好补丁的过程中，这在我看来说得通；现在我知道这不 OK。』他还说要去学着理解别人的情绪、恰当回应。而内核照样拒绝不合格的补丁。」}{「守质量，用不着看低人。」}{①「Torvalds 因为厌蠢道歉」②「他至今觉得骂人说得通」③「他赶走了多少人」}{Torvalds 2018（\#139、\#186、\#172）；Linux 内核行为准则 2018（\#221、\#189）}
\canonrow{朋霍费尔（只引一句）}{「朋霍费尔 1942 年底写了〈论愚蠢〉，紧接着下一节的标题就是『蔑视人？』，开头一句：被推向蔑视人的危险很大。」}{对方引中文流传本「愚蠢是道德上的缺陷」时纠正：「原文说的是人的缺陷，不是智力上的缺陷。」}{①「朋霍费尔只是暂停说服」②「朋霍费尔是我们改法的榜样」③「狱中写下」④「愚蠢是一种道德上的缺陷」⑤再引同一节别的句子⑥拿他说日常同事}{Bonhoeffer〈蔑视人？〉（\#211）；Bonhoeffer〈论愚蠢〉（\#247、\#75）；群学书院（中文流传本）2021（\#133）}
\canonrow{孔子与子路（对方案例，我方口径）}{「孔子骂了子路『野哉』，还把正名的道理讲完。这是圣人、是老师，正是我们要常人学的改法，不是常人厌蠢平常的样子。同一部书里，门人听孔子批评子路一句，就不敬子路了；拦住他们的是孔子，朱熹注：『未可以一事之失而遽忽之也』。」}{「常人没有老师的职分，也没有圣人的修养。」}{①「孔子也厌蠢」②「子路鼓瑟是蠢」}{《论语·子路》；朱熹集注（\#226）；《论语·先进》；朱熹集注（\#245）}
\canonrow{莎莉文与海伦·凯勒（对方案例，我方口径）}{「莎莉文教海伦·凯勒，绝望的时候把这个话题暂时放下，一有机会又重新开始。绝望之后是重新开始，这正是我们要的改法；她是老师，海伦是聋盲，不是蠢。」}{}{①「莎莉文从不发火」②拿它证常人平常就会回头}{Keller 1903（\#142）}
\canonrow{不屑之教诲（对方案例，我方口径）}{「孟子讲的是老师的教法；朱熹的注是条件句：对方若能有所感触、回去自我反省，才算是教诲了他。」}{}{「孟子说不教也是教，所以走开没错」}{《孟子·告子下》；朱熹集注（\#232）}
\canonrow{郑人买履（对方案例，我方口径）}{「郑人那一问，是在『及反，市罷』之后：集市散了，没有鞋可试。」}{「原文是『寧信度，無自信也』。」}{「宁信尺码，也不信自己的脚」}{《韩非子》郑人买履（\#202、\#233）}
\canonrow{苏格拉底的自知}{「苏格拉底说：他什么都不知道却以为自己知道；我既不知道，也不以为自己知道。自以为知的，也可以是那个觉得『这都不会』的人。」}{}{「苏格拉底厌恶无知的人」}{Plato《申辩》（\#146）}
\canonrow{伪名言与读全的古句（防守）}{「这几句先查出处：『别和傻子争论』找不到出自马克·吐温的实质证据；『宇宙和人类的愚蠢』大概不是爱因斯坦说的；颜回『三八二十三』《论语》里没有。《箴言》下一节就是『要照愚昧人的愚妄话回答他，免得他自以为有智慧』；席勒那句，是剧中濒死败将的诅咒。」}{①「虚假论据按章程判负，我方可以当场申诉。」②「英语里常说的汉隆剃刀：能用愚蠢解释的，别归于恶意。」③「suffer fools gladly 源自保罗的反讽，今天常被用来开脱暴躁难相处的行为。」}{①对「可以原谅坏人，不能原谅蠢人」说「这句是假的」或「是某某说的」②报维基当出处（口头只说「英语里常说的」）}{Quote Investigator 2023（\#127）；Quote Investigator 2010（\#132）；《论语》全文；台北市孔庙 2019（\#135）；《箴言》26:4–5（\#128）；《箴言》14:7（\#129）；Schiller 1801（\#131）；《哥林多后书》11:19；Wikipedia（\#136）；Wikipedia「Hanlon's razor」（\#130）}
\canonrow{《奇葩说》2016}{「2016 年《奇葩说》那期辩『跟蠢货做朋友是不是傻』，高晓松说：『坏人不是蠢人。』」}{对方用「蠢人就是坏人」时用它切割：「蠢不是坏。」}{①「蔡康永说过蠢人常常是坏人」②把「朋友不超过 150 人」当论据③说成正式赛事、报胜负}{澎湃新闻·有戏 2016-03-08（\#111）}
\canonrow{「智商税」的口气}{「爱范儿 2022 年那篇里，受访者说：别人用开玩笑的语气说『智商税』，『我会觉得非常不礼貌且傲慢』。说的人觉得轻，听的人未必。」}{}{「智商税一词 2013 年起源于知乎」}{爱范儿 2022-04-10（\#125）}
\canonrow{媒体批评引句（只说「某某文章写道」）}{「有间大学的文章写：『无论怎样把它调侃为一种疾病，都不能掩盖背后浓浓的傲慢和尖刻』；『我只厌打扰到我的蠢』这个前缀，是为了不显得傲慢才加的。」}{①新京报书评周刊认为：「『厌蠢症』这个标签的出现给了我们合理化贬低他人的契机」②央视网评论：「看不得别人犯蠢、没常识、太笨」③中国网的文章写：「有的人或许不会用语言来表达，但会用行动离对方远远的」}{①「新周刊说」②「新京报书评周刊证明精英最容易厌蠢」③报央视网文中无出处的人数等数字④「心理学证明厌蠢是投射」}{有间大学 2023-12-25（\#14）；新京报书评周刊 2023-11-01（\#13）；央视网 2023-02-23（\#10）；中国网·心理中国 2026-02-27（\#16）}
\canonrow{对方会引的「功能说」（我方口径）}{「明确写成医生原话的只有一句：『厌蠢背后其实是对潜在风险的厌烦』。风险是真的，防范也该有；可讨厌麻烦叫怕麻烦，用不着『厌蠢』这个词。」}{「咨询圈那篇是自媒体，原句前面还有『很多情况下』。」}{①把每日商报其他段落说成医生原话②省掉「很多情况下」}{每日商报 2024-01-13（\#15）；张琨 2025-08-11（\#17）；汉典「厌烦」（\#6）}
\canonrow{修辞与话语样本（不当论据）}{「穆齐尔说，骂一句『蠢』，是回应的开始，也是它的结束。」}{只作修辞，结辩可用一次。}{①把穆齐尔当论证②把奇波拉「蠢人比强盗更危险」当研究③把自媒体「蠢就是一种天生无可救药的坏」当论据}{Musil 1937（\#77）；Cipolla 1987（\#73）；有画说（自媒体）2025-03-02（\#124）}
\canongroup{不进稿件的材料（有把握不进文稿；【待核实】文稿、速查都不进；已核实但不上场的照理由处理）}
\canonrow{有把握·Aczél 2015 的数}{只说三类情境与「责任越大、后果越重，越显得蠢」（\#93），不报任何数。}{}{二手报道里的样本数、三类情境各占多少、「自信的无知」评分最高}{Aczél 等 2015（\#69，有把握）}
\canonrow{有把握·基本归因错误定义句}{改用对应偏差（\#61）。}{}{Ross 1977 的定义原句}{Ross 1977（\#90，有把握）}
\canonrow{有把握·医疗差错死亡数}{改用美国医学研究所 1999：多数医疗差错不是「坏苹果」问题（\#104）；专业领域只作机制例证。}{}{「美国每年死于医疗差错的人数」「第三大死因」}{Johns Hopkins 2016（\#103，有把握）}
\canonrow{有把握·Strawson 原刊卷页}{只说「Strawson 1962 年的《自由与怨恨》」，引句照全文转载页。}{}{报原刊卷页}{\#32、\#33、\#34、\#174、\#197 的卷页部分}
\canonrow{有把握·Darwall 卷期页}{只说「据《斯坦福哲学百科》」，承认尊重与评价尊重的区分照 SEP 原文。}{}{报 Ethics 卷期页}{\#42}
\canonrow{有把握·亚里士多德译者与页码}{只说「《尼各马可伦理学》第三卷 / 第四卷」，引句照 Gutenberg 原文。}{}{报译者与 Bekker 页码}{\#39、\#41、\#230 的译者与页码部分}
\canonrow{有把握·蔡康永那一句}{用同条已核实的高晓松「坏人不是蠢人」。}{}{「蔡康永说过蠢人常常是坏人」}{\#111 局限栏复核补记（有把握）}
\canonrow{待核实·《现代汉语词典》「厌」「蠢」}{「汉典释『厌』为嫌恶、憎恶」；汉典「蠢」「蠢人」（\#1、\#66、\#67）。}{}{「《现代汉语词典》释……」}{\#7、\#68（【待核实】）}
\canonrow{待核实·「错」的词典释义}{只说「『我错了』就是认错、该改，不是自认坏人」，当日常语感说。}{}{「词典释『错』为……」}{调研未收}
\canonrow{待核实·朴素实在论的经典出处}{需要时说「Pronin、Gilovich 与 Ross 2004 年的综述」（\#209）；本方主论证不用朴素实在论。}{}{「朴素实在论由 Ross 与 Ward 1996 年提出」}{\#178（【待核实】）}
\canonrow{待核实·厌蠢比例}{任何「百分之几的人厌蠢」都不说；对方报比例时，要求当场出示来源。}{}{任何厌蠢比例}{\#102（检索未果）}
\canonrow{待核实·其他}{本方不用 Hong 与 Page；对「可以原谅坏人，不能原谅蠢人」只说「请出示出处」；对方引 Mason 2003 时答 Bell：轻蔑该回应的是傲慢、伪善这类优越之恶（\#35）。}{}{Thompson 2014 的批评；说「这句是假的」或把它安到任何人名下；Mason 2003 为针对人的轻蔑辩护}{\#96、\#134、\#166（【待核实】）}
\canonrow{待核实·局限栏里的数}{〔智〕材料本方不用，对方报系数时要求出处并要求搭桥；只说新闻稿里「坏苹果」的定义是不尽本分、情绪不稳、欺凌；评论只引已核实的定义句；不说「智商税」的词源。}{}{陌生人智力评分的相关系数；网传「坏苹果」拖低绩效的比例；央视网、每日商报文中无出处的人数、城市数与未核的老人用手机比例；「智商税」起源于某网站的说法}{\#88、\#65、\#10、\#15、\#125 的局限栏}
\canonrow{不上场·比例类}{比例不上场，容易被听成厌蠢比例；对方报出时回一句「那不是厌蠢的比例」。}{}{市场力量约减偏 50\%（\#79）；事前组四分之三认为不值得防（\#206）；87.6\%、51.8\% 的职场新人（\#121、\#173）；91.36\% 子女以为父母学不会等（\#192、\#175）；本科及以上 46.20\%、小学及以下 3.11\%（\#83）；65\% 美国人自认比一般人聪明（\#91）；35\% 与 55\% 的表情识别（\#24）；约九成质量失误被判无可指责（\#105）}{同左各条（已核实）}
\canonrow{不上场·〔智〕组}{测的是智力或认知能力，本方主论证不用；对方用时答：「你们自己的定义说，蠢不等于智商；请先搭桥。」}{}{智力评分效度、自评与实测相关、个体 IQ 解释群体差异、遗传率、教育提高 IQ 分、能力最差者自估的百分位}{\#88、\#89、\#92、\#94、\#95、\#80、\#81、\#82、\#57、\#62}
\canonrow{不上场·其余}{Kuppens 只在把「我」拉近台下时作一句旁证，不打头阵；「认定他听不进」已淘汰；新国辩结辩标题是账号编辑起的，复赛票数说明不了持方优劣；原谅研究的数值不复述；「150 人」不能当论据。}{}{「研究证明厌蠢的人有偏见」；把意见不合看成有偏（\#208）当论据；「厌蠢的人大多被无知伤害过」当统计；报新国辩票型当论据；「两个对数优势单位」或换算成比例}{\#54、\#55、\#56、\#208、\#114、\#113、\#228、\#111}
\canonrow{赛前待办}{康德两句换成李秋零译本之前，场上说「据德文原文的意思」；朋霍费尔同样据德文自译，只引「被推向蔑视人的危险很大」；赛前读一遍 Hom 2023 的方法部分，被问到「设想同伴经历相同」能当场复述。}{}{「李秋零译本写道」}{\#48、\#45、\#211、\#75、\#176}
\end{canon}
\begin{fields}
\field{滑坡清单}{听到队友说出口，下一个发言的人当场拉回。\sub{相关说成导致、可能说成必然}{「轻蔑从愤怒里长出来」只说可能；预测因素不说成一定；Edmondson、Pennycook 与 Rand 是相关研究；后见之明、知识诅咒的实验迁移到「这都不会？」要说成推论。}\sub{智力说成蠢}{〔智〕组一条不用，「智商」「IQ」不出现在我方论证里。}\sub{一时走开说成从此走开}{「小陈从此不理小刘」「冰点到现在」都不说；「一走开就是越线」也不说。}\sub{有时说成每次、偏一截说成判错}{不说「厌蠢每一次都判错」「研究证明厌蠢偏了零点三九」「屡次也救不了」。}\sub{描述说成统计}{「许多」「总有」「满屏」只说「某某文章写道」；任何百分比都不报。}\sub{把有把握、待核实的说出口}{「现汉释……」「医疗差错死亡数」「Aczél 的样本数」「Strawson 原刊卷页」「蔡康永说蠢人常常是坏人」都不说。}\sub{打成霸凌、要求圣人}{不说「厌蠢就是霸凌」「你们第一秒就不该厌」「别有感觉」「压住不厌」「要一直耐心教下去」。}\sub{媒体说成证明}{不说「新京报证明」「新京报认为厌蠢错了」「新京报是央媒」「专家说厌蠢是慕强」。}\sub{推论说成来源}{不说「报道写了回头不讲理」「Strawson 说厌蠢就是客观态度」「后见之明研究测的就是厌蠢」；被问「是不是推论」，直接认，然后给推论的依据。}}
\end{fields}

\consubsection{10.2 反方口径表}
\begin{canon}{反方}
\canongroup{定义}
\canonrow{定1 厌蠢（中立定义）}{「厌蠢，是觉得别人犯蠢时生出的嫌恶、反感，常带着想躲开的心；它可以只冲着那件事，也可以冲着这个人、看低他；它是一种态度，不等于骂人或欺负人。」}{被问「是不是把鄙视塞进了厌」：「两头都写了，都用『可以』，没替任何一方选；平常落在哪一头、落在那里错不错，是今天要辩的。」；正方把「看低」写成厌的全部时：「汉典里，厌是嫌恶、憎恶；看不起是鄙视的释义。中立定义两头都写了，平常落在哪头要证，不能在定义里先定。」}{「厌蠢是看到别人犯蠢时……」；「厌只是生理反应」「厌是自然生起的烦」；「厌蠢就是厌麻烦」「就是不耐烦」；「厌蠢只讨厌蠢事，和人无关」；「题目里有『症』字，所以由不得我」；「《现代汉语词典》释厌为……」}{第三节；汉典「厌」（\#1）；汉典「厌恶」（\#2）；汉典「嫌弃」（\#4）；汉典「鄙视」（\#5）；《重編國語辭典修訂本》2021（\#3）；中国网·心理中国 2026-02-27（\#16）}
\canonrow{定2 蠢（中立定义）}{「蠢，是在一般人看来本该想得到、做得对的事上，明显想错、做错；可能出于不会，也可能出于不肯，原因不在定义里预设。」}{被问「一般人是谁」：「了解他处境、讲道理的一般人，不是评论区的多数票。」（口播不加字，被问才答）；短版：「本该会，却明显做错。」；「明显」二字：「排除鸡毛蒜皮。」}{把「明知故犯」「拒不改正」「理直气壮」「又菜又犟」「让别人买单」写进定义；「不知道叫无知，不肯想才叫蠢」；「蠢就是智商低」，拿 IQ、遗传、教育年限讲蠢；「不会的不算蠢」（改说：「第一次、没人教的不会，一般人不觉得他本该会。」）}{第三节；汉典「蠢」（\#66）；汉典「蠢人」（\#67）；Aczél 等 2015（\#93）；Stanovich 2009（\#70）；Bonhoeffer〈论愚蠢〉（\#75）}
\canonrow{定3 错了 / 没错（与判准同一句，全队只背这一句）}{「判断这份厌对不对，就看一个清醒的我问完三句：看准了吗？对他公道吗？划得来吗？问完觉得这份厌平常的样子就该改，就是错了；站得住，只在个别场合调调分寸，就是没错。」}{「我没错，不等于我完美。」；「平常的样子该改，是错了；平常站得住、个别场合调分寸，是没错。」}{「错就是该受谴责」「伤了人才算错」；「我没错，就是厌蠢是美德」；「情绪无所谓对错」；「词典里『错』的意思是……」}{第三节；第四节}
\canonrow{定4 「我」是谁}{「一个普通的、会说出这句话的我。」}{「不是网暴者，也不是圣人。」}{「我们这种人」「聪明人」「我们中科大的」之类的精英腔；「网友自己都会划界限，所以常人都这样」}{第三节；第四节；上游新闻 2023-02-24（\#168）}
\canonrow{定5 只在心里厌，算不算厌蠢（双方统一答）}{「算，心里那份就是厌蠢本身；它错不错，看判准。」}{「算。心里那份我们也放进三问里称。」}{「只在心里就不算厌蠢」「心里的厌谈不上对错」}{第三节}
\canonrow{定6 第一次不会，算不算蠢（正二最可能问）}{「不算。定义里的一般人，是了解他处境、讲道理的一般人；第一次、没人教，这样的人不觉得他本该会。我们的定义不冤枉人。」}{追问「那评论区为什么说『这不是谁都会吗』」：「评论区的多数票不是定义里的一般人，那是评论区看错了。他们那份也算厌蠢，看错是第一问要称的；但这是个别场合的误判，不是这份厌平常的样子。」}{「评论区那种不算厌蠢」；「第一次不会也该被厌」}{第三节；工人日报 2023-03-03（\#120）}
\canonrow{定7 把骂人切出去，剩下的还有什么可评}{「切出去的只是霸凌。想躲开的心、冲着人的看低都还在定义里，欢迎来评。」}{「定义只切掉骂人、欺负人；态度本身整个留在题目里。」}{「厌不伤人就没错」；「我们只辩护不说出口的厌」}{第三节}
\canongroup{判准}
\canonrow{判1 中立判准}{「判断这份厌对不对，就看一个清醒的我问完三句：看准了吗？对他公道吗？划得来吗？问完觉得这份厌平常的样子就该改，就是错了；站得住，只在个别场合调调分寸，就是没错。」}{自由辩短版：「问完三句：看准了吗，对他公道吗，划得来吗。」；「准、公、值」只在队内说，场上说全称。；「清醒」被问时：「清醒只是说，双方摆出来的事实他都知道：蠢的成因和代价，厌的后果和功能。」}{「只看对不对得上，看准了就没错」；「我有权」「只看伤没伤人」「情绪无对错」；「没证明错就是没错」}{第四节；Frankfurt 1971（\#157）}
\canonrow{判2 反一开场：正方用了中立版}{「定义与判准，和对方大体一致。」（随即接判3 门槛）}{「定义判准我们不争，我们争的是门槛。」}{自造新版定义或判准；重复念一遍正方已念过的原文，浪费时间}{第四节；第九节}
\canonrow{判2b 一辩开头合并版：正方用了偏向版时重述}{「厌蠢，是觉得别人犯蠢时的嫌恶、反感，常带着想躲开的心；可以只冲着那件事，也可以冲着这个人、看低他；它是态度，不等于骂人或欺负人。蠢，是在一般人看来本该想到、做对的事上明显想错、做错，不论出于不会还是不肯。错没错，看清醒的我问完三句：看准了吗？对他公道吗？划得来吗？这份厌平常的样子该改，就是错了；只在个别场合调调分寸，就是没错。」}{先用一句话点出偏在哪：「对方把看低写死在厌里了。」或「对方把门槛压成了『有任何毛病』。」再重述。}{删掉「也可以冲着这个人、看低他」半句；改写成「看到别人犯蠢时」}{第四节}
\canonrow{判3 门槛分界规则}{「要改的如果是这份厌平常的样子，就是这份厌本身该改，哪怕改完还剩一点反感；要改的如果只在个别场合冒出来，那叫调分寸。平常就这样、只是改得掉，照样算错了。」}{「只在个别场合冒出来的，调分寸就好；要改的是这份厌平常的样子，才叫错。」；「门槛不看改不改得掉，看平常还是个别。」}{「改得掉，所以只是调分寸」；「有点毛病也不算错」；「那一份我们认错」「厌蠢有时候是错的」}{第三节；第四节}
\canonrow{判4 三问各问什么}{「看准了吗：对方真的蠢吗，我了解他的处境吗，从这件事推到这个人推得准吗，证据够不够。对他公道吗：这份蠢该不该由他负责，厌了以后还给不给他作为人应得的尊重。划得来吗：这份厌让事情、关系和我自己，离我在乎的东西更近还是更远。」}{「公道问两件事：该不该由他负责；厌了以后，还给不给他作为人应得的尊重。」}{「公道只问该不该他负责」（漏了尊重小问；「他不肯学」只回答了前一件）；「划得来就是有用就对」}{第四节}
\canonrow{判5 两条约定}{「评价对象是一个普通的、会说『我厌蠢』的人，不是网暴者，也不是圣人。『对他公道吗』既算他没选过的成因，也算不肯学这类可以归责的成因。」}{「拿网暴者当靶子，换了评价对象；拿圣人当标准，也换了。」}{「我们评的是最好的那份厌」}{第四节；有画说（自媒体）2025-03-02（\#124）}
\canonrow{判6 只赢一问；替代态度后登场}{「只赢一问，不自动赢下整场。存在更好的，不等于这个就错了；替代态度，只在对方证出平常的缺陷之后才登场。」}{「没错，不等于最优。」}{「我们的厌就是最好的态度」；「对方说的替代我们都做不到」}{第四节；Scarantino \& de Sousa 2018（SEP）（\#164）；D'Arms \& Jacobson 2023（\#59）}
\canonrow{判7 分歧裁决（全场核心一句）}{「对方只证明了有，没证明平常。按判准，一方证出平常就有的东西，另一方只证出个别场合的东西，前者胜。」}{「有，不等于平常。」}{「对方的研究都是假的」；「偏差很小，所以没事」}{第四节；第九节}
\canonrow{判8 分量：归责小问最重}{「常人判断一份恼火对不对，本来就看可不可控、该不该他负责。有记录时，看准和归责一起过关。」}{「该不该他负责，是常人称恼火的第一把秤。」}{「看准了就自动公道」（尊重小问还要另答，见论2）}{第四节；第九节；Weiner 1993（\#58）；Rudolph 等 2004（\#219）；Aczél 等 2015（\#93）}
\canonrow{判9 修正不等于压抑}{「问完三句之后的改，指的是换一种看法，不是别有感觉。我们不拿『控制不了』免于评价，对方也不要求我们压着不厌。」}{「我们不说管不住，我们说站得住。」}{「情绪控制不了，所以无所谓对错」；对方没这样说时，指责对方「要我们压抑」}{第四节；Webb 等 2012（\#155）}
\canonrow{判10 「对自己」的代价分量偏低}{「普通人不会因为生气伤身，就说一份该有的义愤是错的。」}{「累，能解释情绪为什么发生，不说明它错。」}{「生气对身体没坏处」（数17 是相关研究，我们不否认，只说分量低）}{第四节；Chida \& Steptoe 2009（\#63、\#97）}
\canonrow{判11 推定}{「我没错，也是一句正面的自我辩护；我们正面举证：问完三句，这份厌平常的样子站得住。」}{「我们不靠对方举证不足赢。」}{「对方没证明错，我们就没错」}{第四节}
\canonrow{判12 被问必答：看准了就够了吗}{「不够，所以我们同意三问都问；我们要证明的是，问完三句，这份厌平常的样子站得住。」}{「看准只是第一问。」}{「看准了就没错」}{第四节}
\canonrow{判13 被问必答：看准了这个人蠢，就能看低他吗}{「看低他的能力可以，拒绝给他作为人的尊重不行；我们的厌停在前一边。」}{「看低他这件事上的表现可以；拒绝他作为人应得的尊重，不行。」}{「看低他这个人可以」；「我们从来不看低任何人」}{第四节；Kant 1797，§39（\#48）；Roberts 2013 书评（引 Bell）（\#37）；SEP「Respect」（\#44）}
\canonrow{判14 被问必答：分寸调到不看低他，不就是改了你平常的厌}{「看低不是这份厌平常的样子，只在个别场合冒出来，所以是调分寸。平常那份反感冲着事，它告诉我标准被踩了，我们守的就是它。」}{「调掉的是个别场合冒出来的那一步，留下的是平常那份冲着事的反感。」}{「看低也是平常的一部分，但改得掉」}{第三节}
\canonrow{判15 被问必答：厌了之后，是去交涉还是走开}{「厌，可以一时走开，但不把他划出讲理的范围：回头照样教、照样合作。」}{「这回先放下，下回照样说、照样听。」}{「从此走开」「眼不见为净」「我有权远离」}{第四节；Strawson 1962（\#34）}
\canonrow{判16 决胜点}{「全场落在一处：常人说『我厌蠢』时，那份厌平常停在事上，还是落在看低人那一头。我们说：平常停在事上，看低和把人划出讲理的范围，只在个别场合冒出来。」}{「平常停在事上，个别场合才越线。」}{「从来不会越线」}{第四节}
\canonrow{判17 全场退守线}{「我们不说情绪无对错；我们说问完三句，那份平常冲着明显犯蠢的事去的反感，仍然站得住。」}{二线：「就算这份厌有时看错、有时过头，那也只在个别场合，不是它平常的样子；该做的是调分寸，不是放下。」}{「就算错了也情有可原」}{第三节；第四节}
\canongroup{论点 · 两个论点与机制}
\canonrow{论1 教过了还犯（主张）}{「我方第一点：教过了还犯。常人的厌蠢，平常冲着教过了还犯：位置要他会的基本功，他被教过、提醒过，仍一再出错、不问不改。这样厌，看得有据，怪得公道。」}{「我们恼的，是教过还犯。」；「被教过、有时间、还不问不改的那一份。」}{「多数厌蠢冲着拖后腿的同事」；「被教过就可以厌」；「我们从不厌不会的人」（绝对化；改说：「对不会的，我们先教、给耐心；有人厌了，那一次是看错了，要调分寸。」）；「研究证明人们厌的都是不肯」}{第六节论点一；回扣「看准了吗」与归责小问}
\canonrow{论1·机制1 情绪按能不能控制分流}{「人对别人的失败，情绪本来就分流：判断他该负责，才生气；判断他控制不了，就同情。」}{「该负责就恼，控制不了就同情。」}{「生气永远是对的」；「这种判断从不出错」（分流靠判断，判断会错，正方会接 \#61、\#56，答：「所以要看记录，见论1·机制2」）}{Weiner 1993（\#58）；Weiner 1986、1995（\#179）}
\canonrow{论1·机制2 「本该」看位置，「本可以」看记录}{「本该，出自位置的公共要求：操作规范、入门培训、带教交代过的，不是我会什么。本可以，看他被教过没有、待了多久、教过之后改不改，这些查得到。新来的、没人教过的，先教；在问、在试、在改还学不会的，给耐心；只凭他一贯做不到，我们不厌。」}{「没有记录，先问；有记录，看记录。」；「分不清不会还是不肯时，先问他卡在哪。」}{「我会的他就该会」；「小刘转眼就把嘱咐抛诸脑后，这是可查的记录」；「本可以有实验证明」}{亚里士多德《尼各马可伦理学》III.5（\#41、\#230）；《孟子·梁惠王上》（\#231）；Frieze \& Weiner 1971（\#241）；新京报 2023-11-01（\#243）}
\canonrow{论1·机制3 「平常」只有一条主线：厌跟着代价走}{「厌跟着代价走。蠢的代价反复落在一起做事的人身上：同一个组、同一间办公室，他被教过没有、改没改，本来就在手里。评论区里对陌生人的那句比较，我们认：算厌蠢，也看错了，是个别场合。」}{「评论区那句最响，不是厌蠢最常住的地方。」；「对方的平常压在评论区，那里没有代价、没有关系，也没有记录。」}{「评论区那种不算厌蠢」；「新京报调查发现……」「现实里多数厌蠢是……」；其余各条另造「平常」的理由}{Hutcherson \& Gross 2011（\#224）；新京报 2023-11-01（\#252、\#243）；每日商报 2024-01-13（\#248）；Brady 等 2021（\#101）}
\canonrow{论2 恼他不弃他（主张）}{「我方第二点：恼他不弃他。厌蠢收回的，是对他这一次的评价，作为人的尊重不动；恼里那句『你本可以』，是对能负责的人提要求，讲理由它带着：这回放下，下回照样说、照样听。」}{「恼的是这件事，不弃这个人。」；「看低他这件事上的表现，不拒绝他作为人的尊重。」}{「恼还在，所以里面没有轻蔑」「轻蔑要三个前因凑齐才出现」；「恼就是尊重」；「我们从来不看低人」}{第六节论点二；回扣「对他公道吗」的尊重小问}
\canonrow{论2·机制1 尊重线（康德、SEP 整句引）}{「平等的尊严，不排除按品质、成就、功绩区别评价人，只是这种评价始终要受承认尊重的约束。康德说，比较之下在心里看低别人，有时免不了，但表现出来就是冒犯；拒绝给人作为人应得的尊重，无论如何违背义务。我们的厌，停在前一边。」}{「看低他这件事上的表现，可以；拒绝他作为人应得的尊重，不行。」}{只引康德前半句「心里看低有时免不了」；「康德说心里看低没关系」；只引 SEP 前半句；「这是李秋零译本的原文」}{SEP「Respect」（\#44）；Kant 1797，§39（\#48）；Roberts 2013 书评（引 Bell）（\#37）}
\canonrow{论2·机制2 恼是把他当能负责的人}{「Strawson 说，怨恨这类反应态度，回应的是对方意志里的善意、恶意，或者冷漠、缺乏关切，是把他当成要为自己负责的参与者；态度完全客观时，可以跟他谈，甚至跟他谈判，却不能跟他讲理。所以恼他不上心，是把他当能负责的人。恼里带着讲理，这一步是我们的推论。」}{「完全客观，才不能讲理；推不出没讲理就是完全客观。」；「客观态度装不下愤怒，混着恼的厌，到不了『完全』。」；「去人性化的，是不掺愤怒的纯粹轻蔑。」}{「一躲开就不能讲理」；「Strawson 说厌蠢就是反应态度」；「还在恼就一定还在讲理」不标推论；「Strawson 1962 年发表在《英国科学院院刊》第 48 卷」}{Strawson 1962（\#32、\#174、\#197）；Roberts 2013 书评（\#38）}
\canonrow{论2·机制3 越界按时间划}{「越界按时间判：这回先不劝、一时走开、不再主动提醒，不越线；他开口也不听、合作只下指令不给理由、认定他这个人说什么都没用，越线，我们认要调分寸，那是个别场合。Strawson 也说，正常人可以一时借客观态度躲开投入之累，但作为人，通常不能长久这样做。」}{「这回放下，下回照样说、照样听。」；「这件事上从此不听他，也算越线，我们认要调。」}{把「一时」换成「一事」；「从此不理他也没错」；「不再主动提醒就是不讲理」}{Strawson 1962（\#34、\#197）；Kant 1797，§39 附释（\#45）}
\canongroup{论点 · 全队共守的线、价值层与场上用语}
\canonrow{论3 记录线}{「没有记录，先问；新来的、没人教过的，先教；在问、在试、在改还学不会的，给耐心；只凭一贯做不到，不厌。教过、被指出、又做错这串记录，在看准了吗这一问上，双方都按它判。」}{「他学过没有，查得到，用不着估。」}{「只要他错过一次就有记录了」}{Frieze \& Weiner 1971（\#241）}
\canonrow{论4 分格问题（一辩稿抛出，后场反复追问）}{「这个例子里，他被教过没有？」}{对方答「没有」：「那是我们先问、先教的一种。」；对方答「有」：「那正是我们守的一种。」}{换成开放式问法「你怎么知道他是不会还是不肯」}{第九节}
\canonrow{论5 条件模块「换不掉的那点反感」（只在正方摆出替代态度后用）}{「对方的替代，一半本来就在我们这份厌里：没有记录先问、对事批评、回头还讲；另一半管不到我们恼的那一类：纠正治的是不知道，治不了教过还不上心。存在更好的，不等于这个就错了。」}{「先开口问，本来就在我们的厌里。」；「对一个共事两年、劝过多次的同事，先开口问，问什么？」}{主动开「划得来吗」的进攻；「厌蠢守住了社会标准」}{反驳 2、22；Roberts 2013 书评（\#160）；Wood \& Porter 2019（\#109）}
\canonrow{论6 主动让掉的}{「评论区里对陌生人的比较：算厌蠢，看错了，个别场合。第一次、没人教的不会：先教、先问。越线的行为：认要调分寸。专业系统里的失误：那是流程和职业要求的事。把一个群体的处境算成活该：我们和对方一样反对。」}{「这几样我们不争，争的是熟人之间、教过还犯的那一份。」}{「这些都不算厌蠢」}{第三节；第九节}
\canonrow{论7 对方已让出的（反三小结清单；只在正方台上确实认了才说）}{「对方已经认了：凭记录的，在看准了吗上不争；走神、手滑这类，不争看准；对那件事本身的恼火，不反对；这回先不说、一时走开，不算越线。剩下的，只有评论区和一则没写结尾的报道。」}{「对方的战场，只剩没有记录的比较，和回头不讲理。」}{正方没认的也说成认了}{第四节；第五节论点一；第九节}
\canonrow{论8 后场追问（奇袭链首问）}{「你们说的平常，除了评论区和一则没写结尾的报道，还有什么？」「回头不讲理，有直接材料吗，还是推论？」「对一个共事两年、劝过多次的同事，先开口问，问什么？」}{备用：「恼火你们不反对；去掉看低之后剩下的那份，错没错？」}{问题超过 25 字；开放式提问}{第五节论点二；第九节}
\canonrow{论9 价值层（结辩固定段）}{「我没错，不是我比你强，是我把你当成能负责的人：你本可以，所以我恼。在意标准，也在意一起做事的人；为在意这件事，我不必认错。」}{「你本可以的恼，是一种看得起。」；「对教过还犯的人只给怜悯、只管不问，反倒是不再跟他讲理。」（这一句是推论，照说推论）}{「厌蠢是觉醒」「厌蠢推动社会进步」；任何精英腔}{第九节；Strawson 1962（\#32、\#33）；Graham 1984（\#250）}
\canonrow{论10 不上场的论点（禁止栏为主）}{「我们不说情绪无对错，不说我有权远离，不说厌蠢让人进步，不说厌蠢守住了社会标准。」}{被对方安上这些立场时：「这几句我们一句都没说。」}{「情绪控制不了，所以无所谓对错」；「我有权远离谁」「远离蠢人是我的权利」；「被厌的人表现会提升」；「对搭便车者的愤怒维持合作」；「愚蠢比恶更危险」；「根在效率社会、优绩主义」；「我们拿要求自己的标准要求别人」作主论证}{第六节反方被淘汰的候选论点}
\canongroup{论点 · 学理与典籍（论点里引用的原话）}
\canonrow{学1 罪与病（Weiner）}{「Weiner 1993 年的归因理论概括：因为不努力而失败，比因为能力不足而失败，招来更多惩罚。」}{「缺努力招来恼火，缺能力招来同情。」}{「Weiner 证明这样分是对的」}{Weiner 1993（\#58）}
\canonrow{学2 怨恨与轻蔑（Bell 经 Roberts）}{「Roberts 概括 Bell 的区分：怨恨针对行为、促使交涉；轻蔑针对人、促使退出。我们的恼冲着这件事、回头还讲，在怨恨这一格。」}{「冲着事、促使交涉的，是批评；冲着人、促使退出的，才是轻蔑。」}{「Bell 为厌蠢辩护」}{Roberts 2013 书评（\#160）；Bell 2013（\#35）}
\canonrow{学3 亚里士多德：盲人与醉酒致盲}{「亚里士多德说，没人责备天生的盲人，人人责备醉酒致盲的人；对应当知道、又不难知道的事无知，也要担，因为他本可以留意。」}{「天生的不怪，自己招来的才怪。」}{「亚里士多德说可以厌蠢人」；报 Bekker 页码}{亚里士多德《尼各马可伦理学》III.5（\#41、\#230）}
\canonrow{学4 孟子：不为与不能}{「孟子说，为长者折枝，跟人说我不能，是不为也，非不能也。」}{「挟泰山超北海是真不能；折一根树枝说不能，是不肯。」}{「孟子说的是日常的蠢」}{《孟子·梁惠王上》（\#231）}
\canonrow{学5 孔子：困而不学}{「孔子说，困而不学，民斯为下矣：困不怪，不学才怪。」}{「困了不学，才排在最下等。」}{「孔子说要厌不学的人」}{《论语·季氏》（\#49）}
\canonrow{学6 子张：嘉善而矜不能（对方引「人将拒我」时用）}{「同一章里，子张说的是嘉善而矜不能：矜的是不能。对不会的，我们给耐心；我们不拒人。」}{「子张也没说不分不能与不肯。」}{「《论语》裁决了子夏与子张谁对」}{《论语·子张》（\#50、\#182）}
\canonrow{学7 亚里士多德：该怒而不怒}{「亚里士多德说，对该生气的人、在该生气的时候生气，受称赞；该生气而不生气，被看作愚蠢。」}{「该恼的不恼，在亚里士多德看来是愚蠢。」}{「亚里士多德主张厌蠢」；报译者名}{亚里士多德《尼各马可伦理学》IV.5（\#39）}
\canonrow{学8 中节}{「《中庸》说，发而皆中节，谓之和；亚里士多德说，在恰当的时候、对恰当的对象、以恰当的方式去感受，才是德性。中节的厌：对象是本该会、被教过、还明显做错的事；方式是冲着这件事，回头还讲理；时候是这回放下、下回再来。」}{「情绪可以评，所以我们才问完三句。」}{「中节是我们这份厌天生就有的」}{《礼记·中庸》（\#158）；亚里士多德《尼各马可伦理学》II.6（\#152）}
\canonrow{学9 亚里士多德：多数人的轻视是随意的（对方引时）}{「整句是：骄傲的人轻视得正当，因为他想得对；多数人的轻视则是随意的。分界在想得对不对，不在多数少数；我们的厌，根据是记录：被教过没有、改没改。」}{「那句讲的是轻视人，我们不为轻视人辩护。」；「随意，是说没根据。」}{「亚里士多德说可以轻视人」「我们就是大度之人」；「亚里士多德说多数人看不准」；拿前半句立论}{反驳 30；亚里士多德《尼各马可伦理学》IV.3（\#162）}
\canonrow{学10 康德的尊重义务（据 SEP，对方引时；两句一起念）}{「据《斯坦福哲学百科》的概括，康德说我们没有感到尊重的义务，这些义务同时包括不持有某些态度，比如不轻蔑、不嘲笑。百科说的轻蔑，是否认他有任何价值；我们看低的是他这件事上的表现，不否认他作为人的价值。」}{「心里那份也要称，我们不躲。」；「康德自己分了两件事：比较之下的看低有时免不了，拒绝基本尊重无论如何违背义务。」}{只引「没有感到尊重的义务」；「康德说心里怎么想都行」；把「想自我感觉良好的欲望是不尊重的主因，『我比你聪明』就是这种欲望」说成百科原话}{反驳 33；SEP「Respect」（\#43）}
\canongroup{论点 · 对方最可能打的地方，我方怎么回（二辩、四辩、自由辩共用）}
\canonrow{回1 打论1：「新京报读全，同学同事那段就是『你能不能有点常识』『我都可以做到』」}{「这几句我们全认。只凭『我都可以』、不看他学过没有的，是该先问的场合。新京报那段我们只取场景，记录由每日商报那一案补。」}{「比较的口气说的是本该；没说出口的是他被教过没有，熟人之间那一步用不着说，手里有。」（推论）}{「新京报那段写了当事人被教过」；「同学同事那段不是比较」}{第六节论点一；反驳 1；新京报 2023-11-01（\#212、\#243）}
\canonrow{回2 打论1：「教爸妈付款，教三遍还按错，是不会还是不肯？」}{「看教过之后问不问、改不改：在问、在改还学不会的，给耐心；分不清时，先问他卡在哪。」}{「教三遍还在问，是不会；教三遍不问也不改，才是我们恼的那一种。」}{「教三遍还按错就是不肯」；「老人学不会就该被厌」}{第六节论点一}
\canonrow{回3 打论1：「会了的人估不准、自己看不见、提醒压不下，所以要改成先开口问」}{「这些实验估的，都是一个你一无所知的人。被教过、劝过多次的同事不是不知情者，他学过没有，查得到，用不着估。对方的药，先开口问，只治得到没有记录的那一种；没有记录的，我们也认该先问。」}{「偏差长在对他一无所知的地方。」}{「后见之明研究是错的」；「我们从不估别人」；「偏差很小」}{反驳 2；Fischhoff 1975（\#204）；Hom 2023（\#176）；Birch 等 2017（\#215）；Harley 等 2004（\#205）}
\canonrow{回4 打论2：「轻蔑的预测因素在你们的场景里现成：劝了还犯是管不了，『还是那么蠢』是品性归因」}{「预测因素不是必要条件。劝了不听，这回放下、下回照样说，就是还觉得话管得到他；冲着失败生气，对方读出的是没用心，没用心是下回改得了的事。真硬化成『他这个人说什么都没用』，我们认越线、要调分寸，那是个别场合。」}{「熟人之间亲密在，他别处做对的记录，逆着品性归因。」（推论）}{「恼还在，所以里面没有轻蔑」「轻蔑要三个前因凑齐才出现」}{第六节论点二；反驳 3；Fischer \& Roseman 2007（\#185、\#86）；Graham 1984（\#250）；Miceli \& Castelfranchi 2018（\#27）}
\canonrow{回5 打论2：「讲理挂在他的表现上；你们把『一时』换成了『一事』」}{「讲理不等他做对：子路最终没领会，孔子照样把道理讲完。这回先不劝，下回照说；这件事上从此不听他，我们认越线。」}{「讲理不是奖励，是我们本来就带着的。」}{「这件事上从此不听也不算越线」}{第六节论点二；《论语·子路》；朱熹集注（\#226）；Kant 1797，§39 附释（\#45）}
\canonrow{回6 打论2：每日商报一案「劝了多次，关系降至冰点，接着就是『说了那么多遍，怎么还是那么蠢』」}{「三句一起读：多次劝告，讲理没等他改好；冰点是一度，是降到过；现在是很微妙，不是断了往来。那句话没标说话人；之后讲不讲理，报道没写，对方自己标了推论。」}{被追问「怎么还是那么蠢，宾语是人」：「『怎么还是』还在要求他（推论）；蠢字落在人身上，下调的是对他这一方面的评价，平等尊严不排除区别评价。」}{「小陈后来照样教他」「两人早就和好了」；「这句是小陈说的」「这句是记者编的」}{反驳 3；每日商报 2024-01-13（\#248）；SEP「Respect」（\#44）}
\canongroup{数据 · 我方立论数据}
\canonrow{数1 归因分流元分析（论1 主数据）}{「2004 年，一项汇总六十四项研究、一万二千多名被试的元分析发现：判断对方该负责，人就生气；判断他控制不了，人就同情；不同文化里结果高度一致。」}{「一项汇总六十四项研究的元分析发现，该负责就恼，控制不了就同情。」（数字要说就照原样说，不改约数）}{「一万二千人证明厌蠢没错」；「研究对象就是蠢人」；报效应量；「生气就该攻击」}{Rudolph 等 2004（\#219）；64 项研究，12,000 余名被试}
\canonrow{数2 愤怒与轻蔑的社会功能（论2 数据；正方也引，整句认）}{「2007 年的三项研究发现：生气当下冲上去，长远会和好；轻蔑才把人推出去，短期长期都是。」}{被追问时整句认：「这项研究也说，愤怒和轻蔑常常一起出现，轻蔑可能从先前的生气里长出来；可能，不是平常的走向，这些是预测因素，不是必要条件。」}{「生气不会变成轻蔑」；「生气促进关系」「生气导致和解」；报样本量、效应量；「轻蔑常常从生气里长出来」}{Fischer \& Roseman 2007（\#86、\#170、\#185）；3 项研究；摘要原句 “anger and contempt often occur together”、“contempt may develop out of previously experienced anger”}
\canonrow{数3 冲着失败生气，对方读出「没用心」}{「1984 年一项实验里，孩子做砸之后，大人冲他生气，孩子更多想的是『我没用心』；大人表示同情，孩子反倒更多觉得『我能力不行』。」}{「冲着失败生气，对方读出的更多是没用心；没用心，是下回改得了的事。」}{「被生气反而更有信心」（摘要只说同情之后期待下降得更多）；「生气让人进步」；「成年人也一样」不标推论；报效应量}{Graham 1984（\#250）}
\canonrow{数4 与自己相关，推高的是愤怒（整句引）}{「2011 年的一组研究发现：事情与自己相关，推高的是愤怒，不推高别的情绪；轻蔑，在评判无能的行为时起作用。」}{「轻蔑起作用的对象，摘要写的是无能的行为。」}{只引前半句；「研究发现愤怒和轻蔑分不开」；「厌蠢里没有轻蔑」}{Hutcherson \& Gross 2011（\#224、\#26）}
\canongroup{数据 · 反驳库用到的研究}
\canonrow{数5 网上的愤怒被点赞放大}{「2021 年一项对推特的研究发现：网上的道德愤怒得到点赞，以后就更常这样表达。评论区里，是被放大过的声音。」}{「被点赞的表达会被放大；好批、好转，不等于常见。」（后半句是推论）}{「评论区的厌蠢都是表演」「网上的厌蠢是点赞刷出来的」；「微博上也一样」}{Brady 等 2021（\#101）；7,331 名用户、1,270 万条推文，另有行为实验 N=240}
\canonrow{数6 求助的人被看作更有能力}{「2015 年的研究发现：人怕求助显得无能，这种担心放错了地方；求助的人反而被看作更有能力，任务难的时候尤其如此。」}{「开口问，不会被看低。」}{「问简单的问题也加分」（摘要只说任务难时更明显；容易任务见数7）；报样本量}{Brooks 等 2015（\#227）}
\canonrow{数7 任务容易时，求助不损害能力评价（报道）}{「据《MIT 斯隆管理评论》2015 年的报道，任务被看作容易时，求助既不损害、也不提升别人对他能力的看法。」}{「问容易的事，也不会被看低。」}{「论文发现……」}{MIT 斯隆管理评论 2015（\#181）}
\canonrow{数8 心理安全的原文定义}{「Edmondson 1999 年研究一家制造企业的五十一个工作团队，她给心理安全下的定义是：相信团队不会因为某人开口而羞辱、排斥或惩罚他；她特意说，这不是随意的放任。」}{「怕显得无能、不愿认错，Edmondson 说是组织里常见的倾向，没有厌蠢者也在。」；「对方没有一条数据测过厌蠢。」}{「心理安全决定三成学习」；报回归系数；「量表是反向计分」}{Edmondson 1999（\#222、\#184）；51 个工作团队、个体问卷 N=427；定义句在 p. 354；量表第 1 题 “If you make a mistake on this team, it is often held against you.”（p. 382）}
\canonrow{数9 给了归因信息，情绪跟着变}{「1988 年的实验发现：给了归因信息，人的情绪和帮不帮的判断就跟着变。」}{「问到了，情绪就跟着变。」；被追问时整句认：「同一研究也说，心理和行为一类的污名，默认被看作可控；有些归因改不动。」}{举具体病名；「人天生会把蠢看成不肯」}{Weiner 等 1988（\#188）}
\canonrow{数10 一贯做不到，旁人读出的是能力与难度}{「1971 年的实验发现：观察者分原因，会看这个人自己过去做得怎样、别人做得怎样；和他自己过去一致的结果，读成能力和任务难度。所以只凭他一贯做不到，我们不厌。」}{「一贯做不到，读出来的多是不会，不是不肯。」}{「人对别人的失败特别宽厚」；把这里的「能力」说成智力}{Frieze \& Weiner 1971（\#241）}
\canonrow{数11 说的人高估自己说清楚了，旁听者没有}{「2002 年的实验发现：说话的人总高估自己说清楚了；可知道本意的旁听者，没有系统性地高估，偏差出在说话的人评估自己的话时。所以判断不靠『我说过了』，靠他之后改没改。」}{「被指出、又做错，『以为懂了』已经被那次失败推翻；分不清时，问他卡在哪。」}{「研究证明我们教得清楚」；只引前半句}{Keysar \& Henly 2002（\#234、\#191）}
\canonrow{数12 经常接触的人，坏印象更容易改（整句认）}{「2005 年的一篇模型推导写着：坏经历让人少来往，所以坏印象更稳定；可对经常接触的人，坏的初始印象更容易改。厌蠢的典型对象，正是天天见面的同事。」}{「接触在，信息就在。」；「这个机制里没有情绪，只有当前印象；对方要定罪的不是厌，是有坏印象本身。」（后半是据摘要措辞的推论）}{「研究发现同事之间坏印象都会改」；「接触回来，讲理就回来了」}{Denrell 2005（\#193、\#235）}
\canonrow{数13 「偏见」的定义前提是看错了}{「2009 年那项计算机模拟里，『偏见』的定义，就是对其实是好的一类对象抱着错误的负面预期：它咬人的前提是看错了。判对了，少来往不丢任何东西。」}{「锁死的前提是判错。」}{「实验发现人会……」}{Eiser 等 2009（\#236）}
\canonrow{数14 人们把什么叫作蠢}{「2015 年一项研究收集人们把哪些行为叫作蠢的故事，归出三类：信心与能力失衡、注意失误、缺乏控制；行为人被认为责任越大、后果越重，就显得越蠢。」}{「人们叫作蠢的，是具体的行为，而且跟着责任和后果走。」；「不懂又自以为懂，就落在第一类。」（这是我们举的例子）}{报三类的比例、故事数、评分人数；说出「缺乏控制」具体指什么；「研究证明这些人真的蠢」}{Aczél 等 2015（\#93、\#239）}
\canonrow{数15 压住感受几乎没用，换位效果最大}{「2012 年一项比较了三百多个实验的元分析发现：压住感受本身几乎没有效果；在各种调节办法里，换位思考效果最大。」}{「我们不压，也没人要我们压。」；对方问「为什么不换位」：「对没有记录的人先问，就是在换位。」}{念 d 值（场上说「几乎没有效果」「效果最大」，精确值在速查）；「研究证明情绪控制不了」；「换位没用」}{Webb 等 2012（\#155、\#177）；306 个实验比较；压住感受 d(+) = −0.04，换位思考 d(+) = 0.45}
\canonrow{数16 婚姻研究里的轻蔑（对方引时）}{「那项婚姻研究追了七十九对夫妻十四年，测的是夫妻冲突讨论里表达出来的轻蔑；轻蔑只是几个负面预测项之一。表达出来的看低，我们本来就不辩护。」}{「那是对伴侣表达的轻蔑，不是同事填错表时的那一下恼。」；对方说「轻蔑是离婚最强预测因子」：「那是 Gottman 研究所博客的说法；论文里防御性的效应比轻蔑大。」}{「Gottman 2002 年研究发现轻蔑是离婚最佳预测因子」；报准确率}{Gottman \& Levenson 2000（\#84）；Gottman 研究所博客 2013（\#85）；79 对夫妻，1983 年起追踪 14 年，22 对离婚}
\canonrow{数17 愤怒、敌意与冠心病（对方引时）}{「冠心病那项元分析测的是特质性的愤怒与敌意，是相关，不是因果；不是同事教过还犯时冲着那件事的恼。」}{「普通人不会因为生气伤身，就说该有的义愤是错的。」}{「效应很小，所以无所谓」；「生气不伤身」}{Chida \& Steptoe 2009（\#97、\#63）；健康人群 25 项研究，合并 HR 1.19（95\% CI 1.05–1.35）}
\canonrow{数18 系统取向也要分流（Reason，整句引）}{「Reason 2000 年在《英国医学杂志》上写：报告文化靠信任，信任要求公正文化，也就是大家对无可指责与应受指责的界线有共识；他也写，有些不安全行为确实恶劣，但绝大多数不是。系统取向要的是划线，不是什么都不怪。」}{「没教过的不怪，教过还犯的那一份怪：这就是我们划的线。」}{只引「有些确实恶劣」前半句；「Reason 主张追责个人」；报「九成质量差错被判为无可指责」}{Reason 2000（\#223、\#169、\#183）}
\canonrow{数19 轻蔑把人看成愚蠢、无能（推不出反过来）}{「综述说，轻蔑会把对方看成愚蠢、无能，贬的是品性多于行为；推不出反过来：把一件事看成蠢，就是轻蔑。」}{「轻蔑里有蠢，不等于对蠢的反感都是轻蔑。」}{「轻蔑和蠢没关系」}{Fischer \& Giner-Sorolla 2016（\#20）}
\canonrow{数20 轻蔑不看好对方改善}{「理论分析里，轻蔑的特点之一，是不看好对方改善。恼里那句『你本可以』，方向正好相反。」}{「还在说『你本可以』，就是还看好他改。」（推论）}{「研究证明恼的人都看好对方」}{Miceli \& Castelfranchi 2018（\#27）}
\canonrow{数21 不会但态度好的，对应怜悯}{「刻板印象内容模型认为，不会但态度好的，对应怜悯，不对应轻蔑。」}{「对真不会的，被读出不会本来就是准的；耐心是先教，不是摆出同情。」}{「研究发现……」（摘要里这是模型的假设，只说「该模型认为」）}{Fiske 等 2002（\#29）}
\canonrow{数22 学历偏见研究（对方引时）}{「那组研究测的是对『学历』这个群体身份的态度，西方样本，不是对一个本该会却不改的人的反应；其中研究六只报了愤怒由责任归因中介，没有报轻蔑的中介。」}{「对一个群体做责任归因，才是错配；我们归的是一件被教过、又做错的事。」}{报回归系数；「Kuppens 证明高学历者不厌蠢」；「我们这些高学历的人」}{Kuppens 等 2018（\#54、\#55、\#56、\#171）}
\canonrow{数23 事实纠正有效（对方引时）}{「美国政治议题上的几个实验发现，事实纠正有效，没有引发逆火。纠正治的是不知道，我们本来先教；它治不了教过还不上心。」}{「纠正有效，正说明先教是对的。」}{「纠正没用」}{Wood \& Porter 2019（\#109）；5 个实验、10,100 余名被试、52 个议题}
\canonrow{数24 老人与智能手机（对方引「子女以为父母学不会」时）}{「同一篇报道写着：受访老人里多数在用智能手机；『太贵』那组回答，只来自少数不用的老人；两代人都说难在操作太复杂；年轻人表示愿意主动教父母。读偏了，生出来的是教，不是厌。」}{「两代人说的难处是同一个：操作太复杂。」}{报任何百分数；「九成子女看错了父母」；说成厌蠢比例}{川观新闻 2020-10-21（\#175、\#238、\#192）}
\canonrow{数25 后见之明元分析（备用，只在对方夸大偏差时用）}{「2004 年一项汇总九十五项研究的元分析，算出后见之明偏差的平均效应，约零点四个标准差。我们不评它大小，只说：它说明有偏差，没说明这份厌平常判错。」}{对方质疑我方虚报时报全：「零点三九，置信区间零点三六到零点四二。」}{「偏得不多」「才零点四」「效应很小」「说明厌蠢大多看得准」；换算成「百分之几的人」；不进一辩稿与立论；引它时要预判对方补出：常识问答题效应更大；减偏操纵没让效应变小}{Guilbault 等 2004（\#246）；95 项研究、252 个独立效应量，Md = .39，95\% CI .36–.42}
\canonrow{数26 刑法看「应当预见」（对方说「法律也是事后判疏忽」时）}{「刑法第十五条写的是：应当预见自己的行为可能发生危害社会的结果，因为疏忽大意而没有预见。法律看的是事前该不该预见，不是事后的结果。」}{「我们只借这条判别线，不借它的严厉。」}{漏掉同条「过失犯罪，法律有规定的才负刑事责任」（被问到要认）；「法律规定可以厌蠢」}{《刑法》第十五条（\#180）}
\canongroup{数据 · 对方研究：我方统一说法}
\canonrow{数27 后见之明经典实验（Fischhoff）}{「Fischhoff 的被试，估的是不知道结果的『别人』：一个他一无所知的人。」}{「这些实验估的，都是一个你一无所知的人。」}{「Fischhoff 的研究是假的」；报样本量、描述实验材料}{Fischhoff 1975（\#204）}
\canonrow{数28 替陌生同学估字谜（Hom）}{「Hom 的被试替一个不认识的同学估，还被要求设想他和自己经历相同：这条指示本身就推着人往自己身上投射；何况对方自己的研究说，设身处地替同伴想，偏差照旧。」}{「它模拟的是对陌生人的默认；对一个共事两年的人，我们用不着默认。」}{「Hom 的结果不成立」；「被试很少」}{Hom 2023（\#176、\#217）}
\canonrow{数29 知识有多普遍（Birch）}{「Birch 估的是：一道知识题，同龄人里有多少人知道。」}{「位置上的基本功是写得出、教得出的，不是我会什么；教过、交代过的，是记录，用不着估。」}{「Birch 说人总是高估」}{Birch 等 2017（\#215）}
\canonrow{数30 以己度人（Nickerson，整句）}{「Nickerson 同一句前半说，拿自己估别人，总体上对人有用；出错在不加批判。手里有他被教过、又做错的记录，就是批判的依据。」}{「起点是自己，不等于终点是自己。」}{「Nickerson 说以己度人是对的」}{Nickerson 1999（\#216）}
\canonrow{数31 认模糊名人脸（Harley）}{「Harley 的任务，是辨认由模糊变清晰的名人面孔，是知觉，不是做事。」}{「认脸不是填表。」}{「提醒压不下，说明厌蠢改不了」}{Harley 等 2004（\#205）}
\canonrow{数32 防洪过失认定（Kamin 与 Rachlinski）}{「他们判的是一场洪水当初该不该防：事前本来就是不确定的风险，又是法律上的过失认定。常人厌蠢平常判的，是教过就会、不需要预测的基本功。」}{「事后诸葛只管想得到，管不了做得对。」}{报「四分之三」「多数」}{Kamin \& Rachlinski 1995（\#206）}
\canonrow{数33 后见之明综述（Roese 与 Vohs）}{「它管的是结果的可预见性，不管一个人被教过之后改没改。」}{「记录在结果之前，就不是事后诸葛。」}{「综述说后见之明压不下」}{Roese \& Vohs 2012（\#207）}
\canonrow{数34 经理信不信人能改（Heslin）}{「那项研究测的是经理对人性能不能改的一般看法，不是厌，也不是对一个人的判断；辅导下属是经理的职务；也没测听不听他的理由。它反倒替我们说话：决定教不教的，是『人能不能改』，恼里的『你本可以』，站在相信他改得了的那一边。」}{「它的变量是能不能改，恼说的正是你能。」（推论）}{报样本量、效应量；「经理研究证明厌蠢者会继续辅导」}{Heslin 等 2006（\#218）}
\canonrow{数35 觉得对方有偏、冲突螺旋（Kennedy 与 Pronin；Pronin 等）}{「这组研究的自变量是觉得对方有偏，管的是意见分歧，不问有没有厌。同一处基本功错了又错，没有立场之争。」}{「意见分歧不是犯蠢。」}{「朴素实在论是 Ross 与 Ward 1996 年提出的」；描述「冲突升级的做法」}{Kennedy \& Pronin 2008（\#208）；Pronin 等 2004（\#209）}
\canonrow{数36 〔智〕材料（对方引时的统一回应）}{「那条测的是智力。双方共用的中立定义写明：蠢不等于智商低。」}{「两项集体智力研究结论相反，学界有争议，谁也不该拿来定论。」（\#94 对 \#95）}{我方主动引「陌生人看录像能判断智力，所以厌蠢看得准」；「智力八成由基因决定」；拿遗传、教育年限讲蠢}{Borkenau \& Liebler 1993（\#88）；Borkenau 等 2004（\#89）；Heck 等 2018（\#91）；Freund \& Kasten 2012（\#92）；Bates \& Gupta 2017（\#94）；Riedl 等 2021（\#95）；Neisser 等 1996（\#80）；Turkheimer 等 2003（\#81）；Ritchie \& Tucker-Drob 2018（\#82）；Plomin \& Deary 2015（\#57）；Stanovich 2009（\#70）；Bonhoeffer〈论愚蠢〉（\#75）}
\canonrow{数37 职场新人调查（对方报百分比时）}{「那是新人遇到尴尬的情形，不是厌蠢的比例；不成文规定没人教过，正是我们先教的一类。」}{「那项调查没有测厌蠢，也不是随机样本。」；对方说新人头号办法是请教：「那篇报道里闷在工位上的新人，说的原因是同事忙，不是怕被厌。」}{我方复述任何百分数}{中国青年报 2023-05-24（\#121、\#187、\#173）}
\canonrow{数38 Adams「不自愿的罪」（对方引时）}{「Adams 举的是无端的愤怒、自以为是的轻蔑，是哲学论证，不是实证；不自愿也可以被责备，是一派立场，不是共识。我们的厌有端可循，记录在手。」}{「我们管住表现，守住尊重。」}{「Adams 是错的」}{Adams 1985（\#156）}
\canonrow{数39 对应偏差（Gilbert 与 Malone，对方引时）}{「对应偏差，是从完全可以由情境解释的行为，推出一个人持久的特质；摘要列的第一个机制，就是没觉察到情境的约束。它咬人的前提，是处境没被看见；教过、被指出、又做错，处境已经摆出来一回，我们看的是这串记录。」}{「所以要看记录；分不清时，先问他卡在哪。」；「我们下调的是他这件事上的评价；推成『他这个人就是蠢』，那一步越线了，要调分寸。」}{「对应偏差是西方现象」「中国人不这样推」；「对应偏差是错的」；「屡次就一定不是处境」}{反驳 26；Gilbert \& Malone 1995（\#61）}
\canonrow{数40 信特质固定的人推得更狠（Chiu、Hong 与 Dweck，对方引时）}{「那项研究说的是人和人不一样：相信特质固定的人，从行为里推出的特质判断更强；推得狠不狠因人而异。恼里那句『你本可以』，站的是相信他改得了的那一边。」（后半是推论）}{「它是调节变量，打不了头阵。」}{「我们都是渐进论者」；「研究的是大学生」；「研究证明了因果」}{反驳 26、41；Chiu, Hong \& Dweck 1997（\#256）}
\canonrow{数41 知识的诅咒（Birch 与 Bloom，对方引时）}{「知识的诅咒测的是错误信念推理：自己知道事情的结果，再推断别人对这件事以为什么，会被带偏。对一个教过三遍的同事，我们用不着猜他以为什么，他学过没有，记录在手。」}{「它解释得了把信息差当成蠢，那一种我们认是看错了；解释不了教过还犯。」；「知识的诅咒，正是我们没记录先问、新来的先教的理由。」}{「知识的诅咒是假的」；描述实验材料；和 Birch 2017混成一篇}{反驳 31；Birch \& Bloom 2007（\#253）}
\canonrow{数42 问答者实验（Ross、Amabile 与 Steinmetz，对方引时）}{「出题人是随机分的，题是自己编的，可以拿自己知道的冷门知识来出；测的是谁显得更有学识。我们说的本该会，是位置上教过、交代过的基本功，不是我出的题。」}{「校正角色偏差，就是把角色说破：题是谁出的，他学过没有。」；被说「答题者自己都贬低自己」时：「同一篇摘要写明，偏差在答题者和旁观者身上更强。」（这一句在该条方法栏，被要求出示时指方法栏）}{「研究说出题人都瞧不起答题人」；报样本量；「这项研究测的是蠢」}{反驳 32；Ross, Amabile \& Steinmetz 1977（\#254）}
\canonrow{数43 专家之咒（Hinds，对方引时）}{「Hinds 测的是内行预测新手完成一项任务要多久；同一句后半还说，中等水平的人预测新手反而可能更准。没人教过的新手，我们先教，不拿内行的手感判他。」}{「压不下的是估计；有记录时，我们不靠估计。」}{「同学同事都是中等水平」；「内行判断都不准」；报样本量、描述任务内容}{反驳 36；Hinds 1999（\#203）}
\canongroup{案例}
\canonrow{案1 每日商报小陈与小刘（论1 主案例；三句一起说）}{「每日商报 2024 年 1 月写过一位程序员：同办公室的同事工作超过两年，还常把表格数据填错、资料不分类。他多次劝告，同样的错接二连三再犯，两人的关系一度降至冰点，现在很微妙。」}{正方引「说了那么多遍，怎么还是那么蠢」时：「报道紧接着这个案例写了这一句，没有标说话人。」；「劝了多次，讲理没等他改好；冰点是一度，现在是很微妙，不是断了往来。」}{「小陈说：说了那么多遍……」；「小刘就是无所谓」「小刘转眼就把嘱咐忘了」当事实说（改说：「小陈一方说他一副无所谓的样子」）；「这是事后可查的记录」；「职场里厌蠢多半是这种」「报道说厌蠢在职场用得极为频繁」；「小陈后来照样教他」「两人和好了」；只说「多次劝告」，不说「一度降至冰点」「很微妙」；引医生「大部分人讨厌明知故犯」}{每日商报 2024-01-13（\#248）}
\canonrow{案2 孔子与子路论正名（论2 主案例）}{「《论语·子路》里，子路当面说老师迂，孔子先斥一句『野哉，由也』，接着把正名的道理一层层讲完。朱熹的集注里，胡氏说，孔子讲得这样详细，子路最终还是没领会。恼和讲理在同一个当场，讲理也不等他做对。」}{「孔子恼了，也把道理讲完了。」；「子路终不喻，孔子也没把他划出门外。」（后半是推论）}{「野只是对事不对人」；「孔子从不背后定性人」（同卷下一章樊迟出去后孔子说「小人哉」）；「圣人这样，所以常人平常也这样」；「讲完之后子路就懂了」}{《论语·子路》；朱熹集注（\#226）}
\canonrow{案3 莎莉文与海伦}{「莎莉文教海伦分清杯子和水，教到绝望，暂时放下这个话题，又在下一个机会重新开始。有绝望，有放下，有再来。」}{「这回放下，下回再来。」}{「海伦是蠢」；「莎莉文从不发火、一味温和」；「这证明厌蠢没错」}{Keller 1903（\#142）}
\canonrow{案4 新京报帕孜丽娅一文（只说「新京报的文章写道」）}{「新京报的文章写完评论区，下一段就是『当然，在现实生活中』，写的是学习、工作中拖后腿的同学、同事；文章说，现实里说的蠢，是『缺乏所处位置需要的基本知识和技能』。同一篇还写，人会想到『生活中有这样一个人要怎么办』，对蠢的厌恶更多源自它可能造成的后果。」}{「这几句我们全认：只凭『我都可以』、不看他学过没有的，是该先问的场合。新京报那段我们只取场景，记录由每日商报那一案补。」}{「新京报认为厌蠢没错」；「新京报调查发现……」「多数人……」；只引前半（同篇后文「合理化贬低他人」「恐弱」要认）；「新京报是央媒」；拿「后果」一句定义厌蠢}{新京报 2023-11-01（\#243、\#252、\#212、\#213）；新京报书评周刊 2023-11-01（\#12、\#13）}
\canonrow{案5 上游新闻里的网友}{「上游新闻引的一位早期自称厌蠢症的网友说：遇见不会但努力的人，并不会激发她的厌蠢症，她甚至可以耐着性子教对方。」}{「触发厌的，不是开口问。」}{「厌蠢的人都这样」；「网友说远离蠢人是自己的权利」；「我有权远离」}{上游新闻 2023-02-24（\#168、\#9）}
\canonrow{案6 这个词的出生与出圈}{「上游新闻写着，这个词元旦前后就出现了；小张事件是它出圈，不是它出生；让它彻底出圈的，是央视网的报道。」}{「出圈靠的是好批，好批不等于常见。」（推论）}{「评论区那批人都不算厌蠢」；「这个词是小张事件造出来的」}{上游新闻 2023-02-24（\#8、\#167）}
\canonrow{案7 Torvalds 道歉（对方引时）}{「2018 年 Torvalds 道歉的，是邮件里轻率的攻击，尤其是人身攻击。内核社区随后的准则写着：贡献会被仔细审查、常招来批评，拒绝不合格的贡献不算违规；禁止的是辱骂和人身攻击。骂人去掉，标准留下，这叫调分寸。」}{「他说要学会理解别人的情绪、恰当回应，改的是怎么回应。」}{「Torvalds 至今认为骂人说得通」；「他是被迫道歉的」；「Torvalds 证明厌蠢没错」}{Torvalds 2018（\#139、\#186、\#172）；Linux 内核行为准则 2018（\#221）}
\canonrow{案8 朋霍费尔（对方引「蔑视人？」时）}{「朋霍费尔 1942 年底写的〈论愚蠢〉说，愚蠢不是智力上的缺陷，而是人的缺陷，人会任由自己被弄蠢。他下一节警告的是蔑视人：我们无权这样。我们同意，不蔑视人正是我们守的线；他说蔑视一个人就永远不能使他有所成就，而恼还在指望他成。」}{被问「他说再也不说服蠢人」：「那是政权之下对群体愚蠢的结论，期限是外部解放；那一步我们不学，我们回头还讲。」}{「狱中写下」；「道德缺陷」；「朋霍费尔支持厌蠢」；「愚蠢比恶更危险」作论据}{Bonhoeffer〈论愚蠢〉（\#74、\#75、\#247）；Bonhoeffer〈蔑视人？〉（\#76、\#211）；群学书院（中文流传本）2021（\#133）}
\canonrow{案9 何不食肉糜（对方引时）}{「惠帝错在不知道饥民的处境；我们没记录先问、新人先教，问的正是处境。」}{「何不食肉糜的毛病是不问，我们先问。」}{「惠帝不算蠢」}{《晋书·惠帝纪》（\#137）}
\canonrow{案10 汉隆剃刀、「坏人不是蠢人」（对方说我们偷换成坏时）}{「英语里常说的汉隆剃刀，说的是能用愚蠢解释的，就不要归因于恶意。我们同意，所以不把他当坏人，只恼他不上心，回头照样讲理。坏是存心损人；不上心是对该会的事没用心，用不着恶意。」}{「不上心不是恶意。」}{「维基百科说……」（说「英语里常说的」）；引蔡康永「交朋友需要付出心力」}{Wikipedia「Hanlon's razor」（\#130）；澎湃新闻·有戏 2016-03-08（\#111）；Strawson 1962（\#32）}
\canonrow{案11 冯小刚「垃圾观众」（对方引时）}{「冯小刚错在把一个群体公开骂成垃圾：那是群体加羞辱，我们两样都不做。结构让很多人没机会会，我们先教；教过之后还不改的那一部分，才是我们恼的。」}{「结构解释和个人那一份可以并存。」}{「戴锦华同台反驳」}{界面新闻、每日经济新闻 2017（\#138）}
\canonrow{案12 教师准则、大医精诚（对方说「身份要求耐心」时）}{「职业角色不当本题主战场。就算看，教师准则第五条同一句写着严慈相济、严格要求学生，禁止的是歧视、侮辱学生；大医精诚管的是救不救、怎么救，没有要求医者心里不恼。」}{「职业把门槛抬高，恰恰说明普通人的门槛不在那里。」}{「教师可以厌学生」；「高校教师准则也有这一句」}{教育部 2018（\#147）；孙思邈〈大医精诚〉（\#148）}
\canonrow{案13 流行语与自媒体（对方说「厌蠢话语自己写下了处方」时）}{「挨个看来源：『别和蠢人争论』，考证网站找不到它出自马克·吐温的实质证据，而且说的是别争论；『蠢就是天生无可救药的坏』出自自媒体，只能当有人这样认为；中国网那句原话带着『有的人』『或许』。越线的那几句，我们和它切割。」}{「网上最响的几句，不是常人说『我厌蠢』时平常的样子。」}{「马克·吐温说过」；「中国网说厌蠢的人都会离对方远远的」}{Quote Investigator 2023（\#127）；有画说（自媒体）2025-03-02（\#124）；中国网·心理中国 2026-02-27（\#16）}
\canonrow{案14 主流舆论的批评（对方说「舆论一边倒」时）}{「舆论批的，是把别人的不了解当成蠢，是表现出来的贬低；这两样我们都不辩护。国旗下讲话自己也说，没有经历过而不会做，并不是蠢，我们同意。」}{「评有高下，不等于人格不平等。」}{「新京报是央媒」；「舆论都错了」}{央视网 2023-02-23（\#10）；红网 2023-02-25（\#11）；有间大学 2023-12-25（\#14）；中大附中国旗下讲话 2023-05-22（\#117）；新京报 2023-02-25（\#118）}
\canonrow{案15 装蠢人设（对方说「厌蠢逼人藏拙」时）}{「同一篇里，装蠢的大学生自述的好处，是小组作业不分给她最难的部分、大家更照顾她，文章自己说本质也算一种自私自利。被看成蠢，换来的是照顾；能装出来的蠢，是态度。」}{「装蠢的动机是少做脏活累活，不是怕被厌。」}{「装蠢的人都是坏人」；「这说明职场不厌蠢」}{真故研究室 2024-07-07（\#237、\#123）}
\canonrow{案16 其他对方引文（一句话接）}{苏格拉底：「那位政治家什么都不知道，却以为自己知道；不懂又自以为懂，正是人们叫作蠢的一类。」；鲁迅：「那是反讽的称呼，傻子不是真蠢。」；穆齐尔：「他说的是『蠢』作骂人话的用法；骂出口的我们不辩护。」；Goffman：「污名是对整个人的判决；我们下调的是这一次、这一方面。」}{费曼「讲不到新生能懂」：「说的是讲不清是讲者自己的问题，不是对笨学生的态度。」；特教老师：「面对学得慢的孩子，我们本来就给耐心。」；老子「无弃人」：「我们不弃人，这回放下、下回照样说。」}{「鲁迅讽刺傻子」；把特教对象说成「蠢人」；硬核读书会与有间大学当两份证据}{Plato《申辩》（\#146）；鲁迅 1925（\#145）；Musil 1937（\#77）；Goffman 1963（\#163）；Goodstein 1989（\#140）；光明日报 2024-06-11（\#143）；《老子》第二十七章（\#51）；硬核读书会 2023-12-23（\#122）；有间大学 2023-12-25（\#14）}
\canonrow{案17 新国辩复赛（只作备赛背景，稿件不用）}{「这道题今年年初在新国辩打过一场，反方赢了。一场比赛，只说明反方打得赢。」}{无（稿件不用；被问到才答）}{「反方更占优」「评委倾向反方」；说对阵方、票型以外的细节}{北大辩协 2026-02-24（\#113）}
\canonrow{案18 有间大学「我只厌打扰到我的蠢」（对方说我们辩护的是最好版本时）}{「我们的限定不是前缀，是双方共用的定义：一般人看来本该会、明显做错；第一次、没人教的，本来不算蠢。我们不靠谁的自我声明，靠记录。」}{「有间大学那篇是评论，没有数据；它批的是表现出来的贬低，我们不辩护。」；「我们说的常人，会恼，会一时走开，下回会再说；这是常人，不是最好版本。」}{「网友自己都会划界限，所以常人都这样」；「我们只厌打扰到我们的蠢」；「有间大学是《新周刊》正刊」；和硬核读书会那篇当两份证据}{反驳 29；有间大学 2023-12-25（\#14）}
\canonrow{案19 每日商报开头两句（对方说「一年后照列」时）}{「每日商报开头那两句没注出处，只能说是报道开头列出的说法；同一篇第 2 页写的，是一位工作超过两年的同事，多次劝告，同样的错接二连三，关系一度降至冰点，现在很微妙。」}{「持续出现在评论区，量的是评论区。」；「开头是记者挑的，第 2 页的案例也是记者挑的；双方各取一半，都不是比例。」}{「那两句是网友原话」「那两句是记者编的」（只说没注出处）；把同页「常识本就是由个体社会经验组成的」说成医生原话}{反驳 35；每日商报 2024-01-13（\#214）}
\canongroup{不进稿件的材料（有把握不进文稿；【待核实】文稿、速查都不进）}
\canonrow{有把握·Strawson 原刊卷页}{只说「Strawson 1962 年的《自由与怨恨》」，不报卷页。}{}{报原刊卷号与页码}{\#32、\#33、\#34、\#174、\#197 的卷页部分}
\canonrow{有把握·亚里士多德页码与译者}{只说「《尼各马可伦理学》第三卷」「第四卷」。}{}{报 Bekker 页码与译者名}{\#39、\#41、\#230 的页码与译者部分}
\canonrow{待核实·《现代汉语词典》}{只说「汉典释……」。}{}{「《现代汉语词典》释厌 / 蠢为……」}{\#7、\#68（【待核实】）}
\canonrow{未调研·「错」的词典释义}{只用中立定义口播版。}{}{「词典里错的意思是……」}{调研未收}
\canonrow{有把握·Aczél 2015 的比例与样本}{只说三类名称和「责任越大、后果越重越显得蠢」（\#93）。}{}{三类的比例、样本数、故事数}{\#69（有把握）}
\canonrow{待核实·厌蠢比例}{「平常」只靠机制和典型情形：「厌跟着代价走」（论1·机制3）。}{}{任何厌蠢比例；查到也不报}{\#102（检索未果）}
\canonrow{待核实·Mason 2003}{不引；我方本来就不为轻蔑辩护，用不着它。}{}{「Mason 为针对人的轻蔑辩护」}{\#166（【待核实】）}
\canonrow{待核实·Gottman「最佳预测因子」}{我方不说。对方说时：「原始摘要里没有出现轻蔑；论文里防御性的效应比轻蔑大。」（数16）}{}{「Gottman 研究发现轻蔑是离婚最佳预测因子」}{\#23 的这一部分（【待核实】）}
\canonrow{待核实·朴素实在论}{对方说时：「我方只核到 Pronin 等 2004 年的综述；请对方报原文。」}{}{「朴素实在论由 Ross 与 Ward 1996 年提出」}{\#178（【待核实】）}
\canonrow{有把握·费曼在巴西、陶行知}{对方引费曼时：「体制造出的不会不该怪学生，我们同意；对方报出处时请说明。」引陶行知时：「讥笑是表达，我们不辩护；这句的原诗请对方说明出处。」}{}{我方引用这两句}{\#141、\#144（有把握）}
\canonrow{有把握·基本归因错误；与本方无关的两条}{基本归因错误我方不用，对方用时不争理论本身，只问「他被教过没有」；医疗差错死亡数、「多样性胜过能力」的批评不用。}{}{引用这三条}{\#90、\#103（有把握）；\#96（【待核实】）}
\canonrow{赛前待办}{康德两句换成李秋零译本后同步改论2·机制1、论2·机制3、判13；刑法第十五条对照国家法律法规数据库；朋霍费尔没有打开校勘本，中文为自译；只核了摘要的研究不报效应量和样本量，查到全文也只写进备赛文档，稿件口径不变。}{}{「这是李秋零译本的原文」}{\#48、\#45、\#180、\#74、\#75、\#76、\#211、\#247、\#219、\#86、\#250、\#224、\#246、\#188}
\end{canon}
\begin{fields}
\field{对方若报出，可以当场要求出示来源}{伪证申诉前先确认对方确实说了数字或原话。\sub{无出处的数}{央视网、新京报评论、每日商报里坐飞机人数、通地铁城市数都没给出处（\#10、\#118、\#15）；老人用智能手机的比例，央视网称出自某研究中心，调研没核到（\#10）。}\sub{没有的数据}{任何厌蠢比例：调研里没有（\#102）。}\sub{没核到的说法}{「Gottman 研究证明轻蔑是离婚最强预测因子」，原始摘要里没有轻蔑（\#23）；「朴素实在论由 Ross 与 Ward 1996 年提出」（\#178）；陶行知诗的年份与篇名（\#144）；费曼在巴西（\#141）。}\sub{有出处但不是厌蠢比例}{职场新人调查（\#121）、子女以为父母学不会（\#192）、防洪材料（\#206）、市场减偏（\#79）：数字有出处，回一句「那不是厌蠢的比例」，不申诉。}}
\field{统一写法}{被试人数统一说「一万二千多名」（数1）；个别场合一律说「那一次是看错了，要调分寸」，「错了」只留给「平常的样子该改」；后见之明元分析不评大小，只说有偏差、没说明平常判错（数25），只备用，不进立论。}
\end{fields}

\section{十一、论据出处清单}
收第五至第十节引用的全部来源，共 257 条，按编号排列；正文与口径表里的 \#N 就是这里的编号。状态照调研原条，一条没有升级：只有「已核实」进文稿与速查；「有把握」只留在备赛文档；【待核实】只在本清单。已核实的每条都有出处（含年份）、链接或 DOI、原文引句和核实日期。读法：论据是一句话概括，带数字的写数字；原文引句每条只录一句（中文过长的用「……」截取，英文过长的用「…」截断）；核实日期后括号里是核实途径，「＊」表示同条有一部分没到已核实或核实范围有限定，见清单末尾。「用在」：正论一 / 反论二＝第五、六节最终论点，正淘 / 反淘＝被淘汰候选；正驳 N＝第七节第 N 条，反驳 N＝第八节第 N 条；正口径 / 反口径＝第十节对应分组（反方写行号，如「定1」「数14」），「禁」＝只在禁止说法栏，「六·…」＝口径表「不进稿件的材料」。标〔智〕的测的是智力或认知能力，用之前先搭桥；研究规模与效应量只当研究本身的规模说。
\begin{sourcelist}
\source{1}{汉典「厌」：嫌恶、憎恶；本义吃饱}{汉典 zdic.net，「厌」字条，2026 年访问 \url{https://zdic.net/hans/\%E5\%8E\%8C}}{「嫌恶，憎恶。」}{2026-10-11}{已核实}{反驳16；正口径·定义；反口径·定1}{}
\source{2}{汉典「厌恶」：对人或事物产生强烈的反感}{汉典 zdic.net，「厌恶」词条，2026 年访问 \url{https://zdic.net/hans/\%E5\%8E\%8C\%E6\%81\%B6}}{「对人或事物产生强烈的反感。」}{2026-10-11（原页）}{已核实}{正口径·定义；反口径·定1}{}
\source{3}{《重編國語辭典》「厭」：憎惡、嫌棄}{中华民国教育部，《重編國語辭典修訂本》臺灣學術網路第六版（2021），「厭」（ID 11068）、「厭惡」（ID 155997） \url{https://dict.revised.moe.edu.tw/dictView.jsp?ID=11068}}{「憎惡、嫌棄。」}{2026-10-11}{已核实}{正口径·定义；反口径·定1}{}
\source{4}{汉典「嫌弃」：厌恶而不愿接近}{汉典 zdic.net，「嫌弃」词条，2026 年访问 \url{https://zdic.net/hans/\%E5\%AB\%8C\%E5\%BC\%83}}{「厌恶而不愿接近。」}{2026-10-11（原页）}{已核实}{正口径·定义；反口径·定1}{}
\source{5}{汉典「鄙视」「轻蔑」：轻视、看不起}{汉典 zdic.net，「鄙视」「轻蔑」词条（含国语辞典），2026 年访问 \url{https://zdic.net/hans/\%E9\%84\%99\%E8\%A7\%86}}{「轻视，看不起。」}{2026-10-11}{已核实}{反驳16；反口径·定1}{}
\source{6}{汉典「厌烦」：厌恶而不耐烦}{汉典 zdic.net，「厌烦」「不耐烦」词条，2026 年访问 \url{https://zdic.net/hans/\%E5\%8E\%8C\%E7\%83\%A6}}{「对某事感到厌恶不耐烦。」}{2026-10-11}{已核实}{正驳33；正口径·案例}{}
\source{7}{现汉第 7 版「厌」的义项（凭记忆）}{中国社会科学院语言研究所词典编辑室编，《现代汉语词典》第 7 版，商务印书馆，2016}{}{}{【待核实】}{正口径·定义禁、六·待核；反口径·定1禁、六}{纸本现汉第 7 版或官方 App，抄义项与页码}
\source{8}{「厌蠢症」约 2023 年元旦前后出现}{上游新闻，2023-02-24，《网红教人如何坐公交高铁引来冷嘲热讽，“厌蠢症”是种什么病？》 \url{https://news.cqnews.net/1/detail/1078733353267396608/app/content\_1078733353267396608.html}}{「“厌蠢症”一词大致出现在2023年元旦前后，来源出处暂不明确」}{2026-10-11}{已核实}{反驳1；反口径·案6}{}
\source{9}{自称厌蠢者自己划界：不厌努力但不会的人}{上游新闻，2023-02-24，《网红教人如何坐公交高铁引来冷嘲热讽，“厌蠢症”是种什么病？》 \url{https://news.cqnews.net/1/detail/1078733353267396608/app/content\_1078733353267396608.html}}{「遇见不会但努力的人，并不会激发她的厌蠢症」}{2026-10-11}{已核实}{正驳29；反驳5、24、29；反口径·案5}{}
\source{10}{央视网：厌蠢症是看不得别人犯蠢、太笨}{央视网（雲成章），2023-02-23，《“厌蠢症”才是一种“真蠢”》，北京日报客户端转载 \url{https://xinwen.bjd.com.cn/content/s63f6b928e4b03a6b6edc84d8.html}}{「最近网上有个新词，叫“厌蠢症”，是说看不得别人犯蠢、没常识、太笨。」}{2026-10-11（＊）}{已核实}{正论一；正驳33；反驳1、15、16；正口径·定义、论点、案例、六·待核；反口径·案14、六·对方报}{}
\source{11}{红网：厌蠢症上热搜，「自以为是、傲慢自大」}{红网（湖南），2023-02-25，《“厌蠢症”上热搜：其实，你我都有“犯蠢”之时》 \url{https://news.qq.com/rain/a/20230225A04L1B00}}{「央视网也在2月23日发表评论，明确表态“厌蠢症”不可取。」}{2026-10-11}{已核实}{反驳15；反口径·案14}{}
\source{12}{新京报书评：厌蠢近乎「厌烦没有常识」「厌麻烦」}{新京报书评周刊，2023-11-01，《为什么我们容易犯“厌蠢症”？》 \url{https://m.bjnews.com.cn/detail/1698812159168656.html}}{「“我的厌蠢症要犯了”如今成了一大网络流行语。」}{2026-10-11}{已核实}{反淘；正驳33；反口径·定1禁、案4}{}
\source{13}{新京报书评：厌蠢标签合理化贬低他人，「恐弱」}{新京报书评周刊，2023-11-01，《为什么我们容易犯“厌蠢症”？》 \url{https://m.bjnews.com.cn/detail/1698812159168656.html}}{「“厌蠢症”这个标签的出现给了我们合理化贬低他人的契机」}{2026-10-11}{已核实}{正驳32、33、43；反驳8、9、15；正口径·定义、案例；反口径·定4禁、案4}{}
\source{14}{有间大学：调侃成病也掩盖不了傲慢}{有间大学，2023-12-25，《“厌蠢症”流行，挺可悲的》 \url{https://news.qq.com/rain/a/20231225A0727R00}}{「无论怎样把它调侃为一种“疾病”，都不能掩盖背后浓浓的傲慢和尖刻」}{2026-10-11}{已核实}{正驳5、32、43；反驳7、15、29；正口径·案例；反口径·案14、案16、案18}{}
\source{15}{每日商报：医生称厌蠢背后是对潜在风险的厌烦}{每日商报，2024-01-13，《“怎么还会有这种人”“我的厌蠢症又犯了”专家：“厌蠢”的实质是一种“慕强”心态》 \url{https://jrsh.hangzhou.com.cn/read/content/2024-01/13/content\_8674352.htm}}{「“厌蠢”背后其实是对潜在风险的厌烦。」}{2026-10-11（＊）}{已核实}{反淘；正驳33；正口径·案例、六·待核；反口径·定1禁、六·对方报}{}
\source{16}{中国网：厌蠢症非医学术语，成因含优越感、投射}{中国网·心理中国（方丽婷），2026-02-27，《“厌蠢症”是对自我的投射？》 \url{http://psy.china.com.cn/2026-02/27/content\_43359463.htm}}{「对于他人相比自己而言显得无知的言行感到烦躁、难以忍受、厌恶的情绪状态」}{2026-10-11}{已核实}{反淘；正驳36；反驳3、8、16、24；正口径·定义、案例；反口径·定1、案13}{}
\source{17}{张琨（有人认为）：厌蠢病是对冗长无效表达的不耐}{张琨，微信公众号「张琨随笔」，2025-08-11，《麦肯锡，厌蠢病》 \url{https://m.huxiu.com/article/4680123.html}}{「它的主要症状，就是对冗长无效的表达深感不耐。」}{2026-10-11}{已核实}{正驳33；正口径·案例}{}
\source{18}{「你只是太累了」（有人认为）；共情疲劳的学术边界}{(a) 滑洋心理（百家号）2025-04-19；(b) 江语童、李红、孟现鑫 2024，《“关怀的代价”》，《心理科学》47(3): 752–759，doi:10.16719/j.cnki.1671-6981.20240330}{「你不是讨厌“蠢”，你只是太累了」}{2026-10-11（摘要＊）}{已核实}{正驳23；正口径·数据}{}
\source{19}{轻蔑的核心是看低对方，想忽视、排斥}{Fischer \& Giner-Sorolla 2016, Contempt: Derogating others while keeping calm. Emot Rev, 8(4), 346–357. doi:10.1177/1754073915610439 \url{https://pure.uva.nl/ws/files/26394166/1754073915610439.pdf}}{“First, contempt is a discrete emotion, different from anger, disgust, or hate; its core feature is that contemptuous people look down on their targets and try to ignore or exclude them.”}{2026-10-11}{已核实}{反驳4、16；正口径·数据}{}
\source{20}{轻蔑贬对方愚蠢、无能，贬品性多于行为}{Fischer \& Giner-Sorolla 2016, Contempt: Derogating others while keeping calm. Emot Rev, 8(4), 346–357. doi:10.1177/1754073915610439 \url{https://pure.uva.nl/ws/files/26394166/1754073915610439.pdf}}{“However, in contempt we despise the other as stupid, incompetent, or immoral (see also Hutcherson \& Gross, 2011)—deriding their disposition more so than their act.”}{2026-10-11}{已核实}{正论二；正驳10、11、13、17、19、25、36、40、44；反驳3、16、26；正口径·论点、数据；反口径·数19}{}
\source{21}{厌恶隔绝污染，轻蔑向下看}{Fischer \& Giner-Sorolla 2016, Contempt: Derogating others while keeping calm. Emot Rev, 8(4), 346–357. doi:10.1177/1754073915610439 \url{https://pure.uva.nl/ws/files/26394166/1754073915610439.pdf}}{“Disgust is often considered as a kindred emotion to contempt, because they both imply avoidance tendencies, though in the case of disgust this tendency is more to shutting oneself off from …”}{2026-10-11}{已核实}{正驳36}{}
\source{22}{针对可改变品性的轻蔑可以正当}{Fischer \& Giner-Sorolla 2016, Contempt: Derogating others while keeping calm. Emot Rev, 8(4), 346–357. doi:10.1177/1754073915610439 \url{https://pure.uva.nl/ws/files/26394166/1754073915610439.pdf}}{“Finally, although contempt is often applied unfairly to people on the basis of prejudice, it does seem to have a place in justified moral condemnation if it targets an undesirable facet of a person's …”}{2026-10-11}{已核实}{正驳40}{}
\source{23}{轻蔑比愤怒更易诱人自以为优越}{Fischer \& Giner-Sorolla 2016, Contempt: Derogating others while keeping calm. Emot Rev, 8(4), 346–357. doi:10.1177/1754073915610439 \url{https://pure.uva.nl/ws/files/26394166/1754073915610439.pdf}}{“Nonetheless, it is plausible that contempt more so than anger or hate risks hypocrisy by enticing our less-than-perfect selves to feel socially or morally superior.”}{2026-10-11（＊）}{已核实}{反驳21；反口径·数16禁、六、六·对方报}{}
\source{24}{表情识别中轻蔑常被认作厌恶（35\%／55\%）}{Fischer \& Giner-Sorolla 2016, Contempt: Derogating others while keeping calm. Emot Rev, 8(4), 346–357. doi:10.1177/1754073915610439 \url{https://pure.uva.nl/ws/files/26394166/1754073915610439.pdf}}{“For example, Haidt and Keltner (1999), comparing American and Indian recognition rates of various emotions, found that only 35\% of the Americans correctly labeled the contempt expression as contempt …”}{2026-10-11}{已核实}{正驳36；正口径·六·不上场}{}
\source{25}{轻蔑导向拒绝排斥，可由反复愤怒演变（3 项研究）}{Fischer \& Roseman 2007, Beat them or ban them: The characteristics and social functions of anger and contempt. JPSP, 93(1), 103–115. doi:10.1037/0022-3514.93.1.103}{“Anger is characterized more by short-term attack responses but long-term reconciliation, whereas contempt is characterized by rejection and social exclusion of the other person, both in the …”}{2026-10-11（摘要）}{已核实}{正口径·数据}{}
\source{26}{评判无能行为时起作用的是轻蔑（5 项研究）}{Hutcherson \& Gross 2011, The moral emotions: A social-functionalist account of anger, disgust, and contempt. JPSP, 100(4), 719–737. doi:10.1037/a0022408}{“… self-relevance on anger but not other emotions, and (c) a role for contempt in judging incompetent actions.”}{2026-10-11（摘要）}{已核实}{正论二；正驳5、11、36、44；反驳16；正口径·论点、数据；反口径·数4}{}
\source{27}{轻蔑只指向人、带优越感、不看好对方改善}{Miceli \& Castelfranchi 2018, Contempt and disgust: The emotions of disrespect. JTSB, 48(2), 205–229. doi:10.1111/jtsb.12159 \url{https://iris.cnr.it/handle/20.500.14243/390145}}{“Contempt is felt exclusively towards human targets, and implies sense of superiority over them, pessimistic feelings about their possibility of betterment, detachment from them, and avoidance driven …”}{2026-10-11（摘要）}{已核实}{正论二；正驳11、15、36、40；反驳3、15、16、39、40；正口径·定义、论点、数据；反口径·回4、数20}{}
\source{28}{特质性轻蔑对低能力目标反应尤其强}{Schriber 等 2017, Dispositional contempt: A first look at the contemptuous person. JPSP, 113(2), 280–309. doi:10.1037/pspp0000101}{“Dispositional contempt predicted contemptuous reactions to eliciting film clips, particularly when targets showed low competence/power.”}{2026-10-11（摘要）}{已核实}{反驳16}{}
\source{29}{刻板印象模型：低能力高温暖对应怜悯，不是轻蔑}{Fiske 等 2002, A model of (often mixed) stereotype content: Competence and warmth respectively follow… JPSP, 82(6), 878–902. \url{https://pubmed.ncbi.nlm.nih.gov/12051578/}}{“(c) distinct emotions (pity, envy, admiration, contempt) differentiate the 4 competence-warmth combinations.”}{2026-10-11（摘要）}{已核实}{反驳3；反口径·数21}{}
\source{30}{BIAS map：能力刻板印象决定被动伤害（忽视）}{Cuddy 等 2007, The BIAS map: Behaviors from intergroup affect and stereotypes. JPSP, 92(4), 631–648. doi:10.1037/0022-3514.92.4.631}{“Competence stereotypes determine passive behavioral tendencies, attenuating passive harm (neglecting) and eliciting passive facilitation (associating).”}{2026-10-11（摘要）}{已核实}{正驳29；正口径·数据}{}
\source{31}{被轻蔑者表现提升，但攻击性升、内隐自尊降}{Melwani \& Barsade 2011, Held in contempt: The psychological, interpersonal, and performance consequences of… JPSP, 101(3), 503–520. doi:10.1037/a0023492}{“… better task performance but also significantly more interpersonal aggressiveness toward their virtual partners compared with recipients of failure, angry, or neutral feedback.”}{2026-10-11（摘要）}{已核实}{正驳31；正口径·数据}{}
\source{32}{Strawson：反应态度回应意志好坏，不是能力}{Strawson 1962, “Freedom and Resentment”, Proc Br Acad 48: 1–25 \url{https://www.informationphilosopher.com/solutions/philosophers/strawson/freedom\_and\_resentment.html}}{“The reactive attitudes I have so far discussed are essentially reactions to the quality of others' wills towards us, as manifested in their behaviour: to their good or ill will or indifference or …”}{2026-10-11（转载页＊）}{已核实}{反论二；反驳17、25；正口径·六·把握；反口径·论2、论9、案10、六}{}
\source{33}{Strawson：客观态度把人当作要管理、回避的对象}{Strawson 1962, “Freedom and Resentment”, Proc Br Acad 48: 1–25 \url{https://www.informationphilosopher.com/solutions/philosophers/strawson/freedom\_and\_resentment.html}}{“...to be managed or handled or cured or trained; perhaps simply to be avoided, though this gerundive is not peculiar to cases of objectivity of attitude.”}{2026-10-11（转载页＊）}{已核实}{反淘；正驳29；正口径·论点；反口径·判15禁、论9、六}{}
\source{34}{Strawson：客观态度是逃离投入之累的避难所}{Strawson 1962, “Freedom and Resentment”, Proc Br Acad 48: 1–25 \url{https://www.informationphilosopher.com/solutions/philosophers/strawson/freedom\_and\_resentment.html}}{“We have this resource and can sometimes use it; as a refuge, say, from the strains of involvement; or as an aid to policy; or simply out of intellectual curiosity.”}{2026-10-11（转载页＊）}{已核实}{反论二；正驳18、23、29；反驳4、7、23、29、42；正口径·判准、论点；反口径·判15、论2、六}{}
\source{35}{Bell：轻蔑有时是回应优越之恶的最佳方式}{Bell 2013, Hard Feelings: The Moral Psychology of Contempt, Oxford University Press；出版社摘要页 \url{https://academic.oup.com/book/12360}}{“Contempt is, in some cases, the best way to respond to arrogance, hypocrisy, and other vices of superiority.”}{2026-10-11（摘要）}{已核实}{正驳40；正口径·判准；反口径·学2}{}
\source{36}{Bell 的七条轻蔑恰当性条件（据书评）}{Robert C. Roberts，2013-09-03，书评 “Hard Feelings: The Moral Psychology of Contempt”，NDPR；转述 Bell 2013 pp. 147–151 \url{https://ndpr.nd.edu/news/hard-feelings-the-moral-psychology-of-contempt/}}{“Bell lays down seven conditions (147-51) that ”must be satisfied if a particular token of contempt is appropriate“ (151).”}{2026-10-11（书评）}{已核实}{正驳40；正口径·判准}{}
\source{37}{Bell：轻蔑与承认尊重相容（据书评）}{Robert C. Roberts，2013-09-03，书评 “Hard Feelings: The Moral Psychology of Contempt”，NDPR；转述 Bell 2013 pp. 147–151 \url{https://ndpr.nd.edu/news/hard-feelings-the-moral-psychology-of-contempt/}；本条：Roberts 2013 书评，引 Bell 2013 p. 171}{“contempt is not incompatible with recognition respect; there is no reason to suppose that the contemnor will, for example, always fail to acknowledge the moral or legal rights of those they con[t]emn”}{2026-10-11（书评）}{已核实}{反口径·判13、论2}{}
\source{38}{Roberts：纯粹的轻蔑是去人性化的}{Robert C. Roberts，2013-09-03，书评 “Hard Feelings: The Moral Psychology of Contempt”，NDPR；转述 Bell 2013 pp. 147–151 \url{https://ndpr.nd.edu/news/hard-feelings-the-moral-psychology-of-contempt/}}{“Pure contempt, unmixed with outrage or anger or indignation, is dehumanizing because it recognizes nothing valuable in its target.”}{2026-10-11}{已核实}{反论二、反淘；反驳3、16、39；反口径·论2}{}
\source{39}{亚里士多德：该怒时怒受称赞，该怒不怒被视为愚蠢}{Aristotle, Nicomachean Ethics IV.5，D. P. Chase 英译，J. A. Smith 导言本，Project Gutenberg \#8438，2026 年访问 \url{https://www.gutenberg.org/ebooks/8438}}{“Now he who is angry at what and with whom he ought, and further, in right manner and time, and for proper length of time, is praised...”}{2026-10-11（全文＊）}{已核实}{正驳40；反驳6、21；正口径·六·把握；反口径·学7、六}{}
\source{40}{亚里士多德：积怨难消者是自己和至交的大麻烦}{Aristotle, Nicomachean Ethics IV.5，D. P. Chase 英译，J. A. Smith 导言本，Project Gutenberg \#8438，2026 年访问 \url{https://www.gutenberg.org/ebooks/8438}}{“Such men are very great nuisances to themselves and to their best friends.”}{2026-10-11}{已核实}{反驳21}{}
\source{41}{亚里士多德：不责天生盲人，责醉酒致盲者}{Aristotle, Nicomachean Ethics IV.5，D. P. Chase 英译，J. A. Smith 导言本，Project Gutenberg \#8438，2026 年访问 \url{https://www.gutenberg.org/ebooks/8438}}{“...no one would think of upbraiding, but would rather compassionate, a man who is blind by nature, or from disease, or from an accident...”}{2026-10-11（全文＊）}{已核实}{正淘、反论一；反驳6、25；正口径·论点、六·把握；反口径·论1、学3、六}{}
\source{42}{SEP：Darwall 区分承认尊重与评价尊重}{Robin S. Dillon, “Respect”, SEP；转述 Stephen Darwall, 1977, “Two Kinds of Respect”, Ethics 88(1): 36–49 \url{https://plato.stanford.edu/entries/respect/}}{“Darwall (1977) distinguishes two kinds of respect: recognition respect and appraisal respect.”}{2026-10-11（SEP＊）}{已核实}{正口径·六·把握}{}
\source{43}{SEP：康德式尊重义务含不轻蔑、不嘲笑}{Robin S. Dillon, “Respect”, SEP；转述 Stephen Darwall, 1977, “Two Kinds of Respect”, Ethics 88(1): 36–49 \url{https://plato.stanford.edu/entries/respect/}}{“We have no moral duty to feel respect for others; rather, the respect we owe others is ”to be understood as the maxim of limiting our self-esteem by the dignity of humanity in another person...“”}{2026-10-11（SEP）}{已核实}{正驳38；反驳33；正口径·论点；反口径·学10}{}
\source{44}{SEP：尊严平等不排除区别评价，但受承认尊重约束}{Robin S. Dillon, “Respect”, SEP；转述 Stephen Darwall, 1977, “Two Kinds of Respect”, Ethics 88(1): 36–49 \url{https://plato.stanford.edu/entries/respect/}}{“However, the equality of all rational beings does not entail that persons cannot also be differentially evaluated and valued in other ways for their particular qualities, accomplishments, merit, or …”}{2026-10-11}{已核实}{反论二；正驳11；反驳3、7、15、26、28、33、40；正口径·论点；反口径·判13、论2、回6}{}
\source{45}{康德 §39 附释：否认他有理智，怎能让他明白错了}{Kant 1797, Die Metaphysik der Sitten, Tugendlehre §39 Anm., AA VI 463, Z. 22–33，korpora.org，2026 年访问 \url{https://korpora.org/kant/aa06/463.html}}{“Denn spricht man seinem Gegner in einem gewissen Urtheile durch jene Ausdrücke allen Verstand ab, wie will man ihn dann darüber verständigen, daß er geirrt habe?”}{2026-10-11（德文原典＊）}{已核实}{正论二；正驳11；反驳4、42；正口径·论点、六·待办；反口径·论2、回5、六·补核}{}
\source{46}{塞涅卡：自然原谅缺陷时，没有人发怒}{Seneca, De Ira (On Anger) Book II, ch. X；Aubrey Stewart 英译，Bohn's Classical Library, 1900；维基文库英文版 \url{https://en.wikisource.org/wiki/Of\_Anger/Book\_II}}{“The wise man will not be angry with sinners. Why not? Because he knows that no one is born wise, but becomes so...”}{2026-10-11（维基文库）}{已核实}{反驳6}{}
\source{47}{奥勒留：恶行出于无知，不能对同类发怒}{Marcus Aurelius, Meditations II.1；George Long 英译 Thoughts of Marcus Aurelius Antoninus，Project Gutenberg \#15877，2026 年访问 \url{https://www.gutenberg.org/ebooks/15877}}{“All these things happen to them by reason of their ignorance of what is good and evil.”}{2026-10-11（全文）}{已核实}{反驳6}{}
\source{48}{康德 §39：心里看低有时难免，拒绝尊重违背义务}{Kant 1797, Die Metaphysik der Sitten, Tugendlehre §39 Anm., AA VI 463, Z. 22–33，korpora.org，2026 年访问 \url{https://korpora.org/kant/aa06/463.html}}{“Andere Verachten (contemnere), d. i. ihnen die dem Menschen überhaupt schuldige Achtung weigern, ist auf alle Fälle pflichtwidrig; denn es sind Menschen.”}{2026-10-11（德文原典＊）}{已核实}{正论二、反论二；正驳11、17、28、38、44；反驳4、5、7、15、20、28、33、42；正口径·判准、论点、六·待办；反口径·判13、论2、六·补核}{}
\source{49}{《论语》：困而不学，民斯为下矣}{《论语·季氏第十六》，十六之九，维基文库本，2026 年访问 \url{https://zh.wikisource.org/wiki/\%E8\%AB\%96\%E8\%AA\%9E/\%E5\%AD\%A3\%E6\%B0\%8F\%E7\%AC\%AC\%E5\%8D\%81\%E5\%85\%AD}}{「困而學之，又其次也。困而不學，民斯爲下矣！」}{2026-10-11}{已核实}{正淘；反驳6；反口径·学5}{}
\source{50}{《论语》：子张「嘉善而矜不能」}{《论语·子张第十九》，十九之三，维基文库本，2026 年访问 \url{https://zh.wikisource.org/wiki/\%E8\%AB\%96\%E8\%AA\%9E/\%E5\%AD\%90\%E5\%BC\%B5\%E7\%AC\%AC\%E5\%8D\%81\%E4\%B9\%9D}}{「君子尊賢而容眾，嘉善而矜不能。」}{2026-10-11}{已核实}{反驳15；反口径·学6}{}
\source{51}{《老子》：常善救人，故无弃人}{《道德经》第二十七章，王弼本，维基文库，2026 年访问 \url{https://zh.wikisource.org/wiki/\%E9\%81\%93\%E5\%BE\%B7\%E7\%B6\%93\_(\%E7\%8E\%8B\%E5\%BC\%BC\%E6\%9C\%AC})}{「是以聖人常善救人，故無棄人」}{2026-10-11}{已核实}{反驳10；反口径·案16}{}
\source{52}{Rawls：没有人应得他更高的自然能力}{Rawls 1971, A Theory of Justice, §17；本次经 Michael J. Sandel, 2009, Justice: What's the Right Thing to Do?, pp. 156–157 的整段引文核实 \url{https://courses.edx.org/c4x/HarvardX/ER22.1x/asset/Chapter\_6\_-\_John\_Rawls\_\_151-157\_.pdf}}{“No one deserves his greater natural capacity nor merits a more favorable starting place in society.”}{2026-10-11（书中引文＊）}{已核实}{反驳6}{}
\source{53}{Sandel 转述：低学历者最不受喜欢}{Michael J. Sandel, 2020-09-02, “Disdain for the Less Educated Is the Last Acceptable Prejudice”, The New York Times \url{https://www.motherjones.com/kevin-drum/2020/09/disdain-for-the-less-educated-needs-to-stop-now/}}{“Of all these groups, the poorly educated were disliked most of all.”}{2026-10-11（博客节选＊）}{已核实}{反驳9}{}
\source{54}{Kuppens：高学历者对低学历者有偏见}{Kuppens 等, 2018, “Educationism and the irony of meritocracy: Negative attitudes of higher educated people…”, JESP 76: 429–447, doi:10.1016/j.jesp.2017.11.001}{“In contrast with popular views of the higher educated as tolerant and morally enlightened, we find that higher educated participants show education-based intergroup bias: They hold more negative …”}{2026-10-11（录用稿）}{已核实}{反驳9；正口径·六·不上场；反口径·定4禁、数22}{}
\source{55}{Kuppens 研究 6：对低学历者的愤怒由责任归因中介}{Kuppens 等, 2018, “Educationism and the irony of meritocracy: Negative attitudes of higher educated people…”, JESP 76: 429–447, doi:10.1016/j.jesp.2017.11.001；本条：Study 6 结果与讨论，录用稿 pp. 44–48}{“...they elicited much less sympathy (b = -1.76, SE = .22, p < .001), much less pity (b = -1.94, SE = .26, p < .001), more anger (b = .69, SE = .19, p < .001), and more contempt (b = .61, SE = .23, p …”}{2026-10-11（录用稿）}{已核实}{正淘；反驳9；正口径·六·不上场；反口径·数22}{}
\source{56}{Kuppens 研究 7：精英主义信念越强越归咎（r=.47）}{Kuppens 等, 2018, “Educationism and the irony of meritocracy: Negative attitudes of higher educated people…”, JESP 76: 429–447, doi:10.1016/j.jesp.2017.11.001；本条：Study 7，录用稿 p. 57}{“As expected, meritocracy beliefs were strongly related to attributions of responsibility (r = .47, p < .001) and blameworthiness (r = .48, p < .001) in relation to the less educated...”}{2026-10-11（录用稿）}{已核实}{反驳9；正口径·六·不上场；反口径·论1禁、数22}{}
\source{57}{〔智〕智力遗传率约 20\% 升到或许 80\%}{Plomin \& Deary 2015, “Genetics and intelligence differences: five special findings”, Mol Psychiatry 20(1): 98–108, doi:10.1038/mp.2014.105（2014-09-16 在线发表） \url{https://pubmed.ncbi.nlm.nih.gov/25224258/}}{“The heritability of intelligence increases from about 20\% in infancy to perhaps 80\% in later adulthood.”}{2026-10-11（摘要＊）}{已核实}{正淘；反驳6；正口径·六·不上场；反口径·数36}{}
\source{58}{Weiner：缺努力比缺能力招来更多惩罚}{Weiner 1993, “On sin versus sickness: A theory of perceived responsibility and social motivation”, Am Psychol 48(9): 957–965, doi:10.1037/0003-066X.48.9.957}{“Lack of effort as a cause of achievement failure evokes more punishment than does lack of ability.”}{2026-10-11（摘要）}{已核实}{反论一；正驳44；反驳27；正口径·论点；反口径·判8、论1、学1}{}
\source{59}{道德主义谬误：觉得不该有不等于不恰当}{D'Arms \& Jacobson 2023, Rational Sentimentalism, OUP, ch. 8（章节摘要）；原始论文 2000, PPR 61: 65–90 \url{https://academic.oup.com/book/45567/chapter/394733317}}{“Philosophers often infer that an emotion is unfitting because it would be wrong to feel. That is the moralistic fallacy...”}{2026-10-11（摘要＊）}{已核实}{反驳23；反口径·判6}{}
\source{60}{SEP：「鄙视不可鄙的人」是错误种类的理由}{Howard 2023, “Fitting Attitude Theories of Value”, SEP（2023-01-24 首发） \url{https://plato.stanford.edu/entries/fitting-attitude-theories/}}{“For example, Crisp's demon might equally threaten harm unless you admire someone who's not admirable, despise someone who's not despicable, love someone who's not lovable, etc.”}{2026-10-11（SEP）}{已核实}{反驳23}{}
\source{61}{对应偏差：从情境可解释的行为推断特质}{Gilbert \& Malone 1995, “The correspondence bias”, Psychol Bull 117(1): 21–38, doi:10.1037/0033-2909.117.1.21}{“The correspondence bias is the tendency to draw inferences about a person's unique and enduring dispositions from behaviors that can be entirely explained by the situations in which they occur.”}{2026-10-11（摘要）}{已核实}{正驳34、39；反驳26；正口径·论点；反口径·论1禁、数39}{}
\source{62}{〔智〕能力最差者自估第 62 百分位，实为第 12}{Kruger \& Dunning 1999, “Unskilled and unaware of it: How difficulties in recognizing one's own incompetence…”, JPSP 77(6): 1121–1134, doi:10.1037/0022-3514.77.6.1121}{“Although their test scores put them in the 12th percentile, they estimated themselves to be in the 62nd.”}{2026-10-11（摘要）}{已核实}{正口径·六·不上场}{}
\source{63}{愤怒敌意与冠心病相关（HR 1.19）}{Chida \& Steptoe 2009, “The association of anger and hostility with future coronary heart disease: a…”, JACC 53(11): 936–946, doi:10.1016/j.jacc.2008.11.044}{“Anger and hostility were associated with increased CHD events in the healthy population studies (combined hazard ratio [HR]: 1.19; 95\% confidence interval [CI]: 1.05 to 1.35, p = 0.008)...”}{2026-10-11（摘要）}{已核实}{反驳21；反口径·判10、数17}{}
\source{64}{心理安全与学习行为相关；怕显得无能不敢求助（51 个团队）}{Edmondson 1999, “Psychological Safety and Learning Behavior in Work Teams”, ASQ 44(2): 350–383（JSTOR 2666999），全文 PDF \url{https://web.mit.edu/curhan/www/docs/Articles/15341\_Readings/Group\_Performance/Edmondson\%20Psychological\%20safety.pdf}}{“...by admitting an error or asking for help, an individual may appear incompetent and thus suffer a blow to his or her image.”}{2026-10-11（全文＊）}{已核实}{正淘；正驳22、24、26；反驳5；正口径·论点、数据}{}
\source{65}{「坏苹果」拖垮团队；应对含促其改变、排斥}{University of Washington News, 2007-02-12, “Rotten to the core: How workplace 'bad apples' spoil barrels of good employees”；介绍 Will Felps, Terence R. Mitchell \& Eliza Byington, 2006, “How, When, and Why Bad Apples Spoil the Barrel: Negative Group Members and…”, Res Organ Behav 27: 175–222, doi:10.1016/S0191-3085(06)27005-9 \url{https://www.washington.edu/news/2007/02/12/rotten-to-the-core-how-workplace-bad-apples-spoil-barrels-of-good-employees/}}{“Felps and Mitchell define negative people as those who don't do their fair share of the work, who are chronically unhappy and emotionally unstable, or who bully or attack others.”}{2026-10-11（新闻稿＊）}{已核实}{正驳22；正口径·数据、六·待核}{}
\source{66}{汉典「蠢」：愚笨、笨拙；古义「不逊」}{汉典 zdic.net「蠢」字条，2026 年访问 \url{https://www.zdic.net/hans/\%E8\%A0\%A2}}{「蠢，不逊也 —— 《尔雅·释训》」}{2026-10-11}{已核实}{正口径·定义；反口径·定2}{}
\source{67}{汉典「蠢人」：缺乏判断力或谋虑的人}{汉典 zdic.net 词语条目「蠢人」「愚蠢」「愚昧」「无知」，2026 年访问 \url{https://www.zdic.net/hans/\%E8\%A0\%A2\%E4\%BA\%BA}}{「缺乏判断力或谋虑的人」}{2026-10-11}{已核实}{反驳25；正口径·定义；反口径·定2}{}
\source{68}{现汉第 7 版「蠢」「愚蠢」「笨」释义（凭记忆）}{中国社会科学院语言研究所词典编辑室编，《现代汉语词典》第 7 版，商务印书馆，2016}{}{}{【待核实】}{正口径·定义禁、六·待核；反口径·六}{同 \#7，抄三条义项与页码}
\source{69}{Aczél 2015 样本数（154）与三类各占比例}{Aczél 等 2015, “What is stupid? People's conception of unintelligent behavior”, Intelligence 53: 51–58, doi:10.1016/j.intell.2015.08.010；比例与样本只见二手报道（BPS Research Digest、444.hu 等）}{“We found that people use the label stupid for three separate types of situation”}{}{有把握}{正口径·数据禁、六·把握；反口径·数14禁、六}{图书馆打开 Intelligence 53: 51–58 全文，核样本与评分表}
\source{70}{Stanovich：理性障碍＝智力足够却不能理性行事}{Stanovich 2009, “Rational and Irrational Thought”, Sci Am Mind 20(6): 34, doi:10.1038/scientificamericanmind1109-34；首提 Stanovich 1993, J Learn Disabil, doi:10.1177/002221949302600803}{“I coined the term “dysrationalia” (analogous to “dyslexia”), meaning the inability to think and behave rationally despite having adequate intelligence, to draw attention to a large domain of …”}{2026-10-11}{已核实}{正驳21、45；正口径·定义、数据；反口径·定2、数36}{}
\source{71}{7 项研究：多数思维偏差与认知能力无关}{Stanovich \& West 2008, On the relative independence of thinking biases and cognitive ability. JPSP, 94(4), 672–695. DOI 10.1037/0022-3514.94.4.672}{“In 7 different studies, the authors observed that a large number of thinking biases are uncorrelated with cognitive ability.”}{2026-10-11}{已核实}{正驳21；反驳8；正口径·数据}{}
\source{72}{我方偏差与智力几乎无关}{Stanovich 等 2013, Myside bias, rational thinking, and intelligence. Curr Dir Psychol Sci, 22(4), 259–264. DOI 10.1177/0963721413480174}{“The magnitude of the myside bias shows very little relation to intelligence.”}{2026-10-11}{已核实}{反驳8}{}
\source{73}{Cipolla（讽刺文）：蠢人损人不利己}{Carlo M. Cipolla, The Basic Laws of Human Stupidity, Whole Earth Review, No. 54, Spring 1987, pp. 2–7；单行本 WH Allen 2019 \url{https://archive.org/details/wholeearthreview00unse\_9}}{“A stupid person is a person who causes losses to another person or to a group of persons while himself deriving no gain and even possibly incurring losses.”}{2026-10-11}{已核实}{正驳27；正口径·案例}{}
\source{74}{朋霍费尔：愚蠢是比恶意更危险的善之敌}{Bonhoeffer〈Nach zehn Jahren〉「Von der Dummheit」，Widerstand und Ergebung, Chr. Kaiser 1951, S. 17–20；Teuffel 全文转录 \url{https://jochenteuffel.com/wp-content/uploads/2018/06/bonhoeffer-nach-zehn-jahren.pdf}}{“Dummheit ist ein gefährlicherer Feind des Guten als Bosheit.”}{2026-10-11（＊）}{已核实}{反驳24；反口径·论10禁、案8、六·补核}{}
\source{75}{朋霍费尔：愚蠢不是智力缺陷，是人的缺陷}{Bonhoeffer〈Nach zehn Jahren〉「Von der Dummheit」，Widerstand und Ergebung, Chr. Kaiser 1951, S. 17–20；Teuffel 全文转录 \url{https://jochenteuffel.com/wp-content/uploads/2018/06/bonhoeffer-nach-zehn-jahren.pdf}}{“Soviel ist sicher, daß sie nicht wesentlich ein intellektueller, sondern ein menschlicher Defekt ist.”}{2026-10-11（＊）}{已核实}{正驳21、27、35、45；反驳6、20；正口径·定义、案例、六·待办；反口径·定2、数36、案8、六·补核}{}
\source{76}{朋霍费尔：被推向蔑视人的危险很大，我们无权如此}{Bonhoeffer〈Nach zehn Jahren〉「Von der Dummheit」，Widerstand und Ergebung, Chr. Kaiser 1951, S. 17–20；Teuffel 全文转录 \url{https://jochenteuffel.com/wp-content/uploads/2018/06/bonhoeffer-nach-zehn-jahren.pdf}}{“Wer einen Menschen verachtet, wird niemals etwas aus ihm machen können.”}{2026-10-11（＊）}{已核实}{反淘；反驳20；反口径·案8、六·补核}{}
\source{77}{穆齐尔：骂「蠢」是回应的开始也是结束}{Robert Musil, Über die Dummheit, Wien 1937 \url{https://projekt-gutenberg.org/authors/robert-musil/books/ueber-die-dummheit/chapter/2}}{“Der Gebrauch dieser Worte ist der schlichteste und der schlechteste abwehrende Ausdruck, der sich finden läßt, er ist der Anfang einer Erwiderung und schon auch ihr Ende.”}{2026-10-11}{已核实}{正淘；反驳7；正口径·案例；反口径·案16}{}
\source{78}{Sternberg 编：聪明人也会做惊人的蠢事}{Sternberg (Ed.) 2002, Why Smart People Can Be So Stupid, Yale University Press，出版社页面 \url{https://yalebooks.yale.edu/book/9780300101706/why-smart-people-can-be-so-stupid/}}{“One need not look far to find breathtaking acts of stupidity committed by people who are smart, or even brilliant.”}{2026-10-11}{已核实}{反驳8}{}
\source{79}{知识的诅咒：市场力量约减偏 50\%}{Camerer 等 1989, The curse of knowledge in economic settings: An experimental analysis. JPE, 97(5), 1232–1254. DOI 10.1086/261651 \url{https://authors.library.caltech.edu/records/zgwcn-vap73/files/1831894\_1\_.pdf?download=1}}{“Better-informed agents are unable to ignore private information even when it is in their interest to do so; more information is not always better.”}{2026-10-11}{已核实}{正口径·六·不上场；反口径·六·对方报}{}
\source{80}{〔智〕可遗传不等于不可改变（遗传度 .45→.75）}{Neisser 等 1996, “Intelligence: Knowns and unknowns”, Am Psychol 51(2): 77–101（引句在第 86 页）, 期刊版扫描 PDF \url{http://differentialclub.wdfiles.com/local--files/definitions-structure-and-measurement/Intelligence-Knowns-and-unknowns.pdf}}{“A common error is to assume that because something is heritable it is necessarily unchangeable. This is wrong. Heritability does not imply immutability.”}{2026-10-11}{已核实}{正淘；正口径·六·不上场；反口径·数36}{}
\source{81}{〔智〕贫困家庭 IQ 差异约 60\% 由共享环境解释}{Turkheimer 等 2003, Socioeconomic status modifies heritability of IQ in young children. Psychol Sci, 14(6), 623–628. DOI 10.1046/j.0956-7976.2003.psci\_1475.x}{“The models suggest that in impoverished families, 60\% of the variance in IQ is accounted for by the shared environment, and the contribution of genes is close to zero; in affluent families, the …”}{2026-10-11}{已核实}{正口径·六·不上场；反口径·数36}{}
\source{82}{〔智〕每多受一年教育约高 1–5 个 IQ 分}{Ritchie \& Tucker-Drob 2018, How much does education improve intelligence? A meta-analysis. Psychol Sci, 29(8), 1358–1369. DOI 10.1177/0956797618774253}{“Across 142 effect sizes from 42 data sets involving over 600,000 participants, we found consistent evidence for beneficial effects of education on cognitive abilities of approximately 1 to 5 IQ …”}{2026-10-11}{已核实}{正淘；反驳6；正口径·六·不上场；反口径·数36}{}
\source{83}{科学素质比例：本科以上 46.20\%，小学以下 3.11\%}{高宏斌、任磊、胡俊平、李秀菊、黄乐乐、曹金，2025，〈我国公民科学素质的发展现状与特征——第十四次中国公民科学素质抽样调查报告〉，《科普研究》20(2): 5–14，DOI 10.19293/j.cnki.1673-8357.2025.02.001 \url{https://castjournals.cast.org.cn/joweb/kpyj/EN/PDF/10.19293/j.cnki.1673-8357.2025.02.001}}{「大学本科及以上人群具备科学素质的比例达到 46.20\%」}{2026-10-11}{已核实}{反驳6；正口径·六·不上场}{}
\source{84}{79 对夫妻追踪 14 年：轻蔑只预测早离婚}{Gottman \& Levenson 2000, The Timing of Divorce: Predicting When a Couple Will Divorce Over a 14-Year Period, J Marriage Fam 62(3): 737–745, DOI 10.1111/j.1741-3737.2000.00737.x \url{https://bpl.studentorg.berkeley.edu/docs/61-Timing\%20of\%20Divorce00.pdf}}{“Thus, the negative codes that contributed to this prediction were the husband's defensiveness, contempt, and stonewalling, and the wife's criticism, contempt, and defensiveness.”}{2026-10-11}{已核实}{反驳21；反口径·数16}{}
\source{85}{Gottman 研究所博客：轻蔑是离婚最强预测因子}{The Gottman Institute 博客，2013-04-24，“The Four Horsemen: Criticism, Contempt, Defensiveness, and Stonewalling” \url{https://www.gottman.com/blog/the-four-horsemen-recognizing-criticism-contempt-defensiveness-and-stonewalling/}}{“Most importantly, contempt is the single greatest predictor of divorce.”}{2026-10-11}{已核实}{反驳21；反口径·数16}{}
\source{86}{愤怒长期和解，轻蔑导向排斥（3 项研究）}{Fischer \& Roseman 2007, Beat them or ban them: The characteristics and social functions of anger and contempt. JPSP, 93(1), 103–115. doi:10.1037/0022-3514.93.1.103}{“Anger is characterized more by short-term attack responses but long-term reconciliation, whereas contempt is characterized by rejection and social exclusion of the other person, both in the …”}{2026-10-11}{已核实}{正论二、反论二；正驳12、14；反驳3、13、21、39；正口径·论点、数据；反口径·回4、数2、六·补核}{}
\source{87}{被表达轻蔑者表现更好，攻击性也更高}{Melwani \& Barsade 2011, Held in contempt: The psychological, interpersonal, and performance consequences of… JPSP, 101(3), 503–520. doi:10.1037/a0023492}{“recipients of contempt had significantly better task performance but also significantly more interpersonal aggressiveness toward their virtual partners compared with recipients of failure, angry, or …”}{2026-10-11}{已核实}{反淘；正驳31；正口径·数据；反口径·论10禁}{}
\source{88}{〔智〕陌生人看有声录像评智力与实测相关}{Borkenau \& Liebler 1993, Convergence of stranger ratings of personality and intelligence with self-ratings,…, JPSP 65(3): 546–553, DOI 10.1037/0022-3514.65.3.546}{“Moreover, ratings of intelligence by strangers were related to targets' measured intelligence, provided that judges had been exposed to a sound film of the targets”}{2026-10-11（＊）}{已核实}{反淘；正驳21、45；正口径·六·待核、六·不上场；反口径·数36}{}
\source{89}{〔智〕600 名被试：行为薄片能预测智力}{Borkenau 等, 2004, Thin slices of behavior as cues of personality and intelligence, JPSP 86(4): 599–614, DOI 10.1037/0022-3514.86.4.599}{“Inferences of intelligence from thin slices of behavior strongly predicted intelligence test scores.”}{2026-10-11}{已核实}{反淘；正驳21、45；正口径·六·不上场；反口径·数36}{}
\source{90}{Ross 1977 命名「基本归因错误」（定义原句）}{Ross 1977, The Intuitive Psychologist And His Shortcomings: Distortions in the Attribution Process, Adv Exp Soc Psychol Vol. 10: 173–220, Academic Press, DOI 10.1016/S0065-2601(08)60357-3}{“The sources of oversight, error, or bias in his assumptions and procedures may have serious consequences, both for the lay psychologist himself and for the society that he builds and perpetuates.”}{}{有把握}{正口径·六·把握；反口径·六}{Adv Exp Soc Psychol Vol. 10 第 183 页附近抄定义原句}
\source{91}{〔智〕65\% 美国人自认比一般人聪明（N=2,821）}{Heck 等, 2018, 65\% of Americans believe they are above average in intelligence: Results of two nationally representative surveys, PLoS ONE 13(7): e0200103, DOI 10.1371/journal.pone.0200103}{“After weighting each sample to match the demographics of U.S. census data, we found that 65\% of Americans believe they are smarter than average, with men more likely to agree than women.”}{2026-10-11}{已核实}{正淘；反驳8；正口径·六·不上场；反口径·数36}{}
\source{92}{〔智〕智力自评与实测相关 r=.33}{Freund \& Kasten 2012, How smart do you think you are? A meta-analysis on the validity of self-estimates of…, Psychol Bull 138(2): 296–321, DOI 10.1037/a0026556}{“The overall relationship is estimated at r = .33.”}{2026-10-11}{已核实}{正淘；反驳8；正口径·六·不上场；反口径·数36}{}
\source{93}{人们叫作蠢的三类情境；责任、后果越重越蠢}{Aczél 等 2015, “What is stupid? People's conception of unintelligent behavior”, Intelligence 53: 51–58, doi:10.1016/j.intell.2015.08.010；比例与样本只见二手报道（BPS Research Digest、444.hu 等）}{“The level of observed stupidity was always amplified by higher responsibility being attributed to the actor and by the severity of the consequences of the action.”}{2026-10-11}{已核实}{反驳8、12、25；正口径·定义、数据；反口径·定2、判8、数14}{}
\source{94}{〔智〕个体 IQ 解释约 80\% 群体 IQ 差异（312 人）}{Bates \& Gupta 2017, Smart groups of smart people: Evidence for IQ as the origin of collective intelligence…, Intelligence 60: 46–56, DOI 10.1016/j.intell.2016.11.004}{“Contrary to prediction, individual IQ accounted for around 80\% of group-IQ differences.”}{2026-10-11}{已核实}{反驳22；正口径·六·不上场；反口径·数36}{}
\source{95}{〔智〕22 项研究：协作过程比个体技能更能预测集体智力}{Riedl 等, 2021, Quantifying collective intelligence in human groups, PNAS 118(21): e2005737118, DOI 10.1073/pnas.2005737118}{“We also find that, overall, group collaboration process is more important in predicting CI than the skill of individual members.”}{2026-10-11}{已核实}{反驳22；正口径·六·不上场；反口径·数36}{}
\source{96}{Thompson 2014 批评「多样性胜过能力」定理}{Thompson 2014, Does Diversity Trump Ability? An Example of the Misuse of Mathematics in the Social…, Notices AMS 61(9): 1024–1030, DOI 10.1090/noti1163}{}{}{【待核实】}{正口径·六·待核；反口径·六}{Notices AMS 61(9): 1024–1030 原文}
\source{97}{健康人群愤怒敌意与冠心病 HR 1.19（25 项研究）}{Chida \& Steptoe 2009, “The association of anger and hostility with future coronary heart disease: a…”, JACC 53(11): 936–946, doi:10.1016/j.jacc.2008.11.044}{“Anger and hostility were associated with increased CHD events in the healthy population studies (combined hazard ratio [HR]: 1.19; 95\% confidence interval [CI]: 1.05 to 1.35, p = 0.008)”}{2026-10-11}{已核实}{反驳21；反口径·判10、数17}{}
\source{98}{11 项研究（N=1,204）：人因费力回避共情}{Cameron 等, 2019, Empathy is hard work: People choose to avoid empathy because of its cognitive costs, JEP: General 148(6): 962–976, DOI 10.1037/xge0000595}{“In each of 11 studies (N = 1,204) and a meta-analysis, we found a robust preference to avoid empathy, which was associated with perceptions of empathy as more effortful and aversive and less efficacious.”}{2026-10-11}{已核实}{反淘；正驳23；正口径·数据}{}
\source{99}{196 项研究：护理人员中等程度共情疲劳}{Yi 等, 2024, Prevalence and related factors of compassion fatigue among registered nurses and…, BMC Nursing 23: 956, DOI 10.1186/s12912-024-02597-2}{“Both registered nurses and nursing students experience moderate compassion fatigue.”}{2026-10-11}{已核实}{反淘；正驳23；正口径·数据}{}
\source{100}{Kuppens：低学历者比其他群体更不受喜欢}{Kuppens 等, 2018, “Educationism and the irony of meritocracy: Negative attitudes of higher educated people…”, JESP 76: 429–447, doi:10.1016/j.jesp.2017.11.001}{“Finally, the less educated were liked less than any other group (all four mean differences > .50 and ps < .031).”}{2026-10-11}{已核实}{反驳9}{}
\source{101}{推特 7,331 人：道德愤怒被点赞后更常表达}{Brady 等, 2021, How social learning amplifies moral outrage expression in online social networks, Sci Adv 7(33): eabe5641, DOI 10.1126/sciadv.abe5641}{“we find that positive social feedback for outrage expressions increases the likelihood of future outrage expressions, consistent with principles of reinforcement learning.”}{2026-10-11}{已核实}{正驳7；反驳1、14、24、35；正口径·数据；反口径·论1、数5}{}
\source{102}{厌蠢问卷或比例数据：检索未果}{检索记录：中国青年报 cyol.com 站内搜索「厌蠢」、「中国青年报社会调查中心 厌蠢」、「厌蠢症 调查 受访者」、「中国青年报社会调查中心 情绪价值 / 耐心」等，2026-10-11}{}{}{【待核实】}{正口径·六·待核；反口径·六、六·对方报}{中青报客户端、「中青报社会调查」公众号、知网检索「厌蠢」}
\source{103}{美国每年约 251,454 人死于医疗差错（推算）}{Johns Hopkins Medicine 新闻稿，2016-05-03，Study Suggests Medical Errors Now Third Leading Cause of Death in the U.S. \url{https://www.hopkinsmedicine.org/news/media/releases/study\_suggests\_medical\_errors\_now\_third\_leading\_cause\_of\_death\_in\_the\_us}}{“Then, using hospital admission rates from 2013, they extrapolated that based on a total of 35,416,020 hospitalizations, 251,454 deaths stemmed from a medical error”}{}{有把握}{正口径·数据禁、六·把握；反口径·六}{2016 年存档快照或 BMJ 2016;353:i2139，核 251,454、35,416,020、9.5\%}
\source{104}{IOM 1999：每年 44,000–98,000 人死于医疗差错，多数不是烂苹果}{Institute of Medicine, 1999-11, To Err Is Human: Building a Safer Health System \url{https://www.kff.org/wp-content/uploads/sites/2/2013/04/to\_err\_is\_human\_1999\_\_report\_brief.pdf}}{“One of the report’s main conclusions is that the majority of medical errors do not result from individual recklessness or the actions of a particular group--this is not a “bad apple” problem.”}{2026-10-11}{已核实}{反驳18；正口径·数据}{}
\source{105}{Reason：系统取向；约 90\% 质量失误被判无可指责}{Reason 2000, Human error: models and management, BMJ 320(7237): 768–770, DOI 10.1136/bmj.320.7237.768}{“The basic premise in the system approach is that humans are fallible and errors are to be expected, even in the best organisations.”}{2026-10-11}{已核实}{反驳18；正口径·数据、六·不上场}{}
\source{106}{火星气候轨道器：英制单位，流程没拦住}{NASA, Mars Climate Orbiter Mishap Investigation Board, 1999-11-10, Phase I Report \url{https://llis.nasa.gov/llis\_lib/pdf/1009464main1\_0641-mr.pdf}}{“The MCO MIB has determined that the root cause for the loss of the MCO spacecraft was the failure to use metric units in the coding of a ground software file, “Small Forces,” used in trajectory models.”}{2026-10-11}{已核实}{正口径·数据}{}
\source{107}{N=3,446：信假新闻是懒得想，不是党派偏见}{Pennycook \& Rand 2019, Lazy, not biased: Susceptibility to partisan fake news is better explained by lack of…, Cognition 188: 39–50, DOI 10.1016/j.cognition.2018.06.011}{“We find that CRT performance is negatively correlated with the perceived accuracy of fake news, and positively correlated with the ability to discern fake news from real news - even for headlines …”}{2026-10-11}{已核实}{正驳34；正口径·数据}{}
\source{108}{轻提醒准确性，分享新闻质量上升}{Pennycook 等, 2021, Shifting attention to accuracy can reduce misinformation online, Nature 592: 590–595, DOI 10.1038/s41586-021-03344-2}{“Together with additional computational analyses, these findings indicate that people often share misinformation because their attention is focused on factors other than accuracy”}{2026-10-11}{已核实}{正驳34；正口径·数据}{}
\source{109}{5 个实验、10,100 余人、52 个议题：纠正无逆火}{Wood \& Porter 2019, The Elusive Backfire Effect: Mass Attitudes’ Steadfast Factual Adherence, Polit Behav 41: 135–163, DOI 10.1007/s11109-018-9443-y}{“Across all experiments, we found no corrections capable of triggering backfire, despite testing precisely the kinds of polarized issues where backfire should be expected.”}{2026-10-11}{已核实}{正淘；正驳22、34；反驳22；正口径·论点、数据；反口径·论5、数23}{}
\source{110}{能惩罚搭便车者时合作繁荣}{Fehr \& Gächter 2002, Altruistic punishment in humans, Nature 415(6868): 137–140, DOI 10.1038/415137a}{“The evidence indicates that negative emotions towards defectors are the proximate mechanism behind altruistic punishment.”}{2026-10-11}{已核实}{反淘；正驳22；正口径·数据；反口径·论5禁、论10禁}{}
\source{111}{《奇葩说》2016：高晓松「坏人不是蠢人」}{澎湃新闻·有戏，2016-03-08，《〈奇葩说〉第三季最令人期待的是……》 \url{https://m.thepaper.cn/newsDetail\_forward\_1440922}}{「第二集‘讨论跟蠢货做朋友是不是傻’」}{2026-10-11（＊）}{已核实}{反淘；正驳23、27、35；反驳25；正口径·案例、六·把握、六·不上场；反口径·案10}{}
\source{112}{2023「鸣辩仙林」辩题：厌蠢症能推动社会进步吗}{辩论吧转载，原载微信公众号「南森校辩论队」，2023-04-16/18，《第十届“鸣辩仙林”高校辩论赛辩题》 \url{https://www.bianlunba.cn/post/12572.html}}{「厌蠢症”能推动社会进步吗」}{2026-10-11（转载页＊）}{已核实}{正驳31}{}
\source{113}{新国辩 2026 复赛：北大持反方，战报记 4:5 胜}{微信公众号「北大辩协」，2026-02-24，《战报丨北大辩协斩获第十三届“新国辩”季军》 \url{https://mp.weixin.qq.com/s?src=11&timestamp=1791739372&ver=7020&signature=PmITnXz60dhNNwDCjbc3Vg3YE6C*vJzQpr0c8g6-} \url{OJR84EHMunoECHRxp9KQNYf6gWOe2stP2jUzHIj09t0NSdDdDa9pvS6kVXWMkNg6c-Uh9YoWKknfENNCrSof9I67&new=1}（链接过长，分两段排，复制时两段直接相连）}{「复赛 辩题 我“厌蠢”，我错了/我没错 持方：反」}{2026-10-11（＊）}{已核实}{正驳34；正口径·六·不上场；反口径·案17}{}
\source{114}{新国辩反方结辩视频标题：厌蠢的人大多被无知伤害过}{微信公众号「辩论」，2026-01-25，视频帖《厌蠢的人，大多真的被无知伤害过》 \url{https://mp.weixin.qq.com/s?src=11&timestamp=1791739372&ver=7020&signature=9d-cVy3O8gDqtyPcaLqSpYZGWgGCmA0uDYdc7F9NLXRn5Jo0CFhUp6*Ij1-} \url{uUa0baaPYydb7-66rKbCnvOHwU6s6VUMglf2C*lLUWvybPWs2ZgOtV7ZzwzPUvLJ3qP84&new=1}（链接过长，分两段排，复制时两段直接相连）}{「对阵：澳门大学VS北京大学 辩题：我厌蠢，我错了/没错 环节：反方四辩陈晓彤结辩」}{2026-10-11（只核标题）}{已核实}{正驳39；反驳17；正口径·六·不上场}{}
\source{115}{吉林工商学院辩论赛：厌蠢是对认知惰性的批判}{微信公众号「吉林工商学院旅游文化传播学院」，2026-03-28，《旅游文化传播学院“以学心听 以公心辩”大学生辩论赛初赛圆满结束》 \url{https://mp.weixin.qq.com/s?src=11&timestamp=1791740257&ver=7020&signature=XIfd00-D24xD0gbjJMMyUHN1BxtLQDR4F53Z7ETYU-aPiPPg51sdHJbpU*} \url{m35q768JK6YD7qKtqWquMGfbWgJSjgVg8i33WcMkaHaphrSWTQoN6pQyXBffc7AOSIp1Uk&new=1}（链接过长，分两段排，复制时两段直接相连）}{「厌蠢是对认知惰性的理性批判，是对自我成长的正向坚守」}{2026-10-11}{已核实}{正驳34}{}
\source{116}{CHD 2026 高中组辩题：我厌蠢我错了/没错}{微信公众号「CHD辩论」，2026-10-01，《云游CHD·地区赛辩题公示！》 \url{https://mp.weixin.qq.com/s?src=11&timestamp=1791739284&ver=7020&signature=0Z5IvgAGXW29m5kqIX7ZRie8ani-lIcd-W5Ljb4*IDtTMrjSIPRIlcIo-} \url{GuerdiYwEo0kqBRCCT3wjbEDxX7qnkkWk5VaaWmT3zo55ejjivqdQDWHC-QRr-*TI2U2Uiw&new=1}（链接过长，分两段排，复制时两段直接相连）}{「第三轮：在当代，我“厌蠢”我错了/我没错」}{2026-10-11}{已核实}{正驳28}{}
\source{117}{中大附中国旗下讲话：蠢没有明确标准}{中山大学附属中学，高二 1 班黄彦灵，2023-05-22，《【国旗下讲话】我不喜欢“厌蠢症”这个词》 \url{https://fz.sysu.edu.cn/article/5323}}{「“蠢”并没有明确的标准」}{2026-10-11}{已核实}{正驳37；反驳15；反口径·案14}{}
\source{118}{新京报 2023-02-25：请不要轻易自称厌蠢}{新京报（官方账号），2023-02-25 \url{https://news.sina.cn/sh/2023-02-25/detail-imyhwhwx1603244.d.html}}{「而所谓的“厌蠢症”，则是指看不惯别人没常识、太笨。」}{2026-10-11（转载页）}{已核实}{正驳37；反驳15；反口径·案14、六·对方报}{}
\source{119}{小张教坐高铁视频，留言「这个也要教？」}{扬子晚报，2023-02-15，《教人坐飞机、一个人看病 女孩拍“社会生存科普视频”涨粉百万》 \url{https://www.chinanews.com.cn/sh/2023/02-15/9953674.shtml}}{「坐高铁不是谁都会吗？这个也要教？」}{2026-10-11（转载页）}{已核实}{正论一；正驳26、37；反驳14；正口径·案例}{}
\source{120}{工人日报：小事可能是进城农民工的知识盲区}{工人日报，2023-03-03，《教人如何坐高铁、乘地铁 融入城市“第一课”火爆网络》 \url{https://www.chinanews.com.cn/sh/2023/03-03/9964177.shtml}}{「这些大众习以为常的“小事”，可能就是进城农民工的知识盲区。」}{2026-10-11（转载页）}{已核实}{反驳14；反口径·定6}{}
\source{121}{中青报 1334 人：87.6\% 新人有尴尬期，51.8\% 不懂潜规则}{中国青年报，2023-05-24，《刚入职没事做没人教怎么办》 \url{https://zqb.cyol.com/html/2023-05/24/nw.D110000zgqnb\_20230524\_1-03.htm}}{「对1334名职场新人进行的一项调查显示，87.6\%的受访职场新人经历过尴尬期。」}{2026-10-11}{已核实}{正驳37；反驳14；正口径·六·不上场；反口径·数37、六·对方报}{}
\source{122}{硬核读书会：厌蠢把焦点从事情转到「愚蠢」}{微信公众号「硬核读书会」，2023-12-23，《厌蠢情绪背后的“傲慢与偏见”》 \url{https://www.huxiu.com/article/2459024.html}}{「在厌蠢的强大逻辑下，焦点马上被转移了，人们从关注事情本身转移到关注“愚蠢”。」}{2026-10-11（转载页）}{已核实}{正驳32、43；反驳7、29；反口径·案16}{}
\source{123}{真故研究室：职场流行「蠢人」人设}{微信公众号「真故研究室」，2024-07-07，《蠢人设，开始侵蚀职场》 \url{https://m.huxiu.com/article/3218080.html}}{「现在上网，总能看到一张张所谓的职场蠢人头像」}{2026-10-11（转载页）}{已核实}{正驳42；反驳5；正口径·案例；反口径·案15}{}
\source{124}{自媒体（有人认为）：蠢人是最坏的恶人}{腾讯新闻·有画说（自媒体），2025-03-02，《为什么说蠢就是坏？蠢人就是恶人，请远离你身边的蠢人！》 \url{https://news.qq.com/rain/a/20250227A089IG00}}{「最危险的人不是坏人，而是你身边的蠢人！」}{2026-10-11（仅话语样本）}{已核实}{正驳27；反驳24；正口径·案例；反口径·判5、案13}{}
\source{125}{爱范儿：「智商税」的傲慢用法}{爱范儿（ifanr），2022-04-10，《「智商税」这个词被用烂了，有人说想做个快乐的「傻瓜」｜ifanrTalllk》 \url{https://www.ifanr.com/1481605}}{「智商税这个词最早来源于 2013 年的知乎，当时就有人说彩票是一种「低智商税」。」}{2026-10-11（＊）}{已核实}{正驳43；反驳9；正口径·案例、六·待核}{}
\source{126}{时代周报：养生骗局，不该嘲笑老人}{时代周报，2025-12-16 \url{https://www.tfcaijing.com/article/page/34336a423769446257536f56334b59465134505344673d3d}}{「老年人沦为诈骗目标，根源从不在自身，而在需求被精准利用、认知被谣言裹挟」}{2026-10-11}{已核实}{正驳35；反驳9、18}{}
\source{127}{「别和蠢人争论」不是马克·吐温说的}{Quote Investigator, 2023-01-29, “Never Argue With Stupid People. They Will Drag You Down To Their Level and Then Beat You With Experience” \url{https://quoteinvestigator.com/2023/01/29/never-argue/}}{“QI has been unable to find substantive evidence crediting this remark to Mark Twain or George Carlin.”}{2026-10-11}{已核实}{正驳41；反驳24；正口径·案例；反口径·案13}{}
\source{128}{《箴言》26:4–5 两句并列}{《圣经·箴言》26:4–5，和合本与英王钦定本（KJV），据 Bible Gateway，2026 年访问 \url{https://www.biblegateway.com/passage/?search=Proverbs\%2026\%3A4-5&version=CUVS;KJV}}{「不要照愚昧人的愚妄话回答他，恐怕你与他一样。」}{2026-10-11}{已核实}{正驳41；正口径·案例}{}
\source{129}{《箴言》14:7 KJV 与和合本}{《圣经·箴言》14:7，CUVS 与 KJV，据 Bible Gateway，2026 年访问 \url{https://www.biblegateway.com/passage/?search=Proverbs\%2014\%3A7&version=CUVS;KJV}}{「到愚昧人面前，不见他嘴中有知识。」}{2026-10-11}{已核实}{正驳41；正口径·案例}{}
\source{130}{汉隆剃刀：能用愚蠢解释的别归于恶意}{Wikipedia，「Hanlon's razor」条目，2026 年访问 \url{https://en.wikipedia.org/wiki/Hanlon\%27s\_razor}}{“Never attribute to malice that which is adequately explained by stupidity.”}{2026-10-11（二手汇编）}{已核实}{正驳20、27、41；反驳25；正口径·案例；反口径·案10}{}
\source{131}{席勒那句是濒死败将的诅咒}{Friedrich Schiller，《Die Jungfrau von Orleans》（1801），第三幕第六场，Project Gutenberg 电子书 \#6383，2026 年访问 \url{https://www.gutenberg.org/cache/epub/6383/pg6383.txt}}{“TALBOT. Unsinn, du siegst und ich muß untergehn! / Mit der Dummheit kämpfen Götter selbst vergebens.”}{2026-10-11}{已核实}{正驳41；正口径·案例}{}
\source{132}{「宇宙和人类的愚蠢」大概不是爱因斯坦说的}{Quote Investigator，「Two Things Are Infinite: the Universe and Human Stupidity」，2010-05-04 \url{https://quoteinvestigator.com/2010/05/04/universe-einstein/}}{“In conclusion, this quotation has evolved over time in the books of Frederick S. Perls.”}{2026-10-11}{已核实}{正驳41；正口径·案例}{}
\source{133}{朋霍费尔〈论愚蠢〉中文流传版}{腾讯新闻·群学书院，2021-09-15，《对于善来说，愚蠢是比恶意更加危险的敌人》 \url{https://news.qq.com/rain/a/20210915A03QIJ00}}{「对于善来说，愚蠢是比恶意更加危险的敌人。」}{2026-10-11（只核流传版＊）}{已核实}{正驳27、41；反驳20、24；正口径·案例；反口径·案8}{}
\source{134}{「可以原谅坏人，不能原谅蠢人」出处}{无可靠出处；样本见 \#124，以及知乎专栏《一个人可以坏，但不能蠢》，2026 年访问 \url{https://zhuanlan.zhihu.com/p/1938633565622473645}}{}{}{【待核实】}{正驳41；正口径·案例禁、六·待核}{毛姆作品全文检索；奇波拉中译本}
\source{135}{颜回「三八二十三」不在《论语》}{(a) 《论语》全二十篇，维基文库，2026 年访问；(b) 台北市孔庙管理委员会，「三八就是二十三」，资料更新于民国 108 年（2019） \url{https://zh.wikisource.org/wiki/\%E8\%AB\%96\%E8\%AA\%9E}}{「若說三八二十四是對的，他輸了，那可是一條人命！你說冠重要還是人命重要？」}{2026-10-11}{已核实}{正驳41；正口径·案例}{}
\source{136}{suffer fools gladly 出自保罗的反讽}{(a) 《圣经·哥林多后书》11:19，Bible Gateway；(b) Wikipedia，「Suffer fools gladly」条目，2026 年访问 \url{https://www.biblegateway.com/passage/?search=2\%20Corinthians\%2011\%3A19&version=CUVS;KJV}}{「你们既是精明人，就能甘心忍耐愚妄人。」}{2026-10-11（二手汇编）}{已核实}{正驳41；正口径·案例}{}
\source{137}{《晋书》：何不食肉糜}{房玄龄等，《晋书》卷四〈惠帝纪〉，维基文库，2026 年访问 \url{https://zh.wikisource.org/wiki/\%E6\%99\%89\%E6\%9B\%B8/\%E5\%8D\%B7004}}{「及天下荒亂，百姓餓死，帝曰：「何不食肉糜？」其蒙蔽皆此類也。」}{2026-10-11}{已核实}{反驳6；反口径·案9}{}
\source{138}{冯小刚「垃圾观众论」与戴锦华反驳}{(a) 界面新闻，任珊，2017-06-21，《冯小刚认为垃圾观众造就垃圾电影，他说得对吗？》；(b) 每日经济新闻，王嘉琦，2017-06-19，《国产烂片多，你我得背锅？冯小刚一席话，把网友们惹急了》 \url{https://www.jiemian.com/article/1412977.html}}{「有很多垃圾观众，才形成了这么多的垃圾电影。」}{2026-10-11}{已核实}{正驳38；反驳18；反口径·案11}{}
\source{139}{Torvalds 2018 道歉：轻率攻击 not OK}{Torvalds 2018-09-16, “Linux 4.19-rc4 released, an apology, and a maintainership note”, LKML（印第安纳大学存档） \url{https://lkml.iu.edu/hypermail/linux/kernel/1809.2/00117.html}}{“In my quest for a better patch, this made sense to me. I know now this was not OK and I am truly sorry.”}{2026-10-11}{已核实}{正驳31、38；反驳19；正口径·案例；反口径·案7}{}
\source{140}{费曼：讲不到新生水平说明我们不理解}{David L. Goodstein，「Richard P. Feynman, Teacher」，Physics Today 42(2):70–75，1989 年 2 月，DOI 10.1063/1.881195；AIP 官方 PDF 第 75 页（PDF 第 6 页） \url{https://physicstoday.aip.org/features/richard-p-feynman-teacher}}{“I couldn't reduce it to the freshman level. That means we really don't understand it.”}{2026-10-11}{已核实}{反驳10；反口径·案16}{}
\source{141}{费曼在巴西：学生背下一切却不懂意思}{Richard P. Feynman，《Surely You're Joking, Mr. Feynman!》。本次只看到 elmerehbi.com 等二手转引，2026 年访问 \url{https://elmerehbi.com/2014/12/20/richard-feynman-on-education-in-brazil/}}{“After a lot of investigation, I finally figured out that the students had memorized everything, but they didn't know what anything meant.”}{}{有把握}{反驳18；反口径·六、六·对方报}{原书「O Americano, Outra Vez!」一章}
\source{142}{莎莉文绝望时放下，一有机会重新开始}{Helen Keller，《The Story of My Life》（1903），第四章及所附 Anne Sullivan 书信，Project Gutenberg \#2397 \url{https://www.gutenberg.org/cache/epub/2397/pg2397.txt}}{“In despair she had dropped the subject for the time, only to renew it at the first opportunity.”}{2026-10-11}{已核实}{反论二；正驳15、30；反驳10；正口径·案例；反口径·案3}{}
\source{143}{特教教师：有的孩子起跑就会慢一点}{光明日报，2024-06-11，《爱与教，为每个特殊孩子量身定制——走近特殊教育青年教师》，光明网 \url{https://news.gmw.cn/2024-06/11/content\_37372449.htm}}{「有的孩子，起跑就是会慢一点的，起跑方式也不一样。值得！」}{2026-10-11}{已核实}{反驳10；反口径·案16}{}
\source{144}{陶行知：你的讥笑中有爱迪生}{陶行知。转引一：张春生，《今天我们为什么纪念陶行知》，原载光明日报 2013-10-18；转引二：搜狐「语文亦国学」2019-03-04 \url{http://theory.people.com.cn/n/2013/1018/c49157-23245246.html}}{「你的教鞭下有瓦特，你的冷眼里有牛顿，你的讥笑中有爱迪生」}{}{有把握}{反驳10；反口径·六、六·对方报}{《陶行知全集》〈糊涂的先生〉；《师范生》1931 年第 2 期；知网}
\source{145}{鲁迅：傻子是唯一动手砸墙开窗的人}{鲁迅，〈聪明人和傻子和奴才〉，1925-12-26，《野草》，维基文库 \url{https://zh.wikisource.org/wiki/\%E8\%81\%B0\%E6\%98\%8E\%E4\%BA\%BA\%E5\%92\%8C\%E5\%82\%BB\%E5\%AD\%90\%E5\%92\%8C\%E5\%A5\%B4\%E6\%89\%8D}}{「儍子跟奴才到他屋外，動手就砸那泥牆。」}{2026-10-11}{已核实}{反驳8；反口径·案16}{}
\source{146}{苏格拉底：不知却以为知}{Plato，《Apology》，Benjamin Jowett 英译，Project Gutenberg \#1656，2026 年访问 \url{https://www.gutenberg.org/cache/epub/1656/pg1656.txt}}{“for he knows nothing, and thinks that he knows; I neither know nor think that I know.”}{2026-10-11}{已核实}{正驳35；反驳8；正口径·案例；反口径·案16}{}
\source{147}{教师十项准则第五条：诲人不倦，不得歧视侮辱学生}{教育部，教师〔2018〕16 号，落款 2018-11-08，《教育部关于印发〈新时代高校教师职业行为十项准则〉〈新时代中小学教师职业行为十…》 \url{http://www.moe.gov.cn/srcsite/A10/s7002/201811/t20181115\_354921.html}}{「五、关心爱护学生。严慈相济，诲人不倦，真心关爱学生，严格要求学生，做学生良师益友……」}{2026-10-11}{已核实}{反驳10；反口径·案12}{}
\source{148}{孙思邈：华夷愚智，普同一等}{孙思邈，《备急千金要方》卷一〈大医精诚〉，维基文库，2026 年访问 \url{https://zh.wikisource.org/wiki/\%E5\%A4\%A7\%E9\%86\%AB\%E7\%B2\%BE\%E8\%AA\%A0}}{「若有疾厄来求救者，不得问其贵贱贫富，长幼妍蚩，怨亲善友，华夷愚智，普同一等，皆如……」}{2026-10-11}{已核实}{反驳10；反口径·案12}{}
\source{149}{2024 年全国居民健康素养水平 31.87\%}{新华网，2025-01-10，《2024年全国居民健康素养水平达到31.87\%》 \url{https://www.news.cn/20250110/c267287bbe7b420fa2740fe6b2b159b2/c.html}}{「2024年我国居民健康素养水平达到31.87\%，比2023年提高2.17个百分点……」}{2026-10-11}{已核实}{正驳35}{}
\source{150}{《论语》：学而不厌，诲人不倦}{《论语·述而》《论语·公冶长》，维基文库，2026 年访问 \url{https://zh.wikisource.org/wiki/\%E8\%AB\%96\%E8\%AA\%9E/\%E8\%BF\%B0\%E8\%80\%8C\%E7\%AC\%AC\%E4\%B8\%83}}{「子曰：「默而識之，學而不厭，誨人不倦，何有於我哉？」」}{2026-10-11}{已核实}{反驳5}{}
\source{151}{亚里士多德：愤怒恐惧不经选择}{Aristotle, Nicomachean Ethics, Book II, ch. 5（1105b–1106a），W. D. Ross 译，MIT Internet Classics Archive，2026 年访问 \url{https://classics.mit.edu/Aristotle/nicomachaen.2.ii.html}}{“we feel anger and fear without choice, but the virtues are modes of choice or involve choice”}{2026-10-11}{已核实}{正驳28；反驳23；正口径·判准}{}
\source{152}{亚里士多德：恰当时候、对象、方式地感受才是德性}{Aristotle, Nicomachean Ethics, Book II, ch. 5（1105b–1106a），W. D. Ross 译，MIT Internet Classics Archive，2026 年访问 \url{https://classics.mit.edu/Aristotle/nicomachaen.2.ii.html}}{“to feel them at the right times, with reference to the right objects, towards the right people, with the right motive, and in the right way”}{2026-10-11}{已核实}{反驳23；反口径·学8}{}
\source{153}{James 1884：身体变化直接跟随知觉}{William James, 1884, “What is an Emotion?”, Mind 9: 188–205，2026 年访问 \url{https://psychclassics.yorku.ca/James/emotion.htm}}{“the bodily changes follow directly the PERCEPTION of the exciting fact”}{2026-10-11（转载页）}{已核实}{正驳28}{}
\source{154}{重评降厌恶体验，压抑升交感激活（N=120）}{Gross 1998, “Antecedent- and response-focused emotion regulation: divergent consequences for…”, JPSP 74(1): 224–237, doi:10.1037//0022-3514.74.1.224}{“However, reappraisal decreased disgust experience, whereas suppression increased sympathetic activation.”}{2026-10-11（摘要）}{已核实}{正驳28、30；正口径·判准、数据}{}
\source{155}{换位效果最大 d=0.45，压抑无效 d=−0.04（306 个比较）}{Webb 等, 2012, “Dealing with feeling: a meta-analysis of the effectiveness of strategies derived from…”, Psychol Bull 138(4): 775–808, doi:10.1037/a0027600}{“… suppressing thoughts of the emotion-eliciting event did not (d(+) = -0.04 and -0.12, respectively)”}{2026-10-11（摘要）}{已核实}{正驳9、28、30；反驳21、22；正口径·判准、论点、数据；反口径·判9、数15}{}
\source{156}{Adams：愤怒不由选择仍可受责，含对他人的轻蔑}{Adams 1985, “Involuntary Sins”, Philos Rev 94(1): 3–31, doi:10.2307/2184713 \url{https://andrewmbailey.com/rma/1985c-Adams-InvoluntarySins.pdf}}{“contempt for other people, and the lack of a hearty concern for their welfare”}{2026-10-11（全文）}{已核实}{正驳11、28、38；反驳4、23；正口径·定义、判准、论点；反口径·数38}{}
\source{157}{Frankfurt：二阶欲望体现反思性自我评价}{Frankfurt 1971, “Freedom of the Will and the Concept of a Person”, J Philos 68(1): 5–20, doi:10.2307/2024717 \url{https://www.yorku.ca/dimock/documents/FrankfurtFreedomoftheWillandtheConceptofaPerson.pdf}}{“Besides wanting and choosing and being moved to do this or that, men may also want to have (or not to have) certain desires and motives.”}{2026-10-11（全文）}{已核实}{正驳30、32；正口径·定义、判准；反口径·判1}{}
\source{158}{《中庸》：发而皆中节谓之和}{《礼记·中庸》，维基文库全文，2026 年访问 \url{https://zh.wikisource.org/wiki/\%E7\%A6\%AE\%E8\%A8\%98/\%E4\%B8\%AD\%E5\%BA\%B8}}{「喜怒哀樂之未發，謂之中；發而皆中節，謂之和。」}{2026-10-11}{已核实}{反驳23；反口径·学8}{}
\source{159}{Roberts：轻蔑的四特征（整体、轻慢、比较、退出）}{Robert C. Roberts，2013-09-03，书评 “Hard Feelings: The Moral Psychology of Contempt”，NDPR；评 Macalester Bell, 2013, Hard Feelings, Oxford University Press \url{https://ndpr.nd.edu/reviews/hard-feelings-the-moral-psychology-of-contempt/}}{“It takes whole persons (rather than persons' actions or character traits) as its object; thus it is a ”globalist“ or ”totalizing“ evaluative perception of its target”}{2026-10-11（书评）}{已核实}{反驳16}{}
\source{160}{Roberts：轻蔑针对人促退出，怨恨针对行为促交涉}{Robert C. Roberts，2013-09-03，书评 “Hard Feelings: The Moral Psychology of Contempt”，NDPR；评 Macalester Bell, 2013, Hard Feelings, Oxford University Press \url{https://ndpr.nd.edu/reviews/hard-feelings-the-moral-psychology-of-contempt/}}{“Whereas contempt focuses on a person, attributes ”badbeing,“ and motivates withdrawal, resentment focuses on an act, attributes wrongdoing, and motivates engagement with the target.”}{2026-10-11}{已核实}{正驳13、18、22、25；反驳22；正口径·定义、论点；反口径·判15禁、论5、学2}{}
\source{161}{Roberts：轻蔑四特征各引一条道德反对}{Robert C. Roberts，2013-09-03，书评 “Hard Feelings: The Moral Psychology of Contempt”，NDPR；评 Macalester Bell, 2013, Hard Feelings, Oxford University Press \url{https://ndpr.nd.edu/reviews/hard-feelings-the-moral-psychology-of-contempt/}}{“Since persons are never, or almost never, contemptible in every respect, contempt's globalist character condemns it to fail to fit its object: it (nearly) always misrepresents its target.”}{2026-10-11}{已核实}{反驳16；反口径·判15禁}{}
\source{162}{亚里士多德：大度者轻视有据，多数人随意}{Aristotle, Nicomachean Ethics, Book IV, ch. 3（1124b），W. D. Ross 译，MIT Internet Classics Archive，2026 年访问 \url{https://classics.mit.edu/Aristotle/nicomachaen.4.iv.html}}{“For the proud man despises justly (since he thinks truly), but the many do so at random.”}{2026-10-11}{已核实}{正驳40；反驳30；正口径·定义；反口径·学9}{}
\source{163}{Goffman：污名把人贬为被打折扣的人}{Goffman 1963, Stigma: Notes on the Management of Spoiled Identity, ch. 1, p. 3 \url{https://moodle.brynmawr.edu/pluginfile.php/24292/mod\_resource/content/2/Goffman.\%20Chapter\%201\%20Stigma\%20and\%20Social\%20Identity.pdf}}{“He is thus reduced in our minds from a whole and usual person to a tainted, discounted one.”}{2026-10-11}{已核实}{反驳7；反口径·案16}{}
\source{164}{SEP：契合不同于道德或策略上的恰当}{Andrea Scarantino \& Ronald de Sousa, 2018, “Emotion”, SEP §10 \url{https://plato.stanford.edu/entries/emotion/}}{“As D'Arms and Jacobson (2000) have emphasized, fittingness is importantly different from forms of appropriateness that are moral or strategic.”}{2026-10-11}{已核实}{正驳40；正口径·定义、判准；反口径·判6}{}
\source{165}{Mason 编：轻蔑该培养还是远离是开放问题}{Michelle Mason (ed.), 2018, The Moral Psychology of Contempt, Rowman \& Littlefield International, doi:10.5040/9798881816971}{“It invites the general reader to reflect on whether contempt is something to be embraced and cultivated as an emotional safeguard of valued norms or, rather, an emotion from which we have good reason …”}{2026-10-11（出版商简介）}{已核实}{正驳22；正口径·论点}{}
\source{166}{Mason 2003 为针对人的轻蔑辩护}{Mason 2003, “Contempt as a Moral Attitude”, Ethics 113(2): 234–272, doi:10.1086/342860}{}{}{【待核实】}{正口径·六·待核；反口径·六}{PhilPapers「MASCAA」；Ethics 113(2): 234–272}
\source{167}{小张评论区「我厌蠢症犯了」；央视网报道让它出圈}{上游新闻，2023-02-24，《网红教人如何坐公交高铁引来冷嘲热讽，“厌蠢症”是种什么病？》 \url{https://news.cqnews.net/1/detail/1078733353267396608/app/content\_1078733353267396608.html}}{「……“怎么会有人不会坐飞机啊？我厌蠢症犯了。”」}{2026-10-11（原页）}{已核实}{正论一；正驳9；反驳1、35；正口径·论点、案例；反口径·案6}{}
\source{168}{那位网友原话还有「可以耐着性子教对方」}{上游新闻，2023-02-24，《网红教人如何坐公交高铁引来冷嘲热讽，“厌蠢症”是种什么病？》 \url{https://news.cqnews.net/1/detail/1078733353267396608/app/content\_1078733353267396608.html}}{「A反击说，遇见不会但努力的人，并不会激发她的厌蠢症。她甚至可以耐着性子教对方，可……」}{2026-10-11（原页）}{已核实}{正驳29；反驳5、24、29；正口径·判准；反口径·定4、案5}{}
\source{169}{Reason：有些不安全行为恶劣，但绝大多数不是}{Reason 2000, Human error: models and management, BMJ 320(7237): 768–770, DOI 10.1136/bmj.320.7237.768}{“Although some unsafe acts in any sphere are egregious, the vast majority are not.”}{2026-10-11（全文）}{已核实}{正驳39；反驳18；正口径·数据；反口径·数18}{}
\source{170}{愤怒和轻蔑常一起出现，但区分清楚（3 项研究）}{Fischer \& Roseman 2007, Beat them or ban them: The characteristics and social functions of anger and contempt. JPSP, 93(1), 103–115. doi:10.1037/0022-3514.93.1.103}{“The results of the 3 studies show that anger and contempt often occur together but that there are clear distinctions between the 2 emotions”}{2026-10-11（摘要）}{已核实}{正论二、反论二；正驳12、14；反驳3、16、39；正口径·论点、数据；反口径·数2}{}
\source{171}{Kuppens 研究 6 只报愤怒的中介，未报轻蔑}{Kuppens 等, 2018, “Educationism and the irony of meritocracy: Negative attitudes of higher educated people…”, JESP 76: 429–447, doi:10.1016/j.jesp.2017.11.001}{“Consistent with Weiner (1988), the effect of group (less educated versus poor) on anger (see Figure 6) was mediated by responsibility (indirect effect = .32, 95\% CI = [0.08, 0.60])”}{2026-10-11（录用稿）}{已核实}{反驳9；反口径·数22}{}
\source{172}{Torvalds：要学会理解别人的情绪并恰当回应}{Torvalds 2018-09-16, “Linux 4.19-rc4 released, an apology, and a maintainership note”, LKML（印第安纳大学存档） \url{https://lkml.iu.edu/hypermail/linux/kernel/1809.2/00117.html}}{“This week people in our community confronted me about my lifetime of not understanding emotions.”}{2026-10-11（存档）}{已核实}{正驳31、38；反驳19；正口径·案例；反口径·案7}{}
\source{173}{中青报调查 1334 人：即将入职 30.2\%，已入职 69.8\%}{中国青年报，2023-05-24，《刚入职没事做没人教怎么办》 \url{https://zqb.cyol.com/html/2023-05/24/nw.D110000zgqnb\_20230524\_1-03.htm}}{「……对1334名职场新人进行的一项调查」}{2026-10-11（原页）}{已核实}{正口径·六·不上场；反口径·数37}{}
\source{174}{Strawson：完全客观时能谈判，不能讲理}{Strawson 1962, “Freedom and Resentment”, Proc Br Acad 48: 1–25 \url{https://www.informationphilosopher.com/solutions/philosophers/strawson/freedom\_and\_resentment.html}}{“If your attitude towards someone is wholly objective”}{2026-10-11（转载页＊）}{已核实}{反论二；正驳12；反驳4；正口径·论点；反口径·论2、六}{}
\source{175}{受访老人 95.97\% 用智能手机（2690 份问卷）}{川观新闻，2020-10-21 15:31，《2690份问卷画像：超三成老人遭遇“智能机”困境，网络挂号就医难倒“脱网”老人》 \url{https://cbgc.scol.com.cn/news/416034}}{「直接和间接访问到老人中，有95.97\%的老人使用智能手机。」}{2026-10-11（原页）}{已核实}{反驳14；正口径·六·不上场；反口径·数24}{}
\source{176}{替陌生同龄人估字谜，后见偏差与替自己基本相同（193 人）}{Hom 2023, “Perspective-taking and hindsight bias: When the target is oneself and/or a peer”, Curr Psychol 42(16): 13987–13998, doi:10.1007/s12144-021-02413-z（PMC 全文 PMC8783185）}{“Participants faced with adopting the peer’s perspective showed essentially the same hindsight bias as themselves.”}{2026-10-11（全文）}{已核实}{正论一；正驳9、19；反驳2、22、37、43；正口径·论点、数据、六·待办；反口径·回3、数28}{}
\source{177}{重评反应 d=0.23 不如重评刺激 0.36、换位 0.45}{Webb 等, 2012, “Dealing with feeling: a meta-analysis of the effectiveness of strategies derived from…”, Psychol Bull 138(4): 775–808, doi:10.1037/a0027600}{“Finally, reappraising the emotional response proved less effective (d(+) = 0.23) than reappraising the emotional stimulus (d(+) = 0.36) or using perspective taking (d(+) = 0.45).”}{2026-10-11（摘要）}{已核实}{反口径·数15}{}
\source{178}{Ross 与 Ward 1996 提出朴素实在论}{Ross \& Ward 1996, “Naive realism in everyday life: Implications for social conflict and misunderstanding”, in E. S. Reed, E. Turiel \& T. Brown (Eds.), Values and Knowledge, Lawrence Erlbaum}{}{}{【待核实】}{反驳13；正口径·六·待核；反口径·数35禁、六、六·对方报}{纸本 Values and Knowledge（Erlbaum 1996）；Robinson 等 1995，JPSP 68(3): 404–417}
\source{179}{Weiner 1986、1995 两部专著（书目）}{Weiner 1986, An Attributional Theory of Motivation and Emotion, Springer, doi:10.1007/978-1-4612-4948-1；Weiner 1995, Judgments of Responsibility, Guilford, ISBN 0-89862-843-1}{“An Attributional Theory of Motivation and Emotion”}{2026-10-11（书目）}{已核实}{反论一；反口径·论1}{}
\source{180}{刑法第十五条：应当预见而没有预见}{《中华人民共和国刑法》（2023 年修正，2024-03-01 施行）第十五、十六条，维基文库全文转录 \url{https://zh.wikisource.org/wiki/\%E4\%B8\%AD\%E5\%8D\%8E\%E4\%BA\%BA\%E6\%B0\%91\%E5\%85\%B1\%E5\%92\%8C\%E5\%9B\%BD\%E5\%88\%91\%E6\%B3\%95\_(2023\%E5\%B9\%B4})}{「应当预见自己的行为可能发生危害社会的结果，因为疏忽大意而没有预见，或者已经预见而……」}{2026-10-11（维基文库＊）}{已核实}{反驳12；反口径·数26、六·补核}{}
\source{181}{MIT 斯隆报道：任务容易时求助不损不增能力评价}{Posner 2015-03-05, “Does Deciding to Seek Advice Signal Weakness?”, MIT Sloan Management Review（报道 Brooks 等 2015） \url{https://sloanreview.mit.edu/article/does-deciding-to-seek-advice-signal-weakness/}}{“When tasks were seen as difficult, the individual seeking advice was actually viewed more competently”}{2026-10-11（原页＊）}{已核实}{正淘、反淘；正驳26；反驳5；正口径·数据；反口径·数7}{}
\source{182}{《论语》子张：人将拒我，如之何其拒人}{《论语·子张第十九》，十九之三，维基文库本，2026 年访问 \url{https://zh.wikisource.org/wiki/\%E8\%AB\%96\%E8\%AA\%9E/\%E5\%AD\%90\%E5\%BC\%B5\%E7\%AC\%AC\%E5\%8D\%81\%E4\%B9\%9D}}{「……眾，嘉善而矜不能。』我之大賢與，於人何所不容？我之不賢與，人將拒我，如之何其拒人……」}{2026-10-11（维基文库）}{已核实}{反驳15；反口径·学6}{}
\source{183}{Reason：怪个人比追究机构痛快；最好的人也犯最糟的错}{Reason 2000, Human error: models and management, BMJ 320(7237): 768–770, DOI 10.1136/bmj.320.7237.768}{“Blaming individuals is emotionally more satisfying than targeting institutions.”}{2026-10-11（全文）}{已核实}{正淘；正驳33、39；反驳18、22；正口径·数据；反口径·数18}{}
\source{184}{心理安全对学习行为 B=.46、R²=.33（51 个团队）}{Edmondson 1999, “Psychological Safety and Learning Behavior in Work Teams”, ASQ 44(2): 350–383（JSTOR 2666999），全文 PDF \url{https://web.mit.edu/curhan/www/docs/Articles/15341\_Readings/Group\_Performance/Edmondson\%20Psychological\%20safety.pdf}}{“Team psychological safety   .46   5.26   <.001   .33”}{2026-10-11（全文）}{已核实}{正驳22、24；反驳5；正口径·论点、数据；反口径·数8}{}
\source{185}{对同一人的多次愤怒可预测轻蔑（研究 2）}{Fischer \& Roseman 2007, Beat them or ban them: The characteristics and social functions of anger and contempt. JPSP, 93(1), 103–115. doi:10.1037/0022-3514.93.1.103}{“in the 2nd study, the authors tested whether contempt could be predicted from previous anger incidents with the same person”}{2026-10-11（摘要）}{已核实}{正论二；正驳5、12–14、17、19；反驳3、39、40；正口径·论点、数据；反口径·回4、数2}{}
\source{186}{Torvalds：尤其在我搞成人身攻击时}{Torvalds 2018-09-16, “Linux 4.19-rc4 released, an apology, and a maintainership note”, LKML（印第安纳大学存档） \url{https://lkml.iu.edu/hypermail/linux/kernel/1809.2/00117.html}}{“Especially at times when I made it personal.”}{2026-10-11（存档）}{已核实}{正驳31；反驳19；正口径·案例；反口径·案7}{}
\source{187}{中青报：63.9\% 新人建议多向周围人请教}{中国青年报，2023-05-24，《刚入职没事做没人教怎么办》 \url{https://zqb.cyol.com/html/2023-05/24/nw.D110000zgqnb\_20230524\_1-03.htm}}{「顺利度过入职尴尬期，63.9\%的受访职场新人建议多向周围人请教，多和老员工交流。」}{2026-10-11}{已核实}{反驳5；反口径·数37}{}
\source{188}{心理行为类污名默认被看作可控，招愤怒（2 项实验）}{Weiner 等, 1988, “An attributional analysis of reactions to stigmas”, JPSP 55(5): 738–748, doi:10.1037/0022-3514.55.5.738}{“mental-behavioral stigmas were perceived as onset-controllable, and elicited little pity, much anger, and judgments to neglect.”}{2026-10-11（摘要）}{已核实}{正论一、反论一；正驳2、24、42、44；反驳6、14、38；正口径·论点、数据；反口径·数9、六·补核}{}
\source{189}{内核行为准则：对其他成员表现出共情}{Linux 内核 code-of-conduct.rst（提交 8a104f8b，2018-09-16），git.kernel.org \url{https://git.kernel.org/pub/scm/linux/kernel/git/torvalds/linux.git/plain/Documentation/process/code-of-conduct.rst?id=8a104f8b5867c682d994ffa7a74093c54469c11f}}{“Showing empathy towards other community members”}{2026-10-11}{已核实}{正淘、反淘；反驳19；正口径·论点、案例}{}
\source{190}{Kerr：搭档能做不做，认真者减少努力}{Kerr 1983, “Motivation losses in small groups: A social dilemma analysis”, JPSP 45(4): 819–828, doi:10.1037/0022-3514.45.4.819}{“An experiment was performed to empirically demonstrate a third motivation loss effect, termed the sucker effect.”}{2026-10-11（摘要）}{已核实}{正驳22；正口径·数据}{}
\source{191}{说话者高估自己说清楚了}{Keysar \& Henly 2002, “Speakers' overestimation of their effectiveness”, Psychol Sci 13(3): 207–212, doi:10.1111/1467-9280.00439}{“When they estimated the addressee's understanding of the intended meaning, they showed a consistent tendency to overestimate their effectiveness.”}{2026-10-11（摘要）}{已核实}{反驳2、12；反口径·数11}{}
\source{192}{子女 91.36\% 以为父母学不会，老人 70.59\% 说太贵}{川观新闻，2020-10-21 15:31，《2690份问卷画像：超三成老人遭遇“智能机”困境，网络挂号就医难倒“脱网”老人》 \url{https://cbgc.scol.com.cn/news/416034}}{「70.59\%老人因智能手机太贵而未使用，91.36\%受访子女则认为父母未使用智能……」}{2026-10-11（原页）}{已核实}{反驳14；正口径·六·不上场；反口径·数24、六·对方报}{}
\source{193}{Denrell：坏印象让人少接触，所以更稳}{Denrell 2005, “Why most people disapprove of me: Experience sampling in impression formation”, Psychol Rev 112(4): 951–978, doi:10.1037/0033-295X.112.4.951}{“Individuals are typically more likely to continue to interact with people if they have a positive impression of them.”}{2026-10-11（摘要）}{已核实}{正淘；正驳16；反驳11；正口径·数据；反口径·数12}{}
\source{194}{态度靠接近形成，坏的学得更牢}{Fazio 等, 2004, “Attitude formation through exploration: Valence asymmetries”, JPSP 87(3): 293–311, doi:10.1037/0022-3514.87.3.293}{“a learning asymmetry involving better learning for negatively valenced than positively valenced objects”}{2026-10-11（摘要）}{已核实}{反驳11}{}
\source{195}{只有接近才得反馈时，坏预期抗拒改变}{Eiser 等, 2009, “Prejudiced learning: A connectionist account”, Br J Psychol 100(Pt 2): 399–413, doi:10.1348/000712608X357849}{“Where feedback was contingent on approach, the false positive expectancies were corrected by experience, but negative prejudices resisted change, since the network avoided objects deemed to be bad, …”}{2026-10-11（摘要）}{已核实}{反驳11}{}
\source{196}{人高估自己在邮件里说清楚的程度}{Kruger 等, 2005, “Egocentrism over e-mail: Can we communicate as well as we think?”, JPSP 89(6): 925–936, doi:10.1037/0022-3514.89.6.925}{“Five experiments suggest that this limitation is often underappreciated, such that people tend to believe that they can communicate over e-mail more effectively than they actually can.”}{2026-10-11（摘要）}{已核实}{反驳2}{}
\source{197}{Strawson：客观态度不含愤怒，原因列着「白痴」}{Strawson 1962, “Freedom and Resentment”, Proc Br Acad 48: 1–25 \url{https://www.informationphilosopher.com/solutions/philosophers/strawson/freedom\_and\_resentment.html}}{“our adoption of the objective attitude is a consequence of our viewing the agent as incapacitated in some or all respects for ordinary interpersonal relationships”}{2026-10-11（转载页＊）}{已核实}{正驳12、15；反驳3、4、39、42；正口径·论点；反口径·论2、六}{}
\source{198}{网上流行「职场蠢人头像」}{微信公众号「真故研究室」，2024-07-07，《蠢人设，开始侵蚀职场》 \url{https://m.huxiu.com/article/3218080.html}}{「现在上网，总能看到一张张所谓的职场蠢人头像，号称“让同事失去沟通欲望，让老板不敢……」}{2026-10-11（原页）}{已核实}{反驳5}{}
\source{199}{道歉补偿加快原谅（337 名大学生）}{McCullough 等, 2014, “Conciliatory gestures promote forgiveness and reduce anger in humans”, PNAS 111(30): 11211–11216, doi:10.1073/pnas.1405072111}{“In a longitudinal study with 337 participants who had recently been harmed by a relationship partner, we found that conciliatory gestures (e.g., apologies, offers of compensation) were associated …”}{2026-10-11（摘要）}{已核实}{反淘；正驳16；正口径·数据}{}
\source{200}{另有一套附属系统修复关系}{McCullough 等, 2013, “Cognitive systems for revenge and forgiveness”, BBS 36(1): 1–15, doi:10.1017/S0140525X11002160}{“We posit, consequently, a subsidiary computational system designed to restore particular relationships after cost-imposing interactions by inhibiting revenge and motivating behaviors that signal …”}{2026-10-11（摘要）}{已核实}{反淘}{}
\source{201}{Rozenblit 研究 8：被试只写步骤，没去做（37 人）}{Rozenblit \& Keil 2002, “The misunderstood limits of folk science: an illusion of explanatory depth”, Cogn Sci 26(5): 521–562, doi:10.1207/s15516709cog2605\_1}{“We asked 37 participants in Study 8 about their knowledge of various real-world procedures”}{2026-10-11（全文）}{已核实}{正驳8；正口径·数据}{}
\source{202}{郑人买履全段：及反，市罢}{《韩非子·外储说左上》，维基文库本，2026 年访问 \url{https://zh.wikisource.org/wiki/\%E9\%9F\%93\%E9\%9D\%9E\%E5\%AD\%90/\%E5\%A4\%96\%E5\%84\%B2\%E8\%AA\%AA\%E5\%B7\%A6\%E4\%B8\%8A}}{「已得履，乃曰：『吾忘持度。』反歸取之。及反，市罷，遂不得履。」}{2026-10-11（维基文库）}{已核实}{反淘；正驳20；正口径·案例}{}
\source{203}{越内行越低估新手难处，去偏压不下（2 项研究）}{Hinds 1999, “The curse of expertise: The effects of expertise and debiasing methods on prediction of…”, JEP: Applied 5(2): 205–221, doi:10.1037/1076-898X.5.2.205}{“Findings from these studies suggest that experts may have a cognitive handicap that leads to underestimating the difficulty novices face”}{2026-10-11（摘要）}{已核实}{正驳3；反驳36；正口径·论点、数据；反口径·数43}{}
\source{204}{知道结果者高估不知情者所知，自己察觉不到（3 项实验）}{Fischhoff 1975, “Hindsight ≠ foresight: The effect of outcome knowledge on judgment under uncertainty”, JEP: HPP 1(3): 288–299, doi:10.1037/0096-1523.1.3.288（2003 年重印于 Qual Saf Health Care 12(4): 304–311）}{“As a result, they overestimated what they would have known without outcome knowledge (Experiment 2), as well as what others (Experiment 3) actually did know without outcome knowledge.”}{2026-10-11（摘要）}{已核实}{正论一；正驳9、19、25、34、37；反驳2、43；正口径·论点、数据；反口径·回3、数27}{}
\source{205}{先认出者高估不知情者，提醒也压不下（3 项实验）}{Harley 等, 2004, “The 'saw-it-all-along' effect: Demonstrations of visual hindsight bias”, JEP: LMC 30(5): 960–968, doi:10.1037/0278-7393.30.5.960}{“In all experiments, participants overestimated identification performance of naive observers.”}{2026-10-11（摘要）}{已核实}{正论一；正驳37；反驳2、36、43；正口径·论点、数据；反口径·回3、数31}{}
\source{206}{防洪材料：事前四分之三认为不值得防，事后判疏忽}{Kamin \& Rachlinski 1995, “Ex post ≠ ex ante: Determining liability in hindsight”, Law Hum Behav 19(1): 89–104, doi:10.1007/BF01499075}{“The debiasing procedure failed to produce any significant differences from the regular hindsight condition.”}{2026-10-11（摘要）}{已核实}{正淘；反驳12；正口径·论点禁、六·不上场；反口径·数32、六·对方报}{}
\source{207}{鼓励考虑其他原因能减少后见之明}{Roese \& Vohs 2012, “Hindsight bias”, Perspect Psychol Sci 7(5): 411–426, doi:10.1177/1745691612454303}{“an intervention that encourages people to consider alternative causal explanations for a given outcome can reduce hindsight bias.”}{2026-10-11（摘要）}{已核实}{反驳12；正口径·数据；反口径·数33}{}
\source{208}{把意见不合者看成有偏，冲突升级}{Kennedy \& Pronin 2008, “When disagreement gets ugly: Perceptions of bias and the escalation of conflict”, PSPB 34(6): 833–848, doi:10.1177/0146167208315158}{“First, these studies show that people perceive those who disagree with them as biased.”}{2026-10-11（摘要）}{已核实}{正淘；反驳13；正口径·六·不上场；反口径·数35}{}
\source{209}{朴素实在论：把看法不同者看成有偏}{Pronin 等, 2004, “Objectivity in the eye of the beholder: Divergent perceptions of bias in self versus…”, Psychol Rev 111(3): 781–799, doi:10.1037/0033-295X.111.3.781}{“it is also a product of the phenomenological stance of naive realism.”}{2026-10-11（摘要）}{已核实}{反驳13；反口径·数35}{}
\source{210}{人以为自己的知觉直接反映现实}{Pronin 2007, “Perception and misperception of bias in human judgment”, TICS 11(1): 37–43, doi:10.1016/j.tics.2006.11.001}{“The other involves people's conviction that their perceptions directly reflect reality, and that those who see things differently are therefore biased.”}{2026-10-11（摘要）}{已核实}{反驳13}{}
\source{211}{朋霍费尔：蔑视人是诱惑，我们无权这样}{Bonhoeffer〈Nach zehn Jahren〉「Von der Dummheit」，Widerstand und Ergebung, Chr. Kaiser 1951, S. 17–20；Teuffel 全文转录 \url{https://jochenteuffel.com/wp-content/uploads/2018/06/bonhoeffer-nach-zehn-jahren.pdf}}{「Die Gefahr, uns in Menschenverachtung hineintreiben zu lassen, ist sehr groß. Wir wissen wohl, daß wir kein Recht dazu haben und daß wir dadurch in das unfruchtbarste Verhältnis zu den Menschen geraten.」}{2026-10-11（转录＊）}{已核实}{正驳27、35；反驳8、20；正口径·案例、六·待办；反口径·案8、六·补核}{}
\source{212}{新京报：同学同事是「我都可以做到你为什么做不到」}{新京报书评周刊，2023-11-01，《为什么我们容易犯“厌蠢症”？》 \url{https://m.bjnews.com.cn/detail/1698812159168656.html}}{「如果再剖析一下这种愤怒背后的思维，更像是一种“我都可以做到你为什么做不到”的抱怨……」}{2026-10-11（原页）}{已核实}{正论一、反论一、反淘；正驳1、4、6、7、9、19、25、35；反驳1、34；正口径·论点、案例；反口径·回1、案4}{}
\source{213}{新京报：评论区总有「怎么会有人不会坐飞机」}{新京报书评周刊，2023-11-01，《为什么我们容易犯“厌蠢症”？》 \url{https://m.bjnews.com.cn/detail/1698812159168656.html}}{「评论区总有“怎么会有人不会坐飞机啊？我厌蠢症犯了”、“不会坐高铁？中国孩子什么时……」}{2026-10-12（原页）}{已核实}{正论一；正驳32；反驳1、14、35；正口径·论点、案例；反口径·案4}{}
\source{214}{每日商报 2024：「怎么会有人不会打车」}{每日商报，2024-01-13，《“怎么还会有这种人”“我的厌蠢症又犯了”专家：“厌蠢”的实质是一种“慕强”心态》 \url{https://jrsh.hangzhou.com.cn/read/content/2024-01/13/content\_8674352.htm}}{「“怎么会有人不会打车啊？我厌蠢症犯了。”」}{2026-10-11}{已核实}{正论一；正驳7、9、19、37；反驳1、14、35；正口径·案例；反口径·案19}{}
\source{215}{熟悉感足以造成知识的诅咒（3 项实验、359 人）}{Birch 等, 2017, “A 'curse of knowledge' in the absence of knowledge? People misattribute fluency when…”, Cognition 166: 447–458, doi:10.1016/j.cognition.2017.04.015}{“Three experiments with a total of 359 undergraduate students provide the first evidence that fluency misattribution processes are sufficient to induce the curse of knowledge bias.”}{2026-10-12（摘要）}{已核实}{正驳3、6、19；反驳2；正口径·论点、数据；反口径·回3、数29}{}
\source{216}{Nickerson：估别人从自己知道的出发，常不加批判}{Nickerson 1999, “How we know—and sometimes misjudge—what others know: Imputing one's own knowledge to…”, Psychol Bull 125(6): 737–759, doi:10.1037/0033-2909.125.6.737}{“An obvious starting point for building a model of what another knows is what one oneself knows, or thinks one knows.”}{2026-10-12（摘要）}{已核实}{正驳1、6、25、37；反驳2；正口径·数据；反口径·数30}{}
\source{217}{Hom 补句：后见者更信自己替同伴估得准（193 人）}{Hom 2023, “Perspective-taking and hindsight bias: When the target is oneself and/or a peer”, Curr Psychol 42(16): 13987–13998, doi:10.1007/s12144-021-02413-z（PMC 全文 PMC8783185）}{“Although participants in hindsight believed their foreseeable predictions for the peer were more accurate or realistic, it was more challenging to predict for the peer than themselves.”}{2026-10-12（全文）}{已核实}{正驳3；反驳37；正口径·数据；反口径·数28}{}
\source{218}{经理信人能改，辅导意愿上升（3 项研究）}{Heslin 等, 2006, “Keen to help? Managers' implicit person theories and their subsequent employee coaching”, Pers Psychol 59(4): 871–902, doi:10.1111/j.1744-6570.2006.00057.x}{“individuals holding an “entity theory” that human attributes are innate and unalterable are disinclined to invest in helping others to develop and improve, relative to individuals who hold the …”}{2026-10-12（摘要）}{已核实}{正驳15；反驳4、41；正口径·论点、数据；反口径·数34}{}
\source{219}{64 项研究、1.2 万余人：该负责生气，不可控同情}{Rudolph 等, 2004, “A meta-analytic review of help giving and aggression from an attributional perspective:…”, Cogn Emot 18(6): 815–848, doi:10.1080/02699930341000248}{“The present review syntheses 64 investigations on the determinants of helping and aggression involving more than 12,000 subjects”}{2026-10-11（摘要）}{已核实}{反论一；正驳2、44；反驳3、6、27；正口径·数据；反口径·判8、数1、六·补核}{}
\source{220}{Kerr 傻瓜效应：搭档能做不做，认真者减力}{Kerr 1983, “Motivation losses in small groups: A social dilemma analysis”, JPSP 45(4): 819–828, doi:10.1037/0022-3514.45.4.819}{“It was hypothesized that group members would reduce their efforts if they had a capable partner who free rode on their efforts, that is, who was capable of contributing to the group but would not. …”}{2026-10-11（摘要）}{已核实}{反淘；正驳22；正口径·数据}{}
\source{221}{内核准则：拒绝不合格贡献不违规，禁辱骂}{Linux 内核 code-of-conduct.rst（2018-09-16）与 code-of-conduct-interpretation.rst（2018-10-22），git.kernel.org \url{https://git.kernel.org/pub/scm/linux/kernel/git/torvalds/linux.git/plain/Documentation/process/code-of-conduct-interpretation.rst?id=79dbeed36f7335ec6432029dbf019ebce28e963e}}{“rejecting unsuitable contributions are not viewed as a violation of the Code of Conduct.”}{2026-10-11}{已核实}{正驳31；反驳19；正口径·论点、案例；反口径·案7}{}
\source{222}{心理安全定义；量表第 1 题「犯错常被记账」（51 个团队）}{Edmondson 1999, “Psychological Safety and Learning Behavior in Work Teams”, ASQ 44(2): 350–383（JSTOR 2666999），全文 PDF \url{https://web.mit.edu/curhan/www/docs/Articles/15341\_Readings/Group\_Performance/Edmondson\%20Psychological\%20safety.pdf}}{“The term is meant to suggest neither a careless sense of permissiveness, nor an unrelentingly positive affect but, rather, a sense of confidence that the team will not embarrass, reject, or punish …”}{2026-10-11（全文）}{已核实}{正驳24；反驳5；正口径·数据；反口径·数8}{}
\source{223}{Reason：公正文化要划清可责与不可责}{Reason 2000, Human error: models and management, BMJ 320(7237): 768–770, DOI 10.1136/bmj.320.7237.768}{“Although some unsafe acts in any sphere are egregious, the vast majority are not.”}{2026-10-11（全文）}{已核实}{反驳18；反口径·数18}{}
\source{224}{与自己相关只推高愤怒（5 项研究）}{Hutcherson \& Gross 2011, The moral emotions: A social-functionalist account of anger, disgust, and contempt. JPSP, 100(4), 719–737. doi:10.1037/a0022408}{“(b) an effect of self-relevance on anger but not other emotions, and (c) a role for contempt in judging incompetent actions.”}{2026-10-11（摘要）}{已核实}{正驳5；反驳1、16；正口径·论点、数据；反口径·论1、数4、六·补核}{}
\source{225}{扬子晚报：那条留言是「调侃」}{扬子晚报，2023-02-15，《教人坐飞机、一个人看病 女孩拍“社会生存科普视频”涨粉百万》 \url{https://www.chinanews.com.cn/sh/2023/02-15/9953674.shtml}}{「发第一条视频时，有人留言调侃称：“坐高铁不是谁都会吗？这个也要教？”」}{2026-10-11（原页）}{已核实}{正驳26；反驳14；正口径·案例}{}
\source{226}{孔子斥子路「野哉」，随即讲完正名}{《论语·子路第十三》十三之三，维基文库本；朱熹《四书章句集注·论语集注》卷七，维基文库本，2026 年访问 \url{https://zh.wikisource.org/wiki/\%E8\%AB\%96\%E8\%AA\%9E/\%E5\%AD\%90\%E8\%B7\%AF\%E7\%AC\%AC\%E5\%8D\%81\%E4\%B8\%89}}{「子路曰：『有是哉？子之迂也！奚其正？』子曰：『野哉，由也！君子於其所不知，蓋闕如……」}{2026-10-11（维基文库）}{已核实}{反论二；正驳18；正口径·案例；反口径·回5、案2}{}
\source{227}{求助者被看作更能干，任务难时尤甚}{Brooks 等, 2015, “Smart People Ask for (My) Advice: Seeking Advice Boosts Perceptions of Competence”, Manag Sci 61(6): 1421–1435, doi:10.1287/mnsc.2014.2054}{“We find that people are reticent to seek advice for fear of appearing incompetent. This fear, however, is misplaced.”}{2026-10-11（摘要）}{已核实}{正淘、反淘；正驳26；反驳5；正口径·数据；反口径·数6}{}
\source{228}{原谅随时间上升，三个月涨两个对数优势单位}{McCullough 等, 2010, “On the form and function of forgiving: Modeling the time-forgiveness relationship and…”, Emotion 10(3): 358–376, doi:10.1037/a0019349}{“By 3 months after their transgressions, the typical participant's forgiveness had increased by two log-odds units.”}{2026-10-11（摘要）}{已核实}{反淘；正驳16；正口径·数据、六·不上场}{}
\source{229}{原谅三参数；起初越归责越难忍让}{McCullough 等, 2003, “Forgiveness, forbearance, and time: The temporal unfolding of transgression-related…”, JPSP 84(3): 540–557, doi:10.1037/0022-3514.84.3.540}{“trend forgiveness (reductions in avoidance and revenge, and increases in benevolence)”}{2026-10-11（摘要）}{已核实}{正驳16；正口径·数据}{}
\source{230}{亚里士多德：应当知道又不难知道的无知也受罚}{Aristotle, Nicomachean Ethics III.5，D. P. Chase 英译，Project Gutenberg \#8438，2026 年访问 \url{https://www.gutenberg.org/cache/epub/8438/pg8438.txt}}{“Again, those also who are ignorant of legal regulations which they are bound to know, and which are not hard to know, they chastise; and similarly in all other cases where neglect is thought to be …”}{2026-10-11（全文＊）}{已核实}{正驳3；反驳6、25；正口径·论点、六·把握；反口径·论1、学3、六}{}
\source{231}{孟子：为长者折枝，是不为也，非不能也}{《孟子·梁惠王上》，维基文库本，2026 年访问 \url{https://zh.wikisource.org/wiki/\%E5\%AD\%9F\%E5\%AD\%90/\%E6\%A2\%81\%E6\%83\%A0\%E7\%8E\%8B\%E4\%B8\%8A}}{「為長者折枝，語人曰：『我不能。』是不為也，非不能也。」}{2026-10-11（维基文库）}{已核实}{反论一；反驳6；反口径·论1、学4}{}
\source{232}{孟子：不屑之教诲，亦教诲之}{《孟子·告子下》，维基文库本；朱熹《四书章句集注·孟子集注》卷十二，维基文库本，2026 年访问 \url{https://zh.wikisource.org/wiki/\%E5\%AD\%9F\%E5\%AD\%90/\%E5\%91\%8A\%E5\%AD\%90\%E4\%B8\%8B}}{「教亦多術矣。予不屑之教誨也者，是亦教誨之而已矣。」}{2026-10-11（维基文库）}{已核实}{正驳15；正口径·案例}{}
\source{233}{郑人：宁信度，无自信也}{《韩非子·外储说左上》，维基文库本，2026 年访问 \url{https://zh.wikisource.org/wiki/\%E9\%9F\%93\%E9\%9D\%9E\%E5\%AD\%90/\%E5\%A4\%96\%E5\%84\%B2\%E8\%AA\%AA\%E5\%B7\%A6\%E4\%B8\%8A}}{「人曰：『何不試之以足？』曰：『寧信度，無自信也。』」}{2026-10-11（维基文库）}{已核实}{正口径·案例}{}
\source{234}{知情旁听者不高估说话者的效果}{Keysar \& Henly 2002, “Speakers' overestimation of their effectiveness”, Psychol Sci 13(3): 207–212, doi:10.1111/1467-9280.00439}{“In contrast, overhearers who were informed about the speakers intention did not systematically overestimate the speakers' effectiveness.”}{2026-10-11（摘要）}{已核实}{反驳2、12、31；反口径·数11}{}
\source{235}{Denrell 模型里没有情绪变量}{Denrell 2005, “Why most people disapprove of me: Experience sampling in impression formation”, Psychol Rev 112(4): 951–978, doi:10.1037/0033-295X.112.4.951}{“The underlying mechanism is the sample bias generated when the probability of interaction depends on current impressions.”}{2026-10-11（摘要）}{已核实}{反驳11；反口径·数12}{}
\source{236}{模拟中「偏见」定义为对好对象的错误负面预期}{Eiser 等, 2009, “Prejudiced learning: A connectionist account”, Br J Psychol 100(Pt 2): 399–413, doi:10.1348/000712608X357849}{“under conditions where learners can only discover the valence of attitude objects through directly experiencing them.”}{2026-10-11（摘要）}{已核实}{反驳11；反口径·数13}{}
\source{237}{扮蠢的好处：最难的部分不分给你}{微信公众号「真故研究室」，2024-07-07，《蠢人设，开始侵蚀职场》 \url{https://m.huxiu.com/article/3218080.html}}{「首先小组作业时，“大家不会把最难的部分分给你，有事也不会首先想到你”；其次是在寝……」}{2026-10-11（原页）}{已核实}{正驳16、42；反驳5；正口径·案例；反口径·案15}{}
\source{238}{两代人都说难在操作太复杂}{川观新闻，2020-10-21 15:31，《2690份问卷画像：超三成老人遭遇“智能机”困境，网络挂号就医难倒“脱网”老人》 \url{https://cbgc.scol.com.cn/news/416034}}{「年轻人和老人均认为老人学习使用智能机的最佳方式是通过家人和朋友的指导，受访年轻人……」}{2026-10-11（原页）}{已核实}{反驳14；反口径·数24}{}
\source{239}{Aczél 摘要：把行为叫作蠢反映理性期待（作者动机）}{Aczél 等 2015, “What is stupid? People's conception of unintelligent behavior”, Intelligence 53: 51–58, doi:10.1016/j.intell.2015.08.010；比例与样本只见二手报道（BPS Research Digest、444.hu 等）}{“but also because it is a robust behavioral reflection of the rationalistic expectations to which people adjust their own behavior and expect others to.”}{2026-10-11（摘要）}{已核实}{反淘；正驳25；反驳8、25；正口径·数据；反口径·论10禁、数14}{}
\source{240}{越少见的行为在评价里分量越重（80 人、23 人）}{Mende-Siedlecki 等, 2013, “Diagnostic value underlies asymmetric updating of impressions in the morality and…”, J Neurosci 33(50): 19406–19415, doi:10.1523/JNEUROSCI.2334-13.2013}{“While positive behavioral information is diagnostic when evaluating a person's abilities, negative information is diagnostic when evaluating morality.”}{2026-10-11（摘要）}{已核实}{正驳16；正口径·数据}{}
\source{241}{与自己过去一致的结果归于能力和难度（2 项实验）}{Frieze \& Weiner 1971, “Cue utilization and attributional judgments for success and failure”, J Pers 39(4): 591–605, doi:10.1111/j.1467-6494.1971.tb00065.x}{“consistency with one's own past performance is ascribed to ability and task difficulty, while inconsistent outcomes give rise to luck and effort attributions”}{2026-10-11（摘要）}{已核实}{正论一、反论一、反淘；正驳20、21、35；反驳6、26、34、38；正口径·数据；反口径·论1、论3、数10}{}
\source{242}{要求解释细节机制削弱解释深度错觉}{Fernbach 等, 2013, “Political extremism is supported by an illusion of understanding”, Psychol Sci 24(6): 939–946, doi:10.1177/0956797612464058}{“Asking people to explain policies in detail both undermined the illusion of explanatory depth and led to attitudes that were more moderate (Experiments 1 and 2).”}{2026-10-11（摘要）}{已核实}{正驳20}{}
\source{243}{新京报：现实中厌蠢更多冲着拖后腿的同学同事}{新京报书评周刊，2023-11-01，《为什么我们容易犯“厌蠢症”？》 \url{https://m.bjnews.com.cn/detail/1698812159168656.html}}{「……“拖后腿”的同学、同事的不满」}{2026-10-11（原页）}{已核实}{反论一；正驳15；反驳1、2、9、14、15、32、34、35；正口径·案例；反口径·论1、回1、案4}{}
\source{244}{愤怒属趋近型动机，焦虑属回避型}{Carver \& Harmon-Jones 2009, “Anger is an approach-related affect: Evidence and implications”, Psychol Bull 135(2): 183–204, doi:10.1037/a0013965}{“The evidence supports the view that anger relates to an appetitive or approach motivational system, whereas anxiety relates to an aversive or avoidance motivational system.”}{2026-10-11（摘要）}{已核实}{正驳14；正口径·数据}{}
\source{245}{朱熹：未可以一事之失而遽忽之也}{《论语·先进第十一》十一之十四，维基文库本；朱熹《四书章句集注·论语集注》卷六，维基文库本，2026 年访问 \url{https://zh.wikisource.org/wiki/\%E8\%AB\%96\%E8\%AA\%9E/\%E5\%85\%88\%E9\%80\%B2\%E7\%AC\%AC\%E5\%8D\%81\%E4\%B8\%80}}{「未可以一事之失而遽忽之也」}{2026-10-11（维基文库）}{已核实}{正驳18；正口径·案例}{}
\source{246}{后见之明元分析：95 项研究，Md=.39}{Guilbault 等, 2004, “A meta-analysis of research on hindsight bias”, BASP 26(2–3): 103–117, doi:10.1080/01973533.2004.9646399}{“When excluding missing effect sizes, the overall mean effect size was Md = .39 with a 95\% confidence interval of .36 to .42.”}{2026-10-11（摘要）}{已核实}{正论一；正驳19；反驳43；正口径·论点、数据；反口径·数25、六·补核、七}{}
\source{247}{朋霍费尔：理由不起作用}{Bonhoeffer〈Nach zehn Jahren〉「Von der Dummheit」，Widerstand und Ergebung, Chr. Kaiser 1951, S. 17–20；Teuffel 全文转录 \url{https://jochenteuffel.com/wp-content/uploads/2018/06/bonhoeffer-nach-zehn-jahren.pdf}}{「Gründe verfangen nicht」}{2026-10-11（转录＊）}{已核实}{正驳27；反驳20；正口径·案例；反口径·案8、六·补核}{}
\source{248}{小刘工作超过 2 年仍常犯错，劝告后再犯，关系一度冰点}{每日商报，2024-01-13 09:28，《“怎么还会有这种人”“我的厌蠢症又犯了”专家：“厌蠢”的实质是一种“慕强”心态》，第 2 页 \url{https://jrsh.hangzhou.com.cn/read/content/2024-01/13/content\_8674352\_2.htm}}{「小刘已经工作超过2年了，但是他还是经常犯一些常规基础性的错误。」}{2026-10-11（原页）}{已核实}{正论二、反论一；正驳4–6、8、10、11、13–15、17、19；反驳1、3、4、29、34、35、40；正口径·论点、案例；反口径·论1、回6、案1}{}
\source{249}{新京报：社交媒体最常见的吐槽对象}{新京报书评周刊，2023-11-01，《为什么我们容易犯“厌蠢症”？》 \url{https://m.bjnews.com.cn/detail/1698812159168656.html}}{「诸如此类，都是社交媒体中最常见的“厌蠢症”吐槽对象」}{2026-10-12（原页）}{已核实}{反论一；正驳1；反驳1；正口径·案例}{}
\source{250}{大人生气，孩子更多归因没用心；同情归因能力低}{Graham 1984, “Communicating sympathy and anger to Black and White children: The cognitive…”, JPSP 47(1): 40–54, doi:10.1037/0022-3514.47.1.40}{“In addition, a greater decline in expectancy for success followed the emotional cue of sympathy.”}{2026-10-12（摘要）}{已核实}{正驳13；反驳3、4、16、20、39、41；正口径·数据；反口径·论9、回4、数3、六·补核}{}
\source{251}{提高技能后才认出自己能力的局限}{Kruger \& Dunning 1999, “Unskilled and unaware of it: How difficulties in recognizing one's own incompetence…”, JPSP 77(6): 1121–1134, doi:10.1037/0022-3514.77.6.1121}{“Several analyses linked this miscalibration to deficits in metacognitive skill, or the capacity to distinguish accuracy from error.”}{2026-10-12（摘要）}{已核实}{正驳8；正口径·数据}{}
\source{252}{新京报：对蠢的厌恶更多源自可能的后果}{新京报书评周刊，2023-11-01，《为什么我们容易犯“厌蠢症”？》 \url{https://m.bjnews.com.cn/detail/1698812159168656.html}}{「……人要怎么办”，很多时候，我们对“蠢”的厌恶和批判，更多源自于“蠢”可能造成的后果……」}{2026-10-12（原页）}{已核实}{反驳1、7、16；反口径·论1、案4}{}
\source{253}{知识的诅咒：自己知道结果，推断别人信念受损}{Birch \& Bloom 2007, “The curse of knowledge in reasoning about false beliefs”, Psychol Sci 18(5): 382–386, doi:10.1111/j.1467-9280.2007.01909.x}{“That is, adults' own knowledge of an event's outcome can compromise their ability to reason about another person's beliefs about that event.”}{2026-10-11（摘要）}{已核实}{反驳31；反口径·数41}{}
\source{254}{问答者实验：出题者因角色优势被评为更有学识（2 个实验）}{Ross, Amabile \& Steinmetz 1977, Social roles, social control, and biases in social-perception processes, J Pers Soc Psychol 35(7): 485–494, DOI 10.1037/0022-3514.35.7.485}{“Questioners, allowed to display their personal store of esoteric knowledge in composing questions, were consistently rated superior to their partners, who attempted to answer the questions.”}{2026-10-11（摘要）}{已核实}{反驳32；反口径·数42}{}
\source{255}{教师期望预测学生表现，可能更多因为期望准确}{Jussim \& Harber 2005, Teacher expectations and self-fulfilling prophecies: Knowns and unknowns, resolved and unresolved controversies, Pers Soc Psychol Rev 9(2): 131–155, DOI 10.1207/s15327957pspr0902\_3}{“(d) teacher expectations may predict student outcomes more because these expectations are accurate than because they are self-fulfilling.”}{2026-10-11（摘要）}{已核实}{正驳45；正口径·数据}{}
\source{256}{实体论者从行为推出更强的特质判断（5 项研究）}{Chiu, Hong \& Dweck 1997, “Lay dispositionism and implicit theories of personality”, JPSP 73(1): 19–30, doi:10.1037//0022-3514.73.1.19}{“those who believed in fixed traits (entity theorists) used traits or trait-relevant information to make stronger future behavioral predictions (Studies 1 and 2) and made stronger trait inferences from behavior (Study 3)”}{2026-10-11（摘要）}{已核实}{反驳26、41；反口径·数40}{}
\source{257}{课堂自证预言通常很小、不累积；受污名群体可能更强}{Jussim \& Harber 2005, Teacher expectations and self-fulfilling prophecies: Knowns and unknowns, resolved and unresolved controversies, Pers Soc Psychol Rev 9(2): 131–155, DOI 10.1207/s15327957pspr0902\_3}{“Self-fulfilling prophecies in the classroom do occur, but these effects are typically small, they do not accumulate greatly across perceivers or over time, and they may be more likely to dissipate than accumulate”}{2026-10-11（摘要）}{已核实}{正驳45；正口径·数据}{}
\end{sourcelist}
\begin{fields}
\field{三类计数}{已核实 245 条，有把握 5 条，【待核实】7 条。}
\field{整条有把握（5 条）}{\#69（Aczél 2015 的样本与比例）、\#90（基本归因错误定义原句）、\#103（医疗差错死亡数）、\#141（费曼在巴西）、\#144（陶行知诗）。文稿、速查都不进，查证方向见各条。}
\field{整条【待核实】（7 条）}{\#7、\#68（《现代汉语词典》）、\#96（Thompson 2014）、\#102（厌蠢比例检索未果）、\#134（「可以原谅坏人」的出处）、\#166（Mason 2003）、\#178（Ross 与 Ward 1996）。文稿、速查都不进。}
\field{带「＊」的条目}{\sub{有把握的部分}{Strawson 原刊卷页（\#32、\#33、\#34、\#174、\#197），只说「Strawson 1962 年的《自由与怨恨》」；亚里士多德 Gutenberg 本译者与 Bekker 页码（\#39、\#40、\#41、\#230），只说卷次；Darwall 1977 卷期页（\#42），只说「据《斯坦福哲学百科》」；Rawls 的节号（\#52）；蔡康永那一句（\#111 局限栏），改用高晓松「坏人不是蠢人」。}\sub{【待核实】的部分}{「Gottman 发现轻蔑是离婚最佳预测因子」（\#23）；陌生人智力评分的相关系数（\#88 局限栏）；网传「坏苹果」拖低绩效的比例（\#65）；评论里坐飞机人数、通地铁城市数与老人用手机比例（\#10、\#15、\#118）；「智商税」词源（\#125）；Rawls「努力的意愿也受家庭环境影响」（\#52）；Edmondson「心理安全与高标准并存」（\#64）。都不说。}\sub{核实范围的限定}{康德两句（\#45、\#48）和朋霍费尔各句（\#74、\#75、\#76、\#211、\#247）中文为自译，赛前对照李秋零译本与校勘本，之前只说「据德文原文的意思」；\#53 只核到博客节选；\#181 核实的是「报道这样写」；\#133 只核「流传版这样写」；\#18 的自媒体部分只作「有人认为」；\#180 赛前对照国家法律法规数据库；\#59 的 2000 年论文只核书目。}}
\end{fields}

\section{十二、备赛建议}
\begin{fields}
\begin{timeline}
\when{周一}{辩题公布当天：全队对齐三句口播原话（厌蠢、蠢、判准）和门槛分界规则；正反两边按口径表分工，赛前待办分人认领（康德李秋零译本、《现代汉语词典》纸本、Hom 2023 方法部分）。}
\when{周二}{稿件成型：正一立论定稿背熟，其余成段讲稿按套件写齐条件模块；速查与场上记录卡同步；所有数据对照出处清单，只用已核实的。}
\when{周三}{第一次模辩，正反各打一遍，复盘「平常」之争谁描述得更像；24 时前确认参赛（弃赛与改期截止）。}
\when{周四至周五}{第二次模辩，专练奇袭窗口④与四辩对辩；赛前待办落实后同步改口径表与稿件；向组委会确认盘问时被盘问方回答是否计时。}
\when{赛前一天}{全队背同一句；照禁止说法与滑坡清单互相抽查；每人能当场报出所用数据的机构、年份、报告名。}
\end{timeline}
\field{模辩重点检验}{\sub{①「平常」之争}{正方能否在没有比例数据的前提下，用机制与典型情形撑住「平常」，并在被问「回头不讲理是不是推论」时坦白认、给依据；反方能否把「平常」钉在熟人之间、记录在手，同时不把「先问」说成立论许的愿。}\sub{②口径滑坡}{正方会不会滑到「厌蠢就是霸凌」「别有感觉」「屡次也救不了」；反方会不会滑到「情绪无对错」「我有权远离」「那一份我们认错」；双方会不会把推论说成来源、把描述说成统计、把〔智〕材料直接用在「蠢」上。}\sub{③奇袭与末位}{反方在自由辩之后发动奇袭质询，正四能否在结辩里正面回答「平常还有什么材料」「回头不讲理是不是推论」，又不被拖掉最后一发的价值收束；反四封路式结辩能否在 210 秒里封完正方最可能的四发。}}
\begin{checklist}
\item 厌蠢中立定义口播原话：「厌蠢，是觉得别人犯蠢时生出的嫌恶、反感，常带着想躲开的心；它可以只冲着那件事，也可以冲着这个人、看低他；它是一种态度，不等于骂人或欺负人。」
\item 蠢中立定义口播原话：「蠢，是在一般人看来本该想得到、做得对的事上，明显想错、做错；可能出于不会，也可能出于不肯，原因不在定义里预设。」
\item 判准口播原话（与「错」同一句）：「判断这份厌对不对，就看一个清醒的我问完三句：看准了吗？对他公道吗？划得来吗？问完觉得这份厌平常的样子就该改，就是错了；站得住，只在个别场合调调分寸，就是没错。」
\item 门槛分界规则：「要改的如果是这份厌平常的样子，就是这份厌本身该改，哪怕改完还剩一点反感；要改的如果只在个别场合冒出来，那叫调分寸。平常就这样、只是改得掉，照样算错了。」
\item 「你说的是哪个我」：「一个普通的、会说出这句话的我。」「只在心里厌，算不算」：「算，心里那份就是厌蠢本身；它错不错，看判准。」
\item 越线按时间划：这回先不说、一时走开、不再主动提醒，不越线；回头他开口也不听、合作只下指令不给理由、认定「他这个人说什么都没用」，越线。
\item 正方：两个论点「拿我当一般人」「说了也没用」各一句话；「对屡教不改的怎么办」的答法；反二取证五问答「不争；不争；不反对；不算；是」，再补「我们争的是两处：没有记录、凭比较就判的那一下，和看准之后回头还讲不讲理」。
\item 反方：两个论点「教过了还犯」「恼他不弃他」各一句话；分格问题「这个例子里，他被教过没有？」；「厌了之后交涉还是走开」的答法；看准了这个人蠢能不能看低他：「看低他的能力可以，拒绝给他作为人的尊重不行；我们的厌停在前一边。」
\item 双方：不报任何厌蠢比例；〔智〕先搭桥；每日商报一案三句一起说、不标说话人；康德整句引；推论被问就认。
\end{checklist}
\end{fields}
\end{document}
